% !TeX root = ../Book1.tex
% 纲要标题：# 第二编　环、整除与多项式
% 纲要位置：计划纲要.md，第 391 行。

\BookPart{环、整除与多项式}{b1:part:02}
\BookPartGuide
在 \(\mathbb Z/6\mathbb Z\) 中，非零元素 \(\bar2,\bar3\) 的乘积却是 \(\bar0\)；
在 \(M_2(\mathbb Q)\) 中，两矩阵的乘积可能依赖相乘的次序。
把乘法加入加法结构以后，零因子、单位与非交换性便成为新的问题。
理想使同余关系同时与加法和乘法相容，整除理论研究哪些因子分解规律仍然成立；
多项式环与局部化则分别通过加入变量和指定逆元构造新环。

\ProseParagraphBreak
\chapref{b1:ch:10}从整数环、多项式环和矩阵环出发；矩阵乘法提醒我们，
环的乘法不必交换，形成商环所需的理想必须吸收两侧的乘法。
在交换环中，素理想和极大理想可分别通过整环商和域商识别；
中国剩余定理则把互素理想给出的同余条件，转化为商环之间的直积同构。

\ProseParagraphBreak
\chapref{b1:ch:11}考察整数中的整除理论能够推广多远。
Euclid 整环的带余除法给出计算最大公因子的算法；
主理想整环把每个理想归结为一个生成元，并使最大公因子能够写成线性组合；
唯一分解整环则控制元素的分解及不可约因子的重数。
这些性质有强弱之分：能够唯一分解，并不意味着能够带余除法；
在一般整环中，不可约元也未必是素元。

\ProseParagraphBreak
\chapref{b1:ch:12}把除法、导数和插值用于多项式计算。
Gauss 引理\footnote{高斯（Johann Carl Friedrich Gauss，1777-1855），德国数学家，在数论、分析与微分几何中作出重要贡献，著有《算术研究》。}把唯一分解整环上的多项式分解与其分式域上的分解联系起来，
Eisenstein 判据\footnote{艾森斯坦（Ferdinand Gotthold Max Eisenstein，1823-1852），德国数学家，研究二次、三次形式与高次互反律。}和模素数约化提供具体的不可约性检验。
多元情形中，本章讨论对称多项式，并用 Hilbert 基定理\footnote{希尔伯特（David Hilbert，1862-1943），德国数学家，在不变量论、几何基础与数学基础等领域有重要成果。}证明：
若系数环是 Noether 环\footnote{诺特（Amalie Emmy Noether，1882-1935），德国数学家，发展抽象环论和理想理论，也揭示了对称性与守恒律之间的联系。}、变量只有有限个，所得多项式环的每个理想
都可由有限多个多项式生成。
因此，即使一个理想收集了无限多个多项式关系，也可用有限多个关系生成它。
这一有限性为后续交换代数与\appref{b1:app:F}的 Gröbner 基计算提供结构依据，
但 Hilbert 基定理本身还没有给出寻找生成元的算法。

\ProseParagraphBreak
第十二章为 Gauss 引理构造了整环的分式域，使所有非零元都可逆。
如果在交换环 \(R\) 中只想让一个元素 \(f\) 可逆，或者只想让某个素理想 \(\mathfrak p\)
之外的元素可逆，该怎样构造？\chapref{b1:ch:13}用局部化分别得到
\(R[1/f]\) 与 \(R_{\mathfrak p}\)，并把分式域纳入同一种构造。
若原环有零因子，
分式的相等条件便不能简单地通过交叉相乘表达，
原环到局部化环的映射也可能把非零元送到零。
素理想处的局部环使分母避开该素理想，
理想的扩张与收缩则说明局部化前后的理想如何对应。

\begin{center}
\begin{tikzcd}[column sep=large,row sep=large]
\begin{gathered}\text{第10章}\\\text{环、理想与商}\end{gathered}
\arrow[r]\arrow[dr]
&\begin{gathered}\text{第11章}\\\text{整除与唯一分解}\end{gathered}\arrow[r]
&\begin{gathered}\text{第12章}\\\text{多项式与有限性}\end{gathered}\arrow[dl,dashed]\\
&\begin{gathered}\text{第13章}\\\text{分式与局部化}\end{gathered}&
\end{tikzcd}
\end{center}

实线表示主要依赖，虚线表示第13章回顾第12章为 Gauss 引理构造的分式域。
初读第12章时，可以把代数基本定理作为已知结论使用；
完成第19章的有限 Galois 理论以后，再回读\appref{b1:app:E}中结合 Sylow 理论给出的证明。

\ProseParagraphBreak
若希望先学习有限 Galois 理论，需掌握第10章的环与商环、素理想与极大理想，
第11章的唯一分解与不可约元，以及第12章的域上多项式、Gauss 引理和不可约性检验。
分式域的首次完整构造是\subsecref{b1:subsec:1202:02}中的\cref{b1:prop:fraction-field-basic}，
这也是第三编所需的基础。第14章的函数域说明还用到\cref{b1:prop:fraction-in-ambient-field}，
把分式域识别为已有域中包含原环的最小子域，可从\secref{b1:sec:1301}补读。
第13章的一般局部化、理想对应与局部判据可以暂缓。
学习交换环及后续交换代数时，则应继续研究这些局部性质。

\ProseParagraphBreak
本编使用交换群、同态、商结构、整数算术与基本线性代数。
交换性会影响理想和因子分解的结论，
而原环、商环与局部化环中的等式也不能混为一谈。


% !TeX root = ../../Book1.tex
% 纲要标题：## 第10章　环、理想与中国剩余定理
% 纲要位置：计划纲要.md，第 386 行。

\chapter{环、理想与中国剩余定理}
\label{b1:ch:10}
\BookChapterNumber{b1:ch:10}

\begin{chapterabstract}
环把加法、乘法及分配律置于同一结构中；理想使这两种运算能够一起下降到商环。本章研究同态与理想、素理想和极大理想，以及理想间的运算，最后以中国剩余定理说明互素条件下的商环分解。
\end{chapterabstract}

\begin{learninggoals}
\begin{itemize}
\item 明确含幺环、保单位同态与子环的约定，区别单位、零因子和幂零元。
\item 判断左右理想与双边理想，构造商环并准确计算同态的核和像。
\item 用商环判断素理想与极大理想，比较理想的和、积、交、根与理想商。
\item 显式求解同余方程，并由正交幂等元识别环的直积分解。
\end{itemize}
\end{learninggoals}

在整数中，\(2\cdot3=6\) 不是零；把计算放到模 \(6\) 的剩余类里，它却变成零。这并不妨碍剩余类继续相加、相乘，但它说明从整数继承运算之后，乘法的性质可能改变。若要系统地问哪些元素能取逆、哪些乘积会成为零，以及哪些同余关系允许同时做加法和乘法，就需要把两种运算放进同一个结构中。理想控制可用于取商的同余关系，中国剩余定理则说明某些模数上的计算可以怎样拆开并重新合成。

\ProseParagraphBreak

本章使用第一编的交换群、同态、商群与整数子群的分类；矩阵和线性映射只需基本线性代数。环的乘法起初可以不交换，进入交换环的结论时会另行说明。一般极大理想的存在性证明使用 Zorn 引理\footnote{佐恩（Max August Zorn，1906-1993），德国数学家，研究代数并提出与选择公理等价的极大元引理。}。涉及准素理想的习题使用\cref{b1:def:primary-ideal}。

\ProseParagraphBreak
环的定义、基本例子与同态可对照席南华《基础代数》第一卷\cite[第5章第5.3节]{Xi2016BasicAlgebraI}。交换环中的理想、素理想、极大理想与商构造，可继续读 Atiyah 与 Macdonald 的《Introduction to Commutative Algebra》\cite[Chapter 1]{AtiyahMacdonald1969CommutativeAlgebra}；中国剩余定理的整数计算则可对照 Shoup 的《A Computational Introduction to Number Theory and Algebra》\cite[Section 4.3]{Shoup2008NumberTheoryAlgebra}。

% !TeX root = ../../Book1.tex
% 纲要标题：### 10.1 环的约定与例子
% 纲要位置：计划纲要.md，第 388 行。

\section{环的约定与例子}
\label{b1:sec:1001}

群只规定一种运算。整数与矩阵却同时允许加法和乘法，而且乘法能够分配到加法中。环的定义把这两种运算放在一起；它也使零因子、方程的可解性和同余关系进入同一套语言。本节先固定单位元及同态的约定，再比较几类常见例子。

% 纲要内容（仅作写作计划，尚非教材正文）：
%
% * 含幺环、交换环及同态的统一约定
% * 矩阵环、函数环、有限直积环、多项式环与群环
% * 零因子、单位与幂零元
%

\subsection{含幺环、交换环与同态}
\label{b1:subsec:1001:01}
\begin{definition}\label{b1:def:ring}
给定集合 \(R\) 及加法、乘法两个从 \(R\times R\) 到 \(R\) 的运算。如果它们满足下列条件，就称 \(R\) 连同这两种运算为\term[huan]{环}[ring]：
\begin{enumerate}
\item 加法结合且交换，存在零元 \(0\)，每个 \(a\in R\) 有加法逆元 \(-a\)；
\item 乘法结合，存在单位元 \(1\)，使 \(1a=a1=a\)；
\item 对任意 \(a,b,c\in R\)，有 \(a(b+c)=ab+ac\)、\((a+b)c=ac+bc\)。
\end{enumerate}
如果还对任意 \(a,b\in R\) 有 \(ab=ba\)，就称 \(R\) 为\term[jiaohuan-huan]{交换环}[commutative ring]。本书环总含幺，但允许 \(1=0\)；这时 \(R\) 只有一个元素，称为\term[linghuan]{零环}[zero ring]。
\end{definition}
加法公理正是说 \((R,+)\) 为交换群。由 \(a0=a(0+0)=a0+a0\)，在加法群中消去 \(a0\)，得到 \(a0=0\)；同样 \(0a=0\)。再由 \(a(b-b)=0\)，得 \(a(-b)=-(ab)\)，继而 \((-a)(-b)=ab\)。这些等式来自分配律，而不是额外公理。若 \(1=0\)，每个元素 \(a=a1=a0=0\)，说明零环的确只有一个元素。

\ProseParagraphBreak
允许零环会影响后面的边界条件：下面定义的整环、除环、域都要求非零，即 \(1\ne0\)。\secref{b1:sec:1002}将定义理想与真理想；零环的唯一理想就是整个环，因而没有真理想，极大理想的存在性也须在非零环中讨论。

\begin{definition}\label{b1:def:ring-homomorphism}
设 \(R,S\) 为环，\(f:R\rightarrow S\) 为映射。如果对任意 \(a,b\in R\) 都有
\[
f(a+b)=f(a)+f(b),\qquad f(ab)=f(a)f(b),\qquad f(1_R)=1_S,
\]
就称 \(f\) 为从 \(R\) 到 \(S\) 的\term[huan-tongtai]{环同态}[ring homomorphism]。如果环同态 \(f\) 还是双射，就称它为\term[huan-tonggou]{环同构}[ring isomorphism]。如果子集 \(T\subseteq R\) 含 \(1_R\)，并且对减法和乘法封闭，就称 \(T\) 为本书约定下的\term[zihuan]{子环}[subring]。
\end{definition}
同态保持零元和负元，因为它首先是加法群同态。双射环同态的逆映射也保持两种运算与单位：将等式写在原像上再应用 \(f\)，便逐项得到这些性质。有些文献允许环同态不保持单位；本书始终要求 \(f(1_R)=1_S\)，子环也须含母环的同一个单位。按这一约定，任意环到零环的唯一映射都是环同态，而零环到非零环没有环同态，因为它的同一个元素 \(0=1\) 不可能同时映到不同的 \(0\) 与 \(1\)。

\ProseParagraphBreak
要求保持单位使 \(\mathbb Z\rightarrow R\)、\(n\mapsto n1_R\) 成为唯一的环同态；这里正整数倍由重复加法定义，负数倍取其加法逆元。它说明环中的整数从何而来：像由全部 \(n1_R\) 组成，对减法和乘法封闭，并且每个含 \(1_R\) 的子环都必须包含它，所以它是由单位生成的最小子环。

\begin{definition}\label{b1:def:ring-characteristic}
设 \(R\) 为环。若存在正整数 \(n\) 使 \(n1_R=0\)，其中最小者称为 \(R\) 的\term[tezheng]{特征}[characteristic]；若不存在，规定特征为零。环 \(R\) 的特征记为 \(\operatorname{char}R\)。
\end{definition}
\symidx{\(\operatorname{char}R\)：环的特征}
同态 \(\mathbb Z\rightarrow R\) 的核是加法子群，\cref{b1:lem:integer-subgroups}说明它唯一等于 \(n\mathbb Z\)，其中 \(n\geq0\)；这与特征定义中的 \(n\) 相同。特征为零时，这个同态单射，所以由单位生成的子环与 \(\mathbb Z\) 同构；特征为正整数 \(n\) 时，对整数 \(a,b\) 有 \(a1_R=b1_R\) 当且仅当 \(a-b\in n\mathbb Z\)，故 \([a]_n\mapsto a1_R\) 给出从整数剩余类环 \(\mathbb Z/n\mathbb Z\) 到该子环的同构。它保持乘法是因为 \((a1_R)(b1_R)=(ab)1_R\)。特征不是集合的元素个数：例如 \(\mathbb Z\) 和 \(\mathbb Q\) 都具有特征零。零环的特征为一，对应 \(\mathbb Z/\mathbb Z\)。

\begin{definition}\label{b1:def:domain-field}
如果非零交换环 \(R\) 对任意 \(a,b\in R\) 都满足 \(ab=0\) 必有 \(a=0\) 或 \(b=0\)，就称 \(R\) 为\term[zhenghuan]{整环}[integral domain]。如果非零环 \(R\) 的每个非零元素都有双边乘法逆元，就称 \(R\) 为\term[chuhuan]{除环}[division ring]；乘法交换的除环称为\term[yu]{域}[field]。
\end{definition}
域是整环：若 \(a\ne0\) 且 \(ab=0\)，左乘 \(a^{-1}\) 得 \(b=0\)。\(\mathbb Z\) 是整环而不是域，\(\mathbb Q,\mathbb R,\mathbb C\) 是域。整环的特征只能为零或素数：若正特征 \(n=uv\)，\(1<u,v<n\)，则 \((u1_R)(v1_R)=0\)，而由正特征的最小性，两个因子都非零，与整环无零因子矛盾。

\ProseParagraphBreak
这几类环的条件可以按下面的蕴含关系比较：
\[
\begin{tikzcd}[column sep=2em,row sep=1.8em]
\text{域}\arrow[r,Rightarrow]\arrow[d,Rightarrow]
&\text{整环}\arrow[r,Rightarrow]
&\text{交换非零环}\arrow[d,Rightarrow]\\
\text{除环}\arrow[rr,Rightarrow]&&\text{非零环}
\end{tikzcd}
\]
域恰好是乘法交换的除环。整环只排除两个非零元素相乘为零的情形，不要求非零元素都可逆；除环要求每个非零元素都有双边逆，但不要求乘法交换。

\subsection{矩阵环、函数环、多项式环与群环}
\label{b1:subsec:1001:02}
设 \(R\) 为环，\(n\geq1\) 为整数。所有 \(n\times n\) 矩阵在逐项加法与通常的矩阵乘法下组成\term[juzhen-huan]{矩阵环}[matrix ring] \(M_n(R)\)，单位为单位矩阵。乘法结合律可由有限和核验：两种结合方式的第 \((i,j)\) 项都是 \(\sum_{k,\ell}a_{ik}b_{k\ell}c_{\ell j}\)。若 \(R\ne0,n\geq2\)，则矩阵单位满足 \(E_{12}E_{21}=E_{11}\)、\(E_{21}E_{12}=E_{22}\)，故即使 \(R\) 交换，矩阵环也不交换。这里 \(n\ge2\) 不可省略：取唯一矩阵项给出 \(M_1(R)\cong R\)，所以 \(1\times1\) 矩阵环保留系数环原来的交换性。
\symidx{\(M_n(R)\)：环上的 \(n\) 阶矩阵环}

\ProseParagraphBreak
设 \(R\) 为环，\(X\) 为集合。映射集合 \(R^X\) 在逐点加法和乘法下构成\term[hanshu-huan]{函数环}[function ring]，其零元与单位分别为常值函数 \(0,1\)。每条环公理都可在每个 \(x\in X\) 处验证。若 \(R\) 为域而 \(X\) 有两个不同的点 \(x,y\)，定义
\[
f(z)=\begin{cases}1,&z=x,\\0,&z\ne x\end{cases},
\qquad
g(z)=\begin{cases}1,&z=y,\\0,&z\ne y\end{cases}.
\]
两函数分别只在一个点非零，所以 \(f\ne0\)、\(g\ne0\)；因为 \(x\ne y\)，每个点处至少有一个函数取零，故 \(fg=0\)。因此，系数环是域也不能保证函数环是整环。

\ProseParagraphBreak
设 \(R_1,\ldots,R_r\) 为环，其中 \(r\geq1\)。在笛卡尔积 \(\prod_{i=1}^rR_i\) 上逐坐标定义加法与乘法：
\[
(a_i)_i+(b_i)_i=(a_i+b_i)_i,\qquad
(a_i)_i(b_i)_i=(a_ib_i)_i.
\]
零元和单位分别是 \((0_{R_i})_i\) 与 \((1_{R_i})_i\)。环公理逐坐标成立，所得环称为 \(R_1,\ldots,R_r\) 的\term[huan-zhiji]{直积环}[direct product of rings]；对每个 \(1\leq j\leq r\)，坐标投影 \(\prod_iR_i\to R_j\) 都是保单位的环同态。前述函数环 \(R^X\) 则是各坐标均取 \(R\)、指标集可任意选取的情形。

\ProseParagraphBreak
本章关于多项式只需建立形式和的运算、次数以及整环上的次数加法。除法算法、重根与不可约性将在\chapref{b1:ch:12}系统讨论；这里先把多项式环作为一种可以直接核验的环来构造。

\begin{definition}\label{b1:def:polynomial-ring}
设 \(R\) 为环，\(x\) 为与 \(R\) 中每个元素交换的形式符号。考虑系数 \(a_i\in R\) 且只有有限项非零的形式和
\(f=\sum_{i\geq0}a_ix^i\)，并规定两个形式和相等当且仅当对应系数逐项相等。再规定运算
\[
\sum_i a_ix^i+\sum_i b_ix^i=\sum_i(a_i+b_i)x^i,
\qquad
\Bigl(\sum_i a_ix^i\Bigr)\Bigl(\sum_j b_jx^j\Bigr)
=\sum_k\Bigl(\sum_{i+j=k}a_ib_j\Bigr)x^k.
\]
则称这些形式和组成的环为 \(R\) 上的\term[duoxiangshi-huan]{多项式环}[polynomial ring]，记为 \(R[x]\)。如果 \(R\) 不交换，乘积中的系数仍按所示次序相乘。
\end{definition}
\symidx{\(R[x]\)：一元多项式环}
这里的 \(x\) 是与全部系数交换的中心变量；这只规定 \(ax=xa\)，没有规定任意系数 \(a,b\) 也交换。例如
\[
(ax)b=a(xb)=a(bx)=(ab)x,
\]
所得系数仍按 \(a,b\) 的顺序相乘，不能改为 \(ba\)。

\ProseParagraphBreak
系数支集有限使上述各和有限。两次乘法的第 \(k\) 个系数无论如何加括号，都等于 \(\sum_{i+j+\ell=k}a_ib_jc_\ell\)，故乘法结合；分配律来自系数环，常数多项式 \(1\) 为单位。这验证了定义确实给出环。设 \(R\) 为环，\(f=\sum_{i\geq0}a_ix^i\in R[x]\) 为非零多项式，并令 \(n\) 为满足 \(a_n\ne0\) 的最大指标。整数 \(n\) 称为 \(f\) 的\term[duoxiangshi-cishu]{次数}[degree]，记为 \(\deg f\)；系数 \(a_n\) 称为 \(f\) 的\term[shouxiang-xishu]{首项系数}[leading coefficient]。如果 \(a_n=1\)，就称 \(f\) 为\term[shouyi-duoxiangshi]{首一多项式}[monic polynomial]。零多项式的次数记为 \(-\infty\)。
\symidx{\(\deg f\)：多项式的次数}

\begin{proposition}\label{b1:prop:polynomial-degree}
若 \(R\) 是整环，\(f,g\in R[x]\) 非零，则 \(fg\ne0\)，且
\(\deg(fg)=\deg f+\deg g\)。特别地，\(R[x]\) 是整环。
\end{proposition}
\begin{proof}
设 \(f,g\) 的次数为 \(m,n\)，首项系数为 \(a_m,b_n\)。乘积没有次数超过 \(m+n\) 的项，\(x^{m+n}\) 的系数是非零的 \(a_mb_n\)，所以最高次数恰为 \(m+n\)。
\end{proof}
形式多项式与多项式函数必须区别。例如在模 \(2\) 的剩余类中，\(x^2-x\) 对每个输入都取零值，却不是零多项式。若底环存在零因子，次数公式也可能失效：在模 \(4\) 的系数中，非零多项式 \(2x\) 与自身的乘积 \((2x)(2x)=4x^2\) 为零。

\ProseParagraphBreak
设 \(R\) 是交换环，\(G\) 为群。对系数 \(a_g\in R\)，若只有有限多个 \(a_g\) 非零，则称 \(\sum_{g\in G}a_gg\) 为有限形式和。全体这类形式和在逐项加法及
\[
\Bigl(\sum_g a_gg\Bigr)\Bigl(\sum_h b_hh\Bigr)
=\sum_{g,h}a_gb_h(gh)
\]
下组成\term[qunhuan]{群环}[group ring] \(R[G]\)。写在括号中的 \(gh\) 是群乘积，相同群元素的系数须相加。有限支集保证乘积有限，结合律由系数乘法与群乘法的结合律给出；单位是群单位处系数为 \(1\)、其余系数为零的形式和。这一构造把群的乘法保留下来，又允许对群元素作线性组合。例如取 \(G=C_2=\{1,g\}\)、\(g^2=1\)，对任意 \(a,b,c,d\in R\) 有
\[
(a+bg)(c+dg)
=ac+adg+bcg+bdg^2
=(ac+bd)+(ad+bc)g.
\]
群关系 \(g^2=1\) 把最后一项并入常数项。取 \(a=c=d=1\)、\(b=-1\)，便得到 \((1-g)(1+g)=0\)。
\symidx{\(R[G]\)：群环}

\subsection{零因子、单位与幂零元}
\label{b1:subsec:1001:03}
\begin{definition}\label{b1:def:units-zero-nil}
设 \(a\) 为环 \(R\) 的元素。如果存在 \(b\in R\) 使 \(ab=ba=1\)，就称 \(a\) 为\term[danwei]{单位}[unit]，所有单位组成的集合记为 \(R^\times\)。如果 \(a\ne0\)，且存在非零元素 \(b\in R\) 使 \(ab=0\) 或 \(ba=0\)，就称 \(a\) 为\term[lingyinzi]{零因子}[zero divisor]；如果存在正整数 \(m\) 使 \(a^m=0\)，就称 \(a\) 为\term[milingyuan]{幂零元}[nilpotent element]。
\end{definition}
\symidx{\(R^\times\)：环的单位群}
非交换环中，\(ab=0\) 与 \(ba=0\) 不能互相推出，文献中也会分别讨论左零因子与右零因子。本书在这里把这两种情形合称为零因子，只要求存在一个非零元素在至少一侧将它零化。

\ProseParagraphBreak
单位在乘法下组成群：乘积 \(ab\) 的逆是 \(b^{-1}a^{-1}\)，其余群公理从环继承。单位不可能为零因子，因为可以乘逆元消去。非零幂零元则一定是零因子：取使 \(a^m=0\) 的最小正整数 \(m\)，此时 \(m\geq2\)，最小性保证 \(a^{m-1}\ne0\)，而 \(a\cdot a^{m-1}=a^m=0\)。因此确实找到了一个非零的零化因子。注意零本身是幂零元，却按本书定义不叫零因子。

\ProseParagraphBreak
对照 Atiyah 与 Macdonald 的《Introduction to Commutative Algebra》时还须留意：该书在非零环中把零也计为零因子\cite[Chapter 1, p.~2]{AtiyahMacdonald1969CommutativeAlgebra}；本书的“零因子”始终要求元素非零。

\ProseParagraphBreak
有限性会把消去性质转化为可逆性，但非交换情形还须检查得到的是双边逆元。

\begin{proposition}\label{b1:prop:finite-domain-field}
有限非零环中，不是零因子的非零元素一定是单位。特别地，每个有限整环都是域。
\end{proposition}
\begin{proof}
设 \(a\ne0\) 且不是零因子。映射 \(L_a:r\mapsto ar\) 单射，因为 \(ar=as\) 给出 \(a(r-s)=0\)，从而 \(r=s\)。有限集上的单射为满射，故存在 \(b\) 使 \(ab=1\)。再比较
\[
L_a(ba)=a(ba)=(ab)a=a=L_a(1).
\]
由 \(L_a\) 的单射性，\(ba=1\)。这一步只用乘法结合律与单射性，没有使用交换性。整环乘法交换且每个非零元素都满足上述条件，故为域。
\end{proof}
有限性在这里不能删去：整数 \(2\) 既不是单位，也不是零因子。一般的非交换环中，仅有一侧逆元也不够。取域 \(k\) 上以 \(e_0,e_1,\ldots\) 为基的向量空间，令 \(T(e_i)=e_{i+1}\)，\(U(e_0)=0\)、\(U(e_{i+1})=e_i\)。在线性自映射组成的环中，\(UT=1\) 而 \(TU(e_0)=0\ne e_0\)，所以 \(TU\ne1\)。这说明定义中要求双边逆元确有必要。

\ProseParagraphBreak
幂零元还能产生单位。若 \(a^m=0\)，则
\[
(1-a)(1+a+\cdots+a^{m-1})
=(1+a+\cdots+a^{m-1})(1-a)=1.
\]
这里只涉及同一元素的幂，不需要环交换。与之相比，两个幂零元的和在非交换环中未必幂零。例如在任意域 \(k\) 上的 \(M_2(k)\) 中，取 \(a=E_{12}\)、\(b=E_{21}\)，则 \(a^2=b^2=0\)，而
\[
(a+b)^2=E_{11}+E_{22}=I_2.
\]
所以 \(a+b\) 不幂零。交换性怎样保证幂零条件对加法稳定，将在根理想的讨论中出现。

% !TeX root = ../../Book1.tex
% 纲要标题：### 10.2 理想与商环
% 纲要位置：计划纲要.md，第 394 行。

\section{理想与商环}
\label{b1:sec:1002}

在加法群上按一个子群取商，只能保证加法良定义。若还想让乘法通过商映射保留下来，这个子群必须能够吸收环元素的左右乘法，理想由此出现。商环中的零因子与可逆性又将分别引出素理想和极大理想。

% 纲要内容（仅作写作计划，尚非教材正文）：
%
% * 左理想、右理想与双边理想
% * 商环及同构定理
% * 素理想与极大理想的首次出现
%

\subsection{左理想、右理想与双边理想}
\label{b1:subsec:1002:01}
设 \(R\) 为环，先取加法子群 \(I\leq(R,+)\)，希望陪集上的乘法由代表元的乘法给出。若把第一项的代表 \(a\) 换成 \(a+\ell\)，其中 \(\ell\in I\)，所得乘积与原乘积的差是
\[
(a+\ell)b-ab=\ell b.
\]
要使输出陪集不变，必须有 \(\ell b\in I\)。若换的是第二项的代表，同样由 \(a(b+\ell)-ab=a\ell\) 得到 \(a\ell\in I\)。所以 \(I\) 既要吸收右乘，也要吸收左乘。与\cref{b1:thm:quotient-group}中由换代表元引出正规性一样，这里的条件也来自商运算的良定义要求。

\begin{definition}\label{b1:def:ring-ideal}
设 \(R\) 为环，\(I\subseteq R\) 是包含零元且对减法封闭的子集。若对每个 \(r\in R,a\in I\) 都有 \(ra\in I\)，称 \(I\) 为\term[zuo-lixiang]{左理想}[left ideal]；若总有 \(ar\in I\)，称为\term[you-lixiang]{右理想}[right ideal]；同时满足两者时称为\term[shuangbian-lixiang]{双边理想}[two-sided ideal]。未加限定的“理想”均指双边理想。若 \(I\ne R\)，称为真理想。
\end{definition}
理想不要求含环的单位；事实上理想一旦含有单位元或任一单位 \(u\)，就包含 \(u^{-1}u=1\)，继而包含每个 \(r=r1\)，所以等于 \(R\)。因此，一个非零环的真理想不是本书约定下的含同一单位的子环。这个限定不排除它自身的运算：理想仍对加法、减法和乘法封闭，可以作为不要求单位元的环来研究，有时也具有自己的单位元；这里排除的是它含有母环单位 \(1_R\)。

\ProseParagraphBreak
在交换环中左右条件等价。\(n\mathbb Z\) 是 \(\mathbb Z\) 的理想；反过来，整数环的每个理想首先是加法子群，由\cref{b1:lem:integer-subgroups}都为这种形式。

\ProseParagraphBreak
非交换情形确须区分左右。设 \(k\) 为域，在 \(M_2(k)\) 中令
\[
I_\ell=\left\{\,\begin{pmatrix}a&0\\b&0\end{pmatrix}
\;\middle|\;a,b\in k\,\right\},\qquad
I_r=\left\{\,\begin{pmatrix}a&b\\0&0\end{pmatrix}
\;\middle|\;a,b\in k\,\right\}.
\]
它们分别由第二列为零、第二行为零的矩阵组成，都对减法封闭。对任意 \(a,b,r,s,t,u\in k\)，
\[
\begin{aligned}
\begin{pmatrix}r&s\\t&u\end{pmatrix}
\begin{pmatrix}a&0\\b&0\end{pmatrix}
&=\begin{pmatrix}ra+sb&0\\ta+ub&0\end{pmatrix},\\
\begin{pmatrix}a&b\\0&0\end{pmatrix}
\begin{pmatrix}r&s\\t&u\end{pmatrix}
&=\begin{pmatrix}ar+bt&as+bu\\0&0\end{pmatrix}.
\end{aligned}
\]
第一式说明 \(I_\ell\) 吸收左乘，第二式说明 \(I_r\) 吸收右乘。但 \(E_{11}\in I_\ell\) 而 \(E_{11}E_{12}=E_{12}\notin I_\ell\)，所以 \(I_\ell\) 不是右理想；\(E_{11}\in I_r\) 而 \(E_{21}E_{11}=E_{21}\notin I_r\)，所以 \(I_r\) 不是左理想。

\ProseParagraphBreak
若 \(f:R\rightarrow S\) 为环同态，其加法核 \(\ker f=\{a\in R:f(a)=0\}\) 为双边理想：对 \(a\in\ker f\)，有 \(f(ra)=f(r)f(a)=0=f(ar)\)。这里保持零元的条件由加法保证，吸收乘法的条件来自乘法保持性。

\subsection{商环及同构定理}
\label{b1:subsec:1002:02}
\begin{proposition}\label{b1:prop:quotient-ring}
设 \(R\) 为环，\(L\leq(R,+)\) 为加法子群。在加法陪集集合 \(R/L\) 上按
\[
(a+L)+(b+L)=(a+b)+L,\qquad
(a+L)(b+L)=ab+L
\]
规定运算，则乘法与代表元无关当且仅当 \(L\) 是 \(R\) 的双边理想。在此情形下，\(R/L\) 构成\term[shanghuan]{商环}[quotient ring]，单位为 \(1+L\)；自然映射 \(q:R\rightarrow R/L\) 是满环同态，核为 \(L\)。
\end{proposition}
\begin{proof}
加法良定义来自加法群的商构造。若 \(L\) 是双边理想，而 \(a'=a+\ell,b'=b+m\)，其中 \(\ell,m\in L\)，则
\[
a'b'-ab=am+\ell b+\ell m\in L.
\]
其中 \(am\in L\) 用左吸收，因为 \(m\in L\)、\(a\in R\)；\(\ell b\in L\) 用右吸收，因为 \(\ell\in L\)、\(b\in R\)；\(\ell m\) 则用任一侧吸收都可得到。故乘法与代表元无关。结合律和分配律由代表元上的对应等式下降得到，\(1+L\) 在两侧乘任一陪集都保持它。自然映射按定义保持两种运算与单位，陪集都有代表元，且 \(a+L=L\) 恰好等价于 \(a\in L\)。

反过来，设所给陪集乘法良定义。对任意 \(r\in R\) 与 \(\ell\in L\)，由 \(\ell+L=0+L\) 得
\[
\begin{aligned}
r\ell+L&=(r+L)(\ell+L)=(r+L)(0+L)=L,\\
\ell r+L&=(\ell+L)(r+L)=(0+L)(r+L)=L.
\end{aligned}
\]
于是 \(r\ell,\ell r\in L\)。\(L\) 已是加法子群，故为双边理想。
\end{proof}
\symidx{\(R/I\)：环关于双边理想的商环}

把环同态 \(f:R\to S\) 先看作加法群同态，\cref{b1:thm:first-isomorphism}已经给出加法群同构 \((R,+)/\ker f\cong\operatorname{im}f\)。环论需要补的检查是：核吸收左右乘法，使商乘法有意义；诱导映射除保持加法外，还保持乘法和单位。前面已经证明核是双边理想，像则对两种运算封闭并含 \(1_S=f(1_R)\)，是 \(S\) 的子环。因而可以沿用同一个陪集映射得到下面的环同构。

\begin{theorem}[环的第一同构定理]\label{b1:thm:ring-first-iso}
设 \(R,S\) 为环。若 \(f:R\rightarrow S\) 是环同态，则
\[
\overline f:R/\ker f\longrightarrow\operatorname{im}f,
\qquad a+\ker f\longmapsto f(a)
\]
为环同构。更一般地，若 \(I\) 为 \(R\) 的双边理想且 \(I\subseteq\ker f\)，并以 \(q:R\rightarrow R/I\) 表示自然映射，则存在唯一环同态 \(g:R/I\rightarrow S\) 满足 \(g\circ q=f\)。
\end{theorem}
\begin{proof}
若 \(a-a'\in I\subseteq\ker f\)，则 \(f(a)=f(a')\)，所以 \(g(a+I)=f(a)\) 良定义。加法相容性由加法群的商构造给出，乘法与单位则分别满足
\[
\begin{aligned}
g\bigl((a+I)(b+I)\bigr)&=f(ab)=f(a)f(b),\\
g(1_R+I)&=f(1_R)=1_S.
\end{aligned}
\]
所以 \(g\) 是保单位的环同态。由定义有 \(g\circ q=f\)；由于 \(q\) 满射，该公式也是满足此等式的同态的唯一可能取值。取 \(I=\ker f\)，并将 \(g\) 的目标限制为其实际像 \(\operatorname{im}f\)，便得到定理中的 \(\overline f\)。它按像的定义满射；若 \(\overline f(a+\ker f)=0\)，则 \(a\in\ker f\)，所以其核只有零陪集，因而 \(\overline f\) 也是单射。
\end{proof}
条件 \(I\subseteq\ker f\) 使下面的三角形交换，虚线表示唯一满足 \(g\circ q=f\) 的环同态。
\[
\begin{tikzcd}[column sep=3em,row sep=2.5em]
R\arrow[r,"q"]\arrow[dr,"f"']&R/I\arrow[d,dashed,"\exists!\,g"]\\
&S
\end{tikzcd}
\]
这里的泛性质告诉我们：要从商环定义映射，先在原环定义，并检查需要抹去的理想确实映成零。商环不只是一个同构类型，还连同指定的商映射一起发挥作用。

\begin{proposition}[理想对应与商的同构]\label{b1:prop:ring-ideal-correspondence}
设 \(R\) 为环，\(I\) 为 \(R\) 的双边理想，\(q:R\rightarrow R/I\) 为自然映射。则：
\begin{enumerate}
\item 理想对应：\(q\) 给出包含 \(I\) 的 \(R\) 的双边理想与 \(R/I\) 的双边理想之间的保序双射。令 \(J\) 遍历前一类理想，\(L\) 遍历后一类理想，双向公式为
\begin{equation}\label{b1:eq:ring-ideal-correspondence}
J\longmapsto J/I,\qquad L\longmapsto q^{-1}(L).
\end{equation}
\item 环的第三同构定理：若 \(J\) 为 \(R\) 的双边理想且 \(I\subseteq J\)，则
\begin{equation}\label{b1:eq:ring-third-iso}
(R/I)/(J/I)\cong R/J.
\end{equation}
\item 环的第二同构定理：若 \(A\subseteq R\) 为含同一单位的子环，并记 \(A+I=\{a+i:a\in A,\ i\in I\}\)，则 \(A+I\) 是 \(R\) 的子环，\(I\) 是 \(A+I\) 的双边理想，且
\begin{equation}\label{b1:eq:ring-second-iso}
A/(A\cap I)\cong(A+I)/I.
\end{equation}
\end{enumerate}
\end{proposition}
\begin{proof}
\begin{enumerate}
\item 设 \(L\) 是 \(R/I\) 的双边理想。其原像 \(q^{-1}(L)\) 对减法及左右乘法封闭，且包含 \(I\)。含 \(I\) 的理想 \(J\) 的像 \(q(J)=J/I\) 对 \(R/I\) 中元素的左右乘法封闭，因为这些元素都可提升到 \(R\)。有 \(q^{-1}(q(J))=J\)：若 \(q(a)=q(j)\)，则 \(a-j\in I\subseteq J\)。又 \(q(q^{-1}(L))=L\)，因为 \(q\) 满射。两种操作均保持包含关系，得到\cref{b1:eq:ring-ideal-correspondence}。
\item 因 \(I\subseteq J\)，映射 \(R/I\rightarrow R/J\)、\(a+I\mapsto a+J\) 良定义且满射，核为 \(J/I\)。对该映射应用\cref{b1:thm:ring-first-iso}，即得\cref{b1:eq:ring-third-iso}。
\item 对 \(a,b\in A\) 及 \(i,j\in I\)，两元素 \(a+i,b+j\) 的差属于 \(A+I\)，乘积
\[
(a+i)(b+j)=ab+aj+ib+ij\in A+I.
\]
假设 \(A\) 含母环的同一个单位正用于这里：\(1_R\in A\subseteq A+I\)，配合上述封闭性，使 \(A+I\) 成为本书约定下的子环；只有加法与乘法封闭还不足以保证这一点。\(I\subseteq A+I\) 且吸收其左右乘法，因此是它的双边理想。\(A\cap I\) 同样是 \(A\) 的双边理想。映射 \(A\rightarrow(A+I)/I\)、\(a\mapsto a+I\) 保持单位且满射，因为 \(1_R\mapsto1_R+I\)、\(a+i+I=a+I\)，其核为 \(A\cap I\)。应用\cref{b1:thm:ring-first-iso}，即得\cref{b1:eq:ring-second-iso}。
\end{enumerate}
\end{proof}
上面的理想对应和两个商同构分别来自不同的映射。表中 \(I\subseteq J\) 是 \(R\) 的双边理想，\(A\) 是含 \(1_R\) 的子环；把出发映射及其核写出，就能看清每条公式抹去了什么。

\begin{center}
\begin{tabularx}{\linewidth}{@{}l>{\raggedright\arraybackslash}Xl@{}}
\toprule
结论 & 出发映射 & 核 \\
\midrule
理想对应 & \(q:R\to R/I\)，\(a\mapsto a+I\) & \(I\) \\
第三同构 & \(R/I\to R/J\)，\(a+I\mapsto a+J\) & \(J/I\) \\
第二同构 & \(A\to(A+I)/I\)，\(a\mapsto a+I\) & \(A\cap I\) \\
\bottomrule
\end{tabularx}
\end{center}

环同态的像未必是目标环的理想，例如 \(\mathbb Z\hookrightarrow\mathbb Q\) 的像不是 \(\mathbb Q\) 的理想。上述理想对应使用商映射的满射性，不能省去这个条件。

\subsection{素理想与极大理想}
\label{b1:subsec:1002:03}
本小节只讨论交换环。非交换环的素性与极大理想需要区分不同版本，不把下面的乘积判别直接当作其一般定义。

\ProseParagraphBreak
取真理想 \(I\) 后，商环 \(R/I\) 非零。接下来有两个不同的问题：什么时候它仍像整数环一样，两个非零元素的乘积不会为零？什么时候它像域一样，每个非零元素都能取逆？下面的素理想与极大理想将分别刻画这两种商环性质。

\begin{definition}\label{b1:def:prime-maximal-ideal}
设 \(R\) 为交换环，\(\mathfrak p\subsetneq R\) 为真理想。若对任意 \(a,b\in R\)，\(ab\in\mathfrak p\) 都蕴含 \(a\in\mathfrak p\) 或 \(b\in\mathfrak p\)，则称 \(\mathfrak p\) 为\term[su-lixiang]{素理想}[prime ideal]。再设 \(\mathfrak m\subsetneq R\) 为真理想。若满足 \(\mathfrak m\subseteq J\subseteq R\) 的理想只有 \(J=\mathfrak m\) 与 \(J=R\)，则称 \(\mathfrak m\) 为\term[jida-lixiang]{极大理想}[maximal ideal]。
\end{definition}
“真”不能漏掉；否则整个环会满足乘积条件，却不能对应非零整环。\((0)\) 是整环中的素理想，但整数环中它不极大，因为 \((0)\subsetneq2\mathbb Z\subsetneq\mathbb Z\)。

\begin{proposition}\label{b1:prop:prime-max-quotient}
交换环 \(R\) 的理想 \(I\) 为素理想，当且仅当 \(R/I\) 是整环；为极大理想，当且仅当 \(R/I\) 是域。因此每个极大理想都是素理想。
\end{proposition}
\begin{proof}
\(I\ne R\) 恰好保证商环非零；在商环中，\((a+I)(b+I)=0\) 恰好等价于 \(ab\in I\)，故第一项直接由定义得到。

若 \(I\) 极大且 \(a\notin I\)，要使非零类 \(a+I\) 有逆，就要找到 \(r\in R\) 使 \(ra-1\in I\)。这等价于 \(1\in I+Ra\)，所以考虑包含 \(I\) 并加入 \(a\) 的集合
\[
I+Ra=\{i+ra:i\in I,\ r\in R\}.
\]
此集合对减法封闭；由于 \(R\) 交换，对 \(s\in R\) 还有 \(s(i+ra)=si+(sr)a\in I+Ra\)，故它是理想。它含 \(I\) 及 \(a\)，而 \(a\notin I\)，所以严格包含 \(I\)，由极大性知 \(I+Ra=R\)。写 \(1=i+ra\)，便有 \((r+I)(a+I)=1+I\)，所以商中每个非零元素可逆。反过来，若 \(R/I\) 为域，而理想 \(J\) 严格包含 \(I\)，取 \(a\in J\setminus I\)；商中 \(a+I\) 的逆有代表 \(r\)，于是 \(1-ra\in I\subseteq J\)，给出 \(1\in J\)，故 \(J=R\)。最后，域无非零零因子，所以极大理想为素理想。
\end{proof}
素理想不一定极大，正如整环不一定是域：\(\mathbb Z/(0)\cong\mathbb Z\) 是整环，却不是域，所以 \((0)\) 是 \(\mathbb Z\) 的素理想而不是极大理想。这与前面的严格包含链 \((0)\subsetneq2\mathbb Z\subsetneq\mathbb Z\) 给出同一个反例。

\ProseParagraphBreak
整数环的理想为 \(n\mathbb Z\)。当 \(p\) 是素数，任一 \(a\notin p\mathbb Z\) 满足 \(\gcd(a,p)=1\)，整数 Bézout 等式给出 \(ua+vp=1\)，故模 \(p\) 的非零类均有逆。所得域记为 \(\mathbb F_p=\mathbb Z/p\mathbb Z\)。正合数 \(n=ab\)、\(1<a,b<n\) 则产生两个非零类之积为零。因此 \(\mathbb Z\) 的素理想恰为 \((0)\) 及 \(p\mathbb Z\)，极大理想恰为后者。
\symidx{\(\mathbb F_p\)：含 \(p\) 个元素的素数域}

\ProseParagraphBreak
极大理想的存在性寻找的是包含偏序中的极大元，即不能再严格扩大的真理想，并非包含所有真理想的最大元。例如 \(2\mathbb Z\) 与 \(3\mathbb Z\) 都是极大理想，却互不包含；没有哪一个同时包含这两个真理想并仍为真理想。下面用 Zorn 引理证明所需极大元存在。

\begin{proposition}\label{b1:prop:maximal-ideal-exists}
非零交换环中的每个真理想都包含在某个极大理想中。
\end{proposition}
\begin{proof}
固定真理想 \(I\)，考虑包含 \(I\) 的真理想按包含关系构成的偏序集。它非空。非空链的并对减法封闭，因为任意两个元素都落在链上某个共同的理想中；它也吸收乘法。这个并不含 \(1\)，否则链中的某个真理想已含 \(1\)，矛盾。因而链的并仍是此偏序集的上界；空链可用 \(I\) 作上界。应用 Zorn 引理（\cref{b1:thm:A-zorn}），即“每条链有上界的非空偏序集有极大元”，得到 \(\mathfrak m\)。任何严格包含它的真理想仍包含 \(I\)，违背极大性，故 \(\mathfrak m\) 是极大理想。
\end{proof}
这里明确使用选择公理的 Zorn 形式，不给出寻找极大理想的算法。在具体环中仍应尽可能写出商映射，例如对交换环 \(R\)，常数项映射 \(R[x]\rightarrow R\) 满射，核为 \(xR[x]\)。\cref{b1:thm:ring-first-iso}给出 \(R[x]/xR[x]\cong R\)，所以当 \(R\) 为整环或域时，这个理想分别为素理想或极大理想。

% !TeX root = ../../Book1.tex
% 纲要标题：### 10.3 理想运算
% 纲要位置：计划纲要.md，第 400 行。

\section{理想运算}
\label{b1:sec:1003}

一个理想记录一组允许被视为零的元素。设 \(I,J\) 为环 \(R\) 的双边理想。若在同一个商环中令 \(I\) 和 \(J\) 的元素都成为零，就要商去 \(I+J\)；同时属于 \(I\) 和 \(J\) 的元素则组成 \(I\cap J\)，它们在两个商环中都映为零。具体地，自然同态 \(r\mapsto(r+I,r+J)\) 满足
\[
\ker\!\left(R\longrightarrow R/I\times R/J\right)=I\cap J.
\]
理想积 \(IJ\) 由 \(I\) 中元素与 \(J\) 中元素的乘积的有限和组成。从包含关系看，\(I+J\) 是同时包含 \(I,J\) 的最小理想，\(I\cap J\) 是同时包含于 \(I,J\) 的最大理想；\(IJ\) 则由环的乘法产生，并且包含于 \(I\cap J\)。和、交由包含关系确定，积回答的则是乘法问题。本节先写出它们的元素形式，再在交换环中讨论取幂、取根，以及寻找哪些乘子能把一个理想送入另一个理想。

% 纲要内容（仅作写作计划，尚非教材正文）：
%
% * 和、积、交与生成理想
% * 理想的幂与根、准素理想的初步性质
% * 不同理想运算之间的包含关系
% * 理想商、零化子与基本运算性质
%

\subsection{理想的和、积、交与生成}
\label{b1:subsec:1003:01}
\begin{definition}\label{b1:def:ideal-operations}
设 \(I,J\) 为环 \(R\) 的双边理想。它们的\term[lixiang-de-he]{和}[sum of ideals]与\term[lixiang-de-ji]{积}[product of ideals]分别为
\[
I+J=\{\,a+b\mid a\in I,\ b\in J\,\},\qquad
IJ=\Bigl\{\,\sum_{\nu=1}^m a_\nu b_\nu\mid
m\geq0,\ a_\nu\in I,\ b_\nu\in J\,\Bigr\}.
\]
空和表示零。理想的交采用通常的集合交。对任意子集 \(S\subseteq R\)，包含 \(S\) 的所有双边理想之交称为 \(S\) 的\term[shengcheng-lixiang]{生成理想}[generated ideal]，记为 \((S)\)。
\end{definition}
\symidx{\(I+J,\ IJ,\ I\cap J\)：理想的和、积与交}
\symidx{\((S)\)：由子集 \(S\) 生成的理想}

生成双边理想时，两侧系数来自左右吸收条件：一个包含 \(s\in S\) 的双边理想必须包含所有 \(rs\) 和 \(st\)，继而包含所有 \(rst\)，其中 \(r,t\in R\)。再把这些元素的有限和收入其中，就得到下面的具体公式；非交换环中通常不能把两侧系数越过 \(s\) 合并。

\begin{proposition}\label{b1:prop:ideal-operations}
设 \(I,J\) 为环 \(R\) 的双边理想，\(S\subseteq R\)。则 \(I+J\)、\(IJ\) 与 \(I\cap J\) 都是 \(R\) 的双边理想；\(I+J\) 是包含 \(I,J\) 的最小双边理想。由 \(S\) 生成的双边理想有具体表达式
\[
(S)=\Bigl\{\,\sum_{\nu=1}^m r_\nu s_\nu t_\nu
\mid m\geq0,\ r_\nu,t_\nu\in R,\ s_\nu\in S\,\Bigr\}.
\]
若 \(R\) 交换，可把每个 \(r_\nu t_\nu\) 合成一个系数，得到 \(\sum r_\nu s_\nu\) 的形式。
\end{proposition}
\begin{proof}
理想的交保持减法及左右乘法封闭。对 \(I+J\)，两个和相减仍可分别归入 \(I,J\)，且左右乘法分配到两项中，故是理想；任何同时包含 \(I,J\) 的理想都包含这些和。

\(IJ\) 对加法封闭来自有限和的合并，取负可把负号吸收到 \(a_\nu\) 中。对任意 \(r\in R\)，左乘将 \(a_\nu\) 换成 \(ra_\nu\in I\)，右乘将 \(b_\nu\) 换成 \(b_\nu r\in J\)，所以也是理想。最后一个显示式右边同样对这些运算封闭，取 \(r_\nu=t_\nu=1\) 说明它包含 \(S\)；任何包含 \(S\) 的双边理想都包含这些有限和，故恰为 \((S)\)。
\end{proof}
即使 \(S\) 无限，每个生成理想中的元素也只用有限多个生成元。写 \((a_1,\ldots,a_m)\) 表示有限集合的生成理想；由一个元素生成的理想称为\term[zhulixiang]{主理想}[principal ideal]。交换环中 \((a)=Ra\)，非交换环中的双边主理想则一般须保留两侧系数。

\ProseParagraphBreak
理想积取有限和，是为了保证加法封闭。一个具体例子在 \(\mathbb Q[x,y]\) 中可以看清。这里取迭代多项式环 \((\mathbb Q[x])[y]\)，其中 \(y\) 与系数交换，每个元素都唯一写成有限和 \(\sum_{i,j}a_{ij}x^iy^j\)。令 \(I=J=(x,y)\)。单个乘积的集合含有 \(x^2=x\cdot x\)、\(y^2=y\cdot y\)，却不含它们的和 \(x^2+y^2\)。事实上，若 \(x^2+y^2=fg\)，其中 \(f,g\in(x,y)\)，则 \(f,g\) 的常数项都为零。记它们的一次项分别为 \(\alpha x+\beta y\)、\(\gamma x+\delta y\)，其中系数均为有理数。乘积的二次项只能来自这两组一次项相乘，比较系数便有
\[
\alpha\gamma=1,\qquad
\alpha\delta+\beta\gamma=0,\qquad
\beta\delta=1.
\]
首尾两式使四个系数都非零，且 \(\gamma=\dfrac1\alpha\)、\(\delta=\dfrac1\beta\)。代入中间一式并乘以 \(\alpha\beta\)，得到 \(\alpha^2+\beta^2=0\)，与非零有理数的平方为正矛盾。因此 \(x^2+y^2\) 属于 \(IJ\)，但不属于单个乘积的集合。

\ProseParagraphBreak
有限和的定义还使
\[
I(J+K)=IJ+IK,\qquad (I+J)K=IK+JK,\qquad (IJ)K=I(JK).
\]
前两式把一个乘积的因子展开，再合并有限和；末式两边都是元素 \(abc\)、\(a\in I,b\in J,c\in K\) 的有限和。这些论证只用结合律和分配律，不要求交换性。

\ProseParagraphBreak
再设 \(k\) 为任意域，仍取迭代多项式环 \(k[x,y]=(k[x])[y]\)。在这个环中取 \(I=(x,y)\)、\(J=(x,y^2)\)。把两组生成元两两相乘，得到
\[
IJ=(x^2,xy^2,xy,y^3)=(x^2,xy,y^3).
\]
其中 \(xy\) 是容易漏掉的混合乘积；\(xy^2\) 已由 \(xy\) 生成，因此可从生成元列表中删去。

\subsection{理想的幂与根}
\label{b1:subsec:1003:02}
设 \(I\) 为环 \(R\) 的双边理想。约定 \(I^0=R\)、\(I^1=I\)，并对整数 \(n\geq1\) 递归定义 \(I^{n+1}=I^nI\)。对 \(n\geq1\)，\(I^n\) 由 \(n\) 个 \(I\) 中元素的乘积的有限和组成，而不只由某个元素的 \(n\) 次幂组成。交换环的主理想满足 \((a)^n=(a^n)\)，但多生成元时各个混合乘积也必须计算。

\begin{definition}\label{b1:def:radical-ideal}
设 \(R\) 为交换环，\(I\) 为 \(R\) 的理想。集合
\[
\sqrt I\coloneqq\{\,a\in R\mid a^n\in I\;\text{对某个正整数}\;n\;\text{成立}\,\}
\]
称为 \(I\) 的\term[lixiang-de-gen]{根}[radical of an ideal]。若 \(I=\sqrt I\)，则称 \(I\) 为\term[genlixiang]{根理想}[radical ideal]。\(\sqrt{(0)}\) 称为环的\term[milinggen]{幂零根}[nilradical]。
\end{definition}
\symidx{\(\sqrt{(0)}\)：交换环的幂零根}
\symidx{\(\sqrt I\)：理想的根}

取根可以直接从商环来理解。设 \(q:R\to R/I\) 为商映射，则
\[
a\in\sqrt I
\quad\Longleftrightarrow\quad
(a+I)^n=0\;\text{对某个正整数}\;n\;\text{成立}.
\]
因此 \(\sqrt I\) 收集商环中所有幂零元的原像：一个元素未必已经属于 \(I\)，但它在商环中的像经过某次幂会变为零。

\begin{proposition}\label{b1:prop:radical-properties}
设 \(R\) 为交换环，\(I\) 为 \(R\) 的理想。则 \(\sqrt I\) 是包含 \(I\) 的理想，且 \(\sqrt{\sqrt I}=\sqrt I\)。它是包含 \(I\) 的最小根理想。商环 \(R/I\) 中没有非零幂零元，当且仅当 \(I=\sqrt I\)。
\end{proposition}
\begin{proof}
若 \(a^m,b^n\in I\)，则交换性允许对 \((a-b)^{m+n-1}\) 用二项式展开。指数 \(m+n-1\) 的选择保证：若某项中 \(a\) 的指数小于 \(m\)，则至多为 \(m-1\)，因而 \(b\) 的指数至少为 \((m+n-1)-(m-1)=n\)。所以每项或者含 \(a^m\)，或者含 \(b^n\)，都属于 \(I\)。又 \((ra)^m=r^ma^m\in I\)，所以 \(\sqrt I\) 对减法及乘法吸收封闭。\(I\subseteq\sqrt I\) 可取指数一。

若 \(a^r\in\sqrt I\)，再取 \(s\) 使 \((a^r)^s\in I\)，得到 \(a\in\sqrt I\)，证明取两次根不变。若根理想 \(J\) 包含 \(I\)，则 \(a^n\in I\) 蕴含 \(a^n\in J\)，进而 \(a\in J\)，故 \(\sqrt I\subseteq J\)。最后，\((a+I)^n=0\) 等价于 \(a^n\in I\)，而 \(a+I=0\) 等价于 \(a\in I\)，给出商环的判别。
\end{proof}
没有非零幂零元的交换环称为\term[yuehua-huan]{约化环}[reduced ring]。对交换环 \(R\) 及其理想 \(I\)，\cref{b1:prop:radical-properties}给出
\[
\begin{aligned}
R\;\text{约化}&\quad\Longleftrightarrow\quad(0)\;\text{是根理想},\\
R/I\;\text{约化}&\quad\Longleftrightarrow\quad I\;\text{是根理想}.
\end{aligned}
\]
这里 \(\sqrt I/I\) 正是 \(R/I\) 中所有幂零元组成的理想。整环必约化，反之不成立：域的直积 \(k\times k\) 有零因子 \((1,0),(0,1)\)，但若 \((a,b)^n=(0,0)\)，域中取幂不产生非零幂零元，故 \(a=b=0\)。

\ProseParagraphBreak
每个素理想都是根理想，因为 \(a^n\in\mathfrak p\) 可反复用素性推出 \(a\in\mathfrak p\)。\(\mathbb Z\) 中，\(\sqrt{12\mathbb Z}=6\mathbb Z\)：若 \(12\mid a^n\)，素数 \(2,3\) 均整除 \(a\)；反过来 \(6\mid a\) 时 \(12\mid a^2\)。这也说明根理想的作用是忘掉某些幂次信息，而不是忘掉全部整除信息。

\ProseParagraphBreak
若希望分解理想而保留幂次信息，就会用到下述概念。与素理想要求商环没有零因子相比，准素条件允许出现零因子，但要求它们经过某次幂变为零。先在原环中写出这个条件，再翻译到商环。
\begin{definition}\label{b1:def:primary-ideal}
设 \(Q\) 为交换环 \(R\) 的真理想。若对任意 \(a,b\in R\)，只要 \(ab\in Q\) 且 \(a\notin Q\)，就存在正整数 \(m\) 使 \(b^m\in Q\)，则称 \(Q\) 为\term[zhunsu-lixiang]{准素理想}[primary ideal]。
\end{definition}
设 \(R\) 为交换环，\(Q\subsetneq R\) 为理想。准素性的商环判别为
\[
Q\;\text{是准素理想}
\quad\Longleftrightarrow\quad
R/Q\;\text{中每个非零零因子都是幂零元}.
\]
若 \(Q\) 准素，取商环中的一个非零零因子 \(b+Q\)。存在非零元素 \(a+Q\) 使 \((a+Q)(b+Q)=0\)，即 \(ab\in Q\) 且 \(a\notin Q\)。准素性给出正整数 \(m\) 使 \(b^m\in Q\)，所以 \((b+Q)^m=0\)，这个零因子幂零。

\ProseParagraphBreak
反过来，若 \(R/Q\) 的每个非零零因子都幂零，设 \(ab\in Q\) 且 \(a\notin Q\)。若 \(b\in Q\)，取 \(m=1\) 就满足准素条件；若 \(b\notin Q\)，则 \(a+Q,b+Q\) 都非零，而它们的乘积为零，故 \(b+Q\) 是非零零因子。由假设存在正整数 \(m\) 使 \((b+Q)^m=0\)，也就是 \(b^m\in Q\)。因此 \(Q\) 是准素理想。两种情况都要保留，因为本书的零因子定义不把零本身计入其中。

\begin{proposition}\label{b1:prop:primary-radical}
设 \(R\) 为交换含幺环。\(R\) 的每个素理想都是准素理想；若 \(Q\) 是 \(R\) 的准素理想，则 \(\sqrt Q\) 是素理想。因此，准素理想 \(Q\) 是素理想，当且仅当它是根理想。
\end{proposition}
\begin{proof}
若 \(P\) 为素理想，\(ab\in P\) 且 \(a\notin P\) 时有 \(b\in P\)，所以 \(P\) 满足准素条件。设 \(Q\) 为准素理想。由 \(Q\neq R\) 可知 \(1\notin\sqrt Q\)，故 \(\sqrt Q\) 仍是真理想。若 \(ab\in\sqrt Q\) 而 \(a\notin\sqrt Q\)，取正整数 \(n\) 使 \(a^nb^n\in Q\)。此时 \(a^n\notin Q\)，准素性给出正整数 \(m\)，使 \(b^{nm}\in Q\)；于是 \(b\in\sqrt Q\)。这证明 \(\sqrt Q\) 为素理想。若 \(Q\) 还是根理想，则 \(Q=\sqrt Q\) 为素理想；反之，素理想是根理想。
\end{proof}
例如整数环中的 \(p^r\mathbb Z\)（\(p\) 为素数，\(r\geq1\)）是准素理想。若 \(p^r\mid ab\) 而 \(p^r\nmid a\)，整数的素数分解说明 \(p\mid b\)。写 \(b=pc\)，便有 \(b^r=p^rc^r\)，所以 \(p^r\mid b^r\)；通过取幂，就补足了原来可能不足的素数因子个数。特别地，
\[
12\mathbb Z=4\mathbb Z\cap3\mathbb Z
\]
把一个理想写成两个准素理想的交，仍保留因子二所带的平方信息。这里仅用这个例子说明分解问题；一般的准素分解理论，包括 Noether 环中的存在性与唯一性范围，安排在第三卷第4.2--4.4节。

\subsection{理想运算的包含关系}
\label{b1:subsec:1003:03}
\begin{proposition}\label{b1:prop:ideal-containments}
设 \(I,J\) 为环 \(R\) 的双边理想，\(n\geq1\) 为整数。则
\[
IJ\subseteq I\cap J\subseteq I+J,\qquad I^{n+1}\subseteq I^n.
\]
如果 \(R\) 交换，还有
\[
\sqrt{IJ}=\sqrt{I\cap J}=\sqrt I\cap\sqrt J,
\qquad \sqrt{I^n}=\sqrt I\quad(n\geq1).
\]
\end{proposition}
\begin{proof}
每个 \(ab\)、\(a\in I,b\in J\) 既在 \(I\) 中也在 \(J\) 中，因此有限和也在交中；交中的元素可写为自身加零，故在和中。把 \(J\) 换成 \(I\) 的幂给出递降关系。

由包含关系及根的定义，\(\sqrt{IJ}\subseteq\sqrt{I\cap J}\subseteq\sqrt I\cap\sqrt J\)。若 \(a^r\in I\)、\(a^s\in J\)，则 \(a^{r+s}=a^ra^s\in IJ\)，故反向包含也成立。\(I^n\subseteq I\) 给出一个方向；若 \(a^r\in I\)，则 \(a^{rn}\in I^n\)，给出另一个方向。
\end{proof}
\cref{b1:fig:ideal-operation-containments}把和、交、积的包含关系放在一起。箭头沿包含方向由较小的理想指向较大的理想，图中的包含一般都可能严格；和与交分别是共同的最小上界与最大下界，而积位于交之内。

\begin{figure}[H]
\centering
\begin{tikzpicture}[x=1cm,y=1cm,>=stealth]
\node (sum) at (0,3) {\(I+J\)};
\node (leftideal) at (-1.5,2) {\(I\)};
\node (rightideal) at (1.5,2) {\(J\)};
\node (intersection) at (0,1) {\(I\cap J\)};
\node (product) at (0,0) {\(IJ\)};
\draw[->] (product) -- (intersection);
\draw[->] (intersection) -- (leftideal);
\draw[->] (intersection) -- (rightideal);
\draw[->] (leftideal) -- (sum);
\draw[->] (rightideal) -- (sum);
\end{tikzpicture}
\caption[理想运算的包含关系]{理想运算的包含关系。\(I,J\) 为同一环的双边理想，箭头表示包含。}
\label{b1:fig:ideal-operation-containments}
\end{figure}

等号需要额外理由。在 \(\mathbb Z\) 中取 \(I=J=2\mathbb Z\)，则 \(IJ=4\mathbb Z\subsetneq I\cap J=2\mathbb Z\)；取 \(I=2\mathbb Z,J=3\mathbb Z\)，则交为 \(6\mathbb Z\)，和为 \(\mathbb Z\)，也不相等。另一方面 \(I\cup J\) 通常不是理想，例如上述并包含 \(2,3\)，却不包含 \(3-2=1\)。

\ProseParagraphBreak
主理想的包含关系对应于生成元之间的整除关系，但方向相反。整数环中 \((a)\subseteq(b)\) 恰好等价于 \(b\mid a\)。例如 \(6\mathbb Z\subseteq2\mathbb Z\)，不是因为 \(6\mid2\)，而是每个 \(6\) 的倍数都是 \(2\) 的倍数。下一章将把这个简单现象推广到整环。下面先从乘法吸收条件出发，再定义一种理想运算。

\subsection{理想商与零化子}
\label{b1:subsec:1003:04}

本小节始终设 \(R\) 为交换含幺环。给定理想 \(I,J\)，要寻找满足 \(KJ\subseteq I\) 的最大理想 \(K\)。逐元素看，就是考虑哪些 \(r\in R\) 能使 \(J\) 中每个元素乘以 \(r\) 后都进入 \(I\)。把所有这些乘子收集起来，就得到所求的最大理想；下面先给出集合公式，再验证它的理想性质与最大性。

\begin{definition}\label{b1:def:ideal-quotient}
设 \(R\) 为交换含幺环，\(I,J\) 为 \(R\) 的理想。集合
\[
(I:J)\coloneqq\{\,r\in R\mid rJ\subseteq I\,\}
\]
称为\term[lixiang-shang]{理想商}[ideal quotient]，也称商理想。对单个元素 \(a\in R\)，简记
\[
I:a\coloneqq\bigl(I:(a)\bigr)
=\{\,r\in R\mid ra\in I\,\}.
\]
取 \(I=(0)\)，所得理想称为 \(J\) 的\term[linghuazi]{零化子}[annihilator]，也称湮灭子，记为 \(\operatorname{Ann}_R(J)=(0:J)\)；元素的零化子相应记为 \(\operatorname{Ann}_R(a)=0:a\)。把乘以 \(r\) 后使 \(a\) 变为零称为 \(r\) 对 \(a\) 的\term[linghua]{零化}[annihilation]，条件就是 \(ra=0\)。
\end{definition}
\termidx[shang-lixiang]{商理想}
\termidx[yanmiezi]{湮灭子}
\symidx{\((I:J)\)：理想商}
\symidx{\(\operatorname{Ann}_R(J)\)：理想的零化子}
这里的 \((I:J)\) 是 \(R\) 的一个子集，不是商环 \(R/I\)，也不是对两个理想作加法群的商。它的定义不要求 \(I\subseteq J\) 或 \(J\subseteq I\)。单元素简记中的等号成立，是因为 \(ra\in I\) 蕴含对每个 \(s\in R\) 有 \(rsa\in I\)，而反向可取 \(s=1\)。

\begin{proposition}\label{b1:prop:ideal-quotient-rules}
设 \(I,J\) 为交换含幺环 \(R\) 的理想。则 \((I:J)\) 是包含 \(I\) 的理想。对 \(R\) 的任一理想 \(K\)，有
\begin{equation}\label{b1:eq:ideal-quotient-adjunction}
KJ\subseteq I\quad\Longleftrightarrow\quad K\subseteq(I:J).
\end{equation}
因此 \((I:J)\) 是满足左侧包含关系的最大理想，且
\[
(I:R)=I,\qquad (I:0)=R,\qquad
(I:J)=R\ \Longleftrightarrow\ J\subseteq I.
\]
\end{proposition}
\begin{proof}
零在 \((I:J)\) 中。若 \(r,r'\) 在其中，对任意 \(j\in J\) 有 \((r-r')j=rj-r'j\in I\)；若 \(s\in R\)，则 \((sr)j=s(rj)\in I\)。所以它是理想。若 \(r\in I\)，理想的吸收性给出 \(rJ\subseteq I\)，故 \(I\subseteq(I:J)\)。

条件 \(KJ\subseteq I\) 意味着每个 \(k\in K\) 都满足 \(kJ\subseteq I\)，故 \(k\in(I:J)\)。反过来，若所有这些乘积均在 \(I\) 中，其有限和也在 \(I\) 中，正好得到 \(KJ\subseteq I\)。取 \(K=I:J\) 及任一满足条件的 \(K\)，分别得到它满足该条件和它包含全部这样的理想。最后，\(rR\subseteq I\) 等价于 \(r\in I\)，因为可令乘数为一；乘以零没有限制；\(1\in(I:J)\) 则等价于 \(J\subseteq I\)。这证明全部公式。
\end{proof}
例如在 \(\mathbb Z\) 中，
\[
\bigl(12\mathbb Z:8\mathbb Z\bigr)=3\mathbb Z,
\]
因为 \(12\mid8r\) 当且仅当 \(3\mid r\)。这里计算的是满足 \(K(8\mathbb Z)\subseteq12\mathbb Z\) 的最大理想 \(K\)，普通数字除法的结果 \(\dfrac{12}{8}=\dfrac32\) 不参与这个乘子条件。事实上，\((3\mathbb Z)(8\mathbb Z)=24\mathbb Z\subsetneq12\mathbb Z\)，所以把理想商再乘回去，一般只能得到包含而非等号。在 \(\mathbb Z/12\mathbb Z\) 中，\(\operatorname{Ann}_R(\bar4)=(\bar3)\)：条件 \(4r\equiv0\pmod{12}\) 正好要求 \(3\mid r\)。对非零元素，零化子非零恰好说明它是零因子；零元素本身的零化子则是全环。

\begin{proposition}\label{b1:prop:ideal-quotient-identities}
设 \(I,J,K\) 为交换含幺环 \(R\) 的理想。则理想商满足
\[
\begin{aligned}
(I\cap J):K&=(I:K)\cap(J:K),\\
I:(J+K)&=(I:J)\cap(I:K),\\
(I:J):K&=I:(JK).
\end{aligned}
\]
它对第一项保持包含关系，对第二项反转包含关系。若 \(m\geq0\) 为整数，且 \(J=(a_1,\ldots,a_m)\)，则
\[
(I:J)=\bigcap_{i=1}^m(I:a_i).
\]
\end{proposition}
\begin{proof}
一个元素 \(r\) 属于第一式左边，当且仅当 \(rK\) 同时包含于 \(I,J\)，这正是右边的两个条件。第二式中，\(r(J+K)\subseteq I\) 蕴含 \(rJ,rK\subseteq I\)，因为 \(J,K\subseteq J+K\)；反向由加法封闭性得到。

第三式左边的条件是对所有 \(k\in K\) 都有 \(rk\in(I:J)\)，即对所有 \(k\in K,j\in J\) 都有 \(rkj\in I\)。这些乘积的有限和仍在 \(I\) 中，故等价于 \(r(JK)\subseteq I\)。包含方向则由定义可见：扩大允许落入的理想放宽条件，扩大被乘的理想增加条件。最后，\(rJ\subseteq I\) 必须使每个 \(ra_i\in I\)；若每个生成元都满足此条件，则 \(r\sum_i s_i a_i=\sum_i s_i ra_i\in I\)，所以也足够。空生成集对应零理想，此时空交取为 \(R\)，公式仍成立。
\end{proof}
由\cref{b1:prop:ideal-quotient-rules}，\(I\subseteq(I:a)\)。所以商映射下的像 \((I:a)/I=\{\,r+I\mid r\in(I:a)\,\}\) 是 \(R/I\) 中的理想，这也符合\cref{b1:prop:ring-ideal-correspondence}中的理想对应。在商环中，\(r+I\) 零化 \(a+I\) 恰好意味着 \(ra\in I\)，也就是 \(r\in(I:a)\)。因而
\[
\operatorname{Ann}_{R/I}(a+I)=(I:a)/I.
\]
这把理想商与商环中的零因子联系起来。到\secref{b1:sec:1304}，允许乘以任意高次的分母，还会把这些理想商连接成饱和理想；下一节先继续讨论互素理想与环的直积分解。

% !TeX root = ../../Book1.tex
% 纲要标题：### 10.4 中国剩余定理
% 纲要位置：计划纲要.md，第 406 行。

\section{中国剩余定理}
\label{b1:sec:1004}

整数的两个互素模数可以独立指定余数。对一般交换环，替代“模数互素”的条件是两个理想的和等于整个环。中国剩余定理不仅给出同余方程的解，还把一个商环拆成若干分量；逆向看，这些分量由特殊的幂等元识别。本节除最后的辨析外均在交换含幺环中讨论。

% 纲要内容（仅作写作计划，尚非教材正文）：
%
% * 两两互素理想与有限直积分解
% * 整数同余方程
% * 幂等元与环的分解
%
% ---
%

\subsection{两两互素理想与有限直积分解}
\label{b1:subsec:1004:01}
\begin{definition}\label{b1:def:comaximal-ideals}
设 \(R\) 为交换环，\(I,J\) 为 \(R\) 的理想。若 \(I+J=R\)，则称 \(I,J\) 为\term[husu-lixiang]{互素理想}[comaximal ideals]。这等价于存在 \(a\in I,b\in J\) 使 \(a+b=1\)。
\end{definition}

\begin{lemma}\label{b1:lem:comaximal-product}
设 \(I,J\) 为交换环 \(R\) 的互素理想，则 \(I\cap J=IJ\)。更一般地，若 \(r\) 为正整数，且 \(R\) 的理想 \(I_1,\ldots,I_r\) 两两互素，则
\(\bigcap_{i=1}^r I_i=\prod_{i=1}^r I_i\)。
\end{lemma}
\begin{proof}
\cref{b1:prop:ideal-containments}已给出 \(IJ\subseteq I\cap J\)。若 \(x\in I\cap J\)，取 \(a+b=1\)、\(a\in I,b\in J\)，则 \(x=xa+xb\in IJ\)，这里交换性把 \(xa\) 写成 \(ax\)。

对多个理想，取 \(a_j\in I_j,b_j\in I_r\) 使 \(a_j+b_j=1\)，其中 \(j<r\)。乘积 \(a_1\cdots a_{r-1}\) 属于 \(I_1\cdots I_{r-1}\)，且模 \(I_r\) 同余于 \(1\)。所以 \(I_r+I_1\cdots I_{r-1}=R\)。从 \(r=2\) 起归纳，将两理想结论用于这个乘积理想与 \(I_r\)，即可得到有限情形。
\end{proof}
\term[zhongguo-shengyu-dingli]{中国剩余定理}[Chinese remainder theorem]把这一等式与一个满同态结合起来。

\begin{theorem}[中国剩余定理]\label{b1:thm:crt}
设 \(r\geq1\) 为整数，交换环 \(R\) 的理想 \(I_1,\ldots,I_r\) 两两互素。则映射
\[
\Phi:R\longrightarrow\prod_{i=1}^rR/I_i,\qquad
a\longmapsto(a+I_i)_{i=1}^r
\]
是满环同态，核为 \(\bigcap_iI_i=\prod_iI_i\)，因而诱导环同构
\[
R/\Bigl(\prod_{i=1}^rI_i\Bigr)\cong\prod_{i=1}^rR/I_i.
\]
反过来，若 \(\Phi\) 满射，则这些理想必两两互素。
\end{theorem}
\begin{proof}
逐坐标考察可知 \(\Phi\) 保持运算与单位，其核恰为全部 \(I_i\) 的交。对每个 \(i\ne j\)，取 \(u_{ij}\in I_i,v_{ij}\in I_j\) 使 \(u_{ij}+v_{ij}=1\)，并令
\[
e_i=\prod_{j\ne i}v_{ij}.
\]
空乘积为 \(1\)。于是 \(e_i\equiv1\pmod{I_i}\)，且 \(e_i\in I_j\) 对 \(j\ne i\) 成立。给定任意余数组 \((a_i+I_i)_i\)，元素 \(a=\sum_i a_ie_i\) 的第 \(i\) 个余数恰为 \(a_i+I_i\)，所以 \(\Phi\) 满射。由\cref{b1:lem:comaximal-product}确定核，再对 \(\Phi\) 应用\cref{b1:thm:ring-first-iso}，得到所述同构。

反过来，若 \(\Phi\) 满射，对 \(i\ne j\) 选择一个原像 \(a\)，使 \(a\equiv1\pmod{I_i}\)、\(a\equiv0\pmod{I_j}\)。则 \(1=(1-a)+a\in I_i+I_j\)，故两理想互素。
\end{proof}
有限性使构造中的求和和乘积都在环内有意义。若把定理中的有限族改为无限族，即使理想仍两两互素，满射性也可能失效。例如考虑同余方程组
\[
x\equiv0\pmod2,\qquad
x\equiv1\pmod p\quad\text{对每个奇素数 }p.
\]
其中任意有限多个同余式都由中国剩余定理保证有解。若整数 \(x\) 满足全部同余式，则每个奇素数都整除 \(x-1\)；非零整数只有有限多个素因子，故 \(x-1=0\)。这又与 \(x\equiv0\pmod2\) 矛盾。

\subsection{整数同余方程}
\label{b1:subsec:1004:02}
正整数 \(m,n\) 满足 \(m\mathbb Z+n\mathbb Z=\mathbb Z\)，当且仅当 \(\gcd(m,n)=1\)。若 \(um+vn=1\)，方程
\[
x\equiv a\pmod m,\qquad x\equiv b\pmod n
\]
的解可取 \(x=avn+bum\)。模 \(m\) 第一项同余于 \(a\)、第二项为零；模 \(n\) 恰好相反。\cref{b1:thm:crt}说明全部解组成模 \(mn\) 的唯一剩余类。

\ProseParagraphBreak
例如求 \(x\equiv2\pmod3\)、\(x\equiv3\pmod4\)。由于 \(4-3=1\)，可以取 \(x=2\cdot4+3\cdot(-3)=-1\)，所以答案是 \(x\equiv11\pmod{12}\)。直接核验 \(11\) 的两个余数，可同时检查系数的放置次序。

\begin{proposition}\label{b1:prop:integer-crt-compatible}
设 \(m,n\) 为正整数，\(a,b\in\mathbb Z\)，并令 \(d=\gcd(m,n)\)。同余方程组
\[
x\equiv a\pmod m,\qquad x\equiv b\pmod n
\]
有整数解，当且仅当 \(a\equiv b\pmod d\)。有解时，全部解为模 \(\operatorname{lcm}(m,n)\) 的一个剩余类。
\end{proposition}
\begin{proof}
若有解，\(d\) 同时整除 \(x-a,x-b\)，所以整除 \(a-b\)。反过来，写 \(x=a+mt\)，所求条件成为 \(mt\equiv b-a\pmod n\)。若 \(d\mid b-a\)，约去 \(d\) 后，\(m/d\) 与 \(n/d\) 互素，Bézout 等式使 \(m/d\) 在模 \(n/d\) 下可逆，因而存在 \(t\)。两个解之差必须同时被 \(m,n\) 整除，恰好被它们的最小公倍数整除，且加上这个公倍数保持两个同余式，故得到全部解。
\end{proof}
若 \(n=p_1^{a_1}\cdots p_r^{a_r}\) 为整数的素数幂分解，则
\[
\mathbb Z/n\mathbb Z\cong\prod_{i=1}^r\mathbb Z/p_i^{a_i}\mathbb Z.
\]
这使大模数的乘法、可逆性与幂等元计算分别归结到各个素数幂分量。这里用的是整数的唯一分解；下一章才将元素分解的问题推广到一般整环。

\subsection{幂等元与环的分解}
\label{b1:subsec:1004:03}
\begin{definition}\label{b1:def:idempotent}
设 \(R\) 为环，\(e\in R\)。若 \(e^2=e\)，则称 \(e\) 为\term[midengyuan]{幂等元}[idempotent element]。若一族幂等元 \((e_i)_{i\in I}\) 对任意不同指标 \(i,j\in I\) 都满足 \(e_ie_j=e_je_i=0\)，则称这族幂等元两两正交。
\end{definition}

\begin{proposition}\label{b1:prop:idempotent-product}
设 \(r\geq1\) 为整数。若交换环 \(R\) 有两两正交的幂等元 \(e_1,\ldots,e_r\)，且 \(\sum_i e_i=1\)，则
\[
R\longrightarrow\prod_{i=1}^r e_iR,\qquad
a\longmapsto(e_ia)_i
\]
为环同构。每个 \(e_iR\) 的单位是 \(e_i\)，一般不等于 \(1_R\)。反之，若 \(R_1,\ldots,R_r\) 为交换环，并令
\[
R=\prod_{i=1}^rR_i,
\qquad
e_i=(0,\ldots,0,1_{R_i},0,\ldots,0),
\]
其中 \(1_{R_i}\) 位于第 \(i\) 个坐标，则 \(e_1,\ldots,e_r\) 两两正交、各自幂等，且 \(\sum_i e_i=1_R\)。
\end{proposition}
\begin{proof}
\(e_iR\) 对加法及乘法封闭，\(e_i(e_ia)=e_ia\)，故以 \(e_i\) 为单位。映射保持乘法，因为 \((e_ia)(e_ib)=e_iab\)，并把 \(1\) 送到各分量的单位。若每个 \(e_ia=0\)，则 \(a=\sum_i e_ia=0\)，所以单射。给定 \(x_i\in e_iR\)，令 \(a=\sum_i x_i\)；由正交性，\(e_ia=x_i\)，故满射。逆映射就是各分量求和。有限直积中，第 \(i\) 坐标为单位而其余为零的元素显然满足幂等、正交及和为单位的条件。
\end{proof}
当 \(e_i\ne1_R\) 时，包含映射 \(e_iR\hookrightarrow R\) 不保持单位，因而不是本书约定下的环同态，也不是本书约定下的子环嵌入；\(e_iR\) 在上述分解中以 \(e_i\) 为单位独立成环。

\ProseParagraphBreak
这些幂等元也直接给出中国剩余定理中的理想。对每个 \(i\)，映射
\[
p_i:R\longrightarrow e_iR,\qquad a\longmapsto e_i a
\]
是满环同态，并且 \(p_i(1_R)=e_i=1_{e_iR}\)，所以保持单位。若 \(e_i a=0\)，则 \(a=(1-e_i)a\)；反过来，\(e_i(1-e_i)=0\)。因此
\[
J_i:=\ker p_i=(1-e_i)R,
\qquad e_iR\cong R/J_i.
\]
对于 \(i\ne j\)，由 \(1=(1-e_i)+e_i\) 和 \(e_i=(1-e_j)e_i\in J_j\) 可知 \(J_i+J_j=R\)。若 \(a\in\bigcap_iJ_i\)，则每个 \(e_i a=0\)，于是 \(a=\sum_i e_i a=0\)。故 \(\bigcap_iJ_i=0\)，再由\cref{b1:thm:crt}得到
\[
R\cong\prod_{i=1}^rR/J_i\cong\prod_{i=1}^r e_iR.
\]
这个同构在第 \(i\) 个坐标上的映射就是 \(p_i\)。反过来，从两两互素的理想出发时，\cref{b1:thm:crt}证明中构造的元素 \(e_i\) 未必在原环中幂等，但它们在商环 \(R/\bigcap I_i\) 中的像就是这些坐标幂等元。例如 \(R=\mathbb Z/12\mathbb Z\cong\mathbb Z/3\mathbb Z\times\mathbb Z/4\mathbb Z\) 中，\([4]\) 与 \([9]\) 分别对应 \((1,0)\)、\((0,1)\)：它们的平方各等于自身，乘积为零，和为 \([1]\)。\cref{b1:tab:crt-idempotent-components}列出了相应的分量与投影核。

\begin{table}[htbp]
\centering
\caption{\(\mathbb Z/12\mathbb Z\) 中的幂等元、分量与投影核}
\label{b1:tab:crt-idempotent-components}
\begin{tabularx}{\linewidth}{@{}c>{\centering\arraybackslash}Xcc@{}}
\toprule
幂等元 \(e_i\) & 分量 \(e_iR\)（单位） & 投影核 \(J_i=(1-e_i)R\) & 商环 \(R/J_i\) \\
\midrule
\([4]\) & \(\{[0],[4],[8]\}\)（\([4]\)） & \(([3])\) & \(\mathbb Z/3\mathbb Z\) \\
\([9]\) & \(\{[0],[3],[6],[9]\}\)（\([9]\)） & \(([4])\) & \(\mathbb Z/4\mathbb Z\) \\
\bottomrule
\end{tabularx}
\end{table}

\(e_iR\) 是保留下来的分量，\(J_i=(1-e_i)R\) 则是第 \(i\) 个投影消去的部分；表中两者不是同一个理想，而 \(R/J_i\cong e_iR\)。

\ProseParagraphBreak
整环中 \(e^2=e\) 等价于 \(e(e-1)=0\)，所以只有 \(0,1\) 两个幂等元，不能分成两个非零环的直积。反方向不成立：\(\mathbb Z/4\mathbb Z\) 含零因子，却只有 \([0],[1]\) 两个幂等元。一般非交换环要作上述直积分解，还应要求幂等元与所有元素交换，即为中心幂等元。这样证明中的 \((e_ia)(e_ib)=e_iab\) 才成立。设 \(k\) 为域。在 \(M_2(k)\) 中取 \(e=E_{11}\)、\(f=E_{22}\)，则 \(e^2=e\)、\(f^2=f\)、\(ef=fe=0\) 且 \(e+f=I_2\)，但 \(eE_{12}=E_{12}\ne E_{12}e=0\)。候选坐标映射 \(p(A)=eA\) 也不保持乘法：
\[
p(E_{12}E_{21})=E_{11},\qquad
p(E_{12})p(E_{21})=E_{12}\cdot0=0.
\]
所以正交及和为单位仍不足以取代中心性。


\begin{chaptersummary}
商环要求被抹去的加法子群同时吸收左右乘法，这给出了理想的定义及核的基本性质。交换环的素理想对应整环商，极大理想对应域商；根理想对应没有非零幂零元的商环，准素理想则要求商环中的每个非零零因子都是幂零元。素性与极大性和理想是否非零是不同的问题。理想的和、交、积是三种不同的运算，取根识别哪些元素在商环中成为幂零元，理想商则收集满足指定乘法条件的元素；取零理想作第一项就得到零化子。

\ProseParagraphBreak
中国剩余定理的关键是构造能够分别选择各余数分量的元素；商环的有限直积分解由此转化为正交幂等元的分解。下一章将从主理想的包含关系返回元素整除，考察整数式的分解规律在怎样的环中仍然成立。
\end{chaptersummary}

\BookExercises

\begin{exercise}\label{b1:ex:10-residue-hom}
设 \(m,n\) 是正整数。确定何时存在保单位环同态
\(\mathbb Z/m\mathbb Z\rightarrow\mathbb Z/n\mathbb Z\)，并在存在时给出全部这样的同态、核与像。
\end{exercise}
\begin{solution}
保单位且可加迫使 \([a]_m\mapsto[a]_n\)，所以至多一个同态。由于 \([m]_m=0\)，必须有 \([m]_n=0\)，即 \(n\mid m\)。反过来，若 \(n\mid m\)，同一公式与代表元无关，并保持加法、乘法和单位，故给出同态。它满射，核为
\(\bigl\{\,[kn]_m\mid0\leq k<m/n\,\bigr\}\)，即由 \([n]_m\) 生成的理想，含 \(m/n\) 个元素。\(n=1\) 时目标为零环，公式仍成立；\(m=1\) 时则只能有 \(n=1\)。
\end{solution}

\begin{exercise}\label{b1:ex:10-characteristic-product}
设 \(R=\mathbb Z/4\mathbb Z\times\mathbb Z/6\mathbb Z\)。求 \(R\) 的特征及自然同态 \(\mathbb Z\rightarrow R\) 的核、像，并说明该同态为什么不是满射。
\end{exercise}
\begin{solution}
\(n1_R=0\) 等价于 \(4\mid n\) 且 \(6\mid n\)，最小正解为 \(12\)，故特征为 \(12\)，核为 \(12\mathbb Z\)。\cref{b1:thm:ring-first-iso}给出像同构于 \(\mathbb Z/12\mathbb Z\)，有 \(12\) 个元素；而 \(R\) 有 \(24\) 个元素，故不满射。更具体地，\cref{b1:prop:integer-crt-compatible}说明坐标 \(([a]_4,[b]_6)\) 在像中当且仅当 \(a\equiv b\pmod2\)，因为两模数的最大公因子为二。例如 \(([0]_4,[1]_6)\) 不在像中。环的特征决定这个自然像，却不决定全环的元素个数。
\end{solution}

\begin{exercise}\label{b1:ex:10-matrix-ideals}
设 \(k\) 为域，\(n\geq2\)。证明 \(M_n(k)\) 只有零理想和全环两个双边理想，并说明它既不是整环，也不是除环。
\end{exercise}
\begin{solution}
设非零双边理想 \(I\) 含矩阵 \(A=(a_{ij})\ne0\)，选 \(a_{ij}\ne0\)。对任意 \(r,s\)，矩阵乘法给出
\(E_{ri}AE_{js}=a_{ij}E_{rs}\in I\)。再乘标量矩阵 \(a_{ij}^{-1}I_n\)，得到每个 \(E_{rs}\in I\)。它们的线性组合给出任意矩阵，故 \(I=M_n(k)\)。然而 \(E_{11}E_{22}=0\)，两因子都非零，且 \(E_{11}\) 不可逆。因此在非交换环中，“没有非平凡双边理想”与“每个非零元素都可逆”是两个不同的条件。这里零理想已是极大双边理想，商环仍为 \(M_n(k)\)，却不是除环；\cref{b1:prop:prime-max-quotient}的域商判别确实需要交换性。
\end{solution}

\begin{exercise}\label{b1:ex:10-triangular}
设 \(T\) 为域 \(k\) 上的二阶上三角矩阵环，\(N\) 为严格上三角矩阵的集合。
\begin{enumerate}
\item 求 \(T/N\)、\(T\) 的单位和幂零元，并判断对角矩阵集合是否为 \(T\) 的理想。
\item 令
\[
I=\left\{\begin{pmatrix}a&b\\0&0\end{pmatrix}:a,b\in k\right\},
\qquad
J=\left\{\begin{pmatrix}0&b\\0&c\end{pmatrix}:b,c\in k\right\}.
\]
证明 \(I,J\) 都是 \(T\) 的双边理想，并计算 \(I+J\)、\(I\cap J\)、\(IJ\) 与 \(JI\)。
\end{enumerate}
\end{exercise}
\begin{solution}
1.\ 映射
\(\begin{psmallmatrix}a&b\\0&c\end{psmallmatrix}\mapsto(a,c)\)
是到 \(k\times k\) 的满环同态，核为 \(N\)；\cref{b1:thm:ring-first-iso}给出 \(T/N\cong k\times k\)。当 \(a,c\ne0\) 时，逆矩阵为
\[
\begin{pmatrix}a^{-1}&-a^{-1}bc^{-1}\\0&c^{-1}\end{pmatrix}.
\]
若矩阵可逆，其像 \((a,c)\) 在 \(k\times k\) 中必须可逆，故这也是必要条件。矩阵的 \(r\) 次幂的对角元为 \(a^r,c^r\)，所以幂零必有 \(a=c=0\)；反过来这样的矩阵平方为零。对角矩阵集合含单位但不等于 \(T\)，因而不可能为理想；也可直接计算 \(E_{11}E_{12}=E_{12}\) 看出吸收性失败。

2.\ 取右下对角坐标的同态 \(T\to k\) 的核为 \(I\)，取左上对角坐标的同态之核为 \(J\)，故两者都是双边理想。它们相加含 \(E_{11},E_{12},E_{22}\)，所以 \(I+J=T\)；同时属于两者的矩阵恰为严格上三角矩阵，故 \(I\cap J=N\)。直接相乘可得
\[
\begin{pmatrix}a&b\\0&0\end{pmatrix}
\begin{pmatrix}0&u\\0&c\end{pmatrix}
=\begin{pmatrix}0&au+bc\\0&0\end{pmatrix},
\qquad
\begin{pmatrix}0&u\\0&c\end{pmatrix}
\begin{pmatrix}a&b\\0&0\end{pmatrix}=0.
\]
由于 \(E_{11}E_{12}=E_{12}\)，第一类乘积生成整个 \(N\)，故 \(IJ=N\)，而 \(JI=0\)。这里恰好有 \(IJ=I\cap J\)，但交换两个理想的次序以后，\(J\cap I=N\ne0=JI\)。所以即使两个双边理想之和为全环，它们的乘积仍可能随次序改变，也不能保证乘积等于交；\cref{b1:lem:comaximal-product}中的交换性确实不可省。
\end{solution}

\begin{exercise}\label{b1:ex:10-one-sided-quotient}
设 \(k\) 为域，\(R=M_2(k)\)，\(L\) 是第二列为零的矩阵集合。证明 \(L=RE_{11}\) 是左理想，并用具体代表元说明加法商群 \(R/L\) 不能按 \((A+L)(B+L)=AB+L\) 定义乘法。
\end{exercise}
\begin{solution}
右乘 \(E_{11}\) 保留矩阵第一列并把第二列置零，故 \(L=RE_{11}\)。它对减法封闭，且 \(C(AE_{11})=(CA)E_{11}\in L\)，所以是左理想。\(E_{11}\in L\)，因而 \(E_{11}+L=0+L\)。但是
\(E_{11}E_{12}=E_{12}\notin L\)，而 \(0E_{12}=0\in L\)。同一个左因子陪集的两个代表产生不同的乘积陪集，故乘法不良定义。理想的右吸收性在\cref{b1:prop:quotient-ring}中正是为排除这种现象而要求的。
\end{solution}

\begin{exercise}\label{b1:ex:10-dual-numbers}
设 \(k\) 为域，\(R=k[x]/(x^2)\)，以 \(\varepsilon\) 表示 \(x\) 的像。证明每个元素唯一写成 \(a+b\varepsilon\)，并求 \(R\) 的全部理想、单位和幂零元。利用理想对应，列出 \(k[x]\) 中全部包含 \((x^2)\) 的理想，并判断其中哪些是素理想、哪些是极大理想。
\end{exercise}
\begin{solution}
\((x^2)\) 恰由常数项和一次项均为零的多项式组成，所以每个陪集唯一由 \(a+bx\) 代表，且 \(\varepsilon^2=0\)。若 \(a\ne0\)，则
\[
(a+b\varepsilon)^{-1}=a^{-1}-ba^{-2}\varepsilon.
\]
若 \(a=0\)，元素平方为零，不可能可逆。因此一个真理想中的每个元素均为 \(b\varepsilon\)；只要它含某个 \(b\ne0\) 的元素，乘 \(b^{-1}\) 就含 \(\varepsilon\)，从而等于 \((\varepsilon)\)。全部理想为 \((0),(\varepsilon),R\)。若 \(a+b\varepsilon\) 幂零，其常数项的某次幂为零，故 \(a=0\)；反向已验证。商 \(R/(\varepsilon)\cong k\)，故 \((\varepsilon)\) 是唯一极大理想，且 \((\varepsilon)\ne(0)\)。域的唯一极大理想是 \((0)\)，而这里 \(\varepsilon\ne0\) 且 \(\varepsilon^2=0\)，所以有唯一极大理想并不意味着环是域。

商映射 \(k[x]\to R\) 的核为 \((x^2)\)。由\cref{b1:prop:ring-ideal-correspondence}，包含 \((x^2)\) 的理想恰与上述三个理想对应，依次为 \((x^2),(x),k[x]\)，没有其他情形。其中 \((x)\) 的商环是域 \(k\)，故它既是素理想也是极大理想；\((x^2)\) 不是素理想，因为 \(x^2\in(x^2)\) 而 \(x\notin(x^2)\)。全环不是真理想，故不属于这两类。
\end{solution}

\begin{exercise}\label{b1:ex:10-group-ring-augmentation}
设 \(k\) 为域，\(G\) 为群。在群环 \(k[G]\) 上定义系数和映射
\(\epsilon(\sum_g a_gg)=\sum_g a_g\)。证明它为满环同态，并证明其核作为 \(k[G]\) 的双边理想由全部 \(g-1\)、\(g\in G\) 生成。若 \(G=C_2=\{1,g\}\)，证明 \(k[G]\cong k[t]/(t^2-1)\)。当 \(\operatorname{char}k\ne2\) 时，给出 \(k[G]\cong k\times k\) 的具体同构及其对应的正交幂等元；当 \(\operatorname{char}k=2\) 时，证明增广核的平方为零，并说明群环与 \(k[\varepsilon]/(\varepsilon^2)\) 的关系。
\end{exercise}
\begin{solution}
系数和保持加法与单位；乘积的系数和为
\(\sum_{g,h}a_gb_h=(\sum_ga_g)(\sum_hb_h)\)，所以保持乘法。常数 \(a\cdot1_G\) 映到 \(a\)，故满射。每个 \(g-1\) 属于核，故它们生成的理想包含在核中。反过来，若 \(\sum_g a_g=0\)，则
\[
\sum_g a_gg=\sum_g a_g(g-1),
\]
是这些生成元的有限线性组合，得到反向包含。

现设 \(G=\{1,g\}\)、\(g^2=1\)。代入 \(t\mapsto g\) 给出满环同态 \(k[t]\to k[G]\)，且 \(t^2-1\) 在核中。商环中的关系 \(t^2=1\) 使每个偶次幂化为 \(1\)、每个奇次幂化为 \(t\)，所以每个类都由某个 \(a+bt\) 代表。该表示唯一，因为 \(t^2-1\) 的非零多项式倍数由次数公式具有至少二次，而两个这样的代表之差至多一次。又因 \(1,g\) 是群环的 \(k\)-基，诱导同态 \(k[t]/(t^2-1)\to k[G]\) 也是单射，故为同构。

若 \(\operatorname{char}k\ne2\)，则 \((t-1)+(t+1)=k[t]\)，因为两生成元之差为可逆常数 \(2\)。由\cref{b1:thm:crt}，有具体同构
\[
\begin{aligned}
k[G]&\longrightarrow k\times k,\\
a+bg&\longmapsto(a+b,a-b),
\end{aligned}
\qquad
(u,v)\longmapsto\frac{u+v}{2}+\frac{u-v}{2}g
\]
（右式为其逆）。对应 \((1,0)\)、\((0,1)\) 的正交幂等元分别为 \((1+g)/2\)、\((1-g)/2\)。增广映射是第一个坐标投影，所以其核对应 \(0\times k\)，不是平方为零的理想。

若 \(\operatorname{char}k=2\)，则 \(t+1=t-1\)、\(t^2-1=(t-1)^2\)，两个一次因子不互素，不能作上述直积分解。令 \(\varepsilon=g-1\)，可得 \(k[G]\cong k[\varepsilon]/(\varepsilon^2)\)，其中 \(\varepsilon\ne0\)、\(\varepsilon^2=0\)。增广核为 \(k\varepsilon\)，所以它的平方为零。由\cref{b1:thm:ring-first-iso}，此时的增广商环仍为 \(k\)。这两种情形的分界是域 \(k\) 中的 \(2\) 是否可逆：可逆时两个一次因子互素，商环分裂为两个分量；特征为二时两个因子重合，商环中出现非零幂零元。
\end{solution}

\begin{exercise}\label{b1:ex:10-product-ideals}
设 \(R,S\) 是交换环。证明 \(R\times S\) 的每个理想唯一为 \(I\times J\)，其中 \(I,J\) 分别为两环的理想；进一步确定其全部素理想。
\end{exercise}
\begin{solution}
若 \(K\) 为乘积环的理想，\((a,b)\in K\)，乘 \((1,0)\)、\((0,1)\) 得 \((a,0),(0,b)\in K\)。令
\(I=\bigl\{\,a\mid(a,0)\in K\,\bigr\}\)、\(J=\bigl\{\,b\mid(0,b)\in K\,\bigr\}\)，减法与吸收性直接说明两者为理想，上述拆分及相加给出 \(K=I\times J\)。投影到各坐标又保证唯一性。

若 \(K\) 是素理想，由 \((1,0)(0,1)=0\) 可知至少一个坐标幂等元属于 \(K\)；两者不能同时属于 \(K\)，否则其和 \((1,1)\in K\)。若 \((0,1)\in K\)，则 \(J=S\)，且商环同构于 \(R/I\)，由\cref{b1:prop:prime-max-quotient}知 \(I\) 素；另一种情形给出 \(R\times J\)、\(J\) 素。反过来，这两种形式的商分别是整环，故确为素理想。
\end{solution}

\begin{exercise}\label{b1:ex:10-integer-ideal-operations}
在 \(\mathbb Z\) 中取 \(I=12\mathbb Z,J=18\mathbb Z\)。计算 \(I+J\)、\(I\cap J\)、\(IJ\)、\(I^2\) 及这些理想的根，说明哪些理想不同但根相同。
\end{exercise}
\begin{solution}
因为 \(18-12=6\)，且所有两者的整数线性组合都被六整除，有 \(I+J=6\mathbb Z\)。共同倍数恰为 \(36\) 的倍数，故 \(I\cap J=36\mathbb Z\)。积理想为 \(216\mathbb Z\)，而 \(I^2=144\mathbb Z\)，直接由主理想的乘积定义得到。

所列生成数以及 \(12,18\) 的素因子集合均为 \(\{2,3\}\)。若其中任一个数整除 \(a^r\)，必有 \(2\mid a,3\mid a\)；反过来，若 \(6\mid a\)，取足够大的幂，就能包含该生成数所需的二、三的幂次。因此 \(I,J,I+J,I\cap J,IJ,I^2\) 的根全部为 \(6\mathbb Z\)。这些不同理想保留着不同幂次信息，取根以后该区别消失。
\end{solution}

\begin{exercise}\label{b1:ex:10-integer-colon}
设 \(a,b\) 为正整数，\(d=\gcd(a,b)\)。计算整数环中的理想商 \((a\mathbb Z:b\mathbb Z)\)，并判断何时将它乘以 \(b\mathbb Z\) 能恢复 \(a\mathbb Z\)。再求 \(\mathbb Z/a\mathbb Z\) 中 \([b]_a\) 的零化子及其元素个数。
\end{exercise}
\begin{solution}
写 \(a=da_0,b=db_0\)，其中 \(\gcd(a_0,b_0)=1\)。由\cref{b1:def:ideal-quotient}，\(r\) 属于所求理想当且仅当 \(a\mid rb\)，约去 \(d\) 后等价于 \(a_0\mid rb_0\)。取整数 \(u,v\) 使 \(ua_0+vb_0=1\)，两边乘 \(r\) 可知该条件迫使 \(a_0\mid r\)；反向显然成立。故
\[
(a\mathbb Z:b\mathbb Z)=(a/d)\mathbb Z.
\]
乘回 \(b\mathbb Z\) 后得到 \((ab/d)\mathbb Z\)。由于两个生成数均为正，它等于 \(a\mathbb Z\) 当且仅当 \(b/d=1\)，即 \(b\mid a\)。

在商环中，\([r]_a\) 零化 \([b]_a\) 恰好要求 \(a\mid rb\)，所以零化子由 \([a/d]_a\) 生成，其元素为
\[
\bigl\{\,[j a/d]_a\mid0\leq j<d\,\bigr\}.
\]
这些类互不相同：两者相等要求 \(a\mid(j-j')a/d\)，即 \(d\mid j-j'\)，在给定范围内只能 \(j=j'\)。故零化子恰有 \(d\) 个元素。例如 \(a=24,b=6\) 时，理想商为 \(4\mathbb Z\)，商环中零化子有六个元素。
\end{solution}

\begin{exercise}\label{b1:ex:10-nonprincipal-ideal}
在 \(R=\mathbb Z[x]\) 中令 \(I=(2,x)\)。证明 \(I\) 不是主理想，并计算 \(I^2\)、\(\sqrt{I^2}\) 与 \((I^2:I)\)。
\end{exercise}
\begin{solution}
映射 \(\psi:R\to\mathbb F_2\)、\(f\mapsto[f(0)]_2\) 是满环同态，核为 \(I\)。因此 \(I\) 是真素理想，且 \(\sqrt I=I\)。若 \(I=(f)\)，则由 \(2,x\in I\) 可取 \(g,h\in\mathbb Z[x]\) 使
\[
2=fg,\qquad x=fh.
\]
第一式的两因子均非零，次数公式给出 \(\deg f+\deg g=0\)，所以 \(f,g\) 都是整数常数，且 \(f=\pm1\) 或 \(\pm2\)。再比较第二式的一次项系数，得到 \(1=f h_1\)，其中 \(h_1\in\mathbb Z\) 是 \(h\) 的一次项系数，故只能有 \(f=\pm1\)。于是 \((f)=R\)，与 \(I\ne R\) 矛盾。

由理想积的定义，
\[
I^2=(4,2x,x^2).
\]
其中的多项式恰是常数项被 \(4\) 整除、一次项系数为偶数的多项式。因为 \(I^2\subseteq I=\sqrt I\)，有 \(\sqrt{I^2}\subseteq I\)；反过来，每个 \(f\in I\) 都满足 \(f^2\in I^2\)，故 \(\sqrt{I^2}=I\)。

若 \(f\in(I^2:I)\)，则 \(fx\in I^2\)，由一次项系数的条件得到 \(f(0)\) 为偶数，亦即 \(f\in I\)。反过来，\(I\cdot I=I^2\) 表明每个 \(f\in I\) 都属于 \((I^2:I)\)。所以 \((I^2:I)=I\)。
\end{solution}

\begin{exercise}\label{b1:ex:10-units-lift-nil}
设 \(R\) 为交换环，理想 \(I\) 中每个元素都是幂零元。证明 \(a\in R\) 可逆，当且仅当 \(a+I\) 在 \(R/I\) 中可逆。再给出删去幂零假设后结论失败的例子。
\end{exercise}
\begin{solution}
单位在商同态下仍为单位。反过来，若 \(a+I\) 可逆，选逆的一个代表 \(b\)，则 \(ab=1+n\)，其中 \(n\in I\)。这里只需误差元素 \(n=ab-1\) 自己幂零，不要求存在统一的正整数 \(N\) 使整个理想满足 \(I^N=0\)。对这个 \(n\) 取 \(m\geq1\) 使 \(n^m=0\)，有限几何级数公式给出
\[
(1+n)^{-1}=1-n+n^2-\cdots+(-n)^{m-1}.
\]
于是 \(a\) 的逆为 \(b(1+n)^{-1}\)，左右相乘均为一。若 \(R=\mathbb Z,I=5\mathbb Z,a=2\)，则 \(a+I\) 在商域中可逆，但整数二不可逆；这个理想不满足题设的幂零条件。
\end{solution}

\begin{exercise}\label{b1:ex:10-primary-integers}
对整数 \(n\geq2\)，判定何时 \(n\mathbb Z\) 是准素理想。用 \(n=12\) 写出直接违反准素性条件的两个整数，并用\secref{b1:sec:1003}中的准素商环判别比较 \(\mathbb Z/4\mathbb Z\) 与 \(\mathbb Z/12\mathbb Z\) 中非零零因子的幂零性。
\end{exercise}
\begin{solution}
若 \(n=p^r\) 为素数幂，且 \(ab\in n\mathbb Z\)、\(a\notin n\mathbb Z\)，则 \(p^r\mid ab\) 而 \(p^r\nmid a\)，整数的素数分解迫使 \(p\mid b\)。写 \(b=pc\)，就有 \(b^r=p^rc^r\in n\mathbb Z\)，所以 \(n\mathbb Z\) 准素。

若 \(n\) 不是素数幂，可写 \(n=p^r m\)，其中 \(m>1\)、\(r\geq1\)、\(p\nmid m\)。取 \(a=p^r,b=m\)，则 \(ab=n\in n\mathbb Z\)，但 \(a\notin n\mathbb Z\)，而 \(b\) 的任意正次幂都不被 \(p\) 整除，因而不在 \(n\mathbb Z\) 中。这违反准素性。故 \(n\mathbb Z\) 准素当且仅当 \(n\) 为素数幂。

在 \(n=12\) 时可取 \(a=4,b=3\)：\(4\cdot3\in12\mathbb Z\)，\(4\notin12\mathbb Z\)，且没有 \(3^j\) 是十二的倍数。模四时唯一的非零零因子是 \([2]_4\)，其平方为零；模十二时 \([6]_{12}\) 是平方为零的非零零因子，但 \([3]_{12}\) 被非零的 \([4]_{12}\) 零化而所有正次幂都非零，故零因子并非全都幂零。
\end{solution}

\begin{exercise}\label{b1:ex:10-compatible-congruences}
求全部满足 \(x\equiv3\pmod8\)、\(x\equiv7\pmod{12}\) 的整数，并判断把第二个余数改为 \(5\) 后是否有解。
\end{exercise}
\begin{solution}
两模数的最大公因子为 \(4\)，且 \(3\equiv7\pmod4\)，由\cref{b1:prop:integer-crt-compatible}可解。写 \(x=3+8t\)，代入得 \(8t\equiv4\pmod{12}\)，即 \(2t\equiv1\pmod3\)，所以 \(t\equiv2\pmod3\)。全部解为 \(x=19+24u\)，\(u\in\mathbb Z\)。若第二余数为 \(5\)，则 \(3\not\equiv5\pmod4\)，不满足该命题的必要条件，无解。
\end{solution}

\begin{exercise}\label{b1:ex:10-mod-seventy-two}
求 \(\mathbb Z/72\mathbb Z\) 的全部幂等元与幂零元，并计算其单位个数。
\end{exercise}
\begin{solution}
由\cref{b1:thm:crt}，环同构于 \(\mathbb Z/8\mathbb Z\times\mathbb Z/9\mathbb Z\)，同构为 \([a]_{72}\mapsto([a]_8,[a]_9)\)。若 \(a(a-1)\) 被素数幂 \(p^r\) 整除，因相邻整数互素，该素数幂必全部整除其中一个因子。因此模 \(8\) 和模 \(9\) 各只有 \(0,1\) 两个幂等元；四组坐标提升后恰为
\[
\begin{array}{c|cccc}
\text{模 }72\text{ 的类}&[0]_{72}&[1]_{72}&[9]_{72}&[64]_{72}\\\hline
\text{CRT 坐标}&(0,0)&(1,1)&(1,0)&(0,1)
\end{array}
\]
（各坐标分别在模 \(8\)、模 \(9\) 的环中）。特别地，\([9]_{72}\) 保留模 \(8\) 的分量，\([64]_{72}\) 保留模 \(9\) 的分量。

若 \(72\mid a^m\)，则 \(2,3\) 都整除 \(a\)，故 \(6\mid a\)；反过来 \(6\mid a\) 时 \(72\mid a^3\)。所以幂零元为 \([6j]_{72}\)、\(0\leq j<12\)。单位在模 \(8\) 下有 \(1,3,5,7\) 四个，在模 \(9\) 下有 \(1,2,4,5,7,8\) 六个；直积元素可逆当且仅当两坐标均可逆，因此总数为 \(24\)。
\end{solution}

\begin{exercise}\label{b1:ex:10-comaximal-powers}
设交换环 \(R\) 中 \(I+J=R\)，\(r,s\) 为正整数。证明 \(I^r+J^s=R\)，并给出自然同构
\(R/(I^rJ^s)\cong R/I^r\times R/J^s\)。
\end{exercise}
\begin{solution}
写 \(a+b=1\)，其中 \(a\in I,b\in J\)。展开
\(1=(a+b)^{r+s-1}\)，每一项的 \(a\) 次数至少为 \(r\)，或 \(b\) 次数至少为 \(s\)，所以整个和属于 \(I^r+J^s\)。这证明两幂理想互素。对这两个理想应用\cref{b1:thm:crt}，映射 \(x\mapsto(x+I^r,x+J^s)\) 满射，核为 \(I^r\cap J^s=I^rJ^s\)，故诱导所求同构。此题说明互素关系可以提升到任意正整数幂，而不只是 \(\sqrt{I^r}=\sqrt I\) 所表达的根相同。
\end{solution}

\begin{exercise}\label{b1:ex:10-function-idempotents}
设 \(k\) 为域，\(X\) 为集合。求函数环 \(k^X\) 的全部幂等元。对 \(A\subseteq X\)，利用对应的幂等元解释分解 \(k^X\cong k^A\times k^{X\setminus A}\)，并在 \(X\) 有 \(n\) 个元素时求幂等元个数。
\end{exercise}
\begin{solution}
函数 \(e\) 幂等当且仅当每个 \(x\) 都满足 \(e(x)^2=e(x)\)。域中只有零、一两个幂等元，所以 \(e\) 恰为某个子集 \(A\) 的示性函数：在 \(A\) 上为一，其他点为零。对任意 \(f\in k^X\)，逐点相乘得到
\[
(ef)(x)=\begin{cases}f(x),&x\in A,\\0,&x\notin A,\end{cases}
\qquad
((1-e)f)(x)=\begin{cases}0,&x\in A,\\f(x),&x\notin A.\end{cases}
\]
因此这两个投影分别保留 \(A\) 上和其补集上的函数值。\(e\) 与 \(1-e\) 正交且和为一，\cref{b1:prop:idempotent-product}给出 \(k^X\cong e k^X\times(1-e)k^X\)。两个分量分别通过限制到 \(A\) 与 \(X\setminus A\) 识别为题述函数环；逆映射把两个函数在互不相交的定义域上拼成一个函数。因此有限 \(X\) 的幂等元与其全部子集一一对应，共 \(2^n\) 个。空集时函数环为零环，上述计数仍给出唯一元素。
\end{solution}

% !TeX root = ../../Book1.tex
% 纲要标题：## 第11章　整除、主理想整环与唯一分解
% 纲要位置：计划纲要.md，第 414 行。

\chapter{\texorpdfstring{\mbox{整除、主理想整环}}{整除、主理想整环}\texorpdfstring{\mbox{与唯一分解}}{与唯一分解}}
\label{b1:ch:11}
\BookChapterNumber{b1:ch:11}

\begin{chapterabstract}
整数的素因数分解在一般整环中需要重新检验：不可约元未必为素元，最大公因子也未必给出 Bézout 等式。本章比较 Euclid 整环、主理想整环与唯一分解整环，说明它们的蕴含关系、反例和可执行的整除计算。
\end{chapterabstract}

\begin{learninggoals}
\begin{itemize}
\item 区分不可约元与素元，区分最大公因子与 Bézout 关系。
\item 在整数、域上的一元多项式和 Gauss 整数中执行 Euclid 算法及回代。
\item 证明 Euclid 整环是主理想整环、主理想整环是唯一分解整环，并指出分解存在性与唯一性各自的依据。
\item 用具体理想与范数计算检验主性或唯一分解性的失败，理解理想分解问题的来源。
\end{itemize}
\end{learninggoals}

在整数中，\(2\) 与 \(3\) 的最大公因子是 \(1\)，而且能找到整数 \(u,v\) 使 \(2u+3v=1\)。把同样的话移到 \(\mathbb Z[x]\)，\(2\) 与 \(x\) 仍没有非单位公因子，却不可能有 \(2u(x)+xv(x)=1\)：令 \(x=0\) 就会得到偶数等于 \(1\)。最大公因子、Bézout 等式和理想由一个元素生成，在一般整环中不再是同一回事。本章从整除与 Euclid 算法出发，分辨这些性质之间真正成立的推论，再考察元素何时能唯一分解。

\ProseParagraphBreak

本章使用\chapref{b1:ch:10}建立的整环、单位、理想、商环及多项式次数性质；Gauss 整数和非 Euclid 的例子还用到复数的共轭与模。非 Euclid 主理想整环的例子用于辨析条件强弱。\(\mathbb Z[x]\) 的非主性在本章直接证明，唯一分解性则明确留待下一章的多项式定理。

\ProseParagraphBreak
整除、主理想与唯一分解，可结合 Dummit 与 Foote 的《Abstract Algebra》\cite[Sections 8.1--8.3]{DummitFoote2003}学习。整数的 Euclid 算法、扩展算法与模逆可读 Shoup 的《A Computational Introduction to Number Theory and Algebra》\cite[Sections 4.1--4.3]{Shoup2008NumberTheoryAlgebra}；域上一元多项式的 Euclid 算法、模逆及中国剩余计算见同书\cite[Sections 17.3--17.4]{Shoup2008NumberTheoryAlgebra}。唯一分解整环、Euclid 整环与主理想整环的结构性讨论见该书\cite[Section 16.9.1]{Shoup2008NumberTheoryAlgebra}。

% !TeX root = ../../Book1.tex
% 纲要标题：### 11.1 整除与因子
% 纲要位置：计划纲要.md，第 416 行。

\section{整除与因子}
\label{b1:sec:1101}

% 纲要内容（仅作写作计划，尚非教材正文）：
%
% * 相伴、不可约元与素元
% * 最大公因子与 Bézout 关系
% * 元素分解与理想语言的区别
%

整数的素因数分解同时包含两个断言：每个非零非单位都能分解，所得不可约因子又在适当意义下唯一。在一般整环中，这两个断言必须分别考察。本节先澄清整除、不可约与素性之间的关系；本章的环均为交换整环，除非另有说明。

\subsection{相伴、不可约元与素元}
\label{b1:subsec:1101:01}

\begin{definition}\label{b1:def:divisibility-associates}
设 \(R\) 是整环，\(a,b\in R\)。若存在 \(c\in R\) 使 \(b=ac\)，则称 \(a\) 整除 \(b\)，记为 \(a\mid b\)；这一关系称为\term[zhengchu]{整除关系}[divisibility]。若存在单位 \(u\in R^\times\) 使 \(a=ub\)，则称 \(a,b\) \term[xiangban]{相伴}[associate]。
\end{definition}
\symidx{\(a\mid b\)：\(a\) 整除 \(b\)}

每个元素都整除 \(0\)，而 \(0\) 只整除 \(0\)。单位整除所有元素。相伴是等价关系：单位 \(1\)、\(u^{-1}\) 和单位的乘积分别给出自反性、对称性和传递性。在 \(\mathbb Z\) 中，相伴就是相差一个符号；在域中，任意两个非零元素都相伴。

\begin{proposition}\label{b1:prop:associate-ideals}
设 \(R\) 为整环，\(a,b\in R\)。则 \(a\mid b\) 当且仅当 \((b)\subseteq(a)\)。此外，以下三条件等价：\(a,b\) 相伴；\(a\mid b\) 且 \(b\mid a\)；\((a)=(b)\)。
\end{proposition}
\begin{proof}
若 \(b=ac\)，则 \(br=a(cr)\)，故 \((b)\subseteq(a)\)。反之，该包含使 \(b\in(a)\)，于是 \(b=ac\)。相伴当然给出相互整除。若 \(a=bc\)、\(b=ad\) 且 \(a\ne0\)，则 \(a=acd\)，整环中的消去律给出 \(cd=1\)，故 \(c\) 为单位。若 \(a=0\)，相互整除迫使 \(b=0\)，两者也相伴。主理想相等与相互整除的等价性来自第一句话。
\end{proof}

\begin{definition}\label{b1:def:irreducible-prime}
设 \(p\) 是整环 \(R\) 中的非零非单位。
\begin{enumerate}
\item 若对每个满足 \(p=ab\) 的 \(a,b\in R\)，二者中至少一个为单位，则称 \(p\) 为\term[bukeyueyuan]{不可约元}[irreducible element]。
\item 若对任意 \(a,b\in R\)，\(p\mid ab\) 都蕴含 \(p\mid a\) 或 \(p\mid b\)，则称 \(p\) 为\term[suyuan]{素元}[prime element]。
\end{enumerate}
\end{definition}

不可约性限制 \(p\) 自身的分解，素性涉及 \(p\) 对任意乘积的整除。两者并非同一个定义。单位被排除，是因为在分解中插入或删去单位不应改变因子的实质。

\begin{proposition}\label{b1:prop:prime-implies-irreducible}
设 \(R\) 为整环，\(p\in R\) 为非零非单位。若 \(p\) 是素元，则 \(p\) 是不可约元；而且 \(p\) 是素元，当且仅当主理想 \((p)\) 是 \(R\) 的素理想。
\end{proposition}
\begin{proof}
设 \(p=ab\)。素性给出 \(p\mid a\) 或 \(p\mid b\)。若 \(a=pc\)，则 \(p=pcb\)，由 \(p\ne0\) 消去得到 \(cb=1\)，所以 \(b\) 为单位；另一种情形交换 \(a,b\) 即得。又因 \(p\) 非单位，\((p)\ne R\)；条件 \(ab\in(p)\) 蕴含 \(a\in(p)\) 或 \(b\in(p)\)，逐项翻译后正是素元的定义。
\end{proof}

后面将证明，在主理想整环和唯一分解整环中，不可约元也都是素元；同时会看到 \(\mathbb Z[\sqrt{-5}]\) 中的不可约元 \(2\) 并不是素元。此处不能把整数中的直觉当作一般整环的结论。

\subsection{最大公因子与 Bézout 关系}
\label{b1:subsec:1101:02}

\begin{definition}\label{b1:def:gcd}
设 \(a,b\) 为整环 \(R\) 的两个元素。如果 \(d\in R\) 满足 \(d\mid a\)、\(d\mid b\)，并且每个同时整除 \(a,b\) 的元素 \(c\in R\) 都整除 \(d\)，就称 \(d\) 为 \(a,b\) 的\term[zuidagongyinzi]{最大公因子}[greatest common divisor]。如果 \(1\) 是 \(a,b\) 的最大公因子，就称 \(a,b\) \term[husu-yuansu]{互素}[relatively prime]。
\end{definition}

这里的“最大”依整除关系而言，与实数大小无关。最大公因子若存在，由\cref{b1:prop:associate-ideals}可知它在相伴意义下唯一；未经额外约定，\(\gcd(a,b)\) 表示任选的一个最大公因子。在整数中通常取非负者，在域上的多项式环中通常取首一者；约定 \(\gcd(0,0)=0\)。有限多个元素的最大公因子同样按“公共因子都整除它”定义。
\symidx{\(\gcd(a,b)\)：最大公因子（相伴意义下）}

\begin{proposition}\label{b1:prop:bezout-gcd}
设 \(R\) 为整环，\(a,b,d\in R\)。若 \(d\mid a\)、\(d\mid b\)，且存在 \(u,v\in R\) 使
\begin{equation}\label{b1:eq:bezout-ring}
d=ua+vb,
\end{equation}
则 \(d\) 为 \(a,b\) 的最大公因子。特别地，\((a,b)=R\) 蕴含 \(a,b\) 互素。
\end{proposition}
\begin{proof}
任一公因子 \(c\) 整除 \(a,b\)，于是也整除线性组合 \(ua+vb=d\)。这与已给出的 \(d\mid a,b\) 正好构成最大公因子的两项要求。若 \((a,b)=R\)，则 \(1=ua+vb\)；又 \(1\) 整除一切元素，所以取 \(d=1\) 即得。
\end{proof}

设 \(R\) 为整环，\(a,b\in R\)，并设 \(d\) 是 \(a,b\) 的最大公因子。若存在 \(u,v\in R\) 使 \(d=ua+vb\)，则称这个等式为 \(a,b\) 的\term[bezout-guanxi]{Bézout 关系}[Bézout identity]。它比“\(d\) 是最大公因子”多提供了一个线性组合表达式。在 \(\mathbb Z\) 中，\(\gcd(30,18)=6\)，而 \(6=2\cdot18-30\) 给出相应关系。一般整环中，即使最大公因子存在，也未必能如此表示；下面的例子正说明这一点。

\subsection{元素分解与理想}
\label{b1:subsec:1101:03}

整除关系在主理想之间反向排列，因此“共同整除”与“共同生成”不可混淆。对两个非零元素 \(a,b\)，若 \((a,b)=(d)\)，则 \(d\) 整除 \(a,b\)，且 \(d\) 属于它们生成的理想，\cref{b1:prop:bezout-gcd}遂给出最大公因子和 Bézout 关系。反过来，知道最大公因子并不保证这个二生成理想是主理想。

\begin{proposition}\label{b1:prop:integer-poly-nonprincipal}
在 \(\mathbb Z[x]\) 中，\(2,x\) 的最大公因子为 \(1\)，但 \((2,x)\) 是非主理想，特别地不等于全环。
\end{proposition}
\begin{proof}
若非零多项式 \(c\) 同时整除 \(2,x\)，由\cref{b1:prop:polynomial-degree}的次数相加公式，\(c\mid2\) 迫使 \(c\) 是常数，且为 \(\pm1\) 或 \(\pm2\)。后两者不整除 \(x\)，因为 \(x\) 的一次项系数为 \(1\)。故 \(\gcd(2,x)=1\)。然而若有 \(1=2f(x)+x\cdot g(x)\)，比较常数项就得到 \(1=2f(0)\)，不可能。因此 \((2,x)\ne\mathbb Z[x]\)。若 \((2,x)=(h)\)，则 \(h\) 是 \(2,x\) 的公因子，因而为单位；这又会使该理想等于全环，矛盾。
\end{proof}

互素并不蕴含生成单位理想；一个理想仅用两个元素生成，也可能无法缩减成一个生成元。主理想整环要求每个理想都是主理想，不论是否预先知道它有限生成，因此排除了这种情形。仅要求二生成理想都是主理想，只能进一步推出所有有限生成理想都是主理想：对生成元的个数作归纳，每次将已有的一个生成元与下一个合并即可；这一条件还不足以保证所有理想都是主理想。\cref{b1:ex:11-factorization-nonexistence}将给出一个所有有限生成理想都为主、却不是主理想整环的例子。

\ProseParagraphBreak
不可约性也可以用主理想说明。设 \(R\) 为整环，\(p\in R\) 为非零非单位。由整除与主理想包含的反向对应可知，\(p\) 不可约，当且仅当 \((p)\) 在 \(R\) 的真主理想中为极大元：若 \((p)\subseteq(d)\subsetneq R\)，则 \(p=dc\)，其中 \(d\) 非单位；不可约性迫使 \(c\) 为单位，故 \((d)=(p)\)。反过来，若 \(p=ab\) 且 \(a,b\) 都非单位，则 \((p)\subsetneq(a)\subsetneq R\)，与这一极大性矛盾。这里比较的只是主理想；后面在主理想整环中，所有理想都可纳入比较，才得到 \((p)\) 是极大理想。

\ProseParagraphBreak
后文将证明，在 \(R=\mathbb Z[\sqrt{-5}]\) 中，\(2\) 不可约，而非主理想 \(I=(2,1+\sqrt{-5})\) 满足 \((2)\subsetneq I\subsetneq R\)；这与 \((2)\) 在真主理想中极大并不矛盾。

% !TeX root = ../../Book1.tex
% 纲要标题：### 11.2 Euclid 整环
% 纲要位置：计划纲要.md，第 422 行。

\section{Euclid 整环}
\label{b1:sec:1102}

% 纲要内容（仅作写作计划，尚非教材正文）：
%
% * 除法算法与 Euclid 算法
% * 整数环及一元多项式环
% * Gauss 整数的例子
%

整数中求最大公因子的辗转相除法，依靠的是余数不断变小。推广这个方法时，真正需要的不是元素本身有大小，而是给非零余数配一个不能无限下降的非负整数。本节把这一要求写成公理，并在整数、多项式和复平面格点中实际执行除法。

\subsection{除法算法与 Euclid 算法}
\label{b1:subsec:1102:01}

\begin{definition}\label{b1:def:euclidean-domain}
设 \(R\) 是交换整环。若存在函数
\[
\delta:R\setminus\{0\}\longrightarrow\mathbb N
\]
使任意 \(a,b\in R\)、\(b\ne0\) 都可写成 \(a=bq+r\)，其中 \(q,r\in R\)，并且 \(r=0\) 或 \(\delta(r)<\delta(b)\)，则称 \(R\) 为\term[euclid-zhenghuan]{Euclid 整环}[Euclidean domain]，称 \(\delta\) 为\term[euclid-hanshu]{Euclid 函数}[Euclidean function]。
\end{definition}

本书不额外要求 \(\delta(a)\leq\delta(ab)\)（\(a,b\ne0\)）。有些文献把这个乘法单调性并入定义。两种约定对函数的要求不同，但得到的 Euclid 整环类相同。事实上，给定满足上述除法条件的 \(\delta\)，可对每个 \(a\in R\setminus\{0\}\) 定义
\[
\widetilde\delta(a)\coloneqq\min_{c\in R\setminus\{0\}}\delta(ac).
\]
由于 \(R\) 是整环，取最小值的集合是非空的非负整数集合。若 \(a,b\ne0\)，\(ab\) 的每个非零倍数也是 \(a\) 的非零倍数，故
\(\widetilde\delta(a)\leq\widetilde\delta(ab)\)。又设 \(a\in R\)、\(b\in R\setminus\{0\}\)，取 \(c\ne0\) 使 \(\delta(bc)=\widetilde\delta(b)\)。用 \(\delta\) 作带余除法，得到
\[
a=(bc)q+r=b(cq)+r,
\qquad r=0\ \text{或}\ \delta(r)<\delta(bc).
\]
当 \(r\ne0\) 时，在最小值的定义中取乘子 \(1\)，可知 \(\widetilde\delta(r)\leq\delta(r)<\widetilde\delta(b)\)。因此 \(\widetilde\delta\) 同时满足除法条件与乘法单调性\cite[Theorem 2.1]{ConradRemarksEuclideanDomains}。商和余数通常不唯一，Euclid 函数也未必唯一；对带余除法的存在性证明，还须另行判断是否给出了可执行的求商步骤。

\ProseParagraphBreak
设 \(R\) 为带 Euclid 函数 \(\delta\) 的 Euclid 整环，\(a,b\in R\) 不全为零。从 \(a,b\) 出发，不断以上一次非零余数作除数，直到余数为零；每次所得非零余数的 \(\delta\) 值都严格下降。这个过程称为\term[euclid-suanfa]{Euclid 算法}[Euclidean algorithm]，也称辗转相除法。
\termidx[zhanzhuanxiangchufa]{辗转相除法}

\begin{proposition}[Euclid 算法]\label{b1:prop:euclidean-algorithm}
设 \(R\) 为 Euclid 整环，\(a,b\in R\) 不全为零。依次作带余除法，直到余数为零，则最后一个非零余数 \(d\) 是 \(a,b\) 的最大公因子，并可由回代求得 \(u,v\in R\)，使 \(d=ua+vb\)。
\end{proposition}
\begin{proof}
若 \(b=0\)，取 \(d=a\)、\(u=1\)、\(v=0\)。否则置 \(r_0=a,r_1=b\)，并在 \(r_j\ne0\) 时选取
\[
r_{j-1}=q_jr_j+r_{j+1},
\qquad r_{j+1}=0\ \text{或}\ \delta(r_{j+1})<\delta(r_j).
\]
非负整数 \(\delta(r_1),\delta(r_2),\ldots\) 不可能无限严格下降，所以某一步得到 \(r_{m+1}=0\)、\(r_m\ne0\)。等式表明，一个元素同时整除 \(r_{j-1},r_j\)，当且仅当它同时整除 \(r_j,r_{j+1}\)：两个方向分别使用减法和加法。因此每一对相邻余数的公因子集合相同，最终等于 \(r_m,0\) 的公因子集合，故 \(d=r_m\) 为所求最大公因子。

为明确回代，令 \((u_0,v_0)=(1,0)\)、\((u_1,v_1)=(0,1)\)，并递推
\begin{equation}\label{b1:eq:extended-euclid-recurrence}
u_{j+1}=u_{j-1}-q_ju_j,\qquad
v_{j+1}=v_{j-1}-q_jv_j.
\end{equation}
由 \(r_0=u_0a+v_0b\)、\(r_1=u_1a+v_1b\) 开始归纳，除法等式给出每个 \(r_j=u_ja+v_jb\)。取 \(j=m\) 即得所求关系。
\end{proof}

在求余数的同时按\eqref{b1:eq:extended-euclid-recurrence}更新两组系数，便得到\term[kuozhan-euclid-suanfa]{扩展 Euclid 算法}[extended Euclidean algorithm]。\cref{b1:alg:extended-euclid}同时输出最大公因子与 Bézout 系数。Euclid 整环的定义只保证所需的商和余数存在；要实际运行算法，还须能有效完成环运算和符合下降条件的带余除法。

\begin{algorithm}[htbp]
\caption{扩展 Euclid 算法}
\label{b1:alg:extended-euclid}
\begin{algorithmsteps}{配有 Euclid 函数 \(\delta\) 的 Euclid 整环 \(R\)，及不全为零的 \(a,b\in R\)；其中环运算与符合 \(\delta\) 下降条件的带余除法均可有效计算。}{\(a,b\) 的一个最大公因子 \(d\) 及 \(u,v\in R\)，使 \(d=ua+vb\)。}
\State 置 \((r_0,u_0,v_0)=(a,1,0)\)、\((r_1,u_1,v_1)=(b,0,1)\)，并令 \(j=1\)。
\While{\(r_j\ne0\)}
\State 求 \(q_j,r_{j+1}\)，使 \(r_{j-1}=q_jr_j+r_{j+1}\)，且 \(r_{j+1}=0\) 或 \(\delta(r_{j+1})<\delta(r_j)\)。
\State 置 \(u_{j+1}=u_{j-1}-q_ju_j\)、\(v_{j+1}=v_{j-1}-q_jv_j\)。
\State 令 \(j\leftarrow j+1\)。
\EndWhile
\State \Return \(d=r_{j-1},\ u=u_{j-1},\ v=v_{j-1}\)。
\end{algorithmsteps}
\end{algorithm}

\subsection{整数环及一元多项式环}
\label{b1:subsec:1102:02}

在 \(\mathbb Z\) 中，取 \(\delta(n)=\lvert n\rvert\)。对 \(b\ne0\)，普通整数除法可取 \(0\leq r<\lvert b\rvert\)，故它是 Euclid 函数。以 \(252,198\) 为例，四次除法所取的商依次是 \(q_1=1,q_2=3,q_3=1,q_4=2\)：
\[
252=198+54,\qquad 198=3\cdot54+36,\qquad
54=36+18,\qquad36=2\cdot18.
\]
按\eqref{b1:eq:extended-euclid-recurrence}同步更新系数，得到下表；每行都满足 \(r_j=252u_j+198v_j\)。
\begin{center}
\begin{tabular}{rrrr}
\toprule
\(j\) & \(r_j\) & \(u_j\) & \(v_j\) \\
\midrule
\(0\) & \(252\) & \(1\) & \(0\) \\
\(1\) & \(198\) & \(0\) & \(1\) \\
\(2\) & \(54\) & \(1\) & \(-1\) \\
\(3\) & \(36\) & \(-3\) & \(4\) \\
\(4\) & \(18\) & \(4\) & \(-5\) \\
\(5\) & \(0\) & \(-11\) & \(14\) \\
\bottomrule
\end{tabular}
\end{center}
最后非零余数为 \(18\)；同一结果也可由回代读出：
\begin{equation}\label{b1:eq:integer-bezout-example}
18=54-36=4\cdot54-198=4\cdot252-5\cdot198.
\end{equation}
这里不仅算出最大公因子，也获得了以后解线性同余式所需的整数系数。

\begin{proposition}[多项式带余除法]\label{b1:prop:field-poly-division}
设 \(F\) 为域，\(f,g\in F[x]\)，且 \(g\ne0\)。存在唯一的 \(q,r\in F[x]\)，使
\[
f=gq+r,\qquad r=0\ \text{或}\ \deg r<\deg g.
\]
因此 \(F[x]\) 关于非零多项式上的次数函数为 Euclid 整环。
\end{proposition}
\begin{proof}
若 \(f=0\) 或 \(\deg f<\deg g\)，取 \(q=0,r=f\)。其余情形设 \(\deg f=n\)、\(\deg g=m\leq n\)，首项系数分别为 \(a,b\)。因为 \(F\) 是域且 \(b\ne0\)，可以减去 \(ab^{-1}x^{n-m}g\)，把 \(f\) 的最高次项消去。所得
\[
f_1=f-ab^{-1}x^{n-m}g
\]
为零或次数严格小于 \(n\)。对次数作归纳，写成 \(f_1=gq_1+r\)，其中 \(r\) 满足要求，则
\(q=ab^{-1}x^{n-m}+q_1\) 给出所求分解。次数是非负整数，故消项步骤只能进行有限次。

若又有 \(f=gq'+r'\) 且余数满足相同次数限制，则
\(g(q-q')=r'-r\)。若 \(q-q'\ne0\)，由\cref{b1:prop:polynomial-degree}，非零多项式的次数相加，使左端次数至少为 \(\deg g\)，而右端为零或次数小于 \(\deg g\)，矛盾。因此 \(q=q'\)，再得 \(r=r'\)。
\end{proof}

域假设在首项系数 \(b\) 的求逆处使用。例如在 \(\mathbb Z[x]\) 中，\(x\) 无法除以常数 \(2\) 得到次数小于 \(0\) 的余数：这种余数只能为零，而 \(x\notin2\mathbb Z[x]\)。这不否认某些特殊除数可做除法，例如首一多项式的首项系数就是单位。

\ProseParagraphBreak
在 \(\mathbb Q[x]\) 中取 \(f=x^3-2x+1\)、\(g=x^2-1\)。逐项相减得到
\[
f=xg+(-x+1),\qquad g=(-x-1)(-x+1).
\]
所以首一最大公因子为 \(x-1\)，且 \(x-1=xg-f\)。从最后非零余数 \(-x+1\) 改为首一者，需要同时将 Bézout 系数乘以 \(-1\)。这种归一化不影响生成的主理想。

\subsection{Gauss 整数的例子}
\label{b1:subsec:1102:03}

考虑复数的子环
\[
\mathbb Z[i]=\{\,a+bi\mid a,b\in\mathbb Z\,\},\qquad i^2=-1,
\]
称为\term[gauss-zhengshu]{Gauss 整数环}[ring of Gaussian integers]。它是整环，因为它嵌入 \(\mathbb C\)。定义范数
\[
N(a+bi)\coloneqq a^2+b^2.
\]
由 \(N(z)=z\overline z\) 得 \(N(zw)=N(z)N(w)\)，非零元素的范数为正整数。一个元素是单位当且仅当范数为 \(1\)：必要性由逆元与范数的乘法性得到；充分性由 \(z\overline z=1\) 得到。因此单位恰为 \(\pm1,\pm i\)。
\symidx{\(\mathbb Z[i]\)：Gauss 整数环}

若 \(\alpha\mid\beta\)，则 \(N(\alpha)\mid N(\beta)\)；反向推断一般不成立。即使范数相同，两个元素也未必相伴。例如 \(N(2+i)=N(2-i)=5\)，但
\[
\frac{2-i}{2+i}=\frac{3-4i}{5}\notin\mathbb Z[i],
\]
所以 \(2+i\nmid2-i\)，二者也不相伴。范数为整数素数足以证明一个 Gauss 整数不可约，却不是必要条件。例如 \(3\) 在 \(\mathbb Z[i]\) 中不可约（见\cref{b1:ex:11-gaussian-rational-prime}），但 \(N(3)=9\)。

\begin{proposition}\label{b1:prop:gaussian-euclidean}
设
\[
R=\mathbb Z[i]=\{\,a+bi\mid a,b\in\mathbb Z\,\},\qquad
N(a+bi)=a^2+b^2.
\]
则 Gauss 整数环 \(R\) 关于 \(R\setminus\{0\}\) 上的范数 \(N\) 是 Euclid 整环。
\end{proposition}
\begin{proof}
对 \(\alpha,\beta\in\mathbb Z[i]\)、\(\beta\ne0\)，在复数域中写 \(\alpha/\beta=s+ti\)。分别选最近的整数 \(m,n\)，使
\(\lvert s-m\rvert\leq1/2\)、\(\lvert t-n\rvert\leq1/2\)。置 \(q=m+ni\)、\(r=\alpha-\beta q\)，则
\[
N(r)=N(\beta)\bigl((s-m)^2+(t-n)^2\bigr)
\leq\tfrac12N(\beta)<N(\beta).
\]
因 \(q,r\in\mathbb Z[i]\)，这正是所需带余除法；若 \(r=0\)，也满足定义。
\end{proof}

例如 \(5=(2-i)(2+i)\)，所以 \(2+i\) 整除 \(5\)。又 \(N(2+i)=5\) 为整数素数，任何分解 \(2+i=\alpha\beta\) 都使两个非负整数范数的积为 \(5\)，其中一个必为 \(1\)；故 \(2+i\) 不可约。在 \(\mathbb Z\) 中不可约的 \(5\)，进入 \(\mathbb Z[i]\) 后已经分裂。这提示我们：讨论不可约性之前，工作环必须说清楚。

% !TeX root = ../../Book1.tex
% 纲要标题：### 11.3 主理想整环
% 纲要位置：计划纲要.md，第 428 行。

\section{主理想整环}
\label{b1:sec:1103}

% 纲要内容（仅作写作计划，尚非教材正文）：
%
% * 理想主性与素元性质
% * 主理想整环是唯一分解整环
% * 非 Euclid 的主理想整环作为选读例子
%

Euclid 算法提供了求最大公因子的方法，但它还有一个不依赖具体计算的后果：每个理想都能用一个元素生成。把这一结论单独提出，许多分解论证明便不再需要一个指定的除法算法。最后的例子将说明，这样做确实扩大了研究范围。

\subsection{理想主性与素元性质}
\label{b1:subsec:1103:01}

\begin{definition}\label{b1:def:pid}
设 \(R\) 是交换整环。若对每个理想 \(I\subseteq R\)，都存在 \(d\in R\) 使 \(I=(d)=\{\,dr\mid r\in R\,\}\)，则称 \(R\) 为\term[zhulixiangzhenghuan]{主理想整环}[principal ideal domain]，简称 PID。
\end{definition}

定义包含零理想 \((0)\) 和单位理想 \((1)\)。域只有这两个理想，因而是主理想整环；“整环”这一要求不能删去后仍沿用同一个概念。

\ProseParagraphBreak
要从 Euclid 除法得到理想主性，可以把除法用在理想内部：选定一个非零元素 \(d\) 后，任何不被 \(d\) 整除的理想元素都会产生一个非零余数，而这个余数仍在理想中。因此，若 \(d\) 的 Euclid 值已经最小，就不能再有这样的元素，\(d\) 必须整除整个理想。

\begin{theorem}\label{b1:thm:euclidean-pid}
设 \(R\) 为 Euclid 整环，则 \(R\) 是主理想整环。
\end{theorem}
\begin{proof}
设 \(R\) 带 Euclid 函数 \(\delta\)，\(I\) 是非零理想。从 \(I\setminus\{0\}\) 中取 \(d\)，使 \(\delta(d)\) 最小。对任意 \(a\in I\)，写 \(a=qd+r\)，其中 \(r=0\) 或 \(\delta(r)<\delta(d)\)。但 \(r=a-qd\in I\)，最小性排除了非零余数，所以 \(a\in(d)\)。反向包含由 \(d\in I\) 得出，故 \(I=(d)\)。零理想本来就由 \(0\) 生成。
\end{proof}

\begin{proposition}\label{b1:prop:pid-bezout}
设 \(R\) 为主理想整环，\(a,b,d\in R\)。则 \(a,b\) 有最大公因子，且
\[
(a,b)=(d)\quad\Longleftrightarrow\quad d\text{ 是 }a,b\text{ 的最大公因子}.
\]
此时存在 \(u,v\in R\) 使 \(d=ua+vb\)。
\end{proposition}
\begin{proof}
主性给出 \((a,b)=(d_0)\)。其中 \(a,b\in(d_0)\) 表明 \(d_0\) 是公因子，而 \(d_0\in(a,b)\) 给出线性组合表达式；\cref{b1:prop:bezout-gcd}于是说明它是最大公因子。任一其他最大公因子 \(d\) 都与 \(d_0\) 相伴，故\cref{b1:prop:associate-ideals}给出 \((d)=(d_0)=(a,b)\)，再由 \(d\in(a,b)\) 得到相应 Bézout 系数。这也包含 \(a=b=0\) 的情形。
\end{proof}

\begin{proposition}\label{b1:prop:pid-irreducible-prime}
在主理想整环 \(R\) 中，每个不可约元 \(p\) 都是素元，且 \((p)\) 是极大理想。因此每个非零素理想都是极大理想。
\end{proposition}
\begin{proof}
设 \((p)\subseteq I\subseteq R\)。写 \(I=(d)\)，则 \(p=dc\)。不可约性使 \(d\) 为单位或 \(c\) 为单位；前者给 \(I=R\)，后者给 \(I=(p)\)。又 \(p\) 非单位，所以 \((p)\) 是极大理想。由\cref{b1:prop:prime-max-quotient}，\(R/(p)\) 是域，因而是整环，\((p)\) 为素理想；再用\cref{b1:prop:prime-implies-irreducible}的素理想判别得 \(p\) 为素元。

若 \(P\ne(0)\) 是素理想，写 \(P=(p)\)。因 \(P\) 非零且真，\(p\) 非零非单位；同一判别说明 \(p\) 为素元，再由\cref{b1:prop:prime-implies-irreducible}得它不可约。刚才的论证便说明 \(P\) 极大。
\end{proof}

\subsecref{b1:subsec:1101:03}已把不可约性解释为“在真主理想中极大”。主理想整环中每个理想都为主理想，所以这一极大性涵盖了全部真理想。上述证明中的推理可写为
\[
p\text{ 不可约}
\quad\Longrightarrow\quad
(p)\text{ 为极大理想}
\quad\Longrightarrow\quad
(p)\text{ 为素理想}
\quad\Longrightarrow\quad
p\text{ 为素元}.
\]
“每个非零素理想都是极大理想”中的非零条件也有作用。在整环中，\((0)\) 总是素理想；但在主理想整环 \(\mathbb Z\) 中，有 \((0)\subsetneq(2)\subsetneq\mathbb Z\)，所以 \((0)\) 不是极大理想。

\ProseParagraphBreak
例如对域 \(F\)，\cref{b1:prop:field-poly-division}与\cref{b1:thm:euclidean-pid}给出 \(F[x]\) 是主理想整环。若 \(p(x)\) 不可约，商 \(F[x]/(p)\) 因而是域。这个结论会在构造方程根所在的域时发挥作用；它所依赖的是理想极大性，不是预先知道根的存在。

\subsection{主理想整环的唯一分解性}
\label{b1:subsec:1103:02}

理想主性还将保证元素的不可约分解存在且唯一。这里先定义唯一分解整环，以陈述主理想整环的这一结构性质；\secref{b1:sec:1104}将脱离主理想性，单独研究唯一分解本身及其判别。

\begin{definition}\label{b1:def:ufd}
设 \(R\) 是交换整环。若每个非零非单位都能写成有限个不可约元的乘积，并且同一非零非单位 \(a\in R\) 的任意两种表达式
\[
a=u p_1\cdots p_m=v q_1\cdots q_n
\]
中，\(u,v\) 为单位、\(p_i,q_j\) 为不可约元时，必有 \(m=n\)，且可重新排列 \(q_j\)，使每个 \(p_i\) 与 \(q_i\) 相伴，则称 \(R\) 为\term[weiyifenjiezhenghuan]{唯一分解整环}[unique factorization domain]，简称 UFD。
\end{definition}

单位允许看成空乘积再乘一个单位；零不属于分解定理的对象。唯一性既不固定因子的次序，也不固定各因子相伴类中的代表。因此 \(6=2\cdot3=(-2)(-3)\) 并不是两种不同的素因数分解。

\ProseParagraphBreak
一般理想的并未必是理想，甚至未必对加法封闭。例如 \(2\mathbb Z\cup3\mathbb Z\) 包含 \(2,3\)，却不包含 \(2+3=5\)。理想升链的并有所不同：属于不同项的两个元素总能一起放入其中较大的一项，这使理想的封闭条件得以保留。

\begin{lemma}\label{b1:lem:pid-ascending-chain}
设 \(R\) 为主理想整环。则 \(R\) 的每条理想升链
\(I_1\subseteq I_2\subseteq\cdots\) 都从某项开始恒定。
\end{lemma}
\begin{proof}
令 \(I=\bigcup_{j\geq1}I_j\)。若 \(x\in I_r,y\in I_s\)，升链条件使两者都属于 \(I_{\max(r,s)}\)，故 \(x-y\) 仍在 \(I\)；任一元素乘环中元素也留在原来所属的 \(I_j\)。因此 \(I\) 是理想，写成 \((d)\)。生成元 \(d\) 属于某个 \(I_m\)，所以 \(I=(d)\subseteq I_m\subseteq I\)，从而 \(I_m=I\)，所有 \(j\geq m\) 也满足 \(I_j=I\)。
\end{proof}

分解元素与扩大主理想正好方向相反。若非零元素满足 \(a_1=a_2d_1\)，且 \(d_1\) 非单位，则 \((a_1)\subsetneq(a_2)\)。因此，若一个元素总能继续作非平凡分解，就会产生严格上升的主理想链。下面用升链终止来排除无限继续分解的情形，再用不可约元的素性比较两种分解。

\begin{theorem}\label{b1:thm:pid-ufd}
设 \(R\) 为主理想整环，则 \(R\) 是唯一分解整环。
\end{theorem}
\begin{proof}
先证分解存在。若存在不能分解成有限个不可约元乘积的非零非单位 \(a_1\)，则 \(a_1\) 本身不可为不可约元，故 \(a_1=b_1c_1\)，其中两个因子都非单位。至少有一个因子仍不能分解，否则合并两边分解便得到 \(a_1\) 的分解。把这个因子记为 \(a_2\)，另一个记为 \(d_1\)，于是 \(a_1=a_2d_1\)。有 \((a_1)\subsetneq(a_2)\)：若相等，\cref{b1:prop:associate-ideals}使 \(a_1,a_2\) 相伴，再由非零消去律迫使 \(d_1\) 为单位，矛盾。重复选择会得到严格升链
\((a_1)\subsetneq(a_2)\subsetneq\cdots\)，与\cref{b1:lem:pid-ascending-chain}矛盾。因此分解存在。

再证唯一性。设 \(u p_1\cdots p_m=v q_1\cdots q_n\)。由\cref{b1:prop:pid-irreducible-prime}，不可约元 \(p_1\) 是素元；它不整除单位 \(v\)，对乘积逐次应用素性，必整除某个 \(q_j\)。不可约性使 \(q_j=p_1w\)，其中 \(w\) 为单位。把这个因子移至首位，消去非零的 \(p_1\)，得到两个较短分解的等式，单位因子相应并入 \(v\)。如此归纳，配对因子都相伴。若一边先用尽而另一边仍有不可约因子，就会使若干非单位之积成为单位；但乘积有逆元时，每个因子也有逆元，矛盾。因此两边同时用尽，\(m=n\)，唯一性得证。
\end{proof}

这一证明的两部分使用不同的条件：
\[
\begin{array}{c@{\qquad}c}
\text{分解存在性}&\text{分解唯一性}\\
\text{理想升链终止}&\text{不可约元都是素元}
\end{array}
\]
由此前的具体 Euclid 函数，\(\mathbb Z\)、\(F[x]\) 和 \(\mathbb Z[i]\) 均为唯一分解整环。不过，主性证明中的最小值选择只断言某个生成元存在；若没有可执行的除法步骤，它本身还不是计算该生成元的算法。

\OptionalSubsection{非 Euclid 的主理想整环}{b1:subsec:1103:03}

要说明“Euclid 整环 \(\Rightarrow\) 主理想整环”的逆命题不成立，需要对同一个环完成两件事：先证明每个理想都是主理想，再排除所有可能的 Euclid 函数。前一部分将利用格点估计构造允许倍乘的较弱余数；后一部分则从任意 Euclid 函数推出一个不可能的有限商环。设
\[
\omega=\frac{1+\sqrt{-19}}2,\qquad
R=\mathbb Z[\omega]\subseteq\mathbb C.
\]
因为 \(\omega^2=\omega-5\)，每个元素唯一写成 \(a+b\omega\)，\(a,b\in\mathbb Z\)。定义
\[
N(a+b\omega)=\lvert a+b\omega\rvert^2
=a^2+ab+5b^2
=\Bigl(a+\frac b2\Bigr)^2+\frac{19b^2}4.
\]
范数是非负整数且乘法相容，非零元素的范数是正整数。若 \(u\) 是单位，则 \(N(u)N(u^{-1})=N(1)=1\)，所以 \(N(u)=1\)。由最后一个表达式，范数等于 \(1\) 时只能有 \(b=0,a=\pm1\)，而 \(\pm1\) 确实可逆，因此 \(R\) 的单位恰为 \(\pm1\)。

\ProseParagraphBreak
比较下面两种余数要求，其中 \(\beta\ne0\)。左边是以范数为度量的普通带余除法，右边只针对 \(\alpha\notin(\beta)\)：
\[
\begin{array}{c@{\qquad}c}
\text{普通带余除法}&\text{允许倍乘的较弱条件}\\
\alpha=d\beta+r&c\alpha=d\beta+r\\
r=0\ \text{或}\ N(r)<N(\beta)&0<N(r)<N(\beta)
\end{array}
\]
右边的乘数 \(c\in R\) 可以依赖于 \(\alpha/\beta\)，所以所得 \(r\) 是 \(c\alpha\) 的余数，而不是 \(\alpha\) 的余数。这一条件不提供普通带余除法，却足以证明理想主性：若非零理想 \(I\) 中的 \(\beta\) 具有最小正范数，而 \(\alpha\in I\setminus(\beta)\)，则 \(r=c\alpha-d\beta\) 仍在 \(I\) 中，其正范数严格更小，产生矛盾。令 \(z=\alpha/\beta\notin R\)，右边的不等式就化为 \(0<\lvert cz-d\rvert^2<1\)。下面的格点引理构造这样的 \(c,d\)，并保证余数非零。

\begin{lemma}\label{b1:lem:nineteen-small-combination}
令
\[
\omega=\frac{1+\sqrt{-19}}2,\qquad R=\mathbb Z[\omega]\subseteq\mathbb C.
\]
对任意 \(z\in\mathbb C\setminus R\)，存在 \(c,d\in R\)，使
\(0<\lvert cz-d\rvert^2<1\)。
\end{lemma}
\begin{proof}
由 \(R=\mathbb Z+\mathbb Z\omega\) 及 \(\operatorname{Im}\omega=\sqrt{19}/2\)，减去适当整数倍的 \(\omega\) 可以把虚部移入一个长度为 \(\sqrt{19}/2\)、以零为中心的区间；再减去整数，把实部移入长度为 \(1\)、以零为中心的区间。因此可取 \(\rho\in R\)，使调整后的 \(z_0=z-\rho=x+yi\) 落入基本矩形
\[
\lvert x\rvert\leq\frac12,\qquad
\lvert y\rvert\leq\frac{\sqrt{19}}4.
\]
由于 \(z\notin R\)，调整后的 \(z_0\) 仍不在 \(R\) 中，否则 \(z=z_0+\rho\in R\)。特别地 \(z_0\ne0\)。若 \(\lvert z_0\rvert<1\)，取 \(c=1,d=\rho\)，就有 \(0<\lvert cz-d\rvert^2=\lvert z_0\rvert^2<1\)，已经得到结论。

\cref{b1:fig:nineteen-lattice}中矩形与单位圆之间的上下两片区域，正是基本矩形内仍须处理 \(\lvert z_0\rvert\geq1\) 的点。图中取 \(z_0=\omega/2\)：从 \(z_0\) 到 \(\omega\) 是倍乘，从 \(\omega\) 回到原点才是减去格点 \(\omega\)。

\begin{figure}[htbp]
\centering
\begin{tikzpicture}[x=1.55cm,y=1.55cm,>=stealth]
  \def\h{1.08972474} % sqrt(19)/4
  \def\s{0.86602540} % sqrt(3)/2
  \fill[color=AlgebraBlueLight] (-.5,\s)--(-.5,\h)--(.5,\h)--(.5,\s)
    arc[start angle=60,end angle=120,radius=1]--cycle;
  \fill[color=AlgebraBlueLight] (-.5,-\s)--(-.5,-\h)--(.5,-\h)--(.5,-\s)
    arc[start angle=-60,end angle=-120,radius=1]--cycle;
  \draw[->,color=AlgebraGray] (-1.2,0)--(1.25,0) node[right] {\(\operatorname{Re}z\)};
  \draw[->,color=AlgebraGray] (0,-1.26)--(0,2.51) node[above] {\(\operatorname{Im}z\)};
  \draw[dashed,color=AlgebraGray] (0,0) circle (1);
  \draw[thick,color=AlgebraBlue] (-.5,-\h) rectangle (.5,\h);
  \fill (0,0) circle (1.2pt) node[below left] {\(0\)};
  \fill[color=AlgebraBlue] (.25,\h) circle (1.6pt);
  \node[anchor=east,color=AlgebraBlue] at (.17,1.30) {\(z_0=\omega/2\)};
  \fill[color=AlgebraBlue] (.5,2*\h) circle (1.6pt);
  \node[anchor=west,color=AlgebraBlue] at (.55,2*\h) {\(2z_0=\omega\)};
  \draw[->,thick,color=AlgebraBlue] (.28,\h+.08)--(.48,2*\h-.09)
    node[pos=.53,left] {\(\times2\)};
  \draw[->,thick,color=AlgebraGray] (.55,2*\h-.03)
    .. controls (1.68,2.02) and (1.72,.38) .. (.06,.04);
  \node[color=AlgebraGray] at (1.38,1.16) {\(-\omega\)};
\end{tikzpicture}
\caption[基本矩形与单位圆]{格点引理中的基本矩形与单位圆（虚线）。阴影为矩形内满足 \(\lvert z_0\rvert\geq1\) 的上下两片区域；箭头分别表示倍乘与减去格点，图中展示的是零余数实例。}
\label{b1:fig:nineteen-lattice}
\end{figure}

以下处理基本矩形内仍有 \(\lvert z_0\rvert\geq1\) 的情形。由 \(x^2\leq1/4\)，有 \(\lvert y\rvert\geq\sqrt3/2\)。对 \(2z_0\) 再作格点调整：减去虚部与 \(y\) 同号的 \(\omega\) 或 \(-\omega\)，再减去整数使实部绝对值不超过 \(1/2\)，得到 \(w=2z_0-d_0\)，且
\[
\lvert\operatorname{Im}w\rvert
\leq\frac{\sqrt{19}}2-\sqrt3<\frac12,
\qquad \lvert w\rvert^2<\frac12<1.
\]
其中严格不等式可由 \(\sqrt{19}<22/5\)、\(\sqrt3>17/10\) 核对。若 \(w\ne0\)，取 \(c=2,d=d_0+2\rho\)，便有 \(cz-d=2z_0-d_0=w\)，即得结论。

最后处理 \(w=0\)。此时倍乘 \(c=2\) 恰好把 \(z_0\) 送到格点，所得余数是零，不能用它反驳最小正范数。我们改换乘数 \(c\)，目标是制造余数 \(1/2\)，从而同时保证余数非零与模平方小于 \(1\)。令 \(\alpha=2z_0\in R\)；由于 \(z_0\notin R\)，\(\alpha\notin2R\)。

按 \(a+b\omega\) 的两个整数系数的奇偶性，商环 \(R/2R\) 恰有四个剩余类，分别由 \(0,1,\omega,1+\omega\) 代表。三个非零类的逆元分别由 \(1,1+\omega,\omega\) 代表，因为
\(\omega(1+\omega)=2\omega-5\equiv1\pmod{2R}\)。这说明 \(R/2R\) 是四元素域；当前所需的正是非零类可逆。\(\alpha\) 的类非零，故可取 \(c\in R\) 及 \(d_1\in R\)，使 \(c\alpha=1+2d_1\)。于是 \(cz_0-d_1=1/2\)，其模平方为 \(1/4\)。取 \(d=d_1+c\rho\)，便有 \(cz-d=cz_0-d_1=1/2\)，就得到原命题。

例如 \(z_0=\omega/2\) 位于图中的上方阴影区域，因为 \(\lvert z_0\rvert^2=5/4>1\)，但 \(2z_0-\omega=0\)。这时可直接取 \(c=1+\omega\)、\(d_1=\omega-3\)；由 \(\omega^2=\omega-5\) 得
\[
(1+\omega)\frac{\omega}{2}-(\omega-3)=\frac12.
\]
所以零余数分支也确实产生非零且模平方小于 \(1\) 的线性组合。
\end{proof}

\begin{proposition}\label{b1:prop:noneuclidean-pid}
令 \(\omega=(1+\sqrt{-19})/2\)，并令 \(R=\mathbb Z[\omega]\subseteq\mathbb C\)。则 \(R\) 是主理想整环，但不存在任何使 \(R\) 成为 Euclid 整环的 Euclid 函数。
\end{proposition}
\begin{proof}
设 \(I\ne(0)\) 为理想，取 \(\beta\in I\setminus\{0\}\)，使正整数 \(N(\beta)\) 最小。若某个 \(\alpha\in I\) 不属于 \((\beta)\)，则 \(z=\alpha/\beta\notin R\)。由\cref{b1:lem:nineteen-small-combination}取 \(c,d\in R\)，使 \(0<\lvert cz-d\rvert^2<1\)。于是非零元素 \(c\alpha-d\beta\in I\) 满足
\[
0<N(c\alpha-d\beta)
=N(\beta)\lvert cz-d\rvert^2<N(\beta),
\]
与最小性矛盾。所以 \(I=(\beta)\)，零理想也为主理想，得证主性。

现在排除任意 Euclid 函数。由于 \(N(2)=4\ne1\)，\(2\) 是非零非单位，因此所有非零非单位组成的集合非空。假设存在 Euclid 函数 \(\delta\)，可在这个集合中取 \(b\)，使 \(\delta(b)\) 最小。对任意 \(a\in R\) 作除法 \(a=qb+r\)；若 \(r\ne0\)，则 \(\delta(r)<\delta(b)\)。最小性使 \(r\) 不可能是非单位，因而只能是单位，即 \(1\) 或 \(-1\)。

因此每个模 \((b)\) 的剩余类都由 \(0,1,-1\) 之一代表，\(R/(b)\) 至多有三个元素。另一方面，\(b\) 非单位意味着 \(1\notin(b)\)，所以商环中的 \(1\ne0\)，至少有两个元素。合起来得到
\[
2\leq\lvert R/(b)\rvert\leq3.
\]
故它的大小只能为 \(2\) 或 \(3\)。

令 \(S=R/(b)\)，并令 \(p=\lvert S\rvert\in\{2,3\}\)。由\cref{b1:thm:lagrange}，\(1_S\ne0_S\) 在加法群中的阶是 \(p\)，所以 \(1_S\) 生成整个加法群。映射
\[
\mathbb Z/p\mathbb Z\longrightarrow S,\qquad
[n]\longmapsto n1_S
\]
于是良定义且为双射；它保持加法、乘法和单位，其中乘法相容性来自 \((m1_S)(n1_S)=(mn)1_S\)。因此它是含幺环同构，即 \(S\cong\mathbb F_p\)。

然而 \(\omega\) 的像必须满足 \(t^2-t+5=0\)。模 \(2\) 时，\(t=0,1\) 的值都是 \(1\)；模 \(3\) 时，\(t=0,1,2\) 的值依次为 \(2,2,1\)。两种商中都不可能有这样的元素，矛盾。故不存在 \(\delta\)。
\end{proof}

范数本身不能作 Euclid 函数，还不足以推出环不是 Euclid 整环；必须排除所有可能的 Euclid 函数。上述证明的后半部分正是通过有限商环的必要条件完成这一点。主性证明使用的乘数 \(c\) 则足以产生理想中的更小元素，却不提供对原元素的普通带余除法。

\ProseParagraphBreak
关于这个例子，可参见 Keith Conrad 的讲义\cite[Theorems 5.18, 5.22]{ConradRemarksEuclideanDomains}：定理 5.18 提炼了非域 Euclid 整环中某个非单位模下每个剩余类可由零或单位代表的必要条件；定理 5.22 讨论本小节的二次环，其中主理想性只给出证明提纲，可与此处的格点引理作方法对照。

% !TeX root = ../../Book1.tex
% 纲要标题：### 11.4 唯一分解整环
% 纲要位置：计划纲要.md，第 434 行。

\section{唯一分解整环}
\label{b1:sec:1104}

% 纲要内容（仅作写作计划，尚非教材正文）：
%
% * 分解的存在和唯一性
% * 唯一分解不蕴含主性
% * 非唯一分解例子与理想分解的动机
%
% ---
%

主理想性给出了唯一分解的充分条件，却不是必要条件。本节把分解论自身的要求提炼出来，并考察两个相反方向的例子：\(\mathbb Z[x]\) 的元素仍能唯一分解，但有非主理想；\(\mathbb Z[\sqrt{-5}]\) 连元素的分解唯一性也会失败。前者的唯一分解性将在下一章通过多项式方法证明，后者在这里直接计算。

\subsection{分解的存在性与唯一性}
\label{b1:subsec:1104:01}

\begin{proposition}\label{b1:prop:ufd-prime}
设 \(R\) 为整环。如果 \(R\) 是唯一分解整环，则 \(R\) 中的不可约元都是素元。更一般地，\(R\) 是唯一分解整环，当且仅当它的每个非零非单位都能分解成有限个不可约元的乘积，并且每个不可约元都是素元。
\end{proposition}
\begin{proof}
先设 \(R\) 为唯一分解整环，\(p\) 不可约且 \(p\mid ab\)。若 \(a=0\) 或 \(b=0\)，结论成立；否则写 \(ab=pc\)，其中 \(c\ne0\)。把 \(a,b,c\) 分别分解为单位与不可约元的乘积，允许单位的分解含零个不可约因子。等式右侧有因子 \(p\)，唯一性要求它与左侧某个不可约因子相伴。那个因子来自 \(a\) 或 \(b\)，所以 \(p\mid a\) 或 \(p\mid b\)，即 \(p\) 为素元。

反过来，假定分解存在且不可约元都为素元。\cref{b1:thm:pid-ufd}证明唯一性的后一段，在已有两种分解后只用了不可约元的素性、整环的非零消去律及单位的基本性质，并未再用理想主性。因此同一因子配对论证适用于此，得到因子个数相同且对应因子相伴，正是\cref{b1:def:ufd}的唯一性。
\end{proof}

这个判别不能省略分解存在的条件。证明唯一性时以“取一个不可约分解”起步，只在确知每个非零非单位都能如此分解时才有意义。\cref{b1:thm:pid-ufd}为主理想整环补足存在性时，用到的是理想升链最终稳定，两个论证的职责应当分清。

设 \(a,b\) 为整环 \(R\) 的非零元素。如果 \(m\in R\) 同时被 \(a,b\) 整除，并且 \(m\) 整除 \(a,b\) 的每个公倍数，就称 \(m\) 为 \(a,b\) 的\term[zuixiaogongbeishu]{最小公倍数}[least common multiple]，记为 \(\operatorname{lcm}(a,b)\)。最小公倍数只在相伴意义下确定。
\symidx{\(\operatorname{lcm}(a,b)\)：最小公倍数}

\begin{proposition}[唯一分解整环中的最大公因子与最小公倍数]\label{b1:prop:ufd-gcd-lcm}
设 \(R\) 为唯一分解整环，\(a,b\in R\) 非零。则 \(a,b\) 有最大公因子与最小公倍数。把它们分解中涉及的不同不可约相伴类各选一个代表 \(p_1,\ldots,p_t\)，其中 \(t\geq0\) 为整数，并写成
\[
a=u\prod_{j=1}^{t}p_j^{\alpha_j},\qquad
b=v\prod_{j=1}^{t}p_j^{\beta_j},\qquad \alpha_j,\beta_j\in\mathbb Z_{\geq0},
\]
其中 \(u,v\in R^\times\)。此时最大公因子 \(d\) 和最小公倍数 \(m\) 可分别取为
\[
d=\prod_{j=1}^{t}p_j^{\min(\alpha_j,\beta_j)},\qquad
m=\prod_{j=1}^{t}p_j^{\max(\alpha_j,\beta_j)}.
\]
\end{proposition}
\begin{proof}
唯一分解说明，对于非零元素 \(r,s\)，条件 \(r\mid s\) 等价于 \(r\) 中每个不可约因子的指数不超过它在 \(s\) 中的指数：必要性由 \(s=rc\) 及分解唯一性得到，充分性由指数相减构造商得到。因此公因子中各指数的上界是两个指数的较小者，公倍数中各指数的下界是两个指数的较大者，分别给出所列 \(d,m\)。公倍数为 \(0\) 时，\(m\mid0\) 也成立。
\end{proof}

对上一命题选出的 \(d,m\)，乘积 \(dm\) 与 \(ab\) 相伴；适当调整单位后可令二者相等。零元素的情况另定为 \(\gcd(a,0)\) 与 \(a\) 相伴、\(\operatorname{lcm}(a,0)=0\)。

\ProseParagraphBreak
例如在 \(\mathbb Z[i]\) 中，\(2=-i(1+i)^2\)，所以 \(2\) 已不再是不可约元。若 \(a=(1+i)^3(2+i)\)、\(b=(1+i)(2+i)^2\)，则 \(1+i\) 与 \(2+i\) 的范数分别为整数素数 \(2,5\)，两者不可约且不相伴。因而可取 \(\gcd(a,b)=(1+i)(2+i)\)，最小公倍数则为 \((1+i)^3(2+i)^2\)。

\subsection{唯一分解不蕴含主性}
\label{b1:subsec:1104:02}

\cref{b1:prop:integer-poly-nonprincipal}已经说明：\(\mathbb Z[x]\) 的理想 \((2,x)\) 不为主理想。这里的障碍是 \(2,x\) 虽然没有共同的不可约因子，却不能用多项式系数的线性组合得到 \(1\)。

\ProseParagraphBreak
另一方面，\(\mathbb Z[x]\) 确实是唯一分解整环。这个断言是下一章\cref{b1:thm:polynomial-ufd}的一个应用：那里将证明只要 \(R\) 是唯一分解整环，\(R[x]\) 也是；取本章已证的 \(R=\mathbb Z\) 即得。此处把它作为前瞻例子，不用它证明本章其他结论。与已经完成证明的 \(\mathbb Z[\omega]\) 一起，这两个例子说明三个条件之间是严格蕴含：
\[
\text{Euclid 整环}\ \Longrightarrow\
\text{主理想整环}\ \Longrightarrow\
\text{唯一分解整环}.
\]
左边逆向的反例由\cref{b1:prop:noneuclidean-pid}给出，右边逆向的反例则在下一章完成最后的唯一分解性验证。

\ProseParagraphBreak
在 \(\mathbb Z[x]\) 中，逐项模 \(2\) 的商映射给出 \(\mathbb Z[x]/(2)\cong\mathbb F_2[x]\)。右边是整环而不是域，所以 \(2\) 是素元，\((2)\) 却不是极大理想；事实上 \((2)\subsetneq(2,x)\subsetneq\mathbb Z[x]\)。下一章\cref{b1:thm:polynomial-ufd}还将证明 \(\mathbb Z[x]\) 是唯一分解整环；因而这一例子也表明，不可约元都是素元并不能保证每个非零素理想都极大。

\ProseParagraphBreak
在唯一分解整环中，对非零 \(a,b\)，公倍数恰为最小公倍数的倍数，因此 \((a)\cap(b)=(m)\) 仍是主理想；而和 \((a)+(b)\) 未必为主理想。最大公因子刻画公因子之间的整除关系，却未必能表示为原来两个元素的线性组合。将“有最大公因子”与“能写 Bézout 关系”分开，能避免把主理想整环的结论错误推广到全部唯一分解整环。

\subsection{非唯一分解与理想分解}
\label{b1:subsec:1104:03}

令 \(s=\sqrt{-5}\)，考虑整环 \(R=\mathbb Z[s]\subseteq\mathbb C\)。它的元素唯一写成 \(a+bs\)，范数
\[
N(a+bs)=a^2+5b^2
\]
是乘法相容的非负整数，单位恰为 \(\pm1\)。不存在范数为 \(2\) 或 \(3\) 的元素：若 \(b\ne0\)，范数至少为 \(5\)；若 \(b=0\)，整数平方又不可能等于 \(2,3\)。

\begin{proposition}\label{b1:prop:sqrt-minus-five-factorization}
令 \(s=\sqrt{-5}\)，并令 \(R=\mathbb Z[s]\subseteq\mathbb C\)。则等式
\[
6=2\cdot3=(1+s)(1-s)
\]
给出两种不等价的不可约分解。因此 \(R\) 不是唯一分解整环，且不可约元 \(2\) 不是素元。
\end{proposition}
\begin{proof}
元素 \(2,3,1+s,1-s\) 的范数依次为 \(4,9,6,6\)。若 \(2\) 分解成两个非单位，两因子的范数只能为 \(2,2\)，但范数 \(2\) 不存在，故 \(2\) 不可约。对 \(3\)，唯一的非平凡范数分解为 \(3\cdot3\)，同样不可能。对 \(1\pm s\)，非平凡范数分解只能为 \(2\cdot3\)，也不可能。因此四个因子都不可约。

单位的范数为 \(1\)，相伴元素范数相同，所以范数为 \(4\) 的 \(2\) 不可能与范数为 \(6\) 的任一 \(1\pm s\) 相伴。两种分解因而不等价。此外，\(2\mid(1+s)(1-s)=6\)，却不整除任一 \(1\pm s\)，因为它们的整数常数项为奇数。故 \(2\) 不满足素性。
\end{proof}

元素分解失效以后，可以尝试让理想承担因子的作用。本例的
\[
I=(2,1+s)
\]
不是主理想，却满足 \(I^2=(2)\)。先看非主性。映射
\(R\rightarrow\mathbb F_2\)、\(a+bs\mapsto\overline{a+b}\)
保持乘法，因为 \(s^2=-5\) 在模 \(2\) 下对应 \(1^2=1\)；它把 \(I\) 送到零而把 \(1\) 送到 \(1\)，故 \(I\ne R\)。若 \(I=(d)\)，则 \(d\mid2\)。由 \(2\) 不可约，\(d\) 或为单位，或与 \(2\) 相伴；前者与 \(I\ne R\) 矛盾，后者又使 \(1+s\in(2)\)，也矛盾。

\ProseParagraphBreak
再算理想平方。\(I^2\) 由
\[
4,\quad 2+2s,\quad (1+s)^2=-4+2s
\]
生成，三者都属于 \((2)\)。反向包含由
\[
2=(2+2s)-(-4+2s)-4\in I^2
\]
得到，所以 \(I^2=(2)\)。元素 \(2\) 不可约，但它生成的理想仍有这样的真分解。\cref{b1:ex:11-ideal-splitting-three}将继续核对 \(J=(3,1+s)\)、\(K=(3,1-s)\) 满足 \(JK=(3)\)、\(IJ=(1+s)\)、\(IK=(1-s)\)。于是原来的两种元素分解在理想层面汇合为
\[
(6)=(2)\cdot(3)=I^2JK=(IJ)(IK)
=\bigl((1+s)R\bigr)\cdot\bigl((1-s)R\bigr).
\]
最后一项明确写出两个主理想的乘积，以区别于前面的元素等式 \(6=(1+s)(1-s)\)。这些等式只核对当前例子；要得到一般的理想分解定理，还须对环施加适当条件。这是后续数论与交换代数需要研究的问题。


\begin{chaptersummary}
素元必不可约，反向蕴含需要额外条件。Euclid 除法使连续非零余数的 Euclid 函数值严格下降，故迭代过程必定终止，最后一个非零余数给出最大公因子。若已给出可执行的求商余程序，还可同步更新系数，计算 Bézout 关系。对非零理想选取 Euclid 函数值最小的非零元素，可证明它生成该理想。在主理想整环中，理想升链最终稳定保证不可约分解存在，不可约元的素性保证分解唯一。

\ProseParagraphBreak
唯一分解整环中的最大公因子可以通过因子指数求得，却不一定能写成 Bézout 关系。\(\mathbb Z[x]\) 的理想 \((2,x)\) 展示了这种差别。\(\mathbb Z[\sqrt{-5}]\) 中不可约但不素的 \(2\)，则说明元素的唯一分解可能彻底失效；相应的主理想仍可分解为非主理想的乘积。下一章将通过容度与本原多项式，解决多项式环何时保留唯一分解的问题。
\end{chaptersummary}

\BookExercises

\begin{exercise}\label{b1:ex:11-gaussian-gcd}
在 \(\mathbb Z[i]\) 中求 \(a=7+5i\)、\(b=3+i\) 的一个最大公因子及 Bézout 系数，并列出它们的公因子相伴类。求 \(a,b\) 的一个最小公倍数，并计算理想 \((a)+(b)\) 与 \((a)\cap(b)\)。
\end{exercise}
\begin{solution}
选择商 \(3+i\)，有
\[
a=(3+i)b+(-1-i),\qquad
b=(-2+i)(-1-i).
\]
第一步余数范数为 \(2<N(b)=10\)，第二步余数为零，故\cref{b1:prop:euclidean-algorithm}给出最大公因子 \(-1-i\)。乘以单位 \(-1\)，可取 \(d=1+i\)，相应关系为
\[
1+i=-a+(3+i)b.
\]
范数 \(N(1+i)=2\) 为整数素数，若 \(1+i\) 分解，两因子范数的积为 \(2\)，必有一个范数为 \(1\)，即为单位。所以 \(1+i\) 不可约。公因子恰为 \(d\) 的因子，只有单位类和 \(1+i\) 的相伴类，分别以 \(1,1+i\) 为代表。

由\cref{b1:prop:ufd-gcd-lcm}，可取最小公倍数
\[
m=\frac{ab}{d}=19+3i.
\]
前述 Bézout 关系给出 \((a)+(b)=(d)=(1+i)\)。理想 \((a)\cap(b)\) 中的元素恰为 \(a,b\) 的公倍数；最小公倍数的整除刻画因而给出 \((a)\cap(b)=(m)=(19+3i)\)。
\end{solution}

\begin{exercise}\label{b1:ex:11-diophantine}
求出方程 \(252u+198v=36\) 的全部整数解，并说明右端改为 \(30\) 时是否有解。
\end{exercise}
\begin{solution}
将\eqref{b1:eq:integer-bezout-example}的 Bézout 关系乘以 \(2\)，得到 \(36=8\cdot252-10\cdot198\)，故 \((u,v)=(8,-10)\) 是一个解。任意其他解满足
\[
14(u-8)+11(v+10)=0.
\]
模 \(11\) 后，\(3(u-8)\equiv0\pmod{11}\)。因 \(4\cdot3\equiv1\pmod{11}\)，得到 \(u-8=11t\)，代回得 \(v+10=-14t\)。因此全部解为
\[
u=8+11t,\qquad v=-10-14t,\qquad t\in\mathbb Z.
\]
反向代入可检验每个 \(t\) 均给出解。右端为 \(30\) 时无解，因为左端必被 \(18\) 整除，而 \(18\nmid30\)。
\end{solution}

\begin{exercise}\label{b1:ex:11-polynomial-euclid}
在 \(\mathbb F_2[x]\) 中，求
\[
f=x^4+x^3+x^2+1,\qquad g=x^3+1
\]
的首一最大公因子，并求 \(u,v\in\mathbb F_2[x]\)，使该最大公因子等于 \(uf+vg\)。须列出各次带余除法。
\end{exercise}
\begin{solution}
在特征 \(2\) 中，加法与减法相同。按\cref{b1:prop:field-poly-division}消去首项，得到
\[
\begin{aligned}
f&=(x+1)g+(x^2+x),\\
g&=(x+1)(x^2+x)+(x+1),\\
x^2+x&=x(x+1).
\end{aligned}
\]
非零余数的次数依次为 \(2,1\)，最后余数为零。由\cref{b1:prop:euclidean-algorithm}，首一最大公因子为 \(d=x+1\)。回代为
\[
\begin{aligned}
d&=g+(x+1)(x^2+x)\\
&=g+(x+1)\bigl(f+(x+1)g\bigr)\\
&=(x+1)f+x^2g.
\end{aligned}
\]
最后一步使用 \(1+(x+1)^2=x^2\)，这是当前基域特征为 \(2\) 的结果。因此可取 \(u=x+1\)、\(v=x^2\)；把这些系数误作有理数系数时，不能保留相同的运算步骤。
\end{solution}

\begin{exercise}\label{b1:ex:11-gaussian-rational-prime}
设整数素数 \(p\equiv3\pmod4\)。证明 \(p\) 在 \(\mathbb Z[i]\) 中仍为素元，并解释为什么同一论证不能用于 \(p=2\)。
\end{exercise}
\begin{solution}
若 \(p=\alpha\beta\) 且两个因子都为非单位，则两个正整数 \(N(\alpha),N(\beta)>1\) 的积为 \(p^2\)，只能同时等于 \(p\)。然而整数平方模 \(4\) 只能为 \(0,1\)，两平方之和不可能同余 \(3\)，所以不存在范数为 \(p\) 的 Gauss 整数。这证明 \(p\) 不可约。

由\cref{b1:prop:gaussian-euclidean}和\cref{b1:thm:euclidean-pid}，\(\mathbb Z[i]\) 为主理想整环；\cref{b1:prop:pid-irreducible-prime}因此把这里的不可约性提升为素性。当 \(p=2\) 时，范数 \(2\) 确实存在，例如 \(N(1+i)=2\)，并且 \(2=-i(1+i)^2\)，所以 \(2\) 在此环中并不不可约。
\end{solution}

\begin{exercise}\label{b1:ex:11-quotient-inverse}
设 \(R\) 为主理想整环，\(a\) 为非零非单位。证明 \(b+(a)\) 在 \(R/(a)\) 中可逆，当且仅当 \(a,b\) 互素。应用于 \(R=\mathbb Z[i]\)、\(a=3\)，证明该商为九元素域，并求 \(1+i+(3)\) 的逆元。
\end{exercise}
\begin{solution}
若 \(bc-1\in(a)\)，则存在 \(t\in R\) 使 \(1=bc+at\)。任何公因子都整除 \(1\)，所以 \(a,b\) 互素。反过来，\cref{b1:prop:pid-bezout}使互素条件给出 \(1=ua+vb\)，故 \(v+(a)\) 是 \(b+(a)\) 的逆元。

前题已证明 \(3\) 是 \(\mathbb Z[i]\) 中的素元，所以 \((3)\) 是非零素理想。又由\cref{b1:prop:gaussian-euclidean}和\cref{b1:thm:euclidean-pid}，\(\mathbb Z[i]\) 是主理想整环；\cref{b1:prop:pid-irreducible-prime}于是保证 \((3)\) 为极大理想，商环为域。这里 \(N(3)=9\)，可见在 Gauss 整数中，范数为整数素数只是证明不可约的充分条件，不是必要条件。每个剩余类唯一以 \(u+vi\)、\(u,v\in\{0,1,2\}\) 表示，因为相差 \(3\) 的倍数等价于实部与虚部系数分别相差 \(3\) 的倍数，故商有九个元素。最后
\((1+i)(-1+i)=-2\equiv1\pmod3\)，所求逆元为 \(-1+i+(3)\)。
\end{solution}

\begin{exercise}\label{b1:ex:11-squarefree-quotient}
设 \(R\) 为主理想整环，非零非单位元 \(a\) 的不可约分解为 \(a=u p_1^{e_1}\cdots p_t^{e_t}\)，其中 \(u\in R^\times\)，各 \(p_j\) 为两两不相伴的不可约元，\(e_j\geq1\)。证明以下条件等价：商环 \(R/(a)\) 中，每当 \(z^n=0\)、\(n\geq1\)，都有 \(z=0\)；所有指数 \(e_j\) 都等于 \(1\)。
\end{exercise}
\begin{solution}
若全部 \(e_j=1\)，且 \(b^n\in(a)\)，则每个 \(p_j\mid b^n\)。\cref{b1:prop:pid-irreducible-prime}给出 \(p_j\) 的素性，反复用于 \(b^n=b\cdot b^{n-1}\)，得到 \(p_j\mid b\)。唯一分解使各个不同相伴类的因子同时出现于 \(b\)，于是 \(a\mid b\)，即 \(z=b+(a)=0\)。

若某个 \(e_j\geq2\)，取
\[
b=\prod_{j=1}^{t}p_j^{\lceil e_j/2\rceil}.
\]
每个指数在平方后至少达到 \(e_j\)，故 \(a\mid b^2\)。至少一个指数 \(\lceil e_j/2\rceil\) 严格小于 \(e_j\)，所以由分解唯一性，\(a\nmid b\)。因此非零剩余类 \(b+(a)\) 的平方为零，给出反例。
\end{solution}

\begin{exercise}\label{b1:ex:11-no-gcd}
在 \(R=\mathbb Z[\sqrt{-5}]\) 中证明 \(6\) 与 \(2(1+\sqrt{-5})\) 没有最大公因子。
\end{exercise}
\begin{solution}
记 \(s=\sqrt{-5}\)。\(2\) 是这两个元素的公因子；\(1+s\) 也是公因子，因为 \(6=(1+s)(1-s)\)。若存在最大公因子 \(d\)，它首先必须被每个公因子整除，故 \(2\mid d\)、\(1+s\mid d\)。范数的乘法性于是给出
\[
4\mid N(d),\qquad6\mid N(d),\qquad\text{因而 }12\mid N(d).
\]
另一方面，\(d\) 本身必须是公因子，即 \(d\mid6\)、\(d\mid2(1+s)\)，所以
\[
N(d)\mid36,\qquad N(d)\mid24,\qquad\text{因而 }N(d)\mid12.
\]
合并两个方向，正整数 \(N(d)\) 只能等于 \(12\)。但方程 \(a^2+5b^2=12\) 没有整数解：\(\lvert b\rvert\geq2\) 时左端至少为 \(20\)；\(b=0\) 要求 \(a^2=12\)；\(b=\pm1\) 要求 \(a^2=7\)。矛盾，因此不存在最大公因子。
\end{solution}

\begin{exercise}\label{b1:ex:11-coprime-powers}
设 \(R\) 为唯一分解整环，非零元素 \(a,b\) 互素，且 \(ab=c^n\)，其中 \(n\geq2\)。证明存在 \(d,e\in R\) 与单位 \(u,v\)，使 \(a=ud^n\)、\(b=ve^n\)。说明为什么一般不能删去单位因子。
\end{exercise}
\begin{solution}
互素意味着 \(a,b\) 没有公共不可约因子；否则该因子作为公因子应整除 \(1\)，不可能。由 \(ab=c^n\) 及唯一分解，乘积中每个不可约因子的指数都为 \(n\) 的倍数。一个因子只能出现在 \(a,b\) 的一边，所以它在该边的指数本身就是 \(n\) 的倍数。把这些指数分别除以 \(n\)，得到 \(d,e\)；原分解的单位分别保留为 \(u,v\)，即得结论。若某一元素为单位，该边取相应的空乘积 \(d=1\) 或 \(e=1\)。

单位不能普遍删去。例如在 \(\mathbb Z\) 中取 \(a=-1\)、\(b=-4\)、\(c=2\)、\(n=2\)。两者互素且 \(ab=4=c^2\)，但负整数不是整数的平方；乘上单位 \(-1\) 后才得到上述形式。
\end{solution}

\begin{exercise}\label{b1:ex:11-cusp-subring}
设 \(F\) 为域，令
\(A=F[t^2,t^3]\subseteq F[t]\)，即由 \(F\) 中的常数及 \(t^2,t^3\) 通过有限次加法、减法和乘法生成的子环。证明 \(A\) 不是唯一分解整环，尽管 \(F[t]\) 是。
\end{exercise}
\begin{solution}
由于每个整数 \(m\geq2\) 都可写成 \(2r+3s\)、\(r,s\geq0\)，环 \(A\) 恰由一次项系数为零的多项式组成。它的单位是非零常数：若两个多项式乘积为 \(1\)，次数相加迫使两者都是常数。于是每个非零非单位的次数至少为 \(2\)。因此 \(t^2,t^3\) 都不可约，否则两个非单位因子的乘积次数至少为 \(4\)。它们不相伴，因为乘以非零常数不改变次数。但
\[
t^6=(t^2)^3=(t^3)^2
\]
给出长度不同的两种不可约分解，所以 \(A\) 不是唯一分解整环。另一方面，\cref{b1:prop:field-poly-division}给出 \(F[t]\) 为 Euclid 整环，再由\cref{b1:thm:euclidean-pid}与\cref{b1:thm:pid-ufd}得它是唯一分解整环。同一个元素 \(t^3\) 在 \(A\) 中不可约，在 \(F[t]\) 中却可分解为 \(t\cdot t^2\)，因为 \(t\notin A\)。可用的因子随所在环改变，唯一分解性不能任意传给子环。
\end{solution}

\begin{exercise}\label{b1:ex:11-pid-quotient-ideals}
设 \(R\) 为主理想整环，\(a=u p_1^{e_1}\cdots p_t^{e_t}\) 为非零非单位的不可约分解，其中各 \(p_j\) 两两不相伴、\(e_j\geq1\)。描述 \(R/(a)\) 的全部理想及包含关系，并求理想总数与极大理想总数。
\end{exercise}
\begin{solution}
由\cref{b1:prop:ring-ideal-correspondence}，商环的理想与 \(R\) 中包含 \((a)\) 的理想一一对应。后者均为 \((d)\)，且由\cref{b1:prop:associate-ideals}，\((a)\subseteq(d)\) 等价于 \(d\mid a\)。\cref{b1:thm:pid-ufd}的唯一分解使每个这样的 \(d\) 在相伴意义下唯一写为
\[
d=\prod_{j=1}^t p_j^{f_j},\qquad 0\leq f_j\leq e_j.
\]
不同的指数列给出不同的主理想，因为主理想相等恰好等价于生成元相伴。因此商环全部理想为这些 \((d)/(a)\)，总数为 \(\prod_{j=1}^t(e_j+1)\)。

若 \(d'\) 的指数列为 \(f'_j\)，则
\((d)/(a)\subseteq(d')/(a)\) 当且仅当 \(d'\mid d\)，也就是对每个 \(j\) 都有 \(f'_j\leq f_j\)。全环对应全零指数列，极大真理想恰对应只在一个位置取 \(1\)、其余取 \(0\) 的指数列。它们即 \((p_j)/(a)\)，共有 \(t\) 个；若指数和大于 \(1\)，减去一个正指数就会得到严格位于该理想与全环之间的理想，因而不极大。

包含关系与指数大小的方向相反。例如 \(a=p^3\)、\(p\) 不可约时，商环中的全部理想组成严格升链
\[
(p^3)/(p^3)\subsetneq(p^2)/(p^3)
\subsetneq(p)/(p^3)\subsetneq R/(p^3).
\]
从左到右，生成元的指数依次为 \(3,2,1,0\)，而理想依次变大。
\end{solution}

\begin{optionalexercise}\label{b1:ex:11-nineteen-quotient}
沿用\subsecref{b1:subsec:1103:03}中的 \(R=\mathbb Z[\omega]\)、\(\omega=(1+\sqrt{-19})/2\)，取 \(N(z)=\lvert z\rvert^2\)。求商环 \(R/(2)\)，证明它为四元素域，并求 \(\omega+(2)\) 的逆元。另证不存在 \(q,r\in R\) 使 \(\omega=2q+r\)、\(N(r)<N(2)\)。说明这项除法失败本身为何还不能排除别的 Euclid 函数。
\end{optionalexercise}
\begin{solution}
因为每个元素唯一写为 \(a+b\omega\)，模 \((2)\) 后有四个不同的类
\(0,1,\overline\omega,1+\overline\omega\)。关系 \(\omega^2=\omega-5\) 在商中成为
\(\overline\omega^2=\overline\omega+1\)，所以
\[
\overline\omega(1+\overline\omega)=1.
\]
三个非零类 \(1,\overline\omega,1+\overline\omega\) 的逆元分别为 \(1,1+\overline\omega,\overline\omega\)。故商是四元素域；所求逆元为 \(1+\omega+(2)\)。这也显式给出同构
\(R/(2)\cong\mathbb F_2[t]/(t^2+t+1)\)：把 \(t\) 的类送到 \(\overline\omega\)，两边均以 \(1,t\) 或 \(1,\overline\omega\) 唯一表示元素，关系相同。

若 \(q=a+b\omega\)，则 \(r=\omega-2q=-2a+(1-2b)\omega\)。记 \(v=1-2b\)，它是非零奇数。由 \(\omega=1/2+(\sqrt{19}/2)i\) 直接计算实部与虚部，得
\[
N(r)=\Bigl(-2a+\frac v2\Bigr)^2+\frac{19v^2}{4}
\geq\frac{19}{4}>4=N(2).
\]
所以没有所要求的范数余数。这个计算只否定函数 \(N\) 的 Euclid 条件；环的 Euclid 性要求存在某个合适函数，不能据此否定全部可能性。排除任意 Euclid 函数的结论另由\cref{b1:prop:noneuclidean-pid}给出，其证明使用了最小 Euclid 值的非单位元以及商环大小的限制。
\end{solution}

\begin{exercise}\label{b1:ex:11-nonprincipal-square}
在 \(\mathbb Z[x]\) 中令 \(J=(2,x)\)。对每个整数 \(n\geq1\)，求 \(J^n\) 的一组生成元，证明 \(J^n\) 不是主理想，并求 \(\sqrt{J^n}\)。说明为何发现非主理想不能据此断言该环不是唯一分解整环。
\end{exercise}
\begin{solution}
由理想积的有限和定义，归纳可得
\[
J^n=(2^n,2^{n-1}x,\ldots,2x^{n-1},x^n).
\]
确实，右边在 \(n=1\) 时就是 \(J\)；乘以 \((2,x)\) 后，每个生成元 \(2^{n-j}x^j\) 分别给出 \(2^{n+1-j}x^j\) 与 \(2^{n-j}x^{j+1}\)，恰好覆盖下一次幂的全部生成元。这些生成元及其任意多项式系数组合的常数项都被 \(2^n\) 整除，故 \(1\notin J^n\)。若 \(J^n=(h)\)，则 \(h\mid2^n\)。由\cref{b1:prop:polynomial-degree}，\(h\) 为非零整数常数；再由 \(h\mid x^n\)，比较 \(x^n\) 的系数，得 \(h\mid1\)，所以 \(h=\pm1\)。这将给出 \(J^n=\mathbb Z[x]\)，矛盾。

映射 \(\mathbb Z[x]\rightarrow\mathbb F_2\)、\(f\mapsto\overline{f(0)}\) 是满环同态。它的核恰为常数项为偶数的多项式集合，也就是 \(J=(2,x)\)。由\cref{b1:thm:ring-first-iso}和\cref{b1:prop:prime-max-quotient}，\(J\) 为极大理想，因而为素理想。若 \(f^m\in J\)，反复应用素性就得 \(f\in J\)，故 \(\sqrt J=J\)。再由\cref{b1:prop:ideal-containments}中 \(\sqrt{J^n}=\sqrt J\) 的等式，得到 \(\sqrt{J^n}=J\)。

理想 \(J^n\) 的非主性否定的是主理想整环条件。唯一分解整环的定义要求元素的不可约分解存在且唯一，并不要求全部理想为主理想；因此本题的非主性证明不能否定它。下一章\cref{b1:thm:polynomial-ufd}将说明 \(\mathbb Z[x]\) 实际上是唯一分解整环。
\end{solution}

\begin{exercise}\label{b1:ex:11-factorization-nonexistence}
设 \(F\) 为域。令 \(A\) 由形式有限和 \(\sum_q a_qT^q\) 组成，其中 \(a_q\in F\)，指数 \(q\) 取形如 \(m/2^n\) 的非负数，\(m,n\in\mathbb Z_{\geq0}\)。两个形式和相等指每个指数的系数相等；加法逐项进行，乘法由 \(T^qT^r=T^{q+r}\) 及分配律规定。证明 \(A\) 是整环，且 \(T\) 不能分解为有限个不可约元的乘积。再证明 \(A\) 的每个有限生成理想都是主理想，\(A\) 的每个不可约元都是素元，并说明为何 \(A\) 仍不是主理想整环。由此区分分解不存在与分解不唯一两种失败。
\end{exercise}
\begin{solution}
本题要分别检验两种容易被误用的推论：每个不可约元都是素元，仍不能代替不可约分解的存在性；每个有限生成理想都是主理想，也不能代替所有理想都是主理想。下面先考察元素的分解，再考察理想。

任意有限多个这样的形式和，其涉及的指数都可取共同分母 \(2^n\)；以 \(U=T^{1/2^n}\) 表示时，它们就是通常的 \(F[U]\) 中的多项式。运算结果仍属 \(A\)，加法、乘法的环公理可在这个共同多项式环中验证。因此 \(A\) 是交换含幺环。

对非零 \(f\in A\)，记其最大与最小非零系数指数分别为 \(d(f),v(f)\)。因 \(F\) 为域，最高和最低指数处的系数相乘均非零，故对非零 \(f,g\)，有
\[
d(fg)=d(f)+d(g),\qquad v(fg)=v(f)+v(g).
\]
这首先说明 \(fg\ne0\)，所以 \(A\) 为整环；又说明其单位恰为非零常数，因为 \(fg=1\) 迫使两个非负的最大指数都为零。

若 \(T=fg\)，则
\[
0=d(T)-v(T)
=\bigl(d(f)-v(f)\bigr)+\bigl(d(g)-v(g)\bigr).
\]
两个括号内均为非负数，故各自为零，\(f,g\) 都是单项式。因此 \(T\) 的每个非单位因子都是 \(cT^q\)、\(c\ne0,q>0\)。但
\[
cT^q=(cT^{q/2})T^{q/2}
\]
总是两个非单位的乘积，因为 \(q/2\) 仍是允许的正指数。所以 \(T\) 没有不可约因子，更不可能有有限不可约分解。

有限多个元素 \(f_1,\ldots,f_r\in A\) 同属于某个 \(F[U]\)，其中 \(U=T^{1/2^n}\)。由于 \(F[U]\) 是主理想整环，存在 \(d\in F[U]\) 使 \((f_1,\ldots,f_r)F[U]=dF[U]\)。于是每个 \(f_j\) 都是 \(d\) 在 \(A\) 中的倍数，而 \(d\) 是 \(f_1,\ldots,f_r\) 在 \(F[U]\) 中的线性组合，也属于它们在 \(A\) 中生成的理想。因此 \((f_1,\ldots,f_r)A=dA\)；零理想也由 \(0\) 生成。

设 \(p\in A\) 不可约，且 \(p\mid ab\)。若 \(p\mid a\)，素性所需的结论已经成立；现设 \(p\nmid a\)。由于二生成理想为主理想，可写 \((p,a)=(d)\)。由 \(d\mid p\) 和 \(p\) 的不可约性，\(d\) 或为单位，或与 \(p\) 相伴；后一种情形会使 \(p\mid a\)，故 \((p,a)=A\)。于是存在 \(u,v\in A\)，使 \(1=up+va\)。两边乘 \(b\)，再用 \(p\mid ab\)，即得 \(p\mid b\)。所以 \(p\) 是素元。

现在把不断分解元素与不断增大主理想并排来看：
\[
T=(T^{1/2})^2=(T^{1/4})^4=\cdots,
\qquad
(T)\subsetneq(T^{1/2})\subsetneq(T^{1/4})\subsetneq\cdots.
\]
每个包含都严格，因为 \((T^{1/2^n})\) 的非零元素的最小指数至少为 \(1/2^n\)，而 \(T^{1/2^{n+1}}\) 的最小指数更小。\cref{b1:thm:pid-ufd}证明分解存在性时所用的主理想升链稳定性，在这里失效了。还可直接构造一个非主理想：令
\[
J=\bigcup_{n\geq0}(T^{1/2^n}).
\]
这是理想升链的并，因而是理想。若 \(J\) 有有限个生成元，它们必同时落在某个 \((T^{1/2^N})\) 中，便有 \(J\subseteq(T^{1/2^N})\)；但下一项的生成元属于 \(J\)，却不属于该理想，矛盾。因此 \(J\) 不是有限生成理想，尤其不是主理想，\(A\) 不是主理想整环。这也具体说明了为何前面关于有限生成理想的结论仍然成立。

即使每个不可约元都是素元，也不能省略唯一分解整环所要求的分解存在性。与\cref{b1:prop:sqrt-minus-five-factorization}中 \(6\) 有两种不等价不可约分解相比，本例失败的是分解存在性。
\end{solution}

\begin{exercise}\label{b1:ex:11-ideal-splitting-three}
令 \(s=\sqrt{-5}\)、\(R=\mathbb Z[s]\subseteq\mathbb C\)，考虑映射
\[
\Phi:R\longrightarrow\mathbb F_7\times\mathbb F_7,\qquad
a+bs\longmapsto(\overline{a+3b},\overline{a-3b}).
\]
证明 \(\Phi\) 是保持单位的满环同态，且 \(\ker\Phi=(7)\)。按第一、第二坐标的顺序，将两个坐标映射的核记为 \(P,Q\)。求它们的理想生成元，并证明它们是不同的极大理想，且 \(P+Q=R\)、\(P\cap Q=PQ=(7)\)。确定 \(R/(7)\) 的全部幂等元，说明两个非平凡幂等元分别怎样选取上述直积中的一个分量。最后证明 \(7\) 不可约却不是素元，而 \(P,Q\) 都不是主理想。
\end{exercise}
\begin{solution}
两种模 \(7\) 的求值 \(s\mapsto3\)、\(s\mapsto-3\) 都保持关系 \(s^2=-5\)，因为 \((\pm3)^2=9\equiv-5\pmod7\)。所以 \(\Phi\) 保持加法、乘法和单位。任给 \(u,v\in\mathbb F_7\)，选整数 \(a,b\) 满足
\[
\bar a=4(u+v),\qquad \bar b=6(u-v).
\]
这里 \(4=2^{-1}\)、\(6=6^{-1}\) 在 \(\mathbb F_7\) 中成立，代入得到 \(\Phi(a+bs)=(u,v)\)，证明满射。若 \(\Phi(a+bs)=0\)，两坐标相加、相减给出 \(2\bar a=0\)、\(6\bar b=0\)，故 \(7\mid a\)、\(7\mid b\)。反过来，\(7\) 的倍数都在核中，因此 \(\ker\Phi=(7)\)。由\cref{b1:thm:ring-first-iso}，\(\Phi\) 诱导同构
\[
\bar\Phi:R/(7)\longrightarrow\mathbb F_7\times\mathbb F_7.
\]

第一坐标为零的条件是 \(7\mid a+3b\)，此时 \(a+bs=b(s-3)+(a+3b)\)，所以 \(P=(7,s-3)\)；同理，由 \(a+bs=b(s+3)+(a-3b)\) 得 \(Q=(7,s+3)\)。两个坐标映射都满射到域，由\cref{b1:thm:ring-first-iso}分别得到 \(R/P\cong\mathbb F_7\)、\(R/Q\cong\mathbb F_7\)，再由\cref{b1:prop:prime-max-quotient}知 \(P,Q\) 都是极大理想。元素 \(s-3\) 属于 \(P\)，其第二坐标为 \(-6=1\ne0\)，故 \(P\ne Q\)。又有 \(6=(s+3)-(s-3)\in P+Q\) 及 \(7\in P+Q\)，所以 \(1=7-6\in P+Q\)。由\cref{b1:lem:comaximal-product}，\(PQ=P\cap Q\)；而两坐标核的交就是 \(\ker\Phi\)，故 \(PQ=P\cap Q=(7)\)。

域中的幂等元只有 \(0,1\)，所以直积中的幂等元恰为 \((0,0),(1,1),(1,0),(0,1)\)。它们在 \(R/(7)\) 中的原像分别为
\[
0,\qquad1,\qquad e_+=(4-s)+(7),\qquad e_-=(s-3)+(7).
\]
确实，\(\bar\Phi(e_+)=(1,0)\)、\(\bar\Phi(e_-)=(0,1)\)。若 \(\bar\Phi(z)=(u,v)\)，则 \(\bar\Phi(e_+z)=(u,0)\)、\(\bar\Phi(e_-z)=(0,v)\)：乘以这两个幂等元分别保留第一、第二分量。因而 \(e_++e_-=1\)、\(e_+e_-=0\)，且它们生成的两个理想分别对应直积的两个分量。

\(N(7)=49\)。若 \(7\) 有非平凡分解，两个非单位因子的范数只能都为 \(7\)。但 \(a^2+5b^2=7\) 无整数解：\(b=0\) 要求 \(a^2=7\)，\(|b|=1\) 要求 \(a^2=2\)，\(|b|\ge2\) 则左边至少为 \(20\)。所以 \(7\) 不可约。另一方面，
\[
(s-3)(s+3)=s^2-9=-14
\]
被 \(7\) 整除，两个因子却都不被 \(7\) 整除，因为它们的 \(s\) 系数为 \(1\)。故 \(7\) 不是素元。

若 \(P=(d)\)，由 \(7\in P\) 得 \(d\mid7\)。不可约性说明 \(d\) 或为单位，或与 \(7\) 相伴。前者与 \(P\) 为真理想矛盾；后者使 \(P=(7)\)，与 \(s-3\in P\setminus(7)\) 矛盾。因此 \(P\) 非主，对 \(Q\) 用 \(s+3\) 作同样判断即可。这样，\((7)=PQ\) 是非主理想之间的乘积分解，与元素 \(7\) 的不可约性相容。
\end{solution}

% !TeX root = ../../Book1.tex
% 纲要标题：## 第12章　多项式环、不可约性与有限性
% 纲要位置：计划纲要.md，第 442 行。

\chapter{多项式环、不可约性与有限性}
\label{b1:ch:12}
\BookChapterNumber{b1:ch:12}

\begin{chapterabstract}
系数环影响多项式的根、因子和整除计算。本章以带余除法研究一元多项式，建立不可约性判据及唯一分解结果；转入多元情形后，初等对称多项式处理置换不变式，Hilbert 基定理处理理想的有限生成性。
\end{chapterabstract}

\begin{learninggoals}
\begin{itemize}
\item 能区分多项式、它定义的函数及其在不同域中的根，并用形式导数判别重根。
\item 掌握本原多项式与容度的作用，说明分式域上的分解为何能够下降。
\item 会选择不可约性判据，说明判据的假设与检验失败时能够得出的结论。
\item 会用初等对称多项式表示对称多项式，理解消项过程的终止性和表达式的唯一性，并能用根与系数关系作消元计算。
\item 理解 Noether 性与理想升链条件的等价性，掌握 Hilbert 基定理的消项证明。
\end{itemize}
\end{learninggoals}

同一个整数系数多项式，换一个系数域就可能呈现不同的根与因式。譬如在有理数域上 \(X^4-1=(X-1)(X+1)(X^2+1)\)；在 \(\mathbb F_2\) 上，负号与正号相同，整个多项式竟是 \((X+1)^4\)。根的个数、重数与可分解性因此都离不开系数所在的环。带余除法和形式导数使这些差别可以精确计算；当系数来自整环时，还须弄清在分式域中找到的因式怎样回到原来的环。

\ProseParagraphBreak
根还会把我们带向多个不定元。若 \(\alpha,\beta\) 是 \(X^2-uX+v\) 的根，交换二者不会改变 \(\alpha^2+\beta^2\)，而且
\[
\alpha^2+\beta^2=(\alpha+\beta)^2-2\alpha\beta=u^2-2v.
\]
这提示我们研究一切交换变量后不变的多项式，弄清它们能否由少数基本表达式生成。另一种有限性问题出现在多元多项式环的理想中：理想可能有无限多个元素，是否总能选出有限个生成元控制全部元素？这两问把一元多项式的计算引向本章后半的对称多项式与 Hilbert 基定理。

\ProseParagraphBreak

本章使用\cref{b1:prop:field-poly-division}的域上多项式带余除法、\cref{b1:prop:pid-bezout}的 Bézout 等式，以及\cref{b1:prop:ufd-prime}的唯一分解整环中不可约元的素性。分式域的基本构造由\cref{b1:prop:fraction-field-basic}给出。第一节始终在域上讨论，第二、三节的系数整环须满足所列的唯一分解假设；最后一节证明 Hilbert 基定理时，只要求底环交换且为 Noether 环，不要求它是整环。

\ProseParagraphBreak
\secref{b1:sec:1201}至\secref{b1:sec:1203}的一元多项式理论为有限域和 Galois 理论提供根、分解与不可约性判据。\secref{b1:sec:1204}的对称多项式用于后续一般方程；其中关于理想有限生成性的证明则供交换代数使用。

\ProseParagraphBreak
多项式与根的基础见柯斯特利金的《代数学引论》第一卷：基础代数（张英伯译）\cite[第5--6章]{Kostrikin2006AlgebraIChinese}，或席南华的《基础代数》第一卷\cite[第6--7章]{Xi2016BasicAlgebraI}。本章讨论的对称多项式与不变环，可参看 Cox、Little 与 O’Shea 的《Ideals, Varieties, and Algorithms》第7章第1--2节\cite[Chapter 7, \S\S 1--2]{CoxLittleOShea2015IdealsVarietiesAlgorithms}；Hilbert 基定理与 Gröbner 基的联系见该书第2章第5节\cite[Chapter 2, \S 5]{CoxLittleOShea2015IdealsVarietiesAlgorithms}。

\ProseParagraphBreak
该书第2章第2--4、6--8节分别介绍单项式序、多元除法、单项式理想及 Gröbner 基算法，第3章第1--2节讨论消去及其几何应用\cite[Chapter 2, \S\S 2--4, 6--8; Chapter 3, \S\S 1--2]{CoxLittleOShea2015IdealsVarietiesAlgorithms}。延拓定理与闭包定理的完整证明见第3章第5--6节及第4章第7节\cite[Chapter 3, \S\S 5--6; Chapter 4, \S 7]{CoxLittleOShea2015IdealsVarietiesAlgorithms}。

% !TeX root = ../../Book1.tex
% 纲要标题：### 12.1 一元多项式
% 纲要位置：计划纲要.md，第 444 行。

\section{一元多项式}
\label{b1:sec:1201}

% 纲要内容（仅作写作计划，尚非教材正文）：
%
% * 域上的带余除法；整环上的因式定理与根数估计；根与重数
% * 形式导数与重根判别
% * Lagrange 插值与中国剩余定理
% * 代数闭域中的分解与代数基本定理
% * 代数基本定理在此先行引用；利用实闭域刻画的证明归附录 E.2，须先有第19章的有限 Galois 对应及第一编的 Sylow 理论
%

同一个多项式可以被看作系数的有限序列，也可以在给定元素处求值。前一种观点适合讨论整除，后一种观点适合讨论方程的根。带余除法把这两种观点联系起来：除以 \(x-a\) 的余数，恰好就是在 \(a\) 处的值。本节的系数域记为 \(K\)；涉及更大的域时，会明确说明系数与根分别位于哪里。

\subsection{带余除法、根与重数}
\label{b1:subsec:1201:01}

\cref{b1:prop:field-poly-division}已经证明：若 \(g\in K[x]\) 非零，则任意 \(f\in K[x]\) 唯一写成 \(f=qg+r\)，其中 \(r=0\) 或 \(\deg r<\deg g\)。这里能够不断消去最高次项，是因为 \(g\) 的首项系数在域中可逆。以下先把这一结论用于一次多项式。

\begin{definition}\label{b1:def:polynomial-root}
设 \(K\) 为域，\(f\in K[x]\)，\(a\in K\)。若代入所得的元素 \(f(a)\) 为零，就称 \(a\) 为 \(f\) 的一个\term[gen]{根}[root]。若 \(f\ne0\)，且对某个正整数 \(m\) 和某个 \(g\in K[x]\) 有
\[
f=(x-a)^m g,\qquad g(a)\ne0,
\]
就称 \(m\) 为根 \(a\) 的\term[chongshu]{重数}[multiplicity]。重数为一的根称为\term[dangen]{单根}[simple root]，重数至少为二的根称为\term[chonggen]{重根}[multiple root]。
\end{definition}

若 \(L/K\) 为域扩张，我们将 \(f\) 视为 \(L[x]\) 中的多项式，并用同一定义谈 \(L\) 中的根及其重数。把系数域换成整环 \(R\) 时，也可用上述条件定义 \(R\) 中的根及其重数；下面的因式定理与根数估计在这个范围内仍然成立。

\begin{proposition}[因式定理与根数估计]\label{b1:prop:factor-root}
设 \(R\) 为整环，\(f\in R[x]\) 为非零多项式，\(a\in R\)。则 \(f(a)=0\) 当且仅当 \(x-a\mid f\)。多项式 \(f\) 的每个根都有唯一重数；若 \(\deg f=n\)，则 \(f\) 在 \(R\) 中的根按重数计至多有 \(n\) 个。
\end{proposition}
\begin{proof}
这里的除数 \(x-a\) 首一，不需要对 \(R\) 中的非零系数作除法。对每个 \(i\ge1\)，恒等式
\[
x^i-a^i=(x-a)\sum_{j=0}^{i-1}x^{i-1-j}a^j
\]
说明 \(f-f(a)\) 是 \(x-a\) 的倍数，故 \(f=(x-a)q+f(a)\)。反过来，\(x-a\) 的任何倍数在 \(a\) 处都为零，这就证明了因式定理。若 \(a\) 是根，就可以反复除去 \(x-a\)。每次除法都使非零商的次数减少一，所以至多进行 \(\deg f\) 次。记除去的次数为 \(m\)，则最后得到 \(f=(x-a)^m g\)、\(g(a)\ne0\)。若又有 \(f=(x-a)^r h\)、\(h(a)\ne0\)，设 \(m<r\)，在整环 \(R[x]\) 中消去 \((x-a)^m\)，便有 \(g=(x-a)^{r-m}h\)，与 \(g(a)\ne0\) 矛盾。交换两种分解也排除 \(r<m\)，故重数唯一。

若 \(a\ne b\)，从 \(f=(x-a)^m g\) 可见，\(b\) 是 \(f\) 的根当且仅当它是 \(g\) 的根，因为 \((b-a)^m\ne0\)，且整环中非零元素相乘仍非零。再将 \(g\) 中所有 \(x-b\) 因子除去，所得最后因子在 \(b\) 处非零，乘以 \((x-a)^m\) 后仍然非零，因此 \(b\) 的重数也不变。对任意有限个不同的根，逐一除去它们对应的因子，次数总共下降各重数之和，且不会降为负数，故这一和至多为 \(n\)。因此不同根本身也不可能超过 \(n\) 个，根数估计随之成立。
\end{proof}

例如在 \(\mathbb Q[x]\) 中，\((x-1)^3(x+2)\) 有两个不同的根，但按重数共有四个根。根数估计要求多项式非零，也要求系数处于整环；它不要求所有非零系数可逆，所以也适用于 \(\mathbb Z[x]\)。在 \(\mathbb Z/8\mathbb Z\) 中，\(x^2-1\) 在 \(\bar1,\bar3,\bar5,\bar7\) 处都为零，二次多项式却有四个不同的根。失败来自零因子：不同根之差虽非零，却可能是零因子，因而从 \((b-a)^m g(b)=0\) 不能推出 \(g(b)=0\)。

\subsection{形式导数与重根判别}
\label{b1:subsec:1201:02}

导数在这里由系数定义，不借助极限。它在任意特征下都有意义，但正特征会改变整数系数的取值。

\begin{definition}\label{b1:def:formal-derivative}
设 \(K\) 为域，\(n\geq0\) 为整数，\(a_0,\ldots,a_n\in K\)，并令 \(f=\sum_{i=0}^n a_i x^i\)。定义 \(f\) 的\term[xingshi-daoshu]{形式导数}[formal derivative]为
\[
f'\coloneqq\sum_{i=1}^n i a_i x^{i-1}.
\]
这里 \(i a_i\) 是 \(a_i\) 的 \(i\) 次相加，整数系数在 \(K\) 中解释。
\end{definition}
\symidx{\(f'\)：多项式的形式导数}

形式导数记录的是平移变量后的一次项。设 \(L/K\) 为域扩张，\(a\in L\)，\(T\) 为不定元。按二项式公式展开 \(f(a+T)\)，所得常数项为 \(f(a)\)，一次项为 \(f'(a)T\)，所以存在 \(h\in L[T]\)，使
\[
f(a+T)=f(a)+f'(a)T+T^2h(T).
\]
若 \(f\ne0\) 且 \(f(a)=0\)，常数项已经消失；再要求 \(f'(a)=0\)，便消去一次项，恰好得到 \(T^2\mid f(a+T)\)，也就是 \((x-a)^2\mid f(x)\)。因此 \(a\) 是重根恰好要求这两项同时消失。下面的因式证明是这一现象的另一种表述。这个解释只用多项式恒等式。在特征 \(p>0\) 时，\((a+T)^p=a^p+T^p\)，一次项确实不存在，这便是 \((x^p)'=0\) 的代数意义。

\ProseParagraphBreak
定义直接给出 \((f+g)'=f'+g'\) 及 \((cf)'=cf'\)。对单项式 \(ax^i,bx^j\)，由系数恒等式 \((i+j)ab=iab+jab\) 得到
\[
\bigl(ax^i\cdot bx^j\bigr)'=(ax^i)'bx^j+ax^i(bx^j)'.
\]
对有限项求和，便得到乘积法则 \((fg)'=f'g+fg'\)。因而也有 \(\bigl((x-a)^m\bigr)'=m(x-a)^{m-1}\)。

\begin{proposition}[重根判别]\label{b1:prop:multiple-root}
设 \(K\) 为域，\(f\in K[x]\) 非零，\(L\) 是包含 \(K\) 的域。元素 \(a\in L\) 是 \(f\) 的重根，当且仅当 \(f(a)=f'(a)=0\)。此外，以下条件等价：
\begin{enumerate}
\item \(\gcd(f,f')=1\)，其中最大公因式取首一形式；
\item 在任何包含 \(K\) 的域中，\(f\) 都没有重根。
\end{enumerate}
\end{proposition}
\begin{proof}
若 \(f(a)=0\)，在 \(L[x]\) 中由\cref{b1:prop:factor-root}写成 \(f=(x-a)g\)。求导后代入 \(a\)，得到 \(f'(a)=g(a)\)。因此 \(f'(a)=0\) 等价于 \(x-a\mid g\)，也就等价于 \((x-a)^2\mid f\)。

若 \(\gcd(f,f')=1\)，由于 \(K[x]\) 是主理想整环，\cref{b1:prop:pid-bezout}给出 \(u,v\in K[x]\)，使 \(uf+vf'=1\)。这个等式在 \(L[x]\) 中仍成立，所以不可能有 \(a\in L\) 同时使 \(f(a)\) 和 \(f'(a)\) 为零。

反过来，若最大公因式有正次数，在其中选一个不可约因子 \(q\)。\cref{b1:thm:euclidean-pid}及\cref{b1:prop:pid-bezout}说明 \(K[x]/(q)\) 是域：若 \(q\nmid h\)，则 \(\gcd(q,h)=1\)，某个 Bézout 等式模 \(q\) 后给出 \(h+(q)\) 的逆。常数域 \(K\) 嵌入这个商域，因为非零常数不可能被正次数的 \(q\) 整除。令 \(a=x+(q)\)，便有 \(q(a)=0\)；这个商域的作用是向 \(K\) 中形式地添入 \(q\) 的一个根，相关构造将在\secref{b1:sec:1401}系统研究。由于 \(q\mid f\) 及 \(q\mid f'\)，又有 \(f(a)=f'(a)=0\)，故 \(a\) 是其中的重根。
\end{proof}

这里的“任何包含 \(K\) 的域”不能改成仅检查 \(K\) 本身。例如 \((x^2+1)^2\) 在 \(\mathbb Q\) 中没有根，在 \(\mathbb C\) 中却有重根。判别条件只涉及系数域 \(K\) 中的最大公因式，却能控制所有扩域中的重根，不必先求出各个扩域中的根。

\ProseParagraphBreak
设域 \(K\) 的特征为素数 \(p\)。此时 \(f'=0\) 当且仅当 \(f\) 的所有非零系数项的指数都被 \(p\) 整除，即 \(f=\sum_j a_jx^{pj}\)：若 \(p\nmid i\)，则 \(i\cdot1_K\ne0\)，从 \(ia_i=0\) 可消去它而得 \(a_i=0\)。这是关于指数的条件，并未要求系数 \(a_j\) 具有 \(p\) 次根。因此非常数多项式也可能导数为零。例如在 \(\mathbb F_p[x]\) 中，\(x^p-1=(x-1)^p\)，它的导数为零；等式来自中间二项式系数都被 \(p\) 整除。

\ProseParagraphBreak
要进一步写成 \(f=g^p\)，还须解决系数问题：为每个 \(a_j\) 在 \(K\) 中找到 \(b_j\)，使 \(a_j=b_j^p\)。这个条件既必要也充分，因为
\[
\left(\sum_j b_jx^j\right)^p=\sum_j b_j^px^{pj}.
\]
所以 \(f'=0\) 只限制指数，系数是否有 \(p\) 次根则决定 \(f\) 能否成为 \(K[x]\) 中的 \(p\) 次幂。具体例子见\cref{b1:ex:12-primitive-linear-polynomials}。

\subsection{Lagrange 插值与中国剩余定理}
\label{b1:subsec:1201:03}

根数估计还带来另一种用途：两个次数有界的多项式，在足够多的点处相等，就必定相同。

\begin{theorem}[Lagrange 插值]\label{b1:thm:lagrange-interpolation}
设 \(K\) 为域，\(n\geq0\) 为整数，\(a_0,\ldots,a_n\in K\) 互不相同，\(b_0,\ldots,b_n\in K\)。令
\[
\ell_i(x)=\prod_{\substack{0\le j\le n\\j\ne i}}\frac{x-a_j}{a_i-a_j}
\qquad(0\le i\le n).
\]
空乘积取一。这些多项式满足 \(\ell_i(a_j)=\delta_{ij}\)，其中 \(\delta_{ij}\) 在 \(i=j\) 时为一，否则为零。存在唯一的次数不超过 \(n\) 的多项式 \(f\in K[x]\)，使 \(f(a_i)=b_i\) 对所有 \(0\leq i\leq n\) 成立。它由
\begin{equation}\label{b1:eq:lagrange-interpolation}
f(x)=\sum_{i=0}^n b_i\ell_i(x)
\end{equation}
给出；这里也允许 \(f=0\)。
\end{theorem}
\begin{proof}
由于各 \(a_i\) 互异，每个分母均非零，故 \(\ell_i\) 有意义。它在 \(a_i\) 处每个因子都等于一，在其余 \(a_j\) 处有一个因子为零，所以 \(\ell_i(a_j)=\delta_{ij}\)。每个 \(\ell_i\) 的次数至多为 \(n\)，故所列 \(f\) 有所需取值，次数也至多为 \(n\)。这就是用每次只保留一个指定值的插值基多项式，拼出全部指定值。若 \(g\) 也是这样的多项式，则 \(f-g\) 的次数至多为 \(n\)，却有 \(n+1\) 个不同的根。\cref{b1:prop:factor-root}迫使 \(f-g=0\)。
\end{proof}

这称为\term[lagrange-chazhi]{Lagrange 插值}[Lagrange interpolation]。例如在 \(\mathbb Q\) 中指定 \(f(0)=1,f(1)=2,f(2)=5\)，公式给出
\[
f(x)=\frac{(x-1)(x-2)}2-2x(x-2)+\frac{5x(x-1)}2=x^2+1.
\]
唯一性只对次数至多为二的多项式成立；加上任意 \(x(x-1)(x-2)h(x)\)，这三个值都不改变。

\ProseParagraphBreak
\phantomsection\label{b1:note:polynomial-function-distinction}
形式多项式相等，指对应系数逐项相等；它们在每个元素处的取值当然也相等，即
\[
\boxed{\ f=g\;\text{于}\;K[x]\quad\Longrightarrow\quad
f(a)=g(a)\;\text{对每个}\;a\in K\ }.
\]
反向一般不成立。设 \(p\) 为素数。在 \(\mathbb F_p[x]\) 中，\(x^p-x\) 的首项系数为一，所以不是零多项式，却在每个 \(a\in\mathbb F_p\) 处取值为零：\(a=0\) 时直接成立，\(a\ne0\) 时，\cref{b1:thm:lagrange}用于 \(p-1\) 阶群 \(\mathbb F_p^\times\) 给出 \(a^{p-1}=1\)。因此它与零多项式定义同一个函数。要由取值推出多项式相等，还须用根数估计核对次数界；例如在 \(\mathbb F_p\) 上，两个次数均至多为 \(p-1\) 的多项式若处处取值相同，其差便有 \(p\) 个不同根而次数小于 \(p\)，只能是零多项式。

\ProseParagraphBreak
这个插值公式也可以从\cref{b1:thm:crt}理解。对定理中的互异元素 \(a_0,\ldots,a_n\)，令 \(I_i=(x-a_i)\subseteq K[x]\)。若 \(i\ne j\)，则 \((x-a_i)-(x-a_j)=a_j-a_i\) 是非零常数，故 \(I_i+I_j=K[x]\)。中国剩余定理给出同构
\[
K[x]\Big/\left(\prod_{i=0}^n(x-a_i)\right)
\longrightarrow K^{n+1},\qquad
[g]\longmapsto\bigl(g(a_0),\ldots,g(a_n)\bigr).
\]
定理中的插值基多项式 \(\ell_i\) 在这个商环中对应第 \(i\) 个坐标幂等元，因为它的取值向量只有第 \(i\) 项为一，其余项全为零。中国剩余定理负责实现任意指定值，确定相应的唯一陪集；带余除法再为这个陪集选出唯一的低次代表。具体地，对 \(P=\prod_{i=0}^n(x-a_i)\) 作带余除法，所得余式的次数至多为 \(n\)。若两个这样的代表属于同一陪集，它们的差是 \(P\) 的倍数，而非零倍数的次数至少为 \(n+1\)，所以两者必相等。

\subsection{代数闭域中的分解与代数基本定理}
\label{b1:subsec:1201:04}

根数估计限制已有根的数量，却不保证根一定存在。现在换一个问题：在哪些域中，每个非常数多项式都至少有一个根？

\begin{definition}\label{b1:def:algebraically-closed}
若域 \(K\) 上每个非常数多项式在 \(K\) 中都有根，就称 \(K\) 为\term[daishu-bi-yu]{代数闭域}[algebraically closed field]。
\end{definition}

定义只要求至少存在一个根，却会自动产生完全分解：提出一个一次因子以后，商仍在 \(K[x]\) 中；只要它仍是非常数多项式，就能再次使用同一个存在条件，直到次数降为零。

\begin{proposition}\label{b1:prop:algebraically-closed-splitting}
域 \(K\) 代数闭，当且仅当每个次数为 \(n\ge1\) 的多项式 \(f\in K[x]\) 都能写成
\[
f(x)=a_n\prod_{i=1}^n(x-\alpha_i),\qquad \alpha_i\in K,
\]
其中 \(a_n\) 是 \(f\) 的首项系数。这时按重数计共有 \(n\) 个根，根的多重集合由 \(f\) 唯一确定。
\end{proposition}
\begin{proof}
若 \(K\) 代数闭，选 \(f\) 的一个根 \(\alpha_1\)。\cref{b1:prop:factor-root}给出 \(f=(x-\alpha_1)g\)，而 \(\deg g=n-1\)。对次数归纳，直到商为非零常数，得到所需分解，常数因子由首项系数确定。每个根出现的次数恰为\cref{b1:prop:factor-root}中唯一确定的重数，因此根的多重集合唯一。反过来，所列分解直接给出任一非常数多项式的一个根。
\end{proof}

对通常的复数域，根的存在还要另作论证。以下结论称为\term[daishu-jiben-dingli]{代数基本定理}[fundamental theorem of algebra]；它不是复数的定义。此处先引用这一结论；利用实分析的中间值定理和有限 Galois 理论的证明见\secref{b1:sec:E02}。

\begin{theorem}[代数基本定理]\label{b1:thm:fundamental-theorem-algebra}
复数域 \(\mathbb C\) 是代数闭域。因此，每个次数为 \(n\ge1\) 的复系数多项式都有恰好 \(n\) 个复根，按重数计，并可分解为一次因子的乘积。
\end{theorem}

不可约性讨论必须始终注明系数域；例如 \(x^2+1\) 在 \(\mathbb R\) 中无根，在 \(\mathbb C[x]\) 中却等于 \((x-i)(x+i)\)。

% !TeX root = ../../Book1.tex
% 纲要标题：### 12.2 Gauss 引理与多项式环的唯一分解
% 纲要位置：计划纲要.md，第 450 行。

\section{Gauss 引理与多项式环的唯一分解}
\label{b1:sec:1202}

% 纲要内容（仅作写作计划，尚非教材正文）：
%
% * 本原多项式与容度
% * 首次完整构造分式域，再证明分式域上的分解及其下降
% * 唯一分解整环上的多项式仍为唯一分解整环
%

整数系数多项式在有理数域中分解以后，因子的分母是否真的增加了分解的可能性？答案取决于能否先把系数的公共因子剥离出来。本节始终设 \(R\) 为唯一分解整环。前一章已经证明的“不可约元必为素元”，将使约化到 \(R/(p)\) 的办法适用。

\subsection{本原多项式与容度}
\label{b1:subsec:1202:01}

\begin{definition}\label{b1:def:polynomial-content}
设 \(R\) 为唯一分解整环，\(n\geq0\) 为整数，且 \(f=a_0+\cdots+a_nx^n\in R[x]\) 为非零多项式。各系数 \(a_0,\ldots,a_n\) 的任一最大公因子称为 \(f\) 的\term[rongdu]{容度}[content]，记为 \(c(f)\)。如果 \(c(f)\) 是单位，就称 \(f\) 为\term[benyuan-duoxiangshi]{本原多项式}[primitive polynomial]。容度只在相伴意义下确定；当 \(R=\mathbb Z\) 时，统一选取正的最大公因子。
\end{definition}
\symidx{\(c(f)\)：非零多项式的容度}

最大公因子存在，因为在各非零系数的素因子分解中取指数的最小值即可；零系数不增加限制。将每个系数除以 \(c(f)\)，得到 \(f=c(f)f_0\)，其中 \(f_0\) 本原。例如 \(6x^2+9x+3=3(2x^2+3x+1)\)。本原并不要求某个系数为单位：\(2x+3\) 在 \(\mathbb Z[x]\) 中本原，但两个非零系数都不是单位。

\begin{lemma}[Gauss]\label{b1:lem:gauss-primitive}
设 \(R\) 为唯一分解整环。则 \(R[x]\) 中两个本原多项式的乘积仍为本原多项式。因此，对非零 \(f,g\in R[x]\)，容度 \(c(fg)\) 与 \(c(f)c(g)\) 相伴。
\end{lemma}
\begin{proof}
设 \(f,g\) 本原，而 \(fg\) 不本原。则存在不可约元 \(p\)，整除 \(fg\) 的每个系数。由\cref{b1:prop:ufd-prime}，\(p\) 是素元，所以 \(R/(p)\) 是整环：两个剩余类的积为零意味着 \(p\mid ab\)，进而 \(p\mid a\) 或 \(p\mid b\)。逐系数约化得到
\[
\bar f\bar g=0\qquad\text{于 }\bigl(R/(p)\bigr)[x].
\]
因 \(f,g\) 本原，\(p\) 不可能分别整除它们的全部系数，故 \(\bar f,\bar g\) 都非零。\cref{b1:prop:polynomial-degree}说明整环上的非零多项式之积非零，产生矛盾。

一般地，写 \(f=c(f)f_0\)、\(g=c(g)g_0\)。刚才证明了 \(f_0g_0\) 本原，所以从系数的素因子指数取最小值得到，\(fg=c(f)c(g)f_0g_0\) 的容度与 \(c(f)c(g)\) 相伴。
\end{proof}

这个结论称为\term[gauss-yinli]{Gauss 引理}[Gauss lemma]。证明的关键在于 \(R/(p)\) 没有零因子，而非对多项式系数逐项追踪可能发生的抵消。唯一分解性保证不可约元具有这里所需的素性。

\subsection{分式域及分解的下降}
\label{b1:subsec:1202:02}

为了在更大的域中分解多项式，先完整构造整环的分式域。下面的构造不要求唯一分解性；只有随后的约分与本原多项式论证需要这一假设。\secref{b1:sec:1301}将以这里的结果为基础，研究分式域的泛性质，再由\secref{b1:sec:1302}推广到一般局部化。

\begin{proposition}\label{b1:prop:fraction-field-basic}
每个整环 \(R\) 都嵌入一个域 \(\operatorname{Frac}(R)\)，其中每个元素写成 \(a/b\)，\(a,b\in R\)、\(b\ne0\)。对 \(a,b,c,d\in R\) 且 \(b,d\ne0\)，分式的相等判定、加法和乘法由下列公式给出：
\[
\frac ab=\frac cd\quad\Longleftrightarrow\quad ad=bc,\qquad
\frac ab+\frac cd=\frac{ad+bc}{bd},\qquad
\frac ab\frac cd=\frac{ac}{bd}.
\]
这个域称为 \(R\) 的\term[fenshi-yu]{分式域}[field of fractions]。
\end{proposition}
\symidx{\(\operatorname{Frac}(R)\)：整环 \(R\) 的分式域}
\begin{proof}
在所有有序对 \((a,b)\)、\(b\ne0\) 上定义
\((a,b)\sim(c,d)\) 当且仅当 \(ad=bc\)。自反性与对称性直接成立。若另有 \(ce=df\)，则由 \(ad=bc\) 得
\(ade=bce=bdf\)。因 \(d\ne0\)，整环的消去律给出 \(ae=bf\)，所以关系具有传递性。将等价类记为 \(a/b\)。

要验证运算良定义，设 \(ab'=a'b\)、\(cd'=c'd\)。一方面
\[
(ad+bc)b'd'=a'd'bd+b'c'bd=(a'd'+b'c')bd;
\]
另一方面 \(acb'd'=a'c'bd\)。这恰好分别说明更换两个分数的代表后，所给和与积不变。分母的积非零仍由整环假设保证。

加法交换律和乘法交换律继承自 \(R\)。任意三个分数 \(a/b,c/d,e/f\) 的两种加法结合方式都给出 \((adf+bcf+bde)/(bdf)\)，两种乘法结合方式都给出 \(ace/(bdf)\)。分配律两边分别算得
\[
\frac ab\biggl(\frac cd+\frac ef\biggr)
=\frac{a(cf+de)}{bdf},\qquad
\frac{ac}{bd}+\frac{ae}{bf}
=\frac{ab(cf+de)}{b^2df};
\]
它们按交叉相乘相等。零元为 \(0/1\)，单位元为 \(1/1\)，\(a/b\) 的加法逆元为 \((-a)/b\)。分数 \(a/b\) 非零当且仅当 \(a\ne0\)，这时 \(b/a\) 为乘法逆元。因此所得结构是域。映射 \(a\mapsto a/1\) 保持和、积与单位，且 \(a/1=0/1\) 迫使 \(a=0\)，故为所需嵌入。
\end{proof}

以下记 \(F=\operatorname{Frac}(R)\)，并通过上述嵌入把 \(R[x]\) 看作 \(F[x]\) 的子环。由于 \(R\) 唯一分解，每个非零分数均可约去分子、分母的公共素因子，写成互素的 \(a/b\)。

\begin{lemma}\label{b1:lem:primitive-scaling}
设 \(R\) 为唯一分解整环，\(F=\operatorname{Frac}(R)\)。则每个非零 \(h\in F[x]\) 可写成 \(h=\lambda h_0\)，其中 \(\lambda\in F\) 非零，\(h_0\in R[x]\) 本原。若 \(h_0\) 本原且 \(\lambda h_0\in R[x]\)，则 \(\lambda\in R\)；若 \(\lambda h_0\) 也本原，则 \(\lambda\) 是 \(R\) 的单位。
\end{lemma}
\begin{proof}
取 \(h\) 的有限个系数分母的乘积 \(b\ne0\)，则 \(bh\in R[x]\)。把 \(bh\) 的容度 \(a\) 提出，得到 \(h=(a/b)h_0\)，其中 \(h_0\) 本原。

现在将 \(\lambda=a/b\) 约成互素形式。若 \(\lambda h_0\) 的系数都在 \(R\)，对 \(h_0\) 的每个系数 \(u\) 都有 \(au=bv\)，其中 \(v\in R\)。若 \(b\) 不是单位，选不可约因子 \(p\mid b\)。互素性给出 \(p\nmid a\)，而素性迫使 \(p\mid u\)。这对每个系数成立，与 \(h_0\) 本原矛盾。因此 \(b\) 为单位，即 \(\lambda\in R\)。若 \(\lambda h_0\) 也本原，任何非单位 \(\lambda\) 的素因子都会整除它的全部系数，仍产生矛盾，故 \(\lambda\) 为单位。
\end{proof}

\begin{theorem}[Gauss 分解定理]\label{b1:thm:gauss-descent}
设 \(R\) 为唯一分解整环，\(F=\operatorname{Frac}(R)\)，并设 \(f\in R[x]\) 本原且次数为正。则 \(f\) 在 \(R[x]\) 中不可约，当且仅当它在 \(F[x]\) 中不可约。此外，若 \(h\in R[x]\)，且 \(f\mid h\) 在 \(F[x]\) 中成立，则该整除关系已经在 \(R[x]\) 中成立。
\end{theorem}
\begin{proof}
先证最后的整除断言。\(h=0\) 时取商为零即可；否则写 \(h=fq\)，并由\cref{b1:lem:primitive-scaling}写 \(q=\lambda q_0\)，其中 \(q_0\in R[x]\) 本原。\cref{b1:lem:gauss-primitive}说明 \(fq_0\) 本原，而 \(h=\lambda fq_0\in R[x]\)，再由\cref{b1:lem:primitive-scaling}得 \(\lambda\in R\)。所以 \(q\in R[x]\)。

若 \(f=gh\) 在 \(F[x]\) 中分解，且两因子次数均为正，将它们写成 \(g=\alpha g_0\)、\(h=\beta h_0\)，其中 \(g_0,h_0\) 是 \(R[x]\) 中的本原多项式。于是
\(f=\alpha\beta g_0h_0\)。\cref{b1:lem:gauss-primitive}说明 \(g_0h_0\) 本原，且 \(f\) 也本原，所以\cref{b1:lem:primitive-scaling}给出 \(\alpha\beta\in R^\times\)。因此 \(f=(\alpha\beta g_0)h_0\) 是 \(R[x]\) 中的非平凡分解。

反过来，若本原的 \(f\) 在 \(R[x]\) 中分解为两个非单位，由容度的乘法性质，两因子均本原。整环上多项式环的单位只能是底环的单位：若 \(uv=1\)，次数相加为零，故 \(u,v\) 都是常数。因此本原的非单位因子不能为常数，两因子均有正次数，同一分解在 \(F[x]\) 中仍非平凡。
\end{proof}

对正次数的本原多项式，\(\mathbb Z[x]\) 中的不可约性与 \(\mathbb Q[x]\) 中的不可约性等价；本原性排除了非单位常数因子。若只问能否分解为两个正次数的整系数多项式，则不必假设本原：任意正次数的 \(f\in\mathbb Z[x]\) 在 \(\mathbb Q[x]\) 中可约，当且仅当 \(f=gh\)，其中 \(g,h\in\mathbb Z[x]\) 的次数均为正。事实上，写 \(f=c(f)f_0\)，其中 \(f_0\) 本原。在 \(\mathbb Q[x]\) 中，非零常数 \(c(f)\) 是单位；若 \(f_0\) 在该环中可约，先由定理将它写成两个正次数的整系数因子之积，再把 \(c(f)\) 乘入其中一个因子。逆向蕴含由同一分解直接得到。

\ProseParagraphBreak
例如 \(2x+2=2(x+1)\) 在 \(\mathbb Z[x]\) 中是非平凡分解，但在 \(\mathbb Q[x]\) 中 \(2\) 是单位，故 \(2x+2\) 仍为不可约的一次多项式。常数 \(2\) 在 \(\mathbb Z[x]\) 中是不可约元，在 \(\mathbb Q[x]\) 中则为单位，也说明不能不加条件地比较两个环中的不可约元。

\subsection{多项式环的唯一分解性}
\label{b1:subsec:1202:03}

\begin{theorem}\label{b1:thm:polynomial-ufd}
若 \(R\) 是唯一分解整环，则 \(R[x]\) 也是唯一分解整环。因而，对每个整数 \(n\geq1\)，\(R[x_1,\ldots,x_n]\) 仍为唯一分解整环。
\end{theorem}
\begin{proof}
仍令 \(F=\operatorname{Frac}(R)\)。由\cref{b1:prop:field-poly-division}、\cref{b1:thm:euclidean-pid}及\cref{b1:thm:pid-ufd}，\(F[x]\) 是唯一分解整环。对 \(R[x]\) 中非零的 \(f\)，先写 \(f=c(f)f_0\)，其中 \(f_0\) 本原。容度 \(c(f)\) 在 \(R\) 中有不可约分解。若 \(\deg f_0>0\)，再在 \(F[x]\) 中分解 \(f_0\)，把每个正次数不可约因子乘以适当非零常数，化为 \(R[x]\) 中的本原因子 \(q_1,\ldots,q_s\)。于是
\[
f_0=\lambda q_1\cdots q_s.
\]
右边的多项式乘积本原，故\cref{b1:lem:primitive-scaling}说明 \(\lambda\in R^\times\)。而\cref{b1:thm:gauss-descent}说明每个 \(q_i\) 在 \(R[x]\) 中不可约。若 \(f_0\) 为常数，则本原性已说明它为单位，不需这一步。

还要证明这些分解唯一。\(R\) 的每个不可约元 \(p\)，在 \(R[x]\) 中都是素元：模 \(p\) 约化后的多项式环 \(\bigl(R/(p)\bigr)[x]\) 是整环，故 \(p\mid gh\) 蕴含 \(p\mid g\) 或 \(p\mid h\)。正次数的本原因子 \(q_i\) 在 \(F[x]\) 中不可约，\cref{b1:prop:ufd-prime}说明它在这个唯一分解整环中是素元。若 \(q_i\mid gh\) 于 \(R[x]\)，先在 \(F[x]\) 中得到 \(q_i\mid g\) 或 \(q_i\mid h\)，再由\cref{b1:thm:gauss-descent}的整除下降，得到相应的商已经属于 \(R[x]\)。因此 \(q_i\) 也为素元。

现在任一非零非单位都已分解为素元的乘积。任取 \(R[x]\) 中的不可约元 \(h\)，将它按上述方法分解；由不可约性，该分解只能有一个非单位因子，所以 \(h\) 与某个素元相伴，本身也是素元。比较两个不可约分解，第一个分解中的任一素因子必须整除第二个分解中的某个不可约因子，两者因而相伴；消去它们并重复，就得到因子个数相同且逐一相伴。这就是唯一分解性。有限多个不定元的结论，再用
\(R[x_1,\ldots,x_n]=\bigl(R[x_1,\ldots,x_{n-1}]\bigr)[x_n]\)
对 \(n\) 归纳。
\end{proof}

这个定理说明 \(\mathbb Z[x]\) 和域 \(K\) 上的 \(K[x,y]\) 都具有唯一分解性，却没有说明它们是主理想整环。唯一分解描述元素的因子分解；主理想条件要求一个理想中的全部元素由单个元素生成，是更强的要求。

\ProseParagraphBreak
本原性也不意味着系数生成单位理想。取 \(R=\mathbb Z[y]\)，由上一定理可知它是唯一分解整环。在 \(R[x]\) 中，\(f(x)=2x+y\) 的系数 \(2,y\) 没有共同的非单位因子，所以 \(f\) 本原；但 \((2,y)\ne R\)：先令 \(y=0\)，再将整数模 \(2\) 约化，便得到一个把 \(2,y\) 都送到零、把 \(1\) 送到一的环同态。因此一般不能从本原性推出系数存在以 \(1\) 为结果的 Bézout 线性组合。这也解释了上面的证明为何使用素因子的整除性质。

\ProseParagraphBreak
若底环还满足 Noether 条件，则有限多个不定元的多项式环中每个理想都有限生成；本章最后一节将由 Hilbert 基定理证明这一点。前面列举的 \(\mathbb Z[x]\) 和 \(K[x,y]\)，其底环 \(\mathbb Z\) 与 \(K\) 都满足这个条件。

% !TeX root = ../../Book1.tex
% 纲要标题：### 12.3 不可约性判据
% 纲要位置：计划纲要.md，第 456 行。

\section{不可约性判据}
\label{b1:sec:1203}

% 纲要内容（仅作写作计划，尚非教材正文）：
%
% * Eisenstein 判据及平移使用
% * 模素数约化法
% * 有理根检验的适用范围与失效例子
%

唯一分解定理保证因子分解存在而且唯一，却没有立即告诉我们某个给定多项式是否已经不可约。本节给出几种可以实际使用的判据。它们都只能在相应假设下推出结论；某一种检验没有奏效，并不等于多项式可约。

\subsection{Eisenstein 判据及平移}
\label{b1:subsec:1203:01}

设 \(R\) 为唯一分解整环，\(p\in R\) 为素元。若 \(f=a_nx^n+\cdots+a_0\in R[x]\) 的首项系数不被 \(p\) 整除，而较低次系数都被 \(p\) 整除，则 \(f\) 模 \(p\) 后只剩单项式 \(\bar a_nx^n\)。在整环上的多项式乘积中，最低非零项与最高非零项的指数分别相加；若乘积只含一项，两个非零因子各自也只能含一项。首项系数不被 \(p\) 整除又保证两个因子的次数在约化后不变。因此，若原来有两个正次数因子，它们约化后必是正次数单项式，两个常数项都被 \(p\) 整除，原常数项就被 \(p^2\) 整除。Eisenstein 判据正是用常数项排除这种可能。

\begin{theorem}[Eisenstein]\label{b1:thm:eisenstein}
设 \(R\) 为唯一分解整环，\(n\geq1\) 为整数，\(f=a_nx^n+\cdots+a_1x+a_0\in R[x]\) 本原。若 \(R\) 中存在素元 \(p\)，满足
\[
p\nmid a_n,\qquad p\mid a_i\ (0\le i<n),\qquad p^2\nmid a_0,
\]
则 \(f\) 在 \(R[x]\) 及 \(\operatorname{Frac}(R)[x]\) 中均不可约。
\end{theorem}
\begin{proof}
由\cref{b1:thm:gauss-descent}，只需排除 \(R[x]\) 中两个正次数因子的分解。假设 \(f=gh\)，\(\deg g=r>0\)、\(\deg h=s>0\)，则 \(r+s=n\)。首项系数的积为 \(a_n\)，不被 \(p\) 整除，所以 \(g,h\) 的首项系数模 \(p\) 都非零；约化后次数仍为 \(r,s\)。

在整环 \(R/(p)\) 上，有 \(\bar g\bar h=\bar a_nx^n\)。设 \(\bar g,\bar h\) 中最低非零项的指数分别为 \(u,v\)。最低项系数的积非零，故乘积的最低指数为 \(u+v=n\)。另一方面 \(u\le r,v\le s,r+s=n\)，因此只能有 \(u=r,v=s\)：每个约化因子的最低非零项也是其最高次项，故 \(\bar g,\bar h\) 都是单项式。两者的次数均为正数，所以 \(g,h\) 的常数项都被 \(p\) 整除。这使 \(a_0=g(0)h(0)\) 被 \(p^2\) 整除，与假设矛盾。
\end{proof}

这称为\term[eisenstein-panju]{Eisenstein 判据}[Eisenstein criterion]。例如 \(x^5-6x+3\) 对素数三满足该判据，故在 \(\mathbb Q[x]\) 中不可约。没有出现的系数就是零，当然也被三整除；常数项三不被九整除，是不能遗漏的条件。

\ProseParagraphBreak
有时原多项式的系数没有这种形式，平移后却有。对 \(a\in R\)，代入映射 \(T_a:R[x]\rightarrow R[x]\)、\(f(x)\mapsto f(x+a)\) 保持和、积与单位，而且 \(T_{-a}\circ T_a=T_a\circ T_{-a}=\operatorname{id}_{R[x]}\)。若 \(T_a(f)=gh\) 是非平凡分解，施以逆映射就得到
\[
f=T_{-a}(g)T_{-a}(h).
\]
自同构及其逆映射都保持单位与非单位，所以右边仍是非平凡分解。反过来，对 \(f\) 的分解施以 \(T_a\) 也得到 \(T_a(f)\) 的分解。因此平移保持不可约性，所用的是多项式环内部的可逆代换。

\ProseParagraphBreak
以素数 \(p\) 为例，考虑
\[
\Phi_p(x)=1+x+\cdots+x^{p-1}.
\]
记号 \(\Phi_p\) 的一般推广将在后文介绍；这里仅指所列多项式。\((x+1)^p-1\) 的常数项为零，故在 \(\mathbb Z[x]\) 中已经被 \(x\) 整除。由几何级数恒等式，下面的分式记号表示这个整系数多项式商：
\[
\Phi_p(x+1)=\frac{(x+1)^p-1}{x}
=\sum_{j=1}^p\binom pj x^{j-1}.
\]
右端的有限和写出了这个商。它的首项系数为一，因而本原；常数项为 \(p\)，其余二项式系数都被 \(p\) 整除，因为 \(1\le j<p\) 时 \(j!\) 与 \((p-j)!\) 均不含素因子 \(p\)，而 \(p!\) 含有这个因子。故 \(\Phi_p(x+1)\) 对 \(p\) 满足 Eisenstein 判据，\(\Phi_p(x)\) 也在 \(\mathbb Q[x]\) 中不可约。

\subsection{模素数约化法}
\label{b1:subsec:1203:02}

把系数约化到有限域，可以把有理数域中的问题转成有限次的尝试。这个办法需要保住原多项式的次数。

\begin{proposition}\label{b1:prop:reduction-irreducible}
设 \(f\in\mathbb Z[x]\) 本原，次数为正，\(p\) 是不整除其首项系数的素数。若逐系数约化所得的 \(\bar f\in\mathbb F_p[x]\) 不可约，则 \(f\) 在 \(\mathbb Q[x]\) 及 \(\mathbb Z[x]\) 中都不可约。
\end{proposition}
\begin{proof}
若 \(f\) 在 \(\mathbb Q[x]\) 中可约，\cref{b1:thm:gauss-descent}给出 \(f=gh\)，其中 \(g,h\in\mathbb Z[x]\) 的次数均为正。由于首项系数的积不被 \(p\) 整除，每个因子的首项系数也不被 \(p\) 整除，所以 \(\deg\bar g=\deg g>0\)、\(\deg\bar h=\deg h>0\)。约化后 \(\bar f=\bar g\bar h\) 就是非平凡分解，与假设矛盾。再由同一个 Gauss 分解定理，得到 \(\mathbb Z[x]\) 中的结论。
\end{proof}

上述方法称为\term[mo-sushu-yuehua-fa]{模素数约化法}[reduction modulo a prime]。例如 \(f=x^3+x+1\) 模二后，在 \(\bar0,\bar1\) 处的值都为 \(\bar1\)，所以没有根。域上的二次或三次多项式若可约，两个正次数因子的次数之和至多为三，其中必有一次因子；由\cref{b1:prop:factor-root}，这等价于存在域内的根。因此 \(\bar f\) 不可约，所给整数多项式在 \(\mathbb Q[x]\) 中不可约。

\ProseParagraphBreak
若约化后可约，只能说这次检验没有得出结论。例如 \(x^2+1\) 在 \(\mathbb Q[x]\) 中不可约，因为有理数的平方不可能等于负一；但它模二后等于 \((x+1)^2\)。首项系数条件也不能删去：
\[
2x^2+3x+1=(2x+1)(x+1)
\]
模二后却成为一次不可约多项式 \(x+1\)。失去的因子被约化成了常数，原来的次数分配便无法恢复。

\ProseParagraphBreak
对本原的正次数整系数多项式 \(f\)，记其首项系数为 \(a_n\)。选择素数 \(p\) 后，三种情况应分别处理：
\begin{center}
\begin{tabularx}{\linewidth}{@{}p{0.32\linewidth}X@{}}
\toprule
约化时的情况 & 能够得到的结论 \\
\midrule
\(p\nmid a_n\)，\(\bar f\) 不可约 & \(f\) 在 \(\mathbb Q[x]\) 及 \(\mathbb Z[x]\) 中不可约。 \\
\(p\nmid a_n\)，\(\bar f\) 可约 & 这个素数没有判定 \(f\) 的不可约性，仍可尝试其他素数或其他判据。 \\
\(p\mid a_n\) & 约化后次数下降，不能套用此判据；即使 \(\bar f\) 不可约，也不能据此推出 \(f\) 不可约。 \\
\bottomrule
\end{tabularx}
\end{center}

\subsection{有理根检验的适用范围}
\label{b1:subsec:1203:03}

\begin{proposition}[有理根检验]\label{b1:prop:rational-root}
设 \(f=a_nx^n+\cdots+a_0\in\mathbb Z[x]\)，\(a_n\ne0\)。若非零有理数 \(r/s\) 是 \(f\) 的根，其中 \(r,s\in\mathbb Z\)、\(s>0\)、\(\gcd(r,s)=1\)，则 \(r\mid a_0\)、\(s\mid a_n\)。特别地，首一整系数多项式的每个有理根都是整数。
\end{proposition}
\begin{proof}
把 \(f(r/s)=0\) 乘以 \(s^n\)，得到
\[
a_nr^n+a_{n-1}r^{n-1}s+\cdots+a_1rs^{n-1}+a_0s^n=0.
\]
除最后一项以外，每项都被 \(r\) 整除，因此 \(r\mid a_0s^n\)。由于 \(r\) 与 \(s^n\) 互素，由整数的素因子分解可消去 \(s^n\)，得到 \(r\mid a_0\)。同样，除第一项以外，每项都被 \(s\) 整除，故 \(s\mid a_nr^n\)；互素性给出 \(s\mid a_n\)。若 \(a_n=1\)，正分母只能为一。
\end{proof}

这称为\term[youli-gen-jianyan]{有理根检验}[rational root test]。零是否为根，直接查看 \(a_0\) 是否为零即可。当 \(a_0\ne0\) 时，所有有理根都在下面的候选集合中，其中 \(r,s\) 为整数：
\[
\mathcal C_f=
\left\{\,\frac{r}{s}\ \middle|\quad
r\mid a_0,\quad s\mid a_n,\quad
\gcd(r,s)=1,\quad s>0\,\right\}.
\]
分别列出 \(a_0\) 的正负因子与 \(a_n\) 的正因子，删去不互素的分子分母组合，再逐一代入 \(f\)。由于 \(a_0,a_n\) 都非零，这个集合是有限的。

\ProseParagraphBreak
若 \(a_0=0\)，先提出最低非零项所对应的 \(x\) 的幂，写成 \(f=x^kg\)，其中 \(k\ge1\)、\(g\in\mathbb Z[x]\)、\(g(0)\ne0\)。零是 \(f\) 的根，其他有理根恰是 \(g\) 的有理根。若 \(g\) 为常数，便没有其他根；否则对 \(g\) 使用上述有限候选集合即可。

\ProseParagraphBreak
有理根检验何时足以判定不可约性，取决于多项式的次数。对任意域 \(K\)，有
\[
\boxed{
\begin{gathered}
f\in K[x],\qquad \deg f\in\{2,3\}:\\
f\text{ 可约}\quad\Longleftrightarrow\quad
f\text{ 在 }K\text{ 中有根}.
\end{gathered}}
\]
若这种次数的多项式可约，两个正次数因子的次数之和为二或三，必有一个一次因子；域上一次因子给出域内的根。反过来，若 \(a\in K\) 是根，由\cref{b1:prop:factor-root}可提出 \(x-a\)，剩下的因子仍有正次数，因而得到非平凡分解。

\ProseParagraphBreak
例如 \(x^3-2\) 的候选有理根只有 \(\pm1,\pm2\)，代入均不为零；次数为三，故它在 \(\mathbb Q[x]\) 中不可约。这里也可以直接使用素数二的 Eisenstein 判据。

\ProseParagraphBreak
次数达到四以后，没有有理根只能排除一次因子，不能排除两个二次因子的乘积。例如
\[
x^4+x^2+1=(x^2+x+1)(x^2-x+1).
\]
左边对任何实数都严格为正，因而没有有理根，却在 \(\mathbb Q[x]\) 中可约。选择判据时，先问它排除了哪些可能的因子次数，往往比机械地尝试更多代入值更有效。

% !TeX root = ../../Book1.tex
% 纲要标题：### 12.4 多元多项式与有限性
% 本轮补充的纲要细目：
% * 在任意交换系数环上完整证明对称表达的存在唯一性，明确有界次数下消项终止；用幂和及根乘积说明计算
% * 区分有限域上的形式多项式与点集函数；正文与习题均不靠除以对称群阶证明基本定理
% #### 12.4.1 单项式、次数与消去
% #### 12.4.2 变量置换与不变多项式
% #### 12.4.3 初等对称多项式与基本定理
% #### 12.4.4 Noether 环与升链条件
% #### 12.4.5 Hilbert 基定理
% 纲要位置：计划纲要.md，第 462 行。

\section{多元多项式与有限性}
\label{b1:sec:1204}

% 纲要内容（仅作写作计划，尚非教材正文）：
%
% * 单项式、多项式次数及消去的动机
% * 变量置换、不变多项式与对称多项式基本定理
% * Noether 环与升链条件
% * Hilbert 基定理；第三卷只复用证明并发展模论版本
%
% ---
%

一元多项式的带余除法依靠次数，把一个除法过程不断缩短。多元多项式仍有次数，却不能仅凭全次数把任意多项式除以另一个多项式。例如 \(x\) 和 \(y\) 的全次数都是一，全次数没有规定哪一项应当先处理，而且两项也不能通过乘上单项式互相消去。因此，多元消项还需要规定单项式的比较次序，并检查相应的整除关系。本节从方程的消去转向变量置换下的不变性，研究对称多项式能否用少数几个固定的多项式表示；随后另行讨论理想的有限生成性，证明 Hilbert 基定理。前一个有限生成问题关心特定的子环，后一个则关心整个环的每个理想。

\subsection{单项式、次数与消去}
\label{b1:subsec:1204:01}

设 \(R\) 为交换环，\(n\geq1\) 为整数。对 \(\alpha=(\alpha_1,\ldots,\alpha_n)\in\mathbb Z_{\geq0}^n\)，称 \(x_1^{\alpha_1}\cdots x_n^{\alpha_n}\) 为变量 \(x_1,\ldots,x_n\) 的一个\term[danxiangshi]{单项式}[monomial]，并记
\[
x^\alpha=x_1^{\alpha_1}\cdots x_n^{\alpha_n},\qquad
|\alpha|=\alpha_1+\cdots+\alpha_n
\]
若各系数 \(a_\alpha\in R\) 且只有有限多个非零，则称形式和 \(f=\sum_\alpha a_\alpha x^\alpha\) 为 \(R\) 上的 \(n\) 元多项式。按系数相加，并将乘法由 \(x^\alpha x^\beta=x^{\alpha+\beta}\) 分配展开，所得环记为 \(R[x_1,\ldots,x_n]\)。也可以逐次添加不定元，把它看作 \(\bigl(R[x_1,\ldots,x_{n-1}]\bigr)[x_n]\)；两种描述等价，均以只有有限多个非零项的系数族表示多项式。
\symidx{\(x^\alpha\)：多重指数对应的单项式}

\begin{definition}\label{b1:def:multivariate-degree}
设 \(R\) 为交换环，\(n\geq1\) 为整数，并设
\[
f=\sum_{\alpha\in\mathbb Z_{\geq0}^n}a_\alpha x^\alpha
\in R[x_1,\ldots,x_n]
\]
为非零多项式，其中 \(x^\alpha=x_1^{\alpha_1}\cdots x_n^{\alpha_n}\)，\(|\alpha|=\alpha_1+\cdots+\alpha_n\)，且只有有限多个 \(a_\alpha\) 非零。令
\[
\deg f=\max\bigl\{\,|\alpha|\,\bigm|\,a_\alpha\ne0\,\bigr\}
\]
称整数 \(\deg f\) 为多项式 \(f\) 的\term[quan-cishu]{全次数}[total degree]。
对 \(1\leq i\leq n\)，关于变量 \(x_i\) 的次数是非零项中 \(\alpha_i\) 的最大值。如果某个非负整数 \(d\) 满足所有非零项的全次数都等于 \(d\)，就称 \(f\) 为 \(d\) 次\term[qici-duoxiangshi]{齐次多项式}[homogeneous polynomial]。
\end{definition}
\symidx{\(\lvert\alpha\rvert\)：多重指标的全次数}

例如 \(x^2y+xy^3+1\) 的全次数是四，关于 \(x\) 的次数是二，关于 \(y\) 的次数是三。每个多项式可以唯一拆成不同全次数的齐次部分。若 \(R\) 是整环，反复使用\cref{b1:prop:polynomial-degree}可知多元多项式环也是整环；若 \(f_d,g_e\) 是 \(f,g\) 的最高次齐次部分，则 \(f_dg_e\ne0\)，而它正是 \(fg\) 的 \(d+e\) 次齐次部分。因此非零乘积满足 \(\deg(fg)=\deg f+\deg g\)。

\ProseParagraphBreak
多元多项式常同时描述多个方程。以 \(\mathbb Q[x,y]\) 中的
\[
y-x^2=0,\qquad y-1=0
\]
为例，两式相减便得到 \(x^2-1=0\)，其中已经不含 \(y\)。将方程左端所生成的理想记为 \(I\)，有
\[
I=(y-x^2,y-1)=(y-1,x^2-1),\qquad
I\cap\mathbb Q[x]=(x^2-1).
\]
最后一个等式的反向包含需要说明：若 \(g(x)=A(x,y)(y-1)+B(x,y)(x^2-1)\)，代入 \(y=1\)，便得到 \(g(x)=B(x,1)(x^2-1)\)，所以 \(g\) 被 \(x^2-1\) 整除。交 \(I\cap\mathbb Q[x]\) 收集了理想中只含 \(x\) 的关系，这是\term[xiaoqu]{消去}[elimination]的一个初步例子。此例中，消去方程的两个根 \(x=\pm1\) 都能取 \(y=1\)，得到原方程组的解。

\ProseParagraphBreak
一般地，设 \(K\) 为域，\(I=(f_1,\ldots,f_m)\subseteq K[x,y]\) 由方程左端生成。\(g(x)\in I\cap K[x]\) 意味着 \(g\) 可以写成这些 \(f_j\) 的多项式系数组合，因此每个原方程组解的第一坐标都必须使 \(g\) 为零：
\[
\bigl[f_j(a,b)=0\ (1\le j\le m)\bigr]
\quad\Longrightarrow\quad
g(a)=0\quad(g\in I\cap K[x]).
\]
消去理想收集的是由原方程左端作多项式系数线性组合所得、且只含 \(x\) 的那些关系。但即使 \(a\in K\) 满足全部消去方程，反向提升仍可能失败：
\[
\bigl[g(a)=0\ (\forall g\in I\cap K[x])\bigr]
\quad\not\Longrightarrow\quad
\exists b\in K:\ f_j(a,b)=0\ (1\le j\le m).
\]
例如在 \(K[x,y]\) 中，\((xy-1)\cap K[x]=0\)：非零的 \((xy-1)h\) 关于 \(y\) 的次数至少为一，所以不能只含 \(x\)。消去理想因而不限制 \(x\)，可是 \(x=0\) 满足全部消去方程，却不存在 \(b\in K\) 使 \(0\cdot b=1\)；只有 \(x\ne0\) 时才能取 \(y=x^{-1}\)。\cref{b1:fig:elimination-real-points}只画实数点，直观比较这两种提升情形。\appref{b1:app:F}将给出 Gröbner 基及消去的初步计算方法，几何解释在第三卷第8章继续发展。

\begin{figure}[htbp]
\centering
\begin{tikzpicture}[x=1cm,y=1cm,>=stealth]
  \begin{scope}
    \draw[->,color=AlgebraGray] (-1.8,0)--(1.8,0) node[right] {\(x\)};
    \draw[->,color=AlgebraGray] (0,-.35)--(0,2.1) node[above] {\(y\)};
    \draw[thick,color=AlgebraBlue,domain=-1.4:1.4,samples=50]
      plot (\x,{\x*\x});
    \draw[thick,color=AlgebraGray] (-1.55,1)--(1.55,1);
    \node[color=AlgebraBlue] at (-1.5,1.75) {\(y=x^2\)};
    \node[color=AlgebraGray] at (1.55,1.22) {\(y=1\)};
    \draw[densely dotted,color=AlgebraGray] (-1,1)--(-1,0) (1,1)--(1,0);
    \fill[color=AlgebraBlue] (-1,1) circle (1.5pt) (1,1) circle (1.5pt);
    \fill[color=AlgebraBlue] (-1,0) circle (1.5pt) (1,0) circle (1.5pt);
    \node[below] at (-1,0) {\(-1\)};
    \node[below] at (1,0) {\(1\)};
    \node at (0,-.85) {可提升：\(x=\pm1\)};
  \end{scope}
  \begin{scope}[xshift=5.6cm]
    \draw[->,color=AlgebraGray] (-2.2,0)--(2.2,0) node[right] {\(x\)};
    \draw[->,color=AlgebraGray] (0,-2.05)--(0,2.05) node[above] {\(y\)};
    \draw[thick,color=AlgebraBlue,domain=.55:2,samples=50]
      plot (\x,{1/\x});
    \draw[thick,color=AlgebraBlue,domain=-2:-.55,samples=50]
      plot (\x,{1/\x});
    \node[color=AlgebraBlue] at (1.45,1.45) {\(xy=1\)};
    \draw[very thick,color=AlgebraBlue] (-2,0)--(-.13,0) (.13,0)--(2,0);
    \draw[draw=AlgebraBlue,fill=white,line width=.7pt] (0,0) circle (3.2pt);
    \node[below] at (0,0) {\(0\)};
    \node at (0,-2.45) {可提升：\(x\ne0\)};
  \end{scope}
\end{tikzpicture}
\caption[消去方程与坐标提升]{两个实数点例子。左：消去方程 \(x^2-1=0\) 的两个实根都能提升到交点；右：消去理想为零，故消去方程不排除 \(x=0\)，但 \(xy=1\) 的实数解只有非零的第一坐标。满足消去方程只是原方程组有解的必要条件。}
\label{b1:fig:elimination-real-points}
\end{figure}

\subsection{变量置换与不变多项式}
\label{b1:subsec:1204:invariants}

一元首一多项式若已分解为一次因子的乘积，改变各因子的次序不会改变它的系数。例如
\[
(T-x)(T-y)=T^2-(x+y)T+xy
\]
中的 \(x+y\) 和 \(xy\)，在交换 \(x,y\) 后都不变。暂且把根看作独立的不定元，就得到一个可以在多项式环内讨论的问题；此时无需先求出某个方程的根。

\begin{definition}\label{b1:def:permutation-invariant-polynomial}
设 \(R\) 为交换环，\(n\ge1\)，\(G\) 为对称群 \(S_n\) 的一个子群。对 \(\sigma\in G\) 和 \(f\in R[x_1,\ldots,x_n]\)，规定
\[
\sigma\cdot f=f\bigl(x_{\sigma(1)},\ldots,x_{\sigma(n)}\bigr).
\]
如果 \(\sigma\cdot f=f\) 对所有 \(\sigma\in G\) 都成立，就称 \(f\) 为这个作用下的\term[bubian-duoxiangshi]{不变多项式}[invariant polynomial]。全部不变多项式组成的子环记为 \(R[x_1,\ldots,x_n]^G\)，称为该作用的\term[bubian-huan]{不变环}[invariant ring]。当 \(G=S_n\) 时，称这样的 \(f\) 为\term[duichen-duoxiangshi]{对称多项式}[symmetric polynomial]。
\end{definition}
\symidx{\(R[x_1,\ldots,x_n]^G\)：变量置换作用下的不变环}

这里采用直接置换不定元下标的约定，即 \(\sigma\cdot x_i=x_{\sigma(i)}\)；置换乘积 \(\sigma\tau\) 表示先作用 \(\tau\) 再作用 \(\sigma\)。上述代入确实给出左群作用。它固定系数，并保持加法与乘法；先作用 \(\tau\) 再作用 \(\sigma\)，会把 \(x_i\) 送到 \(x_{\sigma(\tau(i))}\)，与 \(\sigma\tau\) 的作用一致。恒等置换给出恒等映射，逆置换给出逆映射，因此各个作用映射都是环自同构。如果 \(f,g\) 不变，则 \(f-g\) 与 \(fg\) 也不变，常数同样不变，故定义中的集合确为包含 \(R\) 的子环。

\ProseParagraphBreak
检验不变性时，可以比较各单项式的系数。记
\[
\sigma\alpha=
\bigl(\alpha_{\sigma^{-1}(1)},\ldots,\alpha_{\sigma^{-1}(n)}\bigr),
\]
则 \(\sigma\cdot x^\alpha=x^{\sigma\alpha}\)。因此，对 \(f=\sum_\alpha c_\alpha x^\alpha\)，有
\[
f\in R[x_1,\ldots,x_n]^G
\quad\Longleftrightarrow\quad
\boxed{c_{\sigma\alpha}=c_\alpha\quad(\sigma\in G,\ \alpha\in\mathbb Z_{\ge0}^n)}.
\]
也就是说，每个单项式轨道中的系数必须相同。这里指数元组的重排出现 \(\sigma^{-1}\)，只是为了记录代入后每个变量所得的指数；不定元本身仍按上述约定直接置换。例如三元多项式
\[
x^2y+y^2z+z^2x
\]
在循环置换 \(x\mapsto y\mapsto z\mapsto x\) 下不变，但交换 \(x,y\) 后变为 \(xy^2+yz^2+zx^2\)。这两个多项式含有不同的单项式，故在任何非零系数环上都不相等。它是循环子群的不变量，却不是对称多项式。

\ProseParagraphBreak
这里仍须注意\hyperref[b1:note:polynomial-function-distinction]{形式多项式与函数的区别}。在 \(\mathbb F_2[x,y]\) 中，\(x^2-x\) 给出零函数，交换两个坐标当然不改变这个函数；可是交换变量得到的形式多项式是 \(y^2-y\)，与原多项式不同。因而它并不属于 \(\mathbb F_2[x,y]^{S_2}\)。上面的轨道系数判据检验的是形式多项式的不变性。

\subsection{初等对称多项式与基本定理}
\label{b1:subsec:1204:symmetric-polynomials}

\begin{definition}\label{b1:def:elementary-symmetric-polynomials}
设 \(R\) 为交换环，\(x_1,\ldots,x_n\) 为 \(n\ge1\) 个不定元。对 \(1\le j\le n\)，令
\[
e_j=\sum_{1\le i_1<\cdots<i_j\le n}x_{i_1}\cdots x_{i_j}.
\]
称 \(e_j\) 为这 \(n\) 个变量的第 \(j\) 个\term[chudeng-duichen-duoxiangshi]{初等对称多项式}[elementary symmetric polynomial]。另约定 \(e_0=1\)。
\end{definition}
\symidx{\(e_j\)：第 \(j\) 个初等对称多项式}

例如 \(e_1=x_1+\cdots+x_n\)，\(e_n=x_1\cdots x_n\)。变量置换只会重排求和中的各个 \(j\) 元子集，故每个 \(e_j\) 都对称。展开一次因子的乘积得到
\begin{equation}\label{b1:eq:elementary-symmetric-product}
\prod_{i=1}^n(T-x_i)
=T^n-e_1T^{n-1}+e_2T^{n-2}-\cdots+(-1)^ne_n.
\end{equation}
事实上，\(T^{n-j}\) 的系数来自选择 \(j\) 个因子中的 \(-x_i\)、其余因子中的 \(T\)，因而恰为 \((-1)^je_j\)。这就是根与系数关系的多项式恒等式。

为了逐步消去一个确定的单项式，需要先规定多元单项式的大小。我们所选的次序将使对称性限制最大项的指数排列，并使初等对称多项式的乘积能够精确匹配这个最大项。下面使用按 \(x_1\succ\cdots\succ x_n\) 排列的\term[zidian-xu]{字典序}[lexicographic order]：对不同的指数元组 \(\alpha,\beta\in\mathbb Z_{\ge0}^n\)，从第一项起寻找第一个不同的位置 \(i\)；若 \(\alpha_i>\beta_i\)，就规定 \(x^\alpha\succ x^\beta\)。任一非零多项式只有有限项，故有唯一的最大单项式，称它为\term[shou-danxiangshi]{首单项式}[leading monomial]。这个比较与乘法相容：给两组指数都加上同一个元组，第一个不同的位置及大小关系不变。

\begin{theorem}[对称多项式基本定理]\label{b1:thm:fundamental-symmetric-polynomials}
设 \(R\) 为交换环，\(n\ge1\)，\(e_1,\ldots,e_n\in R[x_1,\ldots,x_n]\) 为初等对称多项式。对每个对称多项式 \(f\)，存在唯一的多项式 \(F\in R[t_1,\ldots,t_n]\)，使
\[
f=F(e_1,\ldots,e_n).
\]
换言之，固定系数的代入同态
\[
R[t_1,\ldots,t_n]\longrightarrow R[x_1,\ldots,x_n]^{S_n},
\qquad F\longmapsto F(e_1,\ldots,e_n)
\]
是环同构。
\end{theorem}
\begin{proof}
设非零对称多项式 \(f\) 的首项为 \(cx_1^{\alpha_1}\cdots x_n^{\alpha_n}\)，其中 \(c\ne0\)。必有
\[
\alpha_1\ge\alpha_2\ge\cdots\ge\alpha_n.
\]
否则，若 \(\alpha_i<\alpha_{i+1}\)，交换 \(x_i,x_{i+1}\) 后，新的指数元组在前 \(i-1\) 个位置与原元组相同，第 \(i\) 个位置却由 \(\alpha_i\) 变为较大的 \(\alpha_{i+1}\)。因此新单项式在字典序下严格更大；对称性保证它的系数仍为 \(c\ne0\)，与首单项式的选择矛盾。

\(e_j\) 的首单项式是 \(x_1\cdots x_j\)，系数为一。由比较与乘法的相容性，任意乘积 \(e_1^{b_1}\cdots e_n^{b_n}\) 的首单项式为
\begin{equation}\label{b1:eq:symmetric-leading-monomial}
x_1^{b_1+\cdots+b_n}x_2^{b_2+\cdots+b_n}\cdots x_n^{b_n},
\end{equation}
系数仍为一：只有在每个因子中都选首单项式才能得到这个最大项，其他选择所得的项都更小。例如在三变量情形，要匹配 \(x_1^5x_2^3x_3\)，可以取
\[
b_1=5-3=2,\qquad b_2=3-1=2,\qquad b_3=1.
\]
相应乘积 \(e_1^2e_2^2e_3\) 的首单项式正是
\[
x_1^2(x_1x_2)^2(x_1x_2x_3)=x_1^5x_2^3x_3.
\]
一般地，取
\[
b_i=\alpha_i-\alpha_{i+1}\ (1\le i<n),
\qquad b_n=\alpha_n,
\]
便使 \(e_1^{b_1}\cdots e_n^{b_n}\) 的首单项式恰为 \(x^\alpha\)。从 \(f\) 减去 \(ce_1^{b_1}\cdots e_n^{b_n}\)，首项被消掉，余下的多项式仍然对称，且为零或具有更小的首单项式。

还需证明反复消项一定终止。记原多项式的全次数为 \(d\)。刚才所减的乘积齐次，次数为
\[
\sum_{j=1}^n jb_j=\alpha_1+\cdots+\alpha_n\le d.
\]
因而消项过程中所有多项式的全次数始终不超过 \(d\)。满足 \(\alpha_1+\cdots+\alpha_n\le d\) 的非负整数元组只有有限多个，例如每个坐标都在 \(0,\ldots,d\) 之间。首单项式每步严格变小，不可能在这个有限集合中无限进行，所以经过有限步必得到零。这里有界全次数已经把可能的首项限制在有限集合中，无须先建立一般单项式序的良序理论。把各次减去的初等对称多项式乘积相加，便得到所需的 \(F\)。若 \(f=0\)，可直接取 \(F=0\)。

最后证明唯一性。关键在于，不同的 \(t\)-单项式代入后具有不同的首单项式。具体说，\(t_1^{b_1}\cdots t_n^{b_n}\) 的像以\eqref{b1:eq:symmetric-leading-monomial}为首单项式，而从这个指数元组的相邻差及最后一个指数可以恢复全部 \(b_i\)。因此若非零多项式 \(H\in R[t_1,\ldots,t_n]\) 在代入后成为零，将它写成不同单项式的有限线性组合，并在各像的首单项式中选取最大的一个。这个最大项只由一个 \(H\) 的单项式贡献，其系数就是 \(H\) 中相应的非零系数；其余像的全部单项式都比它小，故不能将它消去。这一步没有使用 \(R\) 为整环：相应乘积的首项系数为一，所保留的系数只是原非零系数乘以一。这与 \(H(e_1,\ldots,e_n)=0\) 矛盾。因此代入同态的核为零，表达式唯一。
\end{proof}

上述证明只作加减和乘法，不需要除以 \(n!\) 或任何整数，所以对整数系数、正特征系数以及有零因子的交换环都成立。它也给出了实际计算的步骤：先固定变量次序，依次消去对称多项式的首项，并记录对应的 \(e_j\) 乘积。这里的唯一性指作为 \(t_1,\ldots,t_n\) 的多项式的表达式唯一；任意加括号或恒等变形当然可以写出不同外观的同一个表达式。

\begin{example}\label{b1:exm:symmetric-cubic-expression}
在 \(\mathbb Z[x,y,z]\) 中，令 \(e_1=x+y+z\)、\(e_2=xy+xz+yz\)、\(e_3=xyz\)，并求 \(f=x^3+y^3+z^3\) 的对称表达式。按 \(x\succ y\succ z\) 的字典序，\(f\) 的首单项式为 \(x^3\)，先减去 \(e_1^3\)，得
\[
f-e_1^3=-3(x^2y+x^2z+xy^2+y^2z+xz^2+yz^2)-6xyz.
\]
余式的首单项式为 \(x^2y\)，系数为 \(-3\)。而
\[
e_1e_2=x^2y+x^2z+xy^2+y^2z+xz^2+yz^2+3xyz,
\]
所以 \(f-e_1^3+3e_1e_2=3e_3\)。这给出
\begin{equation}\label{b1:eq:power-sum-three}
x^3+y^3+z^3=e_1^3-3e_1e_2+3e_3.
\end{equation}
它是整数系数恒等式，因而可以代入任意交换环中的三个元素。特别地，在特征为三的环上，右端后两项消失；计算时不曾除以三，所以这个特例仍有效。
\end{example}

\begin{example}\label{b1:exm:symmetric-elimination}
仍在 \(\mathbb Z[x,y,z]\) 中记 \(e_1=x+y+z\)、\(e_2=xy+xz+yz\)、\(e_3=xyz\)。若要找出三个两两乘积 \(xy,yz,zx\) 共同满足的方程，可以先写
\[
\begin{aligned}
(T-xy)(T-yz)(T-zx)
&=T^3-(xy+yz+zx)T^2\\
&\quad+xyz(x+y+z)T-x^2y^2z^2\\
&=T^3-e_2T^2+e_1e_3T-e_3^2.
\end{aligned}
\]
右端的系数都已经用初等对称多项式表示。于是，若 \(a,b,c\) 是 \(X^3-3X+1\) 在某个包含 \(\mathbb Q\) 的域中的三个根，按重数列出，则根与系数关系给出 \(e_1(a,b,c)=0\)、\(e_2(a,b,c)=-3\)、\(e_3(a,b,c)=-1\)。代入上式便得到
\[
(T-ab)(T-bc)(T-ca)=T^3+3T^2-1.
\]
因此不必分别求出 \(a,b,c\)，就能消去它们，得到两两乘积满足的方程。基本定理中的表达式唯一性是在独立不定元组成的多项式环中成立的；代入具体的根后，例如 \(e_1(a,b,c)=0\)，不同的多项式 \(t_1\) 与 \(0\) 便有相同的代入值。所得方程是否不可约、次数是否最小，仍须另作判断；最小多项式的概念将在第十四章介绍。
\end{example}

对称多项式基本定理告诉我们，不变环 \(R[x_1,\ldots,x_n]^{S_n}=R[e_1,\ldots,e_n]\) 由系数环 \(R\) 中的常数和 \(e_1,\ldots,e_n\) 生成，而且不存在非零的 \(H\in R[t_1,\ldots,t_n]\) 使 \(H(e_1,\ldots,e_n)=0\)。其他群的不变环未必也有后一性质。接下来讨论的 Noether 性则要求一个环的每个理想都由有限个元素以环系数作线性组合生成，不能仅从上述特例直接推出。

\subsection{Noether 环与升链条件}
\label{b1:subsec:1204:02}

\begin{definition}\label{b1:def:noetherian-ring}
设 \(R\) 为交换环。若它的每个理想都能由有限个元素生成，即每个理想 \(I\) 都可写成
\[
I=(a_1,\ldots,a_m)=\{\,r_1a_1+\cdots+r_ma_m\mid r_i\in R\,\},
\]
其中 \(m\geq0\) 为整数、\(a_1,\ldots,a_m\in I\)，则称 \(R\) 为\term[noether-huan]{Noether 环}[Noetherian ring]。这里生成元的数目可以依赖于 \(I\)，零理想允许由空集生成。
\end{definition}

Noether 性要求的是每个理想各有有限生成元，不能把它与底集合有限或环作为代数有有限生成元混同。例如 \(R[x_1,\ldots,x_n]\) 总是由有限个变量作为 \(R\)-代数生成，但要使其中每个理想有限生成，仍需底环的条件与后面的 Hilbert 基定理。另一方面，每个主理想整环都是 Noether 环，故 \(\mathbb Z\) 和 \(K[x]\) 都是；后者即使在有限域上也有无穷多个元素。

\ProseParagraphBreak
上一章在主理想整环中已经用理想升链稳定来证明不可约分解的存在性。下面把这种稳定性与每个理想有限生成直接联系起来，从而得到检验 Noether 性的另一种方式。

\begin{proposition}\label{b1:prop:noetherian-acc}
交换环 \(R\) 是 Noether 环，当且仅当每条理想升链
\[
I_1\subseteq I_2\subseteq I_3\subseteq\cdots
\]
都最终稳定，即存在 \(N\)，使所有 \(n\ge N\) 都有 \(I_n=I_N\)。
\end{proposition}
\begin{proof}
设每个理想有限生成。给定升链，其并集 \(I=\bigcup_{n\ge1}I_n\) 是理想：任取其中两个元素，它们同时属于足够靠后的一个 \(I_n\)，因而它们的差以及任意环元素与它们的乘积都仍在并集中。取有限生成元 \(a_1,\ldots,a_m\)，再取 \(N\) 使它们都属于 \(I_N\)。于是 \(I\subseteq I_N\)，结合相反包含得 \(I=I_N\)，升链自此稳定。

反过来，若某个理想 \(I\) 不有限生成，可以先取 \(a_1\in I\)。对已经取出的任意有限组 \(a_1,\ldots,a_n\)，都有 \(I\ne(a_1,\ldots,a_n)\)，所以总能继续取
\(a_{n+1}\in I\setminus(a_1,\ldots,a_n)\)。不有限生成正好保证这个选择过程不会停止。这给出严格升链
\((a_1)\subsetneq(a_1,a_2)\subsetneq\cdots\)，与假设矛盾。
\end{proof}

设 \(R\) 为交换环。如果 \(R\) 的每条理想升链都最终稳定，就称 \(R\) 满足理想的\term[sheng-lian-tiaojian]{升链条件}[ascending chain condition]。它排除了理想中不断出现新的、无法由此前元素生成的关系。有限生成与升链稳定是同一性质的两种表述，下面的证明将交替使用它们。

\begin{proposition}\label{b1:prop:noetherian-quotient}
设 \(R\) 为 Noether 环，\(J\) 为 \(R\) 的理想。则商环 \(R/J\) 是 Noether 环。
\end{proposition}
\begin{proof}
设 \(R\) 为 Noether 环，\(J\) 为理想，\(\pi:R\rightarrow R/J\) 为商映射。对 \(R/J\) 的任一理想 \(L\)，原像 \(\pi^{-1}(L)\) 是 \(R\) 的理想，所以有限生成，设生成元为 \(a_1,\ldots,a_m\)。\(\pi\) 的满射性给出 \(L=\pi\bigl(\pi^{-1}(L)\bigr)\)，将每个线性组合映到商环，就得到 \(L=\bigl(\pi(a_1),\ldots,\pi(a_m)\bigr)\)。故每个 \(L\) 有限生成。
\end{proof}

无穷多个不定元时，\(K[x_1,x_2,\ldots]\) 不是 Noether 环。这里每个多项式仍只含有限项、有限多个变量，但理想升链 \((x_1)\subsetneq(x_1,x_2)\subsetneq\cdots\) 永不稳定：若 \(x_{n+1}\in(x_1,\ldots,x_n)\)，将前 \(n\) 个变量代为零，会在 \(K[x_{n+1},x_{n+2},\ldots]\) 中得到 \(x_{n+1}=0\)，矛盾。因而下一个定理中的“不定元个数有限”不是形式上的限制。

\subsection{Hilbert 基定理}
\label{b1:subsec:1204:03}

证明将按次数记录理想中多项式的系数。先看 \(I=(2x,x^2)\subseteq\mathbb Z[x]\)：令 \(J_n\) 为 \(I\) 中次数至多为 \(n\) 的多项式的 \(x^n\) 系数所成的集合，则
\[
J_0=0,\qquad J_1=2\mathbb Z,\qquad J_n=\mathbb Z\quad(n\ge2).
\]
第 \(1\) 层可由 \(2x\) 代表，第 \(2\) 层可由 \(x^2\) 代表。虽然从第 \(2\) 层起系数理想已稳定，只用 \(x^2\) 仍不能生成 \(2x\)；低次层的信息也须保留。一般证明将为稳定前的各层各选有限个代表多项式，再用它们逐次消去首项。

\begin{theorem}[Hilbert 基定理]\label{b1:thm:hilbert-basis}
若交换环 \(R\) 是 Noether 环，则 \(R[x]\) 也是 Noether 环。因此，对每个整数 \(n\geq1\)，\(R[x_1,\ldots,x_n]\) 及其任一商环都是 Noether 环。
\end{theorem}
\begin{proof}
给定理想 \(I\subseteq R[x]\)。若 \(I=0\)，它已经有限生成，故设 \(I\ne0\)。对每个整数 \(n\ge0\)，令 \(J_n\subseteq R\) 为 \(I\) 中次数至多为 \(n\) 的多项式的 \(x^n\) 系数所组成的集合，零多项式也包括在内。次数低于 \(n\) 的多项式贡献零。这种定义把每个 \(J_n\) 放在固定次数层上，便于用乘 \(x\) 比较相邻两层。

集合 \(J_n\) 是理想。两个所取多项式的差仍在 \(I\) 中且次数至多为 \(n\)，其 \(x^n\) 系数是相应系数之差；乘以 \(R\) 中任一元素后也如此。若 \(a\in J_n\)，将给出这个系数的多项式乘以 \(x\)，便有 \(a\in J_{n+1}\)。故 \(J_0\subseteq J_1\subseteq\cdots\) 是升链。由\cref{b1:prop:noetherian-acc}，存在 \(N\) 使 \(J_n=J_N\) 对所有 \(n\ge N\) 成立。

对每个 \(0\le n\le N\)，选取 \(J_n\) 的有限个非零生成元 \(a_{n1},\ldots,a_{nr_n}\)；若 \(J_n=0\)，令 \(r_n=0\)。按 \(J_n\) 的定义，选 \(f_{nj}\in I\)，次数至多为 \(n\)，其 \(x^n\) 系数为 \(a_{nj}\)。由于该系数非零，实际上 \(\deg f_{nj}=n\)。我们断言，所有这些有限个 \(f_{nj}\) 生成 \(I\)。

对非零 \(f\in I\) 的次数 \(d\) 归纳。若 \(d\le N\)，就用第 \(d\) 层的 \(f_{dj}\) 消去首项：\(f\) 的首项系数属于 \(J_d\)，可以写成这些 \(a_{dj}\) 的线性组合。若 \(d>N\)，则 \(J_d=J_N\)，可用第 \(N\) 层的代表，再乘上 \(x^{d-N}\) 把它们的次数升到 \(d\)。将两种情形统一写成 \(n=\min(d,N)\)，则 \(f\) 的首项系数 \(a\) 属于 \(J_d=J_n\)，所以
\(a=\sum_{j=1}^{r_n}b_j a_{nj}\)，其中 \(b_j\in R\)。于是
\[
f-\sum_{j=1}^{r_n}b_jx^{d-n}f_{nj}
\]
仍在 \(I\) 中，而最高次项的系数已经消去，故所得多项式为零或次数严格小于 \(d\)。若 \(d=0\)，它必为零，给出归纳起点；若 \(d>0\)，归纳假设说明这个差属于所有 \(f_{nj}\) 生成的理想，因而 \(f\) 也属于该理想。断言得证，所以 \(R[x]\) 是 Noether 环。

反复添加一个不定元，得到有限多个不定元的结论；再用\cref{b1:prop:noetherian-quotient}，其任意商环也为 Noether 环。
\end{proof}

这称为\term[hilbert-ji-dingli]{Hilbert 基定理}[Hilbert basis theorem]。证明中的有限性来自两处：系数理想升链稳定，使我们只需保留有限多层、每层有限个代表；这些代表能够消去任意多项式的最高次项，再由次数归纳完成生成。因此证明可概括为有限控制首项与次数归纳。整个过程没有假定 \(R\) 是整环：即使它有零因子，按固定次数收集系数仍然得到理想升链，最后的消项也只用加法和乘法。这一点使定理能用于很多商环。

\ProseParagraphBreak
定理名称中的“基”不要求理想生成元像自由模的基那样使系数表达唯一。在刚才的例子中，\(2x\) 与 \(x^2\) 之间便有关系 \(x(2x)-2(x^2)=0\)。

\ProseParagraphBreak
特别地，域上有限个不定元的多项式环中，每个理想都由有限个多项式生成。因此任意一族多项式方程所生成的理想，均可由有限个多项式生成；这些生成元各自又是原方程中有限项的组合，所以甚至可以从原方程族中选出有限多个，生成同一个理想。不过，这个存在性证明没有给出从任意输入有效寻找生成元的算法，也没有证明它们组成某种适于消去的标准形式。\appref{b1:app:F}将在\cref{b1:thm:F-buchberger-algorithm}中用这里的升链条件证明算法终止；它从给定有限生成组求出适于成员判定与消去的 Gröbner 基。第三卷第8章进一步讨论计算与几何应用，并在相应模论部分发展有限性条件。这里的 Hilbert 基定理为后面的算法提供有限性依据。


\begin{chaptersummary}
因式定理把根转化为一次因子，形式导数把重根转化为最大公因式问题；插值则用次数界保证有限个取值条件的唯一性。Gauss 引理从素元约化出发，证明容度的乘法性质，使分式域上的非平凡分解下降到本原多项式的系数环，并进一步建立多项式环的唯一分解性。不可约性检验始终依赖所取的基域和多项式次数，尤其不能把“没有有理根”当作高次多项式不可约的充分条件。

\ProseParagraphBreak
交换变量而保持不变的多项式，可以唯一写成初等对称多项式的多项式；首项消去在有界全次数的有限单项式集合中终止，这一证明适用于任意交换系数环。根与系数的关系由此成为处理对称表达式的具体工具。Hilbert 基定理处理的是另一类有限性：有限多个不定元不会破坏 Noether 性，所需的底环条件是交换与 Noether 性。本章已经构造分式域，并用它比较多项式分解；下一章将研究这一构造的映射性质，再讨论只允许指定元素作分母的局部化。
\end{chaptersummary}

\BookExercises

\begin{exercise}\label{b1:ex:12-polynomial-functions}
设 \(p\) 为素数。证明 \(x^p-x\) 在 \(\mathbb F_p\) 的每个元素处取值为零，但不是零多项式。进而证明任意函数 \(\varphi:\mathbb F_p\rightarrow\mathbb F_p\) 都由唯一一个次数至多为 \(p-1\) 的多项式给出，并确定所有表示同一函数的多项式之间的关系。
\end{exercise}
\begin{solution}
若 \(a=0\)，等式 \(a^p-a=0\) 直接成立；若 \(a\ne0\)，它属于 \(p-1\) 阶乘法群 \(\mathbb F_p^\times\)。\cref{b1:thm:lagrange}说明 \(a\) 的阶整除 \(p-1\)，所以 \(a^{p-1}=1\)，仍得 \(a^p=a\)。多项式 \(x^p-x\) 的首项系数为一，故不是零多项式。

将 \(\mathbb F_p\) 中全部 \(p\) 个元素作为插值点，给定值取为 \(\varphi\) 的值，\cref{b1:thm:lagrange-interpolation}就给出次数至多为 \(p-1\) 的唯一代表。若 \(f,g\) 给出同一个函数，对 \(f-g\) 除以 \(x^p-x\)，得到 \(f-g=q(x^p-x)+r\)，其中 \(r=0\) 或 \(\deg r<p\)。前两项在每个点都为零，因此 \(r\) 在 \(p\) 个点都为零。\cref{b1:prop:factor-root}迫使 \(r=0\)，所以 \(x^p-x\mid f-g\)。反过来，该整除关系显然保证两者的全部取值相同。故它们恰好相差理想 \((x^p-x)\) 中的元素。
\end{solution}

\begin{exercise}\label{b1:ex:12-repeated-roots}
在 \(\mathbb Q[x]\) 中求
\[
f=x^4-2x^3+2x^2-2x+1
\]
与 \(f'\) 的首一最大公因式，并给出它在 \(\mathbb C\) 中各根的重数。再将系数模二约化，比较重根情况。
\end{exercise}
\begin{solution}
直接展开可验
\[
f=(x-1)^2(x^2+1),\qquad
f'=2(x-1)(2x^2-x+1).
\]
其中 \(2x^2-x+1\) 在 \(x=1\) 处的值为二，所以不能再提出 \(x-1\)。它与 \(x^2+1\) 的最大公因式也为一：两者相减适当倍数后得到 \(-x-1\)，而 \(x^2+1\) 在负一处的值为二，故不能含有这个一次因子。因此首一最大公因式为 \(x-1\)。在 \(\mathbb C[x]\) 中，\(f=(x-1)^2(x-i)(x+i)\)，三个根互异，所以一的重数为二，\(i,-i\) 的重数均为一。

模二后 \(\bar f=x^4+1=(x+1)^4\)，形式导数为零。因此在任何包含 \(\mathbb F_2\) 的域中，它只有根一，重数为四。约化不仅可能使不同的因子合并，也可能使导数中的系数全部消失，不能照搬特征零下的重数计算。
\end{solution}

\begin{exercise}\label{b1:ex:12-interpolation}
求次数至多为 \(2\)，且满足 \(f(0)=1,f(1)=0,f(2)=3\) 的唯一有理系数多项式。再确定所有满足这些取值条件的 \(\mathbb Q[x]\) 中的多项式。
\end{exercise}
\begin{solution}
在\eqref{b1:eq:lagrange-interpolation}中代入这三个点和值，得到
\[
f_0(x)=\frac{(x-1)(x-2)}2+\frac{3x(x-1)}2
=2x^2-3x+1.
\]
若 \(f\) 满足相同的条件，\(f-f_0\) 在零、一、二处均为零。由\cref{b1:prop:factor-root}，先提出 \(x\)，所得商在一、二处仍为零，因为在这两点 \(x\) 的值非零；再依次提出 \(x-1,x-2\)，得到
\[
f=f_0+x(x-1)(x-2)h,\qquad h\in\mathbb Q[x].
\]
反向代入即可验证每个这样的多项式都符合条件。若还要求次数至多为二，则只能 \(h=0\)，这也直接验证了唯一性。
\end{solution}

\begin{exercise}\label{b1:ex:12-content-factorization}
沿用正文对 \(6x^2+9x+3\) 的容度分解，分别在 \(\mathbb Z[x]\) 与 \(\mathbb Q[x]\) 中验证所列因子确为不可约因子，并比较两个环中的单位部分与常数因子。
\end{exercise}
\begin{solution}
容度为三，本原部分为 \(2x^2+3x+1\)，而
\[
6x^2+9x+3=3(2x+1)(x+1).
\]
在 \(\mathbb Z[x]\) 中，三为素元：逐系数模三约化的核为 \((3)\)，像为整环 \(\mathbb F_3[x]\)，所以若三整除某个乘积，约化后必有一个因子为零，即三整除这个因子的全部系数。两个一次多项式均本原，在 \(\mathbb Q[x]\) 中因次数为一而不可约，再由\cref{b1:thm:gauss-descent}，在 \(\mathbb Z[x]\) 中也不可约。故所列三因子就是不可约分解，单位只能是 \(\pm1\)。

在 \(\mathbb Q[x]\) 中，所有非零有理常数都是单位，可以写成
\(6(x+\tfrac12)(x+1)\)。这里只有两个正次数不可约因子，常数六属于单位部分。唯一分解的比较必须在同一个环中进行；扩大系数环以后，原来的某些非单位会变为单位。
\end{solution}

\begin{exercise}\label{b1:ex:12-monic-descent}
设首一多项式 \(f\in\mathbb Z[x]\) 在 \(\mathbb Q[x]\) 中写成 \(f=gh\)，其中 \(g,h\) 也首一。证明 \(g,h\in\mathbb Z[x]\)。再判断以下两种变式能否保证因子仍有整数系数：只要求 \(f\) 本原，并允许写成 \(f=cgh\)，其中 \(c\in\mathbb Q^\times\)、\(g,h\) 首一；或者保持 \(f=gh\) 与 \(f\) 首一，却取消两个因子的首一要求。
\end{exercise}
\begin{solution}
首一的 \(f\) 本原。由\cref{b1:lem:primitive-scaling}，可以写
\[
g=\alpha g_0,\qquad h=\beta h_0,
\]
其中 \(g_0,h_0\in\mathbb Z[x]\) 本原，\(\alpha,\beta\in\mathbb Q^\times\)。\cref{b1:lem:gauss-primitive}说明 \(g_0h_0\) 本原；而 \(f=\alpha\beta g_0h_0\) 也本原，所以\cref{b1:lem:primitive-scaling}给出 \(\alpha\beta=\pm1\)。设 \(g_0,h_0\) 的首项系数为非零整数 \(a,b\)。比较 \(f\) 的首项系数，得 \(\alpha\beta ab=1\)，因而 \(ab=\pm1\)，所以 \(a,b\) 都等于 \(\pm1\)。再分别比较首一的 \(g,h\) 的首项系数，得 \(\alpha=1/a\)、\(\beta=1/b\)，两者也是整数单位。因此 \(g,h\) 的全部系数都是整数。

第一种变式不能保证结论。例如
\[
2x^2+3x+1=2\Bigl(x+\frac12\Bigr)(x+1).
\]
如果保持 \(f\) 首一，却取消两个因子的首一要求，结论同样失败：\(x^2=(2x)(x/2)\)。因子的非零常数倍可以互相转移，首一条件恰好排除了这种任意缩放。
\end{solution}

\begin{exercise}\label{b1:ex:12-shift-eisenstein}
证明 \(x^4+1\) 在 \(\mathbb Q[x]\) 中不可约，并说明对原多项式直接使用 Eisenstein 判据为何行不通。再证明它对每个素数 \(p\) 的约化在 \(\mathbb F_p[x]\) 中都可约。
\end{exercise}
\begin{solution}
对 \(f=x^4+1\)，常数项为一，不被任何素数整除，所以原式不符合 Eisenstein 判据。作可逆平移得
\[
f(x+1)=x^4+4x^3+6x^2+4x+2.
\]
首项系数为一，其余系数都被二整除，常数项二不被四整除。由\cref{b1:thm:eisenstein}，\(f(x+1)\) 在 \(\mathbb Q[x]\) 中不可约。若 \(f\) 有非平凡分解，逐因子代入 \(x+1\) 会保持各因子的正次数，给出 \(f(x+1)\) 的非平凡分解，矛盾。因此 \(f\) 不可约。

模二时 \(x^4+1=(x+1)^4\)，所以约化式可约。设 \(p\) 为奇素数。非零平方在 \(\mathbb F_p^\times\) 中组成指数为二的子群：每个非零平方恰有两个平方根 \(a,-a\)。若 \(-1\) 是平方，取 \(a^2=-1\)，则
\[
x^4+1=(x^2-a)(x^2+a).
\]
若 \(-1\) 不是平方，则乘以这个非平方元会交换平方子群及其非平方陪集，故 \(2\) 与 \(-2\) 中恰有一个是平方。若 \(a^2=2\)，则
\[
x^4+1=(x^2+ax+1)(x^2-ax+1);
\]
若 \(a^2=-2\)，则
\[
x^4+1=(x^2+ax-1)(x^2-ax-1).
\]
每种情形都给出了两个二次因子。因此此例虽在 \(\mathbb Q[x]\) 中不可约，却找不到任何不可约的模素数约化；逐个尝试素数的办法在此不会成功。
\end{solution}

\begin{exercise}\label{b1:ex:12-quartic-mod-two}
用模二约化证明 \(x^4+x+1\) 在 \(\mathbb Q[x]\) 中不可约。论证中除去一次因子的可能性后，还需要检查什么？
\end{exercise}
\begin{solution}
约化多项式在 \(\bar0,\bar1\) 处均取 \(\bar1\)，所以没有一次因子。若它可约，按不可约因子的次数分拆四，只能是两个不可约二次因子的乘积。\(\mathbb F_2[x]\) 中的首一不可约二次多项式只有 \(x^2+x+1\)：常数项必须为一，以免零为根；对 \(x^2+ax+1\) 代入一得到 \(a\)，所以还必须 \(a=1\)。但是
\[
(x^2+x+1)^2=x^4+x^2+1\ne x^4+x+1
\]
于 \(\mathbb F_2[x]\)。因此二次乘二次的可能性也排除了，约化式不可约。原多项式本原且首项系数为一，\cref{b1:prop:reduction-irreducible}于是给出 \(\mathbb Q[x]\) 中的不可约性。
\end{solution}

\begin{exercise}\label{b1:ex:12-rational-root-limits}
用有理根检验分解 \(6x^3+5x^2-12x+4\) 为 \(\mathbb Q[x]\) 中的不可约因子。再证明 \(x^4+4\) 没有有理根，却在 \(\mathbb Q[x]\) 中可约，并写出其不可约分解。
\end{exercise}
\begin{solution}
由\cref{b1:prop:rational-root}，第一个多项式的非零有理根约成 \(r/s\)、\(s>0\) 后，必须满足 \(r\mid4\)、\(s\mid6\)。从这些候选中代入 \(1/2\)，得到零。除以 \(2x-1\)，有
\[
6x^3+5x^2-12x+4=(2x-1)(3x^2+4x-4).
\]
对剩下的二次多项式，检验 \(2/3\) 得到零，继续相除得
\[
6x^3+5x^2-12x+4=(2x-1)(3x-2)(x+2).
\]
三个因子在 \(\mathbb Q[x]\) 中都是一次多项式，所以均不可约；它们给出了全部三个根 \(1/2,2/3,-2\)。

对任何有理数 \(a\)，都有 \(a^4+4>0\)，所以第二个多项式没有有理根。但将其写成平方差，可得
\[
x^4+4=(x^2+2)^2-(2x)^2
=(x^2-2x+2)(x^2+2x+2).
\]
两个二次因子分别为 \((x-1)^2+1\)、\((x+1)^2+1\)，因此都没有有理根。二次多项式如果可约就必须有一次因子，而\cref{b1:prop:factor-root}会给出一个有理根；所以这两个二次因子均不可约。原式没有有理根，只排除了一次因子，并没有排除这里的二次因子。
\end{solution}

\begin{exercise}\label{b1:ex:12-primitive-linear-polynomials}
设 \(k\) 为任意域。
\begin{enumerate}
\item 分别判断 \(yx+y+1\) 和 \(yx+y\) 在 \(k[x,y]\) 中是否不可约。说明“关于 \(x\) 的次数为一”为什么不足以直接作出判断。
\item 设 \(n\ge2\)，把 \(y\) 看作系数环 \(k[y]\) 中的元素，证明 \(x^n-y\) 在 \(k[x,y]\) 和 \(\operatorname{Frac}(k[y])[x]\) 中都不可约。
\item 若 \(\operatorname{char}k=p>0\) 且 \(n=p\)，求 \(x^p-y\) 关于 \(x\) 的形式导数，并举出一个包含 \(\operatorname{Frac}(k[y])\) 的域，使这个不可约多项式在其中有重根。
\end{enumerate}
\end{exercise}
\begin{solution}
1.\ 把两个多项式看作 \(k[y][x]\) 中的元素。\(k[y]\) 是唯一分解整环，因为它是域上一元多项式环，适用\cref{b1:thm:polynomial-ufd}。第一个多项式的两个系数 \(y,y+1\) 互素：任何公因子都整除它们之差一。因此 \(yx+y+1\) 本原。进入 \(\operatorname{Frac}\bigl(k[y]\bigr)[x]\) 后，它是一次多项式，故不可约；再由\cref{b1:thm:gauss-descent}，它在 \(k[y][x]=k[x,y]\) 中也不可约。

第二个多项式的容度可取 \(y\)，并且 \(yx+y=y(x+1)\)。这里两个因子都是非单位：若它们有多项式逆元，非零乘积的全次数相加将迫使它们的次数为零，矛盾。所以它可约。\(y\) 虽然关于 \(x\) 是常数，在系数环 \(k[y]\) 中却不是单位；这正是不能仅凭关于 \(x\) 的次数来套用域上结论的原因。

2.\ \(k[y]\) 是唯一分解整环，\(y\) 是其中的素元。把 \(x^n-y\) 看作 \(k[y][x]\) 中关于 \(x\) 的多项式：首项系数为一，因此它本原；中间各项系数为零，常数项 \(-y\) 被 \(y\) 整除但不被 \(y^2\) 整除。由\cref{b1:thm:eisenstein}，它在 \(\operatorname{Frac}(k[y])[x]\) 中不可约。再由\cref{b1:thm:gauss-descent}，它在 \(k[x,y]\) 中也不可约。

3.\ 若 \(n=p=\operatorname{char}k\)，记 \(F=\operatorname{Frac}(k[y])\)。此时 \((x^p-y)'=px^{p-1}=0\)。取
\[
L=F[t]/(t^p-y),\qquad \alpha=t+(t^p-y)\in L.
\]
已证 \(t^p-y\) 在 \(F[t]\) 中不可约，故 \(L\) 是域，且 \(\alpha^p=y\)。在 \(L[x]\) 中，特征为 \(p\)，于是
\[
x^p-y=x^p-\alpha^p=(x-\alpha)^p.
\]
所以在 \(F[x]\) 中不可约并不妨碍它在扩大系数域后出现重根。特别地，\((x^p-y)'=0\) 并不意味着 \(x^p-y\) 在原系数域 \(F\) 上就是一个 \(p\) 次幂；这里各非零项的指数符合要求，但常数项 \(-y\) 在 \(F\) 中没有 \(p\) 次根。
\end{solution}

\begin{exercise}\label{b1:ex:12-cyclic-polynomial-invariants}
设 \(k\) 为域，三阶循环群 \(C_3=\langle\sigma\rangle\) 通过 \(\sigma\cdot x=y\)、\(\sigma\cdot y=z\)、\(\sigma\cdot z=x\) 作用在 \(k[x,y,z]\) 上。
\begin{enumerate}
\item 列出全次数至多为二的全部不变多项式，并说明它们是否一定对称。
\item 求全部三次齐次不变多项式，并判断其中哪些对称。结论是否依赖 \(k\) 的特征？
\end{enumerate}
\end{exercise}
\begin{solution}
1.\ 次数至多为二的单项式在这个作用下分为四个轨道：
\[
\{1\},\qquad \{x,y,z\},\qquad
\{x^2,y^2,z^2\},\qquad \{xy,yz,zx\}.
\]
比较单项式系数，不变性恰好要求每个轨道中的系数相等。因此全部所求多项式唯一写成
\[
a+b(x+y+z)+c(x^2+y^2+z^2)+d(xy+yz+zx),
\qquad a,b,c,d\in k.
\]
每个括号内的表达式都在任意变量置换下不变，所以次数至多为二时，循环不变性确实蕴含对称性。

2.\ 三次单项式的四个循环轨道分别为 \(\{x^3,y^3,z^3\}\)、\(\{x^2y,y^2z,z^2x\}\)、\(\{xy^2,yz^2,zx^2\}\) 和 \(\{xyz\}\)。因此全部三次齐次不变多项式唯一写成
\[
\begin{aligned}
&A(x^3+y^3+z^3)+B(x^2y+y^2z+z^2x)\\
&\qquad+C(xy^2+yz^2+zx^2)+Dxyz,
\qquad A,B,C,D\in k.
\end{aligned}
\]
交换 \(x,y\) 时，中间两组轨道和互换，另外两组不变。由单项式系数的唯一性，这个多项式对称，当且仅当 \(B=C\)。特别地，取 \(B=1,C=0,A=D=0\) 就得到正文中的循环不变而不对称的例子。上述判断只比较系数，不涉及除以群阶，因此在特征三时仍然成立。
\end{solution}

\begin{exercise}\label{b1:ex:12-symmetric-fourth-power}
设 \(R\) 为交换环，在 \(R[x,y,z]\) 中记 \(e_1=x+y+z\)、\(e_2=xy+xz+yz\)、\(e_3=xyz\)。把 \(x^4+y^4+z^4\) 写成 \(e_1,e_2,e_3\) 的多项式，不作任何除法。说明结果在特征二时怎样简化，以及为什么这种简化不违背对称表达式的唯一性。
\end{exercise}
\begin{solution}
记 \(p_j=x^j+y^j+z^j\)。展开平方以及\eqref{b1:eq:power-sum-three}分别给出
\[
p_1=e_1,\qquad p_2=e_1^2-2e_2,\qquad
p_3=e_1^3-3e_1e_2+3e_3.
\]
由\eqref{b1:eq:elementary-symmetric-product}，代入 \(T=x\)、\(y\)、\(z\)，每个变量 \(u\) 都满足 \(u^3-e_1u^2+e_2u-e_3=0\)。分别乘以 \(u\) 后相加，得到
\[
\begin{aligned}
p_4&=e_1p_3-e_2p_2+e_3p_1\\
&=e_1^4-4e_1^2e_2+2e_2^2+4e_1e_3.
\end{aligned}
\]
全程只有环中的加法和乘法，所以对所给任意交换环成立。在特征二时，二和四的系数都为零，因而 \(p_4=e_1^4\)。唯一性是在固定系数环 \(R\) 中比较多项式；此时
\(t_1^4-4t_1^2t_2+2t_2^2+4t_1t_3\) 与 \(t_1^4\) 本来就是 \(R[t_1,t_2,t_3]\) 中的同一个多项式，并没有给出两个不同表达式。
\end{solution}

\begin{exercise}\label{b1:ex:12-invariants-fixed-variable}
设 \(k\) 为域，\(S_2\) 在 \(k[x,y,z]\) 上交换 \(x,y\) 而固定 \(z\)。证明每个不变多项式唯一写成 \(F(x+y,xy,z)\)，其中 \(F\in k[s,p,t]\)。进一步，令 \(S_2\times S_2\) 分别交换 \(x_1,x_2\) 和 \(y_1,y_2\)，求它在 \(k[x_1,x_2,y_1,y_2]\) 上的不变环，并证明所取生成元之间没有非零的多项式关系。
\end{exercise}
\begin{solution}
第一问把 \(z\) 放进系数环 \(R=k[z]\)。一个多项式在交换 \(x,y\) 后不变，恰好说明它是 \(R[x,y]\) 中的对称多项式。\cref{b1:thm:fundamental-symmetric-polynomials}给出唯一的 \(H\in R[s,p]\)，满足 \(f=H(x+y,xy)\)。将 \(H\) 的各系数唯一展开为 \(z\) 的多项式，就得到所需的 \(F\in k[s,p,t]\)；这两个唯一性合在一起也证明 \(F\) 唯一。

第二问记
\[
s_x=x_1+x_2,\quad p_x=x_1x_2,\qquad
s_y=y_1+y_2,\quad p_y=y_1y_2.
\]
先要求交换 \(x_1,x_2\) 后不变，应用同一定理，得到唯一表达式
\[
f=\sum_{a,b}c_{ab}(y_1,y_2)s_x^ap_x^b,
\qquad c_{ab}\in k[y_1,y_2],
\]
其中只有有限多个系数非零。再交换 \(y_1,y_2\)。由于关于 \(s_x,p_x\) 的表达式唯一，第二个交换也固定 \(f\)，当且仅当每个 \(c_{ab}\) 都对称。再次使用同一定理，每个系数唯一写成 \(s_y,p_y\) 的多项式。故不变环为
\[
k[x_1,x_2,y_1,y_2]^{S_2\times S_2}
=k[s_x,p_x,s_y,p_y].
\]
若这四个生成元满足一个非零多项式关系，先按前两个变量的幂收集，第一轮唯一性迫使每个关于 \(s_y,p_y\) 的系数为零；第二轮唯一性又迫使这些系数多项式的所有系数为零，与关系非零矛盾。因而它们之间没有非零多项式关系。
\end{solution}

\begin{exercise}\label{b1:ex:12-squared-root-elimination}
设 \(L\) 为包含 \(\mathbb Q\) 的域，并且在 \(L[X]\) 中有分解
\[
X^3-2X^2+3X-4=(X-a)(X-b)(X-c).
\]
不逐个求出 \(a,b,c\)，确定首一多项式 \((T-a^2)(T-b^2)(T-c^2)\)，并证明其系数属于 \(\mathbb Q\)。
\end{exercise}
\begin{solution}
记 \(e_1=a+b+c\)、\(e_2=ab+ac+bc\)、\(e_3=abc\)。比较给定分解的系数得 \(e_1=2\)、\(e_2=3\)、\(e_3=4\)。展开可得
\[
\begin{aligned}
a^2+b^2+c^2&=e_1^2-2e_2=-2,\\
a^2b^2+a^2c^2+b^2c^2&=e_2^2-2e_1e_3=-7,\\
a^2b^2c^2&=e_3^2=16.
\end{aligned}
\]
其中第二式的依据是 \((ab+ac+bc)^2\) 的交叉项之和为 \(2abc(a+b+c)\)。再次使用根与系数关系，得到
\[
(T-a^2)(T-b^2)(T-c^2)=T^3+2T^2-7T-16.
\]
这些系数都是 \(e_1,e_2,e_3\) 的整数系数多项式，而三个 \(e_i\) 都是有理数，故所求多项式属于 \(\mathbb Q[T]\)。这里不需要知道各个根属于哪个具体扩域，也不需要判断平方后的三个根是否互异。
\end{solution}

\begin{exercise}\label{b1:ex:12-noetherian-not-pid}
证明 \(\mathbb Q[x,y]\) 是 Noether 唯一分解整环，却不是主理想整环。再证明 \(\mathbb Q[x,y]/(xy)\) 是 Noether 环，但不是整环。
\end{exercise}
\begin{solution}
域 \(\mathbb Q\) 是 Noether 唯一分解整环，故\cref{b1:thm:polynomial-ufd}与\cref{b1:thm:hilbert-basis}分别给出 \(\mathbb Q[x,y]\) 的唯一分解性和 Noether 性。理想 \((x,y)\) 为真理想，因为其中每个多项式在 \((0,0)\) 处为零。如果 \((x,y)=(g)\)，则 \(g\mid x\) 且 \(g\mid y\)。用全次数的乘法公式，由 \(g\mid x\) 知 \(g\) 是非零常数或与 \(x\) 相伴；后者不可能整除 \(y\)，因为代入 \(x=0\) 会给出 \(y=0\)。因此 \(g\) 为非零有理常数，是单位，与理想为真矛盾。所以该环不是主理想整环。

由\cref{b1:prop:noetherian-quotient}，商环 \(\mathbb Q[x,y]/(xy)\) 为 Noether 环。\(x,y\) 的剩余类都非零：若 \(x\in(xy)\)，代入 \(y=0\) 就得 \(x=0\)，矛盾；对 \(y\) 代入 \(x=0\) 同理。但两者的积为零，故该商环不是整环。
\end{solution}

\begin{exercise}\label{b1:ex:12-elimination}
在 \(\mathbb Q[x,y,z]\) 中令 \(I=(y-x^2,z-xy,z-1)\)。求 \(I\cap\mathbb Q[x]\)，并解释它如何限制原方程组的第一坐标。
\end{exercise}
\begin{solution}
直接计算
\[
x^3-1=(z-1)-(z-xy)-x(y-x^2),
\]
说明 \(x^3-1\in I\)。反过来，
\[
z-xy=(z-1)-(x^3-1)-x(y-x^2),
\]
所以 \(I=(y-x^2,z-1,x^3-1)\)。若 \(g(x)\in I\)，将它表示成这三个生成元的多项式系数组合，再代入 \(y=x^2,z=1\)，得到 \(g(x)=C(x,x^2,1)(x^3-1)\)。故 \(g\) 是 \(x^3-1\) 的倍数。结合此前已经证明的包含，得到
\[
I\cap\mathbb Q[x]=(x^3-1).
\]
原方程组的任一解必满足 \(x^3=1\)。本题中反向提升也成立：在任意包含 \(\mathbb Q\) 的域中，若 \(a^3=1\)，则 \((a,a^2,1)\) 同时满足三个原方程，且后两个坐标由第一个坐标唯一确定。这一充分性来自本题明确的代入，不是消去理想在一般情况下自动保证的结论。
\end{solution}

\begin{exercise}\label{b1:ex:12-noetherian-surjection}
证明 Noether 交换环 \(R\) 的每个满环自同态 \(\varphi:R\rightarrow R\) 都是同构。用无穷多个不定元的多项式环，说明 Noether 条件不能删去。
\end{exercise}
\begin{solution}
记 \(\varphi^n\) 为 \(\varphi\) 的 \(n\) 次复合。各次复合的核组成理想升链
\[
\ker\varphi\subseteq\ker\varphi^2\subseteq\cdots.
\]
因为若 \(\varphi^n(a)=0\)，再作用一次 \(\varphi\) 仍为零，所以这些包含成立。\cref{b1:prop:noetherian-acc}给出某个 \(N\ge1\)，使 \(\ker\varphi^N=\ker\varphi^{N+1}\)。任取 \(a\in\ker\varphi\)，由于 \(\varphi\) 满射，\(\varphi^N\) 也满射，可以取 \(b\in R\) 使 \(\varphi^N(b)=a\)。于是
\[
\varphi^{N+1}(b)=\varphi(a)=0,
\]
故 \(b\in\ker\varphi^{N+1}=\ker\varphi^N\)，从而 \(a=0\)。这证明 \(\varphi\) 单射，结合满射性即为同构。

去掉 Noether 条件，取 \(R=k[x_1,x_2,\ldots]\)，令 \(\varphi\) 固定 \(k\)，把 \(x_1\) 送到零，把 \(x_{i+1}\) 送到 \(x_i\)。每个多项式只含有限多个变量，这些代入确定保持加法、乘法和单位的映射。任意多项式的原像可以把其中每个 \(x_i\) 换成 \(x_{i+1}\) 得到，所以映射满射；非零元素 \(x_1\) 却在其核中，故不单射。相应核升链确实不稳定：\(x_{n+1}\) 在 \(\ker\varphi^{n+1}\) 中，却不在 \(\ker\varphi^n\) 中，因为 \(\varphi^n(x_{n+1})=x_1\ne0\)。
\end{solution}

\begin{exercise}\label{b1:ex:12-nonnoetherian-subring}
设 \(k\) 为域，在 \(k[x,y]\) 中取子环
\[
A=k[x,xy,xy^2,xy^3,\ldots].
\]
证明 \(A=k+xk[x,y]\)，并证明 \(A\) 不是 Noether 环，尽管 \(k[x,y]\) 是。这里的 \(k+xk[x,y]\) 表示所有 \(c+x\cdot h(x,y)\)、\(c\in k\)、\(h\in k[x,y]\) 组成的集合。
\end{exercise}
\begin{solution}
所列每个生成元都属于 \(k+xk[x,y]\)，而后者在加法、取负和乘法下封闭，故 \(A\subseteq k+xk[x,y]\)。反过来，任意 \(x^iy^j\)、\(i\ge1,j\ge0\)，都可写成 \(x^{i-1}(xy^j)\)，属于 \(A\)。每个 \(x\cdot h(x,y)\) 是这种单项式的有限线性组合，所以另一个包含也成立。

在 \(A\) 中令 \(I_n=(x,xy,\ldots,xy^n)\)。这些理想构成升链。如果 \(xy^{n+1}\in I_n\)，则存在 \(a_0,\ldots,a_n\in A\)，使
\[
xy^{n+1}=\sum_{j=0}^n a_jxy^j.
\]
根据刚证的环的描述，写 \(a_j=c_j+x\cdot h_j(x,y)\)，其中 \(c_j\in k\)。比较两边作为 \(k[y][x]\) 中多项式的 \(x\) 的一次项系数，便得到
\[
y^{n+1}=\sum_{j=0}^n c_jy^j,
\]
这与多项式系数的唯一性矛盾。所以每个 \(I_n\subsetneq I_{n+1}\)，由\cref{b1:prop:noetherian-acc}，\(A\) 不是 Noether 环。另一方面，域 \(k\) 是 Noether 环，\cref{b1:thm:hilbert-basis}保证 \(k[x,y]\) 是 Noether 环。因此 Noether 性并不任意传给子环。
\end{solution}

\begin{exercise}\label{b1:ex:12-closed-field-evaluation}
设 \(k\) 为代数闭域，\(f\in k[x]\) 首一且次数为 \(n\ge1\)，其不同的根为 \(a_1,\ldots,a_r\)。证明取值公式给出满环同态
\[
E:k[x]/(f)\longrightarrow k^r,\qquad
g+(f)\longmapsto\bigl(g(a_1),\ldots,g(a_r)\bigr),
\]
并证明 \(\ker E\) 恰为商环中全部幂零元组成的集合。最后证明 \(E\) 是同构，当且仅当 \(f\) 没有重根。
\end{exercise}
\begin{solution}
由\cref{b1:prop:algebraically-closed-splitting}，存在正整数 \(m_1,\ldots,m_r\)，使
\[
f=\prod_{i=1}^r(x-a_i)^{m_i},\qquad
\sum_{i=1}^r m_i=n.
\]
如果 \(g-g'\in(f)\)，那么在每个 \(a_i\) 处两者取值相同，所以 \(E\) 良定义。求值保持加法、乘法和单位，故 \(E\) 是环同态。\cref{b1:thm:lagrange-interpolation}说明任意指定的 \(r\) 个值都可由一个次数至多为 \(r-1\) 的多项式取得，因此 \(E\) 满射。

若 \(z=g+(f)\) 幂零，存在 \(N\ge1\) 使 \(g^N\in(f)\)。代入每个 \(a_i\) 得 \(g(a_i)^N=0\)；域中没有非零幂零元，所以 \(g(a_i)=0\)，即 \(z\in\ker E\)。反过来，若这些取值全为零，\cref{b1:prop:factor-root}保证每个 \(x-a_i\) 都整除 \(g\)。这些一次因子两两不相伴，而\cref{b1:thm:polynomial-ufd}保证 \(k[x]\) 的唯一分解性，所以其乘积
\(h=\prod_{i=1}^r(x-a_i)\) 因而整除 \(g\)。取 \(N=\max_i m_i\)，则 \(f\mid h^N\mid g^N\)，所以 \(z^N=0\)。这就刻画了核。

如果 \(f\) 没有重根，所有 \(m_i=1\)，于是 \(f=h\)。刚才的论证说明 \(E\bigl(g+(f)\bigr)=0\) 就蕴含 \(f\mid g\)，故 \(E\) 单射，从而为同构。如果至少一个 \(m_i>1\)，则 \(r<n\)，多项式 \(h\) 的次数为 \(r\)，小于 \(f\) 的次数，所以 \(h\notin(f)\)；但 \(h\) 在各 \(a_i\) 处均为零。因此 \(h+(f)\) 是核中的非零元素，\(E\) 不是同构。
\end{solution}

% !TeX root = ../../Book1.tex
% 纲要标题：## 第13章　分式域、局部化与局部判据
% 纲要位置：计划纲要.md，第 478 行。

\chapter{分式域、局部化与局部判据}
\label{b1:ch:13}
\BookChapterNumber{b1:ch:13}

\begin{chapterabstract}
局部化通过指定可作分母的元素改变一个环，从而突出选定素理想附近的结构。本章从分式域推广到一般乘法集，建立同态的延拓性质和素理想对应，再以饱和及局部判据比较原环与各个局部环的信息。
\end{chapterabstract}

\begin{learninggoals}
\begin{itemize}
\item 区分形式分式、分式域中的元素与允许在某点求值的表达式。
\item 解释一般局部化等价关系中的额外乘子，并判断元素何时为零、何时可逆。
\item 利用泛性质构造和比较同态，计算素理想处及单个元素处的局部化。
\item 运用扩张、收缩和一般饱和分析理想，理解素理想的对应。
\item 用理想商计算单元素饱和，说明 Noether 条件如何保证相应升链稳定，并理解局部判据为什么需要检查全部极大理想。
\end{itemize}
\end{learninggoals}

从整数出发若只允许把 \(2\) 和 \(3\) 的幂放在分母，便得到 \(\mathbb Z[1/6]\)。在这个环里，\(2\) 与 \(3\) 都可逆，\(5\) 仍不可逆；它保留了整数的许多整除信息，却改变了我们观察这些信息的范围。把“允许哪些分母”写成一套一般构造以后，可以追问原环的理想与素理想哪些仍可辨认，哪些已经因分母变为单位而消失。局部化正是把全局问题限制到指定的部分，再比较这些局部结果能否决定原来的结论。

\ProseParagraphBreak

本章使用\chapref{b1:ch:10}的交换环、理想与商环知识，并使用\cref{b1:prop:fraction-field-basic}构造的分式域。局部化本身不要求原环具有唯一分解性，也不以 \(R\rightarrow S^{-1}R\) 为单射作为先决条件。\chapref{b1:ch:11}和\chapref{b1:ch:12}的整除与多项式计算，为本章提供具体例子；模的局部化留待第三卷第2章在模论基础上展开。

\ProseParagraphBreak
第13.1—13.3节建立局部化的基本构造，第13.4节据此研究理想的扩张、收缩与素理想对应，末小节给出局部—整体判据。单元素饱和及其有限步稳定的计算以这些基本对应为基础，章末关于一般理想饱和的习题也使用相应方法。

\ProseParagraphBreak
局部化的基础可参阅 Atiyah 与 Macdonald 的《Introduction to Commutative Algebra》\cite[Chapter 3]{AtiyahMacdonald1969CommutativeAlgebra}。Matsumura 的《Commutative Ring Theory》\cite[Section 4]{Matsumura1989CommutativeRingTheory}把局部化与环的素谱一起讨论；Eisenbud 的《Commutative Algebra: with a View Toward Algebraic Geometry》\cite[Sections 1.2, 1.4, 1.6, 2.1, 4.1--4.2, 4.4--4.5]{Eisenbud1995CommutativeAlgebra}进一步联系代数几何，其第 4.3 节的解析正规化属于另一方向\cite[Section 4.3]{Eisenbud1995CommutativeAlgebra}。Nagata 的《Local Rings》\cite[Sections 5--6]{Nagata1975LocalRings}把仅有一个极大理想的环称为 quasi-local，把 Noether 的 quasi-local 环称为 local，且第 6 节的乘法集不含零。本书的“局部环”不要求 Noether 性，并允许含零乘法集，此时局部化为零环。查阅这些著作的其他部分还需更多交换环与模的知识。

% !TeX root = ../../Book1.tex
% 纲要标题：### 13.1 分式域
% 纲要位置：计划纲要.md，第 472 行。

\section{分式域}
\label{b1:sec:1301}

\cref{b1:prop:fraction-field-basic}已经完整构造整环的分式域，并验证等价关系、运算和域公理。本节沿用这个对象，研究它与其他域的关系：在已有的域中如何识别分式域，从整环出发的同态又何时能够延拓到分式域？这些问题将引出泛性质，并为下一节只让指定元素取逆的一般局部化作准备。

% 纲要内容（仅作写作计划，尚非教材正文）：
%
% * 分式域的回顾与例子；完整构造已在第12.2节建立
% * 分式相等与分式域的识别
% * 分式域的泛性质，为一般局部化准备同态延拓的观点
%

\subsection{分式域的回顾与例子}
\label{b1:subsec:1301:01}

设 \(R\) 为交换整环。回顾\cref{b1:prop:fraction-field-basic}，域 \(\operatorname{Frac}(R)\) 的元素由 \(a/b\) 表示，其中 \(a,b\in R\)、\(b\neq0\)，相等的判据是
\begin{equation}\label{b1:eq:fraction-equivalence}
\frac ab=\frac cd\quad\Longleftrightarrow\quad ad=bc.
\end{equation}
这里两个分母都非零，整环的消去律使交叉乘积相等成为完整的判据。在有零因子的一般环中，\(u(ad-bc)=0\) 未必允许消去 \(u\)；但若让 \(u\) 变成单位，交叉乘积的差便会变成零。下一节分式等价关系中的额外因子，正是为了记录这种因取逆而消失的差异。

\ProseParagraphBreak
后面使用的运算公式为
\begin{equation}\label{b1:eq:fraction-operations}
\frac ab+\frac cd=\frac{ad+bc}{bd},
\qquad
\frac ab\cdot\frac cd=\frac{ac}{bd}.
\end{equation}
分式域也可以不指定等价类模型，而用它包含原环的方式来刻画。

\begin{definition}\label{b1:def:fraction-field}
设 \(R\) 为交换整环，\(K\) 为域，\(i:R\rightarrow K\) 为保持单位的单射环同态。如果每个 \(x\in K\) 都能写成
\[
x=i(a)i(b)^{-1}\qquad(a,b\in R,\ b\neq0),
\]
就称域 \(K\) 连同嵌入 \(i\) 为 \(R\) 的分式域。
\end{definition}
已构造的 \(\operatorname{Frac}(R)\) 与嵌入 \(i(a)=a/1\) 满足这个定义。通常通过 \(i\) 把 \(R\) 看成分式域的子环，并将 \(i(a)\) 简记为 \(a\)。本节末将证明：不同模型之间有唯一保持这个嵌入的同构。

\ProseParagraphBreak
例如 \(\operatorname{Frac}(\mathbb Z)=\mathbb Q\)。若 \(k\) 为域，则 \(k[x]\) 的分式域记为 \(k(x)\)，称为 \(k\) 上的\term[youli-hanshu-yu]{有理函数域}[rational function field]；元素是多项式的分式 \(f(x)/g(x)\)，其中 \(g\neq0\)。
\symidx{\(k(x)\)：域 \(k\) 上的一元有理函数域}
这里说的是形式分式，而不是在 \(k\) 上处处有定义的函数。域元素 \(h\in k(x)\) 可以有多个分式表示 \(h=f/g\)；对于指定的一个表示，只有 \(g(a)\ne0\) 时才能在 \(a\in k\) 处直接代入。例如 \(1/x\) 是 \(k(x)\) 的元素，但所写表示不能在 \(x=0\) 处直接代入。若 \(R\) 本来就是域，则 \(a/b\mapsto ab^{-1}\) 给出 \(\operatorname{Frac}(R)\cong R\)，构造并没有加入新的元素。

\subsection{分式相等与分式域的识别}
\label{b1:subsec:1301:02}

分式计算的依据是\eqref{b1:eq:fraction-equivalence}，不必先选定所谓最简分式。一般整环中的最大公因子未必存在，因而不能把“约分到最简”作为分式域定义的一部分。即使分子、分母只有单位作为公因子，也未必能保证两种表示只差同时乘一个单位；\cref{b1:ex:13-13}将在 \(\mathbb Z[\sqrt{-5}]\) 中给出这样的既约分式反例。

\ProseParagraphBreak
在有零因子的环中，分母取逆还可能使自然映射不再单射：若 \(sa=0\) 而要求 \(s\) 可逆，那么在新环中 \(a=s^{-1}(sa)=0\)，即使原环中的 \(a\) 非零也一样。因此下一节的局部化既要加入指定元素的逆，也可能使原来的某些元素变成零。

\ProseParagraphBreak
分式相等不表示分子、分母分别相等。例如在 \(k(x)\) 中，
\[
\frac{x^2-1}{x-1}=x+1,
\]
因为 \(x^2-1=(x-1)(x+1)\)，而多项式 \(x-1\) 非零。这是有理函数域内的等式；左端原表示在 \(x=1\) 处不能直接代入，右端则可以。给定域元素 \(h\in k(x)\)，如果它至少有一种表示 \(h=f/g\) 满足 \(g(a)\ne0\)，就说 \(h\) 在 \(a\) 处正则，并用这一表示求值。这个值不依赖所选的可代入表示：若 \(f/g=f_1/g_1\)，且两个分母在 \(a\) 处都非零，则把 \(fg_1=f_1g\) 在 \(a\) 处求值，再消去非零分母，便得到相同的商值。因此，在一点正则是域元素具有某种可代入表示的性质，不要求它的每一种表示都能直接代入；上式中的元素在 \(1\) 处正则。

\ProseParagraphBreak
若 \(R\) 已经放在一个域 \(L\) 中，每个包含 \(R\) 的子域都被迫包含非零 \(b\in R\) 的逆，进而包含所有 \(ab^{-1}\)。所以，要在 \(L\) 中识别分式域，就应收集这些必需的商，并检查它们自身已经组成子域；这时它便是包含 \(R\) 的最小子域。

\begin{proposition}\label{b1:prop:fraction-in-ambient-field}
设 \(R\) 是域 \(L\) 的子环。则
\[
K=\left\{\,ab^{-1}\in L\mid a,b\in R,\ b\neq0\,\right\}
\]
是 \(L\) 中包含 \(R\) 的最小子域，且 \(a/b\mapsto ab^{-1}\) 是从 \(\operatorname{Frac}(R)\) 到 \(K\) 的同构。
\end{proposition}
\begin{proof}
对 \(a,b,c,d\in R\)、\(b,d\neq0\)，域 \(L\) 中有
\[
ab^{-1}+cd^{-1}=(ad+bc)(bd)^{-1},
\qquad
(ab^{-1})(cd^{-1})=(ac)(bd)^{-1}.
\]
这些元素仍在 \(K\) 中，且 \(K\) 含 \(0,1\)，在取负和对非零元素取逆下也封闭，所以是子域。包含 \(R\) 的子域必须包含每个 \(b^{-1}\) 及 \(ab^{-1}\)，故包含 \(K\)。最后，\(ab^{-1}=cd^{-1}\) 在 \(L\) 中等价于 \(ad=bc\)；这个等式也在子环 \(R\) 中成立。因此所给映射良定义且单射，分式运算说明它是环同态，\(K\) 的定义保证满射。
\end{proof}
这个命题允许在已有的域中直接寻找分式域。例如把 \(\mathbb Z[i]\) 放在 \(\mathbb C\) 中，其分式域就是 \(\mathbb Q(i)=\{\,a+bi\mid a,b\in\mathbb Q\,\}\)。后一集合在和、积及取负下封闭；非零元素的逆为
\[
(a+bi)^{-1}=\frac{a-bi}{a^2+b^2},
\]
仍具有有理数的实部与虚部，所以它是包含 \(\mathbb Z[i]\) 的域。每个这样的元素又可通分成一个 Gauss 整数除以非零整数，故确为 \(\mathbb Z[i]\) 的分式域。

\subsection{分式域的泛性质}
\label{b1:subsec:1301:03}

把分式代入映射时，分母必须仍可取逆。这给出分式域的泛性质。目标可以是一般交换含幺环 \(A\)：不要求其中每个非零元素都可逆，只要求原环中将被取逆的非零元素的像可逆。这正是下一节局部化泛性质所保留的条件。

\begin{theorem}\label{b1:thm:fraction-field-universal}
设 \(R\) 为交换整环，\(A\) 为交换含幺环，\(f:R\rightarrow A\) 为保持单位的环同态，记自然嵌入为 \(\iota(a)=a/1\)。存在保持单位的环同态
\(\widetilde f:\operatorname{Frac}(R)\rightarrow A\) 满足 \(\widetilde f(a/1)=f(a)\) 对每个 \(a\in R\) 成立，当且仅当每个 \(b\in R\setminus\{0\}\) 的像 \(f(b)\) 都是 \(A\) 的单位。存在时，此同态唯一；对 \(a,b\in R\)、\(b\ne0\)，有
\[
\widetilde f\left(\frac ab\right)=f(a)f(b)^{-1}.
\]
延拓存在时，下面的图交换：
\[
\begin{tikzcd}[column sep=large,row sep=large]
R\arrow[r,"\iota"]\arrow[dr,"f"']
&\operatorname{Frac}(R)\arrow[d,dashed,"\exists!\,\widetilde f"]\\
&A.
\end{tikzcd}
\]
若 \(A\ne0\) 且上述延拓存在，则 \(f\) 和 \(\widetilde f\) 都单射。特别地，若 \(A=L\) 为域而 \(f\) 单射，则延拓存在，并且是域的单射同态。
\end{theorem}
\begin{proof}
若延拓存在，\(b/1\) 的逆元为 \(1/b\)，故 \(f(b)\widetilde f(1/b)=1\)，说明 \(f(b)\) 必须可逆。反过来，假设每个这样的 \(f(b)\) 可逆。若 \(a/b=c/d\)，则 \(ad=bc\)，从而
\(f(a)f(d)=f(b)f(c)\)。乘以 \(f(b)^{-1}f(d)^{-1}\)，得到
\(f(a)f(b)^{-1}=f(c)f(d)^{-1}\)，所给公式因而良定义。由\eqref{b1:eq:fraction-operations}，
\[
\begin{aligned}
\widetilde f\left(\frac ab+\frac cd\right)
&=\bigl(f(a)f(d)+f(b)f(c)\bigr)f(b)^{-1}f(d)^{-1}\\
&=f(a)f(b)^{-1}+f(c)f(d)^{-1},\\
\widetilde f\left(\frac ab\cdot\frac cd\right)
&=f(a)f(c)f(b)^{-1}f(d)^{-1}.
\end{aligned}
\]
故它保持加法、乘法及单位。任何延拓都必须把 \(a/b=(a/1)(b/1)^{-1}\) 映为 \(f(a)f(b)^{-1}\)，所以唯一。若 \(A\ne0\) 且延拓条件满足，假设有 \(0\ne a\in\ker f\)，则 \(f(a)=0\) 必须是单位，但非零环中的零元不可逆，矛盾。因此 \(f\) 单射。再由 \(\widetilde f(a/b)=0\) 得 \(f(a)=0\)，进而 \(a=0\)，故延拓也单射。若 \(A=L\) 为域且 \(f\) 单射，每个非零 \(b\) 的像都非零，因而可逆，延拓条件自动满足。
\end{proof}

例如单射 \(\mathbb Z\hookrightarrow\mathbb R\) 唯一地延拓为 \(\mathbb Q\hookrightarrow\mathbb R\)。目标不必是域：对角映射 \(\mathbb Z\rightarrow\mathbb Q\times\mathbb Q\)、\(a\mapsto(a,a)\) 也将每个非零整数送到单位，故唯一延拓为 \(a/b\mapsto(a/b,a/b)\)，尽管 \(\mathbb Q\times\mathbb Q\) 有零因子。模素数 \(p\) 的商映射 \(\mathbb Z\rightarrow\mathbb F_p\) 则不能延拓到 \(\mathbb Q\)：它的目标非零而核含有非零整数 \(p\)，不满足上面自动单射的必要条件；等价地说，\(p\) 的像为零，不能作逆。这个障碍不能靠另选分式代表来消除。

\begin{corollary}\label{b1:cor:fraction-field-unique}
设 \(R\) 为交换整环。若 \((K,i)\)、\((K',i')\) 都是 \(R\) 的分式域，则存在唯一保持指定 \(R\) 嵌入的域同构 \(\varphi:K\rightarrow K'\)，即 \(\varphi\circ i=i'\)。这样的同构记作 \((K,i)\cong_R(K',i')\)。
\end{corollary}
\begin{proof}
由\cref{b1:prop:fraction-in-ambient-field}，映射
\[
u:\operatorname{Frac}(R)\xrightarrow{\sim}K,
\qquad
v:\operatorname{Frac}(R)\xrightarrow{\sim}K',
\]
其中 \(u(a/b)=i(a)i(b)^{-1}\)、\(v(a/b)=i'(a)i'(b)^{-1}\)，都是同构：按分式域的定义，\(K\)、\(K'\) 中每个元素都可写成相应的商。取
\[
\varphi=v\circ u^{-1}:K\xrightarrow{\sim}K',
\]
便有 \(\varphi\circ i=i'\)。若另有这样的同态，则它在每个 \(i(a)i(b)^{-1}\) 上的值必为 \(i'(a)i'(b)^{-1}\)，所以与 \(\varphi\) 一致。
\end{proof}
因此，具体等价类只是构造模型；决定分式域的是它对 \(R\) 的包含关系和分母取逆的要求。强调“保持 \(R\) 的嵌入”也很必要：抽象域之间可能有多个同构，唯一性指满足指定条件的那一个。

% !TeX root = ../../Book1.tex
% 纲要标题：### 13.2 环的局部化
% 纲要位置：计划纲要.md，第 478 行。

\section{环的局部化}
\label{b1:sec:1302}

未必每一次都要把所有非零元素取逆。在研究整数的整除性时，我们可能只希望除以 \(2\) 的幂；在研究多项式在一点附近的性质时，则希望允许在该点不为零的多项式作分母。这两种选择都对乘法封闭。以下固定交换含幺环 \(R\)，允许它有零因子；构造仍由分式开始，但必须仔细处理不能消去的因子。

% 纲要内容（仅作写作计划，尚非教材正文）：
%
% * 乘法集与分母
% * 局部化中的零与单位
% * 同态延拓的泛性质
%

\subsection{乘法集与分母}
\label{b1:subsec:1302:01}

\begin{definition}\label{b1:def:multiplicative-set}
设 \(R\) 为交换含幺环，\(S\subseteq R\)。如果 \(1\in S\)，且对任意 \(s,t\in S\) 都有 \(st\in S\)，就称 \(S\) 为 \(R\) 的\term[chengfa-ji]{乘法集}[multiplicative subset]。本书允许 \(0\in S\)，但会明确说明它导致的退化情形。
\end{definition}
典型例子是 \(\{1\}\)、单位群 \(R^\times\)，以及固定元素 \(f\in R\) 的幂组成的集合 \(\{\,1,f,f^2,\ldots\,\}\)。若 \(R\) 为整环，则 \(R\setminus\{0\}\) 也是乘法集；有零因子时，这个结论可能失败。

\ProseParagraphBreak
设 \((a,s),(b,t)\in R\times S\)。若直接照搬整环中的交叉乘积判据，得到的候选关系是
\[
(a,s)\approx(b,t)\quad\Longleftrightarrow\quad at=bs.
\]
它在有零因子时未必满足传递性。例如取 \(R=\mathbb Z/6\mathbb Z\)、\(S=\{\bar1,\bar3\}\)，则
\[
(\bar2,\bar1)\approx(\bar0,\bar3),\qquad
(\bar0,\bar3)\approx(\bar0,\bar1),
\]
但 \((\bar2,\bar1)\not\approx(\bar0,\bar1)\)。失败之处是不能消去中间分母 \(\bar3\)。因此，正式使用的关系允许交叉乘积之差被某个分母零化：
\begin{equation}\label{b1:eq:localization-equivalence}
(a,s)\sim(b,t)
\quad\Longleftrightarrow\quad
\text{存在 }u\in S\text{，使 }u(at-bs)=0.
\end{equation}
下文将证明 \(\sim\) 是等价关系。在这个关系下，取 \(u=\bar3\) 便有 \((\bar2,\bar1)\sim(\bar0,\bar1)\)；上述三个有序对属于同一等价类。额外的 \(u\) 表示：允许的分母取逆后，原先由某个分母零化的差也必须成为零。

\begin{theorem}\label{b1:thm:localization-construction}
设 \(R\) 为交换含幺环，\(S\) 为 \(R\) 的乘法集。在 \(R\times S\) 上定义关系
\[
(a,s)\sim(b,t)
\quad\Longleftrightarrow\quad
\text{存在 }u\in S\text{，使 }u(at-bs)=0.
\]
则 \(\sim\) 是等价关系。将 \((a,s)\) 的等价类记为 \(a/s\)，则这些等价类在运算
\begin{equation}\label{b1:eq:localization-operations}
\frac as+\frac bt=\frac{at+bs}{st},\qquad
\frac as\cdot\frac bt=\frac{ab}{st}
\end{equation}
下组成交换含幺环，零元为 \(0/1\)，单位元为 \(1/1\)。称这个环为 \(R\) 关于 \(S\) 的\term[jubu-hua]{局部化}[localization]，记为 \(S^{-1}R\)。映射
\[
\iota:R\longrightarrow S^{-1}R,\qquad a\longmapsto a/1
\]
是保持单位的环同态；对每个 \(s\in S\)，元素 \(\iota(s)=s/1\) 都可逆，其逆元为 \(1/s\)。
\end{theorem}
\begin{proof}
自反性取 \(u=1\)，对称性把等式取负即可。为证传递性，设
\(u(at-bs)=0\)、\(v(br-ct)=0\)，其中 \(s,t,r,u,v\in S\)。有
\[
uvt(ar-cs)=uv\bigl(r(at-bs)+s(br-ct)\bigr)=0.
\]
因为 \(uvt\in S\)，故 \((a,s)\sim(c,r)\)。

再证运算良定义。若 \(u(as'-a's)=0\)，则与固定分式 \(b/t\) 相加所得两个候选代表的交叉乘积之差是
\[
(at+bs)s't-(a't+bs')st=t^2(as'-a's),
\]
它被 \(u\) 零化；相乘所得差为
\[
ab\,s't-a'b\,st=bt(as'-a's),
\]
也被 \(u\) 零化。因此改变第一个分式的代表不改变结果。由于两个运算的候选公式在交换两个分式后不变，改变第二个代表也不改变结果；先后改变两个代表，再用已证的传递性，便得到良定义。

加法、乘法交换律直接来自 \(R\) 的交换律。三个分式的两种加法结合次序都给出
\[
\frac{atr+bsr+cst}{str},
\]
乘法的两种结合次序都给出 \(abc/(str)\)。分配律左端 \((a/s)(b/t+c/r)\) 表示 \(a(br+ct)/(str)\)；右端 \(ab/(st)+ac/(sr)\) 表示 \(as(br+ct)/(s^2tr)\)。这两个分式的交叉乘积相等，因而代表同一元素。最后，\(0/1\)、\(1/1\) 分别满足加法、乘法单位律，\((-a)/s\) 是 \(a/s\) 的加法逆元。由运算公式，\(\iota\) 保持加法、乘法与单位；等式 \((s/1)(1/s)=1/1\) 说明 \(s/1\) 可逆。
\end{proof}

\symidx{\(S^{-1}R\)：环 \(R\) 关于乘法集 \(S\) 的局部化}
每个分式都满足 \(a/s=(a/1)(1/s)\)。不过自然同态 \(\iota\) 未必单射，所以在证明单射以前，不能把局部化一概理解为给原环“增加一些元素”。

\subsection{局部化中的零与单位}
\label{b1:subsec:1302:02}

\begin{proposition}\label{b1:prop:localization-zero-unit}
设 \(R\) 为交换含幺环，\(S\) 为 \(R\) 的乘法集，\(\iota:R\rightarrow S^{-1}R\) 为自然同态。对 \(a\in R\)、\(s\in S\)，有以下判据：
\begin{enumerate}
\item \(a/s=0\)，当且仅当存在 \(t\in S\) 使 \(ta=0\)。
\item \(a/s\) 是 \(S^{-1}R\) 的单位，当且仅当存在 \(b\in R\) 使 \(ab\in S\)。
\end{enumerate}
特别地，\(S^{-1}R\) 为零环当且仅当 \(0\in S\)；自然映射 \(\iota\) 单射当且仅当对所有 \(s\in S\) 与 \(a\in R\)，都有 \(sa=0\Longrightarrow a=0\)。
\end{proposition}
\begin{proof}
将 \(a/s\) 与 \(0/1\) 代入\eqref{b1:eq:localization-equivalence}，立即得到第一个判据。

若存在 \(b\in R\) 使 \(ab\in S\)，则分式 \(bs/(ab)\) 有意义，且
\[
\frac as\cdot\frac{bs}{ab}=\frac{abs}{abs}=1.
\]
反过来，若 \((a/s)(c/t)=1\)，则存在 \(u\in S\) 使 \(u(ac-st)=0\)。令 \(b=uc\)，即有 \(ab=ust\in S\)，得到第二个判据。

环为零环等价于 \(1/1=0\)。第一个判据把它化为存在 \(t\in S\) 使 \(t\cdot1=0\)，也就是 \(0\in S\)。同一判据给出
\[
\ker\iota=\{\,a\in R\mid ta=0\text{ 对某个 }t\in S\,\},
\]
故 \(\iota\) 单射恰好等价于最后所列条件。
\end{proof}

在整环中，只要 \(0\notin S\)，自然映射就单射，而且 \(S^{-1}R\) 是分式域 \(\operatorname{Frac}(R)\) 的子环：将局部化分式 \(a/s\) 送到分式域中的同名分式即可。局部化等式中出现的 \(u\in S\) 非零，可以消去，因此这个映射的单射性也直接来自两种等价关系的一致。

\ProseParagraphBreak
单位判据中的条件不能缩减为 \(a\in S\)。例如取 \(R=\mathbb Z\)、\(S=\{\,6^n\mid n\geq0\,\}\)，整数 \(2\) 不在 \(S\) 中，却在局部化中可逆，因为 \(2\cdot3=6\in S\)，逆元为 \(3/6\)。需要检查的是分子能否乘上某个环元素落入分母集，而不是分子本身是否被列入分母集。

\subsection{同态延拓的泛性质}
\label{b1:subsec:1302:03}

\begin{theorem}\label{b1:thm:localization-universal}
设 \(R,A\) 为交换含幺环，\(S\) 为 \(R\) 的乘法集，\(\iota:R\rightarrow S^{-1}R\) 为自然同态，\(f:R\rightarrow A\) 为保持单位的环同态。若 \(f(s)\in A^\times\) 对所有 \(s\in S\) 成立，则存在唯一保持单位的环同态 \(\widetilde f:S^{-1}R\rightarrow A\)，使 \(\widetilde f\circ\iota=f\)。对 \(a\in R\)、\(s\in S\)，它由下式给出：
\[
\widetilde f\left(\frac as\right)=f(a)f(s)^{-1}.
\]
反过来，只要这样的延拓存在，每个 \(f(s)\) 就必为单位。
\end{theorem}
\begin{proof}
若 \(a/s=b/t\)，取 \(u\in S\) 使 \(u(at-bs)=0\)。施加 \(f\) 得
\[
f(u)\bigl(f(a)f(t)-f(b)f(s)\bigr)=0.
\]
由于 \(f(u)\)、\(f(s)\)、\(f(t)\) 均可逆，先乘以 \(f(u)^{-1}\)，再乘以两个分母像的逆元，就得到
\(f(a)f(s)^{-1}=f(b)f(t)^{-1}\)。因而候选公式良定义。它保持运算，因为
\[
\begin{aligned}
f(at+bs)f(st)^{-1}
&=f(a)f(s)^{-1}+f(b)f(t)^{-1},\\
f(ab)f(st)^{-1}
&=f(a)f(s)^{-1}f(b)f(t)^{-1}.
\end{aligned}
\]
它也把 \(1/1\) 送到 \(1\)，并延拓 \(f\)。任一延拓都须把 \(1/s\) 送到 \(f(s)^{-1}\)，再利用 \(a/s=(a/1)(1/s)\) 便知唯一。必要性则来自环同态将单位送到单位。
\end{proof}

当 \(f(S)\subseteq A^\times\) 时，下图表示延拓关系：虚线是唯一满足 \(\widetilde f\circ\iota=f\) 的保单位环同态。
\[
\begin{tikzcd}[column sep=3em,row sep=2.5em]
R\arrow[r,"\iota"]\arrow[dr,"f"']&S^{-1}R\arrow[d,dashed,"\exists!\,\widetilde f"]\\
&A
\end{tikzcd}
\]
泛性质也覆盖 \(0\in S\) 的情形。此时条件要求 \(f(0)=0\) 是 \(A\) 的单位，只能有 \(A\) 为零环；而 \(S^{-1}R\) 也是零环，确实有唯一这样的同态。把退化情形保留在定义内，不会破坏定理。

\begin{proposition}\label{b1:prop:localization-comparison}
设 \(R\) 为交换含幺环，\(S\subseteq T\) 为 \(R\) 的两个乘法集。存在唯一保持 \(R\) 中元素像的环同态
\[
S^{-1}R\longrightarrow T^{-1}R,\qquad a/s\longmapsto a/s.
\]
这里 \(a\in R\)、\(s\in S\)。
若 \(T\) 中每个元素的像已经在 \(S^{-1}R\) 中可逆，则这个同态是同构。
\end{proposition}
\begin{proof}
自然映射 \(R\rightarrow T^{-1}R\) 将 \(S\) 中元素送到单位，由\cref{b1:thm:localization-universal}得到所需同态 \(\alpha\)。若 \(T\) 中元素在 \(S^{-1}R\) 中也都可逆，则同一定理对 \(R\rightarrow S^{-1}R\) 给出 \(\beta:T^{-1}R\rightarrow S^{-1}R\)。复合 \(\beta\alpha\) 与恒等映射都延拓 \(R\rightarrow S^{-1}R\)，由唯一性相等；\(\alpha\beta\) 与 \(T^{-1}R\) 的恒等映射也如此。故 \(\alpha\)、\(\beta\) 互逆。
\end{proof}
比较映射不必单射。例如从 \(\mathbb Z/6\mathbb Z\) 出发，先取 \(S=\{1\}\)，再允许 \(3\) 的幂作分母，\(\bar2\) 的像便成为零。分母越多，既可能产生新的分式，也可能使原来不同的元素变得相等。

% !TeX root = ../../Book1.tex
% 纲要标题：### 13.3 素理想处与单元素处的局部化
% 纲要位置：计划纲要.md，第 484 行。

\section{素理想处与单元素处的局部化}
\label{b1:sec:1303}

分母集的选择决定哪些差别会被保留。本节考察两种常用选择：避开一个素理想的全部元素，以及一个固定元素的各次幂。前者把注意力集中于一个素理想，后者只要求一个元素可逆。整数环和多项式环的具体计算会说明，这些构造虽然采用同一套分式运算，却有不同的用途。

% 纲要内容（仅作写作计划，尚非教材正文）：
%
% * 素理想处的局部环
% * 单个元素处的局部化
% * 整数的局部化与 Laurent 多项式环
%

\subsection{素理想处的局部环}
\label{b1:subsec:1303:01}

设 \(\mathfrak p\) 为交换环 \(R\) 的素理想。因为 \(\mathfrak p\neq R\)，\(1\notin\mathfrak p\)；又因为
\(ab\in\mathfrak p\Longrightarrow a\in\mathfrak p\) 或 \(b\in\mathfrak p\)，其补集 \(R\setminus\mathfrak p\) 对乘法封闭。于是可以定义
\[
R_{\mathfrak p}\coloneqq(R\setminus\mathfrak p)^{-1}R.
\]
\symidx{\(R_{\mathfrak p}\)：环 \(R\) 在素理想 \(\mathfrak p\) 处的局部化}
它称为 \(R\) 在 \(\mathfrak p\) 处的局部化。由于 \(0\in\mathfrak p\)，分母集不含零，所以\cref{b1:prop:localization-zero-unit}保证 \(R_{\mathfrak p}\) 不是零环。

\begin{definition}\label{b1:def:local-ring}
设 \(A\) 为非零交换含幺环。如果 \(A\) 恰有一个极大理想 \(\mathfrak m\)，就称 \(A\) 为\term[jubu-huan]{局部环}[local ring]，并常以 \((A,\mathfrak m)\) 同时注明环和它的唯一极大理想。此时商环 \(A/\mathfrak m\) 是域，称为 \(A\) 的\term[shengyu-yu]{剩余域}[residue field]。
\end{definition}
这里商环确为域，是因为\cref{b1:prop:prime-max-quotient}已证明极大理想的商为域。域自身也是局部环，其唯一极大理想为零理想；“局部环”并不意味着一定存在非零非单位元。

\begin{proposition}\label{b1:prop:localization-at-prime}
设 \(R\) 为交换含幺环，\(\mathfrak p\) 为 \(R\) 的素理想。则 \(R_{\mathfrak p}=(R\setminus\mathfrak p)^{-1}R\) 是局部环，唯一极大理想为
\[
\mathfrak pR_{\mathfrak p}
=\left\{\,a/s\mid a\in\mathfrak p,\ s\in R\setminus\mathfrak p\,\right\}.
\]
对 \(a\in R\)、\(s\in R\setminus\mathfrak p\)，分式 \(a/s\) 是单位当且仅当 \(a\notin\mathfrak p\)。
\end{proposition}
\begin{proof}
若 \(a\notin\mathfrak p\)，则 \(a\) 也可作分母，\(s/a\) 就是 \(a/s\) 的逆。若 \(a\in\mathfrak p\) 而 \(a/s\) 为单位，\cref{b1:prop:localization-zero-unit}给出某个 \(b\in R\) 使 \(ab\notin\mathfrak p\)，与理想的吸收性矛盾。

记陈述中的集合为 \(\mathfrak n\)。两个分子在 \(\mathfrak p\) 中的分式相加时，通分后的分子仍在 \(\mathfrak p\) 中；取负及乘以任意分式也保持这一性质，所以 \(\mathfrak n\) 是理想。刚才的单位判据表明它恰为全部非单位组成的集合，因此不含 \(1\)，是真理想。任何真理想都不能含有单位，从而必包含于 \(\mathfrak n\)；这同时说明 \(\mathfrak n\) 极大，而且是唯一极大理想。它也正是由所有 \(a/1\)、\(a\in\mathfrak p\) 生成的理想，故可记为 \(\mathfrak pR_{\mathfrak p}\)。
\end{proof}

进一步可以算出这个局部环的剩余域。记 \(\bar a\) 为 \(a\) 在 \(R/\mathfrak p\) 中的像。这个商是整环，故有分式域。映射
\[
R\longrightarrow\operatorname{Frac}(R/\mathfrak p),\qquad
a\longmapsto\bar a/\bar1
\]
把每个 \(s\notin\mathfrak p\) 送到非零元素，从而送到单位。由\cref{b1:thm:localization-universal}，它唯一延拓为
\[
R_{\mathfrak p}\longrightarrow\operatorname{Frac}(R/\mathfrak p),
\qquad a/s\longmapsto\bar a/\bar s.
\]
任何非零 \(\bar s\) 都有 \(s\notin\mathfrak p\) 的代表，故此映射满射；其核由 \(\bar a=0\) 的分式组成，正是 \(\mathfrak pR_{\mathfrak p}\)。\cref{b1:thm:ring-first-iso}将商环同构到这个实际像，得到
\begin{equation}\label{b1:eq:local-residue-field}
R_{\mathfrak p}/\mathfrak pR_{\mathfrak p}
\cong\operatorname{Frac}(R/\mathfrak p).
\end{equation}
这说明素理想处局部化的剩余域一般不只是 \(R/\mathfrak p\)；当 \(\mathfrak p\) 极大时，后一商本来就是域，才无须再增加分母。

例如设 \(k\) 为域，取 \(R=k[x,y]\)、\(\mathfrak m=(x,y)\) 和 \(\mathfrak p=(x)\)。在 \(\mathfrak m\) 处，允许作分母的多项式 \(g\) 恰好满足 \(g(0,0)\neq0\)。因此原点求值延拓为满同态
\[
R_{\mathfrak m}\longrightarrow k,\qquad
\frac fg\longmapsto\frac{f(0,0)}{g(0,0)}.
\]
它的核为 \(\mathfrak mR_{\mathfrak m}\)，故剩余域是 \(k\)。在 \(\mathfrak p\) 处，分母条件变成 \(g\notin(x)\)，也就是 \(g(0,y)\) 作为 \(k[y]\) 中的多项式不为零。此时有满同态
\[
R_{\mathfrak p}\longrightarrow k(y),\qquad
\frac fg\longmapsto\frac{f(0,y)}{g(0,y)},
\]
其核为 \(\mathfrak pR_{\mathfrak p}\)，所以剩余域是有理函数域 \(k(y)\)。特别地，\(y\notin(x)\)，故 \(y\) 在 \(R_{\mathfrak p}\) 中可逆；但它在原点取值为零，原点求值不能延拓到 \(R_{\mathfrak p}\)。

\ProseParagraphBreak
暂取 \(k=\mathbb R\)，将 \((x)\) 与直线 \(x=0\) 联系起来。条件 \(g\notin(x)\) 只说明 \(g(0,y)\) 不是零多项式，并不要求它在这条直线的每一点都非零。例如 \(g=y\) 满足分母条件，却在原点为零。因此 \(\mathbb R[x,y]_{(x)}\) 允许 \(1/y\)，却不允许 \(1/x\)，不能将它理解为在直线上所有实点都有定义的有理函数环。用 \(\mathbb R[x,y]\) 的唯一分解将分子、分母约去公因子后，属于这个局部环的条件就是分母不含因子 \(x\)。若既约分式的分母含有 \(x\)，分子便不含 \(x\)；任何相等的分式由交叉相乘和 \(x\) 的素性，仍须有含 \(x\) 的分母，因而无法改写成允许的表示。

\ProseParagraphBreak
再取一元多项式环 \(k[x]\) 及 \(a\in k\)，在极大理想 \((x-a)\) 处局部化。由于 \(k[x]\) 是整环，可以把所得局部环看成 \(k(x)\) 的子环。一个有理函数属于这个局部环，当且仅当它可以写成分母在 \(a\) 处非零的分式：
\[
k[x]_{(x-a)}
=\left\{\,h\in k(x)\ \middle|\,
h=\frac fg\text{ 对某些 }f,g\in k[x]\text{ 成立，且 }g(a)\neq0
\right\}.
\]
这不要求有理函数的每一种表示都能直接代入 \(a\)。若 \(f/g=f_1/g_1\) 且 \(g(a),g_1(a)\neq0\)，交叉相乘后代入 \(a\)，便有 \(f(a)g_1(a)=f_1(a)g(a)\)。因此 \(f/g\mapsto f(a)/g(a)\) 与所选的合适表示无关，给出从 \(k[x]_{(x-a)}\) 到 \(k\) 的求值同态。

\ProseParagraphBreak
回到\secref{b1:sec:1301}中的例子，\((x^2-1)/(x-1)=x+1\) 所表示的元素属于 \(k[x]_{(x-1)}\)，可用表示 \(x+1\) 在 \(x=1\) 处求值。相反，\(1/(x-1)\) 不属于这个局部环：若它等于某个分母满足 \(g(1)\neq0\) 的 \(f/g\)，则在 \(k[x]\) 中有 \(g=(x-1)f\)，从而 \(g(1)=0\)，矛盾。

\subsection{单个元素处的局部化}
\label{b1:subsec:1303:02}

设 \(R\) 为交换含幺环，\(f\in R\)。令 \(S=\{\,f^n\mid n\geq0\,\}\)，其中 \(f^0=1\)。称局部化 \(S^{-1}R\) 为 \(R\) 在元素 \(f\) 处的局部化，记为 \(R_f\) 或 \(R[1/f]\)；其中每个元素都能写成 \(a/f^n\)，\(a\in R\)、\(n\geq0\)。
\symidx{\(R_f=R[1/f]\)：使元素 \(f\) 可逆的局部化}
两种记号都很常用；写 \(R[1/f]\) 时，\(1/f\) 表示构造中的形式逆元，不预先假定它已经存在于某个包含 \(R\) 的域中。局部化是一种使指定元素可逆的构造，所得环未必是局部环；上文的素理想处局部化则一定是局部环。若存在非零 \(a\in R\) 使 \(fa=0\)，则 \(a/1=0\) 于 \(R_f\)，所以自然映射 \(R\rightarrow R_f\) 必不单射。

\begin{proposition}\label{b1:prop:single-element-localization}
设 \(R\) 为交换含幺环，\(f\in R\)，并令
\[
R_f=\{\,1,f,f^2,\ldots\,\}^{-1}R.
\]
则有保持 \(R\) 中元素像的同构
\[
R_f\cong R[t]/(ft-1),
\]
其中 \(t\) 为一个多项式变量，\(1/f\) 对应于 \(t\) 的剩余类。此外，\(R_f\) 为零环当且仅当 \(f\) 幂零。
\end{proposition}
\begin{proof}
记 \(B=R[t]/(ft-1)\)，并用上横线表示剩余类。在 \(B\) 中有 \(\bar f\bar t=1\)，故自然映射 \(R\rightarrow B\) 把每个 \(f^n\) 送到单位。\cref{b1:thm:localization-universal}给出唯一同态
\[
\alpha:R_f\longrightarrow B,\qquad a/f^n\longmapsto\bar a\bar t^{\,n}.
\]
另一方面，多项式求值
\[
R[t]\longrightarrow R_f,\qquad
\sum_{j=0}^m a_jt^j\longmapsto\sum_{j=0}^m\frac{a_j}{f^j}
\]
保持加法、乘法与单位，并把 \(ft-1\) 送到零，因此诱导同态 \(\beta:B\rightarrow R_f\)。复合 \(\beta\alpha\) 固定每个 \(a/f^n\)，所以是恒等映射；复合 \(\alpha\beta\) 固定全部 \(\bar a\) 和 \(\bar t\)，故固定它们组成的每个多项式，即固定 \(B\) 的全部元素。两者互逆。

由\cref{b1:prop:localization-zero-unit}，\(R_f=0\) 当且仅当其分母集中含零，即某个 \(f^n=0\)。这恰为 \(f\) 幂零；若 \(n=0\) 已使 \(1=0\)，则 \(R\) 本来就是零环，也仍满足这一结论。
\end{proof}

例如对 \(R=\mathbb Z/6\mathbb Z\)、\(f=\bar3\)，模 \(2\) 的映射将 \(\bar3\) 送到 \(1\)，所以延拓为 \(R_{\bar3}\rightarrow\mathbb F_2\)。又因 \(\bar3\bar2=0\)，在局部化环中有 \(\bar2/1=0\)，从而 \(\bar3/1=\bar1/1\)。每个分母 \(\bar3^n\) 在局部化环内都等于 \(1\)，每个分式因此都等于 \(0\) 或 \(1\)。上述映射把两者分别送到 \(0\)、\(1\)，于是 \(R_{\bar3}\cong\mathbb F_2\)。

\subsection{整数的局部化与 Laurent 多项式环}
\label{b1:subsec:1303:03}

对一个素数 \(p\)，整数环在 \((p)\) 处的局部化为
\[
\mathbb Z_{(p)}
=\bigl\{\,a/b\in\mathbb Q\,\bigm|\,a,b\in\mathbb Z,\ p\nmid b\,\bigr\}.
\]
它的单位是其中分子也不被 \(p\) 整除的分式，唯一极大理想是 \(p\mathbb Z_{(p)}\)，剩余域为 \(\mathbb F_p\)。这些断言分别来自\cref{b1:prop:localization-at-prime}和\eqref{b1:eq:local-residue-field}。例如 \(2/3\) 在 \(\mathbb Z_{(5)}\) 中是单位，\(5/3\) 则不是。分母可含任意不被 \(5\) 整除的整数，不限于某个固定整数的幂。

\ProseParagraphBreak
与此不同，\(\mathbb Z[1/p]\) 只允许 \(p\) 的幂作分母：
\[
\mathbb Z[1/p]=\{\,a/p^n\in\mathbb Q\mid a\in\mathbb Z,\ n\geq0\,\}.
\]
前一个环使 \(p\) 以外的素数可逆，后一个环使 \(p\) 可逆。\(\mathbb Z_{(p)}\) 是局部环，唯一极大理想为 \(p\mathbb Z_{(p)}\)；\(\mathbb Z[1/p]\) 则仍有其他素数所对应的极大理想。以 \(A=\mathbb Z[1/2]\) 为例，因为 \(2\) 在 \(\mathbb F_3\) 和 \(\mathbb F_5\) 中都是单位，模 \(3\)、模 \(5\) 的映射分别延拓为满同态
\[
A\longrightarrow\mathbb F_3,\qquad A\longrightarrow\mathbb F_5.
\]
两者的核分别为不同的极大理想 \(3A\)、\(5A\)，所以 \(A\) 不是局部环。更一般地，对整数 \(n\geq1\)，\(\mathbb Z[1/n]\) 允许 \(n\) 的所有素因子取逆；例如 \(\mathbb Z[1/6]\) 的元素可以写成 \(a/(2^r3^s)\)，其中 \(r,s\geq0\)，因为这类分母都整除足够高次的 \(6\) 的幂。

\ProseParagraphBreak
设 \(k\) 为域。在 \(k[t]\) 中使 \(t\) 可逆，得到\term[Laurent-duoxiangshi-huan]{Laurent 多项式环}[Laurent polynomial ring]：\footnote{洛朗（Pierre Alphonse Laurent，1813-1854），法国数学家，以复分析中的 Laurent 级数著称。}
\[
k[t,t^{-1}]\coloneqq k[t]_t.
\]
\symidx{\(k[t,t^{-1}]\)：Laurent 多项式环}
其中每个元素有形式
\[
\sum_{j=m}^n a_jt^j,\qquad m,n\in\mathbb Z,\quad a_j\in k,
\]
仅有有限多个非零系数；不允许任意无限级数。因为分式 \(f(t)/t^r\) 可把 \(f\) 的每个指数减去 \(r\)，它们恰好给出这些有限和。若两个有限和相等，乘以足够大的 \(t\) 的幂后，就成为 \(k[t]\) 中的多项式等式。自然映射 \(k[t]\rightarrow k[t]_t\) 单射，所以同次幂的系数分别相等，有限和表达由此唯一。

\begin{proposition}\label{b1:prop:laurent-units}
设 \(k\) 为域。Laurent 多项式环 \(k[t,t^{-1}]\) 的单位恰为 \(ct^r\)，其中 \(c\in k^\times\)、\(r\in\mathbb Z\)。
\end{proposition}
\begin{proof}
这样的元素有逆 \(c^{-1}t^{-r}\)。反过来，设非零 Laurent 多项式 \(u\) 的最低、最高非零次数分别为 \(m,n\)，而 \(v\) 的为 \(r,s\)。因为 \(k\) 是域，乘积的最低和最高项系数都不为零，故 \(uv\) 的两个端点次数分别为 \(m+r,n+s\)。若 \(uv=1\)，则二者均为零，因而
\[
(n-m)+(s-r)=0.
\]
两个括号都是非负整数，只能各为零。所以 \(m=n\)，\(u\) 只有一个非零项，必为 \(ct^m\)。
\end{proof}
例如 \(t\) 与 \(t^{-1}\) 都是单位，而 \(t+1\) 不是。由\cref{b1:prop:localization-zero-unit}，若 \(S\) 是交换含幺环 \(R\) 的乘法集，则 \(a\in R\) 的像在 \(S^{-1}R\) 中可逆，当且仅当某个 \(b\in R\) 满足 \(ab\in S\)；也就是说，\(a\) 整除分母集中的某个元素。这个条件不要求所有非零元素可逆，所以 \(k[t,t^{-1}]\) 不是分式域 \(k(t)\)。

% !TeX root = ../../Book1.tex
% 纲要标题：### 13.4 理想的局部化与局部判据
% 纲要位置：计划纲要.md，第 498 行。

\section{理想的局部化与局部判据}
\label{b1:sec:1304}

局部化改变了单位，也就改变了真理想的范围：一个理想只要含有允许的分母，到了局部化后便包含单位，因而成为全环。本节先说明哪些理想能够保留、哪些素理想会对应起来，再计算单元素饱和，并证明一个初等的局部—整体判据。

% 纲要内容（仅作写作计划，尚非教材正文）：
%
% * 理想的扩张、收缩与一般饱和
% * 素理想的对应
% * 单元素饱和与 Noether 情形的有限步稳定
% * 元素与理想的局部—整体判据；模局部化留待第三卷
%
% ---
%

\subsection{理想的扩张、收缩与一般饱和}
\label{b1:subsec:1304:01}

设 \(R\) 为交换含幺环，\(S\) 为 \(R\) 的乘法集，\(\iota:R\rightarrow S^{-1}R\) 为自然同态。对 \(R\) 的理想 \(I\)，把其像生成的理想记为 \(IS^{-1}R\)，也记为 \(S^{-1}I\)，并称其为 \(I\) 的\term[lixiang-kuozhang]{扩张}[extension of an ideal]。对 \(S^{-1}R\) 的理想 \(J\)，原像 \(\iota^{-1}(J)\) 称为 \(J\) 的\term[lixiang-shousuo]{收缩}[contraction of an ideal]。
\symidx{\(S^{-1}I=IS^{-1}R\)：理想的扩张}
不能把扩张误写为像集 \(\iota(I)\)：后者虽能吸收 \(\iota(R)\) 中的元素，却未必能吸收新加入的分母逆元以及一般分式。

\begin{proposition}\label{b1:prop:localization-ideal-correspondence}
设 \(R\) 为交换含幺环，\(S\) 为 \(R\) 的乘法集，\(\iota:R\rightarrow S^{-1}R\) 为自然同态，\(I\) 为 \(R\) 的理想。则有以下结论。
\begin{enumerate}
\item \(S^{-1}I=\{\,a/s\mid a\in I,\ s\in S\,\}\)，而对任意 \(a\in R\)、\(s\in S\)，
\begin{equation}\label{b1:eq:localized-ideal-membership}
a/s\in S^{-1}I
\quad\Longleftrightarrow\quad
ta\in I\text{ 对某个 }t\in S.
\end{equation}
\item 先扩张再收缩得到
\begin{equation}\label{b1:eq:localized-ideal-contraction}
\iota^{-1}(S^{-1}I)
=\{\,a\in R\mid ta\in I\text{ 对某个 }t\in S\,\}.
\end{equation}
\item 对 \(S^{-1}R\) 的每个理想 \(J\)，先收缩再扩张仍得到 \(J\)：
\[
S^{-1}\bigl(\iota^{-1}(J)\bigr)=J.
\]
\item \(S^{-1}I\) 为真理想当且仅当 \(I\cap S=\varnothing\)。
\end{enumerate}
\end{proposition}
\begin{proof}
分子属于 \(I\) 的分式集合在相加、取负及乘以任意分式下封闭，所以是理想。它包含每个 \(a/1\)、\(a\in I\)；另一方面，每个 \(a/s=(a/1)(1/s)\) 都在这些元素生成的理想中，所以该集合恰为 \(S^{-1}I\)。

若 \(a/s=b/u\)，其中 \(b\in I\)、\(u\in S\)，则有 \(v\in S\) 使 \(v(au-bs)=0\)，于是 \((vu)a=vbs\in I\)，且 \(vu\in S\)。反过来，若 \(ta\in I\)、\(t\in S\)，则 \(a/s=(ta)/(ts)\)，右端分子在 \(I\) 中，所以属于 \(S^{-1}I\)。这证明了第一个结论；其中令 \(s=1\)，便得到第二个。

为证第三个结论，注意 \(a/s\in J\) 当且仅当 \(a/1\in J\)：从前者乘以 \(s/1\) 得到后者，从后者乘以 \(1/s\) 得到前者。而 \(a/1\in J\) 正是 \(a\in\iota^{-1}(J)\)。因此按第一个结论，\(J\) 恰为收缩理想的扩张。

最后，\(S^{-1}I\) 不是真理想当且仅当 \(1/1\in S^{-1}I\)。\eqref{b1:eq:localized-ideal-membership}将这个条件化为存在 \(t\in S\) 使 \(t\in I\)，即 \(I\cap S\neq\varnothing\)。
\end{proof}

设 \(R\) 为交换含幺环，\(S\) 为 \(R\) 的乘法集，\(I\) 为 \(R\) 的理想。集合
\[
\{\,a\in R\mid ta\in I\text{ 对某个 }t\in S\,\}
=\iota^{-1}(S^{-1}I)
\]
称为 \(I\) 关于 \(S\) 的\term[chengfa-ji-baohe]{饱和}[saturation with respect to a multiplicative subset]，其中 \(\iota:R\rightarrow S^{-1}R\) 是自然同态。若这个集合等于 \(I\)，则称 \(I\) 关于 \(S\) 饱和；等价地，\(ta\in I\)、\(t\in S\) 总能推出 \(a\in I\)。因此扩张与收缩在 \(S^{-1}R\) 的全部理想和 \(R\) 的全部 \(S\)-饱和理想之间给出保持包含关系的一一对应。满射性、互逆性来自\cref{b1:prop:localization-ideal-correspondence}的第二、三个结论；包含关系的保持则来自生成理想和取原像的定义。

\begin{corollary}\label{b1:cor:noetherian-localization}
设 \(R\) 为 Noether 交换含幺环，\(S\) 为 \(R\) 的乘法集。则局部化环 \(S^{-1}R\) 也是 Noether 环。
\end{corollary}
\begin{proof}
任取 \(S^{-1}R\) 的理想 \(J\)，设 \(\iota:R\rightarrow S^{-1}R\) 为自然同态，并令 \(I=\iota^{-1}(J)\)。由 \(R\) 的 Noether 性，存在 \(a_1,\ldots,a_r\in R\)，使 \(I=(a_1,\ldots,a_r)\)。\cref{b1:prop:localization-ideal-correspondence}给出 \(J=S^{-1}I=(a_1/1,\ldots,a_r/1)\)。因此 \(J\) 有限生成，\(S^{-1}R\) 为 Noether 环。
\end{proof}

由\cref{b1:thm:hilbert-basis}，Noether 环上有限个不定元的多项式环及其商环，在局部化后仍是 Noether 环。局部环的定义本身不附加 Noether 条件；这里的有限性来自原环。

\ProseParagraphBreak
“与 \(S\) 不相交”只保证扩张仍为真理想，不保证扩张后可以恢复原理想。取 \(S=\{\,2^n\mid n\geq0\,\}\)、\(A=S^{-1}\mathbb Z=\mathbb Z[1/2]\)，并以 \(\iota:\mathbb Z\hookrightarrow A\) 表示自然嵌入。

\begin{samepage}
对下列三个与 \(S\) 不相交的理想，像集、扩张及收缩分别为
\begin{center}
\begin{tabular}{@{}cccc@{}}
\toprule
原理想 \(I\) & 像集 \(\iota(I)\) & 扩张 \(IA\) & 收缩 \(\iota^{-1}(IA)\) \\
\midrule
\(3\mathbb Z\) & \(3\mathbb Z\) & \(3A\) & \(3\mathbb Z\) \\
\(6\mathbb Z\) & \(6\mathbb Z\) & \(3A\) & \(3\mathbb Z\) \\
\(12\mathbb Z\) & \(12\mathbb Z\) & \(3A\) & \(3\mathbb Z\) \\
\bottomrule
\end{tabular}
\end{center}
\end{samepage}
第二列把整数视为 \(A\) 的元素。例如 \(6\in\iota(6\mathbb Z)\)，但 \(\dfrac12\cdot6=3\notin\iota(6\mathbb Z)\)，所以像集不是 \(A\) 的理想。三个原理想扩张后都等于 \(3A\)，收缩后也都等于 \(3\mathbb Z\)：在 \(A\) 中，因子 \(2\) 已成为单位。

\ProseParagraphBreak
用\cref{b1:def:ideal-quotient}的理想商记号，\(I:s=\{\,a\in R\mid sa\in I\,\}\)，因而\eqref{b1:eq:localized-ideal-contraction}也可写成
\begin{equation}\label{b1:eq:saturation-colon-union}
\iota^{-1}(S^{-1}I)=\bigcup_{s\in S}(I:s).
\end{equation}
这里的条件是存在一个合适的 \(s\)，并不要求对 \(S\) 中的每个元素都成立。这个并集仍是理想：若两个元素分别属于 \(I:s\) 与 \(I:t\)，则它们同时属于 \(I:st\)，因为 \(st\in S\)；相减与乘以任意环元素便可在这个理想中完成。单个元素的各次幂组成分母集的情形，将在素理想对应之后讨论。

\subsection{素理想的对应}
\label{b1:subsec:1304:02}

素理想恰好使饱和条件变得简单。若 \(\mathfrak p\cap S=\varnothing\)，\(ta\in\mathfrak p\)、\(t\in S\)，则素性与 \(t\notin\mathfrak p\) 迫使 \(a\in\mathfrak p\)。

\begin{theorem}\label{b1:thm:localization-prime-correspondence}
设 \(R\) 为交换含幺环，\(S\) 为 \(R\) 的乘法集，\(\iota:R\rightarrow S^{-1}R\) 为自然同态。则 \(R\) 中与 \(S\) 不相交的素理想，与 \(S^{-1}R\) 的全部素理想之间有保持包含关系的一一对应。两个方向分别由扩张、收缩给出：
\[
\begin{aligned}
\mathfrak p&\longmapsto S^{-1}\mathfrak p,\\
\mathfrak q&\longmapsto\iota^{-1}(\mathfrak q).
\end{aligned}
\]
\end{theorem}
\begin{proof}
先设 \(\mathfrak p\) 为与 \(S\) 不相交的素理想。上一段已说明它饱和，而\cref{b1:prop:localization-ideal-correspondence}说明 \(S^{-1}\mathfrak p\) 为真理想。若
\((a/s)(b/t)=ab/(st)\in S^{-1}\mathfrak p\)，由\eqref{b1:eq:localized-ideal-membership}，有 \(u\in S\) 使 \(uab\in\mathfrak p\)。先用 \(u\notin\mathfrak p\) 推出 \(ab\in\mathfrak p\)，再用素性推出 \(a\in\mathfrak p\) 或 \(b\in\mathfrak p\)。于是 \(a/s\) 或 \(b/t\) 属于 \(S^{-1}\mathfrak p\)，证明它为素理想。

反过来，若 \(\mathfrak q\) 为 \(S^{-1}R\) 的素理想，则其原像 \(\mathfrak p=\iota^{-1}(\mathfrak q)\) 是理想，不含 \(1\)。若 \(ab\in\mathfrak p\)，则 \((a/1)(b/1)\in\mathfrak q\)，从而 \(a/1\) 或 \(b/1\) 属于 \(\mathfrak q\)，即 \(a\) 或 \(b\) 属于 \(\mathfrak p\)，所以 \(\mathfrak p\) 素。每个 \(s/1\)、\(s\in S\) 都可逆，不能属于真理想 \(\mathfrak q\)，故 \(\mathfrak p\cap S=\varnothing\)。

对前一方向，饱和性保证收缩恢复 \(\mathfrak p\)；对后一方向，\cref{b1:prop:localization-ideal-correspondence}保证收缩后的扩张恢复 \(\mathfrak q\)。因此两种操作互逆。它们都保持包含关系。
\end{proof}

例如 \(R_{\mathfrak p}\) 中的素理想对应于 \(R\) 中包含于 \(\mathfrak p\) 的素理想，因为
\(\mathfrak q\cap(R\setminus\mathfrak p)=\varnothing\) 等价于 \(\mathfrak q\subseteq\mathfrak p\)。在 \(\mathbb Z_{(p)}\) 中，只剩零理想和 \(p\mathbb Z_{(p)}\) 两个素理想：整数环中的非零素理想是各 \((q)\)，其中只有 \((p)\) 包含于 \((p)\)。在 \(R_f\) 中，对应的则是满足 \(f\notin\mathfrak p\) 的素理想；因为素理想包含 \(f\) 的某个幂当且仅当包含 \(f\)。

\ProseParagraphBreak
这里的对应针对素理想。局部化中一个极大理想的收缩未必在原环中极大。例如整环 \(\mathbb Z\) 在全部非零元素处局部化为 \(\mathbb Q\)，后者的极大理想 \((0)\) 收缩为 \(\mathbb Z\) 的零理想，却不是 \(\mathbb Z\) 的极大理想。极大性在此应理解为与分母集不相交的素理想之中的极大性。

\subsection{单元素饱和与有限步稳定}
\label{b1:subsec:1304:04}

把分母限制为同一个元素的各次幂，饱和便可写成逐次理想商的并。

\begin{definition}\label{b1:def:element-saturation}
设 \(I\) 为交换含幺环 \(R\) 的理想，\(f\in R\)。定义
\begin{equation}\label{b1:eq:element-saturation}
I:f^\infty\coloneqq\bigcup_{n\geq0}(I:f^n)
=\{\,a\in R\mid f^na\in I\text{ 对某个整数 }n\geq0\,\}.
\end{equation}
这个集合称为 \(I\) 关于 \(f\) 的饱和；这里 \(f^0=1\)，所以 \(I:f^0=I\)。若 \(I:f^\infty=I\)，就称 \(I\) 关于 \(f\) 饱和，简称 \(f\)-饱和。
\end{definition}
\symidx{\(I:f^\infty\)：理想关于元素 \(f\) 的饱和}
\termidx[baohe-lixiang]{饱和理想}

上标 \(\infty\) 表示允许取某个足够大的有限次幂，不是在环中定义元素 \(f^\infty\)。\(I:f\) 只允许乘一次 \(f\)，一般还不是 \(I:f^\infty\)。

\begin{proposition}\label{b1:prop:element-saturation}
设 \(I\) 为交换含幺环 \(R\) 的理想，\(f\in R\)，并令 \(R_f=\{\,1,f,f^2,\ldots\,\}^{-1}R\)，以 \(\iota_f:R\rightarrow R_f\) 表示自然同态。则 \(I:f^\infty\) 是理想，并满足
\begin{equation}\label{b1:eq:element-saturation-contraction}
I:f^\infty=\iota_f^{-1}(IR_f).
\end{equation}
它是包含 \(I\) 的最小 \(f\)-饱和理想。若 \(R\) 为 Noether 环，则存在整数 \(N\geq0\)，使
\begin{equation}\label{b1:eq:element-saturation-stable}
I:f^n=I:f^N=I:f^\infty\qquad(n\geq N).
\end{equation}
\end{proposition}
\begin{proof}
在\eqref{b1:eq:saturation-colon-union}中取 \(S=\{\,f^n\mid n\geq0\,\}\)，便得到\eqref{b1:eq:element-saturation-contraction}；理想的收缩仍为理想，故 \(I:f^\infty\) 是理想。记它为 \(K\)。若 \(f^ma\in K\)，则存在 \(n\geq0\) 使 \(f^{n+m}a\in I\)，所以 \(a\in K\)。因此 \(K:f^\infty=K\)，且取 \(n=0\) 可知 \(I\subseteq K\)。

若 \(L\) 为任意包含 \(I\) 的 \(f\)-饱和理想，对 \(a\in K\) 取 \(n\geq0\) 使 \(f^na\in I\subseteq L\)，则 \(L\) 的饱和性给出 \(a\in L\)。所以 \(K\subseteq L\)，这就证明了最小性。

由于 \(f^na\in I\) 蕴含 \(f^{n+1}a\in I\)，这些理想商构成升链
\[
I=I:f^0\subseteq I:f\subseteq I:f^2\subseteq\cdots.
\]
若 \(R\) 为 Noether 环，\cref{b1:prop:noetherian-acc}说明这条理想升链从某个 \(N\) 起稳定。它的并集因此等于 \(I:f^N\)，即得\eqref{b1:eq:element-saturation-stable}。
\end{proof}

\begin{samepage}
\eqref{b1:eq:element-saturation-contraction}不要求 Noether 性；只有最后的有限步稳定性使用了这一假设。一旦得到 \(I:f^N=I:f^{N+1}\)，以后各项也都相等。依据是\cref{b1:prop:ideal-quotient-identities}的逐次理想商公式
\[
\bigl(I:f^n\bigr):f=I:f^{n+1}.
\]
对相邻两项的等式再取一次关于 \(f\) 的理想商，便得到下一对相邻项相等，归纳即可。
\end{samepage}

要把这个过程用作算法，还须能够计算理想商，并判定相邻两项是否相等；Noether 性保证有限步后稳定，但不单独提供这些计算方法。\cref{b1:ex:13-14}还说明，在固定的多项式环及固定的 \(f\) 下，稳定下标也没有适用于所有理想的统一上界。

\ProseParagraphBreak
两个极端情形也与定义相容。若 \(f\) 已是单位，\(f^na\in I\) 就能乘以 \(f^{-n}\) 得到 \(a\in I\)，因此 \(I:f^\infty=I\)。若 \(f\) 幂零，某个 \(f^n=0\) 把每个元素都送入 \(I\)，所以 \(I:f^\infty=R\)；这也对应于 \(R_f\) 为零环的情形。

\ProseParagraphBreak
例如在 \(R=k[x,y]\) 中取 \(I=(x^2,xy)\)，其中 \(k\) 为域。若 \(xg\in I\)，就有 \(xg=x^2u+xyv\)；多项式环是整环，约去 \(x\) 得 \(g=xu+yv\)。反过来，后一表达式乘以 \(x\) 显然属于 \(I\)。因此
\[
I:x=(x,y),\qquad I:x^2=R,\qquad I:x^\infty=R.
\]
可见只取一次理想商与取饱和确实可能不同。另一方面，\(xy\in I\) 给出 \((x)\subseteq I:y^\infty\)。若 \(y^ng\in I\subseteq(x)\)，令 \(x=0\) 得 \(y^n\cdot g(0,y)=0\)；由于 \(k[y]\) 是整环，\(g(0,y)=0\)，即 \(g\) 的每个非零单项式都含 \(x\)，所以 \(g\in(x)\)。故
\[
I:y^\infty=(x).
\]
在 \(R_x\) 中，\(x^2\) 已可逆，\(IR_x=R_x\)；在 \(R_y\) 中，\(xy\) 与 \(x\) 生成同一主理想，且 \(x^2\) 已被 \(x\) 吸收，所以 \(IR_y=(x)R_y\)。这也从局部化说明了两种饱和为什么不同。关于由多个元素生成的理想的饱和，可进一步参看\cref{b1:ex:13-15}；它的条件与允许这些元素同时作为分母有所区别。

\subsection{元素与理想的局部—整体判据}
\label{b1:subsec:1304:03}

局部化可能使一个非零元素消失，但所有极大理想处的局部化合在一起，能够检测原来的元素。这是局部—整体思想的一个初等形式。

\begin{proposition}\label{b1:prop:local-detection}
设 \(R\) 为非零交换含幺环，\(I,J\) 为 \(R\) 的理想，\(a,b\in R\)。则
\[
a\in I
\quad\Longleftrightarrow\quad
a/1\in IR_{\mathfrak m}
\text{ 对每个极大理想 }\mathfrak m\text{ 成立}.
\]
此外，下列判据成立，其中 \(\mathfrak m\) 遍历 \(R\) 的全部极大理想：
\[
\begin{aligned}
a=0
&\quad\Longleftrightarrow\quad
a/1=0\text{ 在每个 }R_{\mathfrak m}\text{ 中成立},\\
a=b
&\quad\Longleftrightarrow\quad
a/1=b/1\text{ 在每个 }R_{\mathfrak m}\text{ 中成立},\\
I=J
&\quad\Longleftrightarrow\quad
IR_{\mathfrak m}=JR_{\mathfrak m}\text{ 对每个 }\mathfrak m\text{ 成立}.
\end{aligned}
\]
\end{proposition}
\begin{proof}
若 \(a\in I\)，则对每个 \(\mathfrak m\)，其像当然在扩张理想中。反过来，设 \(a\notin I\)，考虑\cref{b1:def:ideal-quotient}中的理想商
\[
K=I:a=\{\,r\in R\mid ra\in I\,\}.
\]
\cref{b1:prop:ideal-quotient-rules}保证 \(K\) 是理想。因为 \(a\notin I\)，\(1\notin K\)，所以 \(K\) 为真理想。由\cref{b1:prop:maximal-ideal-exists}，存在极大理想 \(\mathfrak m\) 包含 \(K\)。\cref{b1:prop:prime-max-quotient}又说明 \(\mathfrak m\) 是素理想，因此可作局部化的中心。

如果 \(a/1\in IR_{\mathfrak m}\)，\eqref{b1:eq:localized-ideal-membership}便给出 \(s\notin\mathfrak m\) 使 \(sa\in I\)。但这意味着 \(s\in K\subseteq\mathfrak m\)，矛盾。因此在这个极大理想处，\(a/1\) 不属于 \(IR_{\mathfrak m}\)。

取 \(I=(0)\) 即得元素为零的判据；把 \(a\) 换成两个元素之差，得到相等的判据。若两个理想在所有 \(R_{\mathfrak m}\) 中的扩张相等，则对第一个理想的任意元素应用刚证明的成员判据，知它属于第二个理想；交换二者后得到反向包含，故两个理想相等。
\end{proof}

回到前面的计算，设 \(k\) 为域，\(R=k[x,y]\)、\(I=(x^2,xy)\)、\(a=x\)，并令 \(\mathfrak m_0=(x,y)\)。由 \(I:x=\mathfrak m_0\) 可知 \(x\notin I\)。若 \(x/1\in IR_{\mathfrak m_0}\)，成员判据给出某个 \(s\notin\mathfrak m_0\) 使 \(sx\in I\)；这又意味着 \(s\in I:x=\mathfrak m_0\)，矛盾。对任意极大理想 \(\mathfrak m\ne\mathfrak m_0\)，必有 \(x\notin\mathfrak m\) 或 \(y\notin\mathfrak m\)。在前一种情形中，\(x^2/1\) 可逆，故 \(IR_{\mathfrak m}=R_{\mathfrak m}\)；在后一种情形中，\(xy\in I\) 且 \(y/1\) 可逆，故 \(x/1\in IR_{\mathfrak m}\)。因此
\[
x/1\notin IR_{\mathfrak m}
\quad\Longleftrightarrow\quad
\mathfrak m=\mathfrak m_0.
\]
这个成员条件只在 \(\mathfrak m_0\) 处失败；证明中使用的理想商 \(I:a\) 在此恰好找到了检出失败的位置。

\ProseParagraphBreak
在上述判据中，必须检查元素在局部环里的像；再取剩余域可能丢失信息。设 \(k\) 为域，\(R=k[t]/(t^2)\)，以 \(\varepsilon\) 表示 \(t\) 在商环中的像，并令 \(\mathfrak m=(\varepsilon)\)。这是 \(R\) 的唯一极大理想，且其外的元素均为单位，故 \(R_{\mathfrak m}\cong R\)。映射链
\[
R\longrightarrow R_{\mathfrak m}
\longrightarrow R_{\mathfrak m}/\mathfrak mR_{\mathfrak m}\cong k
\]
将非零元素 \(\varepsilon\) 先送到局部环中的非零元素，再送到剩余域中的零。因此，“在每个局部环中为零”不能换成“在每个剩余域中为零”。

\begin{corollary}\label{b1:cor:local-unit-detection}
设 \(R\) 为非零交换含幺环，\(a\in R\)。则 \(a\) 为单位，当且仅当它在每个 \(R_{\mathfrak m}\) 中的像均为单位，其中 \(\mathfrak m\) 遍历 \(R\) 的极大理想。
\end{corollary}
\begin{proof}
单位在环同态下仍是单位，故必要性成立。若每个 \(a/1\) 都是单位，则 \(1/1\) 属于由 \(a/1\) 生成的理想，即 \((a)R_{\mathfrak m}\)。对\cref{b1:prop:local-detection}取 \(I=(a)\)、待检测元素为 \(1\)，便得到 \(1\in(a)\)。于是 \(ab=1\) 对某个 \(b\in R\) 成立，故 \(a\) 是单位。
\end{proof}

第三卷第2章将把分母作用推广到模，并在建立模的局部化后研究更一般的局部—整体判据。本章的零判据和理想成员判据将推广到模；相应的定义与证明在第三卷给出。


\begin{chaptersummary}
分式域使整环中所有非零元素可逆；局部化则允许指定一个乘法集。一般环中的等式必须容许再乘一个分母集元素，才能适应零因子。零判据、单位判据和同态延拓的泛性质，分别说明哪些元素变为零、哪些元素成为单位，以及如何把局部化映入其他环。

\ProseParagraphBreak
局部化中的全部理想都来自原环，但扩张后再收缩得到的是原理想的饱和；与分母集不相交的素理想已经饱和，因而与局部化的素理想一一对应。若原环为 Noether 环，局部化后的理想仍有限生成。在单个元素 \(f\) 处，饱和总可写为 \(I:f^\infty=\bigcup_{n\geq0}(I:f^n)\)；若原环为 Noether 环，相应升链最终稳定，从而对这个饱和理想中的所有元素可取同一个有限指数。商环使指定理想中的元素成为零，局部化则使指定乘法集中的元素成为单位。

\ProseParagraphBreak
在一个局部环中成立，尚不能保证全局成立。所有极大理想处的局部化合在一起，能够检测元素相等、理想成员以及可逆性；理想商可以指出某个成员条件失败的具体位置。若 \(f,g\) 生成单位理想，在 \(R_f\) 和 \(R_g\) 中检查理想成员条件也已足够。

\ProseParagraphBreak
第二编关于环、整除与多项式的讨论至此结束。下一编转向域扩张，研究扩大系数域以后，多项式的根与分解如何变化。模的局部化、正合性及更系统的局部性质将在第三卷讨论。
\end{chaptersummary}

\BookExercises

\begin{exercise}\label{b1:ex:13-1}
设 \(k\) 为域，\(R=k[t^2,t^3]\) 为 \(k[t]\) 中由 \(k,t^2,t^3\) 生成的子环。证明 \(t\notin R\)，但 \(\operatorname{Frac}(R)\cong k(t)\)，且可使这一同构保持 \(R\) 中每个元素。
\end{exercise}
\begin{solution}
\(R\) 中每个元素是 \(t^2,t^3\) 的多项式，因此出现的指数都为 \(2a+3b\)，其中 \(a,b\geq0\)。这样的指数不可能等于 \(1\)，所以每个元素的 \(t\) 项系数为零，\(t\notin R\)。

将 \(R\) 视为 \(k(t)\) 的子环。由\cref{b1:prop:fraction-in-ambient-field}，\(\operatorname{Frac}(R)\) 可同构到 \(k(t)\) 中由 \(R\) 生成的最小子域。这个子域包含 \(t=t^3/t^2\)，也包含 \(k\)，故包含 \(k[t]\) 的全部非零分式，也就是 \(k(t)\)。反向包含原本就成立，因而该子域恰为 \(k(t)\)。同构由 \(a/b\mapsto ab^{-1}\) 给出，保持 \(R\) 中元素。
\end{solution}

\begin{exercise}\label{b1:ex:13-9}
设 \(k\) 为域，\(t\) 为多项式变量。证明替换映射
\[
k[t]\longrightarrow k(t),\qquad f(t)\longmapsto f(t^2)
\]
唯一地延拓为域同态 \(\sigma:k(t)\rightarrow k(t)\)。求它在任意分式上的公式，并证明 \(\sigma\) 单射但不满射；结论对特征为 \(2\) 的域也成立。
\end{exercise}
\begin{solution}
多项式替换保持加法、乘法及单位。若 \(f\neq0\)，其最高次非零项 \(a_dt^d\) 变为 \(a_dt^{2d}\)，不会与其他项合并，所以 \(f(t^2)\neq0\)。故给定同态把所有非零多项式送到域 \(k(t)\) 的单位。\cref{b1:thm:fraction-field-universal}给出唯一延拓，公式为
\[
\sigma\left(\frac{f(t)}{g(t)}\right)
=\frac{f(t^2)}{g(t^2)},\qquad g\neq0.
\]
若像为零，则 \(f(t^2)=0\)，由刚才的最高项论证知 \(f=0\)。因此 \(\sigma\) 单射。

若 \(t\) 在其像中，就存在 \(f,g\in k[t]\)、\(g\neq0\)，使 \(t=f(t^2)/g(t^2)\)。分式相等的判据\eqref{b1:eq:fraction-equivalence}给出多项式等式
\[
t\,g(t^2)=f(t^2).
\]
左边只含奇数次幂，且不为零；右边只含偶数次幂。同次幂系数分别相等便迫使 \(g\) 的全部系数为零，矛盾。所以 \(t\) 不在像中。这个论证区分的是形式多项式的指数奇偶，并未把整数 \(2\) 当作系数作除法，因而在特征 \(2\) 时仍然成立。
\end{solution}

\begin{exercise}\label{b1:ex:13-13}
令 \(s=\sqrt{-5}\)，\(R=\mathbb Z[s]\)。证明在 \(\operatorname{Frac}(R)\) 中有
\[
\frac{2}{1+s}=\frac{1-s}{3}.
\]
证明两边各自的分子、分母都只有单位作为公因子，但不存在单位 \(u\in R^\times\) 能把一组分子、分母同时乘为另一组。这说明把分式化为“分子、分母无非单位公因子”的形式后，是否一定具有只差同乘单位的唯一性？
\end{exercise}
\begin{solution}
两个分母非零，而且 \(2\cdot3=6=(1+s)(1-s)\)，所以\eqref{b1:eq:fraction-equivalence}直接给出两个分式相等。

在 \(R\) 中取范数 \(N(a+bs)=a^2+5b^2\)。共轭相乘的表达式说明它对乘法相容；非零元素的范数为正整数，范数为 \(1\) 的元素恰为单位 \(\pm1\)。若 \(d\) 同时整除 \(2\) 和 \(1+s\)，则 \(N(d)\) 同时整除 \(4\) 和 \(6\)，所以只能为 \(1\) 或 \(2\)。但方程 \(a^2+5b^2=2\) 无整数解：\(b\neq0\) 时左边至少为 \(5\)，\(b=0\) 时又要求整数平方为 \(2\)。故 \(N(d)=1\)，\(d\) 为单位。

同样，\(1-s\) 和 \(3\) 的公因子范数同时整除 \(6\) 和 \(9\)，只能为 \(1\) 或 \(3\)；而 \(a^2+5b^2=3\) 也无整数解，故公因子仍只能是单位。若能同乘单位 \(u\) 把第一组分子、分母变成第二组，必有 \(2u=1-s\)。两边取范数却得到 \(4=6\)，矛盾。因此在一般整环中，即使分子、分母已经没有非单位公因子，这样的分式表示仍不一定只差同乘单位。
\end{solution}

\begin{exercise}\label{b1:ex:13-2}
在 \(R=\mathbb Z/12\mathbb Z\) 中取 \(S=\{\,\bar3^n\mid n\geq0\,\}\)。证明 \(S^{-1}R\cong\mathbb Z/4\mathbb Z\)，并求 \(R\rightarrow S^{-1}R\) 的核。
\end{exercise}
\begin{solution}
由\cref{b1:prop:localization-zero-unit}，\(\bar a/1=0\) 当且仅当某个 \(3^na\) 被 \(12\) 整除。这个条件等价于 \(4\mid a\)：必要性来自 \(3^n\) 与 \(4\) 互素，充分性取 \(n=1\)。因此核为 \((\bar4)=\{\,\bar0,\bar4,\bar8\,\}\)。

还需确认自然映射的像就是整个局部化。因为 \(\bar8/1=0\)，局部化中有 \((\bar3/1)^2=\bar9/1=1\)，所以 \(1/\bar3=\bar3/1\)，进而每个分式 \(\bar a/\bar3^n\) 都等于 \(\overline{a3^n}/1\)。自然映射因而满射。\cref{b1:thm:ring-first-iso}给出
\[
S^{-1}R\cong(\mathbb Z/12\mathbb Z)/(\bar4)
\cong\mathbb Z/4\mathbb Z.
\]
最后一个同构由模 \(4\) 的映射诱导；它在模 \(12\) 环上的核也恰为 \((\bar4)\)。
\end{solution}

\begin{exercise}\label{b1:ex:13-3}
设 \(S\) 为交换环 \(R\) 的乘法集，令
\[
T=\{\,a\in R\mid ab\in S\text{ 对某个 }b\in R\,\}.
\]
正文已说明 \(T\) 是包含 \(S\) 的乘法集，且自然映射 \(S^{-1}R\rightarrow T^{-1}R\) 为同构。证明 \(T\) 还满足：若 \(ab\in T\)，则 \(a,b\in T\)；并证明它是包含 \(S\) 且满足这一性质的最小乘法集。
\end{exercise}
\begin{solution}
若 \(ab\in T\)，取 \(c\in R\) 使 \(abc\in S\)，则 \(a(bc)\in S\)、\(b(ac)\in S\)，所以 \(a,b\in T\)。这证明了所要求的乘积性质。

设 \(U\) 是另一个包含 \(S\) 且满足同一性质的乘法集。对任意 \(a\in T\)，可取 \(b\in R\) 使 \(ab\in S\subseteq U\)。由 \(U\) 的乘积性质，\(a\in U\)，所以 \(T\subseteq U\)，证明了最小性。因而 \(T\) 不仅收集了所有像已可逆的元素，还是使分母集对乘积因子的提取封闭所需的最小扩大。
\end{solution}

\begin{exercise}\label{b1:ex:13-4}
设 \(k\) 为域，\(R=k[\varepsilon]/(\varepsilon^2)\)，并仍以 \(\varepsilon\) 表示其剩余类。证明 \(R\) 是局部环，求其单位与剩余域，并比较 \(R_{(\varepsilon)}\) 和 \(R_{\varepsilon}\)。
\end{exercise}
\begin{solution}
每个元素唯一写成 \(a+b\varepsilon\)，其中 \(a,b\in k\)。若 \(a\neq0\)，则
\[
(a+b\varepsilon)(a^{-1}-ba^{-2}\varepsilon)=1,
\]
因为 \(\varepsilon^2=0\)。若 \(a=0\)，则 \(b\varepsilon\) 平方为零，不可能在非零环 \(R\) 中可逆。全部非单位因而组成理想 \((\varepsilon)\)。任何真理想都只含非单位，所以都包含于 \((\varepsilon)\)；这说明它是唯一极大理想。映射 \(a+b\varepsilon\mapsto a\) 满射到 \(k\)，核为 \((\varepsilon)\)，故剩余域为 \(k\)。

在 \((\varepsilon)\) 处局部化只把原环的单位作为分母，逆元已经存在。由\cref{b1:thm:localization-universal}，\(a/s\mapsto as^{-1}\) 与自然映射互逆，所以 \(R_{(\varepsilon)}\cong R\)。单个元素 \(\varepsilon\) 处的局部化则为零环，因为 \(\varepsilon\) 幂零，适用\cref{b1:prop:single-element-localization}。下标是否表示理想，在这个例子中决定了完全不同的结果。
\end{solution}

\begin{exercise}\label{b1:ex:13-5}
设 \(e\) 为交换含幺环 \(R\) 的幂等元，即 \(e^2=e\)。证明自然映射 \(R\rightarrow R_e\) 满射，核为 \((1-e)\)，从而 \(R_e\cong R/(1-e)\)。
\end{exercise}
\begin{solution}
在 \(R_e\) 中，\(e/1\) 是可逆幂等元。把 \((e/1)^2=e/1\) 乘以其逆元，得 \(e/1=1\)。因此 \(a/e^n=a/1\)，自然映射满射。

由\cref{b1:prop:localization-zero-unit}，\(a/1=0\) 等价于存在 \(n\geq0\) 使 \(e^na=0\)。若 \(n=0\)，则 \(a=0\)，当然有 \(ea=0\)；若 \(n\geq1\)，\(e^n=e\)，所以条件仍等价于 \(ea=0\)。而
\[
ea=0\quad\Longleftrightarrow\quad a=(1-e)a
\quad\Longleftrightarrow\quad a\in(1-e).
\]
最后一个反向蕴含用到 \(e(1-e)=0\)。由\cref{b1:thm:ring-first-iso}即得所需同构。\(e=0\) 时两边均为零环，\(e=1\) 时两边均为 \(R\)，也包括在内。

这与\cref{b1:prop:idempotent-product}中的正交幂等元分解相合：\(R\cong eR\times(1-e)R\)，其中 \(eR\) 的单位元为 \(e\)。投影 \(a\mapsto ea\) 的核为 \((1-e)R\)，故 \(R_e\cong eR\)。也就是说，把 \(e\) 变成单位恰好保留 \(eR\) 这一分量，并消去另一个分量。
\end{solution}

\begin{exercise}\label{b1:ex:13-6}
设 \(p\) 为素数。证明 \(\mathbb Z_{(p)}\) 的全部理想为零理想和 \(p^n\mathbb Z_{(p)}\)，其中 \(n=0,1,2,\ldots\)，并证明这些非零理想两两不同。
\end{exercise}
\begin{solution}
设 \(J\neq(0)\) 为 \(\mathbb Z_{(p)}\) 的理想。取非零分式 \(a/b\in J\)，乘以单位 \(b/1\) 得非零整数 \(a/1\in J\)，所以收缩 \(I\subseteq\mathbb Z\) 非零。整数带余除法与\cref{b1:thm:euclidean-pid}已给出 \(\mathbb Z\) 的主理想性，故 \(I=d\mathbb Z\)，可取 \(d>0\)。

\cref{b1:prop:localization-ideal-correspondence}说明 \(J=I\mathbb Z_{(p)}=d\mathbb Z_{(p)}\)。将 \(d\) 写成 \(p^nu\)，其中 \(p\nmid u\)，则 \(u\) 在 \(\mathbb Z_{(p)}\) 中为单位，因此 \(J=p^n\mathbb Z_{(p)}\)。反过来，这些集合当然都是理想。若 \(m>n\) 而 \(p^n\in p^m\mathbb Z_{(p)}\)，则存在 \(a,b\in\mathbb Z\)、\(p\nmid b\)，使 \(p^n=p^ma/b\)，从而 \(b=p^{m-n}a\)，与 \(p\nmid b\) 矛盾，故它们两两不同。
\end{solution}

\begin{exercise}\label{b1:ex:13-10}
设整数 \(n\geq2\)。求存在保持单位的环同态
\[
\mathbb Z[1/6]\longrightarrow\mathbb Z/n\mathbb Z
\]
的充要条件；存在时证明它唯一，并写出其作用于任意元素的公式。当 \(n=5\) 时，分别计算 \(1/2\)、\(1/3\) 的像。
\end{exercise}
\begin{solution}
任何保持单位的同态在 \(\mathbb Z\) 上都必须是 \(a\mapsto[a]_n\)。由于 \(6\) 在定义域中可逆，其像 \([6]_n\) 必须在目标中可逆。这个条件等价于 \(\gcd(6,n)=1\)：若 \(6u\equiv1\pmod n\)，则 \(6u-nv=1\) 对某个 \(v\) 成立，所以公因子只能为 \(1\)；反过来，互素时整数 Bézout 等式给出这样的 \(u,v\)。

当 \(\gcd(6,n)=1\) 时，商映射 \(\mathbb Z\rightarrow\mathbb Z/n\mathbb Z\) 将乘法集 \(\{\,6^r\mid r\geq0\,\}\) 送到单位。\cref{b1:thm:localization-universal}给出唯一延拓
\[
\varphi(a/6^r)=[a]_n[6]_n^{-r}.
\]
该泛性质同时保证公式不依赖分式代表。\(n=5\) 时，\([2]_5[3]_5=[1]_5\)，故 \(\varphi(1/2)=[3]_5\)、\(\varphi(1/3)=[2]_5\)。这里 \(1/2=3/6\)、\(1/3=2/6\)，两者确实都在 \(\mathbb Z[1/6]\) 中。
\end{solution}

\begin{exercise}\label{b1:ex:13-11}
设 \(k\) 为域，\(A=k[t,t^{-1}]\)。求 \(A\) 的所有逐点固定 \(k\) 的环自同态，判断其中哪些单射、哪些为自同构。
\end{exercise}
\begin{solution}
设 \(\varphi:A\rightarrow A\) 为这样的同态。因为 \(t\) 是单位，\(\varphi(t)\) 也是单位。\cref{b1:prop:laurent-units}说明必有
\(\varphi(t)=ct^m\)，其中 \(c\in k^\times\)、\(m\in\mathbb Z\)。又因为 \(\varphi\) 固定系数，且必须将 \(t^{-1}\) 送到 \(c^{-1}t^{-m}\)，它在所有有限 Laurent 和上的值只能是
\[
\varphi_{c,m}\Bigl(\sum_{j=r}^{s}a_jt^j\Bigr)
=\sum_{j=r}^{s}a_jc^jt^{mj}.
\]
反过来，对每一对 \(c,m\)，替换 \(t\mapsto ct^m\) 给出环同态 \(k[t]\rightarrow A\)，并将 \(t\) 的所有非负次幂送到单位。\cref{b1:thm:localization-universal}使它唯一延拓为 \(A=k[t]_t\) 的自同态，恰好具有上式。因此这些公式穷尽全部所求同态。

若 \(m=0\)，非零元素 \(t-c\) 被送到零，同态不单射。若 \(m\neq0\)，不同整数 \(j\) 对应不同的指数 \(mj\)，而 \(c^j\neq0\)。若像为零，给等式两边乘以足够高次的 \(t\)，便得到普通多项式等式，其中不同 \(j\) 对应的项仍有不同次数。比较系数给出每个 \(a_jc^j=0\)，从而 \(a_j=0\)，所以此时同态单射。

当 \(|m|>1\) 时，像中所有非零项的指数都被 \(m\) 整除，故 \(t\) 不在像中，不能满射。\(m=0\) 时像为 \(k\)，也不能满射。\(m=1\) 时，\(\varphi_{c^{-1},1}\) 是逆映射；\(m=-1\) 时，\(\varphi_{c,-1}\) 自己就是逆映射，因为它将 \(t\) 连续两次变换为 \(c(ct^{-1})^{-1}=t\)，并固定 \(k\)。所以自同构恰好对应 \(m=1\) 或 \(m=-1\)。
\end{solution}

\begin{exercise}\label{b1:ex:13-7}
设 \(I\) 为交换环 \(R\) 的理想，\(S\) 为乘法集，\(\bar S\) 为 \(S\) 在 \(R/I\) 中的像。证明
\[
(S^{-1}R)/(S^{-1}I)\cong\bar S^{-1}(R/I),
\]
其中同构由 \(a/s\mapsto\bar a/\bar s\) 诱导。说明 \(I\cap S\neq\varnothing\) 时等式两边是什么环。
\end{exercise}
\begin{solution}
复合 \(R\rightarrow R/I\rightarrow\bar S^{-1}(R/I)\) 把 \(S\) 中每个元素送到单位，故由\cref{b1:thm:localization-universal}延拓为所给同态
\(\varphi:S^{-1}R\rightarrow\bar S^{-1}(R/I)\)。目标中的每个分式都能选取分子、分母的原像，所以 \(\varphi\) 满射。

由\cref{b1:prop:localization-zero-unit}，\(\varphi(a/s)=0\) 当且仅当存在 \(t\in S\) 使 \(\bar t\bar a=0\)，也就是 \(ta\in I\)。\eqref{b1:eq:localized-ideal-membership}说明这又等价于 \(a/s\in S^{-1}I\)。故 \(\ker\varphi=S^{-1}I\)，由\cref{b1:thm:ring-first-iso}得到商环同构。若 \(I\cap S\neq\varnothing\)，则 \(S^{-1}I=S^{-1}R\)，左端是零环；此时 \(0\in\bar S\)，右端也由同一零判据成为零环。
\end{solution}

\begin{exercise}\label{b1:ex:13-localized-finite-intersection}
设 \(R\) 为交换含幺环，\(S\subseteq R\) 为乘法集，\(I,J\subseteq R\) 为理想。证明
\[
S^{-1}(I\cap J)=S^{-1}I\cap S^{-1}J.
\]
举例说明，把有限交改为无限多个理想的交后，等式未必成立。
\end{exercise}
\begin{solution}
左边包含于右边由理想包含直接得到。反过来，若 \(a/s\in S^{-1}I\cap S^{-1}J\)，则\eqref{b1:eq:localized-ideal-membership}分别给出 \(t,u\in S\)，使 \(ta\in I\)、\(ua\in J\)。于是 \(tua\in I\cap J\)；由于 \(tu\in S\)，同一成员判据说明 \(a/s\in S^{-1}(I\cap J)\)。这里能把两个条件合并，是因为只需相乘有限个分母。

对无限交，取 \(R=\mathbb Z\)、\(S=\{\,2^m\mid m\geq0\,\}\)、\(I_n=2^n\mathbb Z\)（\(n\geq1\)）。非零整数不可能被任意高次的 \(2\) 整除，故 \(\bigcap_{n\geq1}I_n=0\)，而 \(S^{-1}0=0\)。另一方面，\(2^n\) 在 \(\mathbb Z[1/2]=S^{-1}\mathbb Z\) 中是单位，所以对每个 \(n\) 都有 \(S^{-1}I_n=\mathbb Z[1/2]\)。因此
\[
S^{-1}\!\left(\bigcap_{n\geq1}I_n\right)=0,
\qquad
\bigcap_{n\geq1}S^{-1}I_n=\mathbb Z[1/2].
\]
即使 \(\mathbb Z\) 为 Noether 环，也不能从无限多个局部成员条件中直接选出共同分母。
\end{solution}

\begin{exercise}\label{b1:ex:13-12}
设 \(k\) 为域，\(f=x(x-1)\)，\(A=k[x]_f\)。求 \(A\) 的全部素理想，并求每个非零素理想的商环。说明原来 \(k[x]\) 中的素理想 \((x)\)、\((x-1)\) 为什么没有留在这一对应中。
\end{exercise}
\begin{solution}
由\cref{b1:prop:field-poly-division}和\cref{b1:thm:euclidean-pid}，\(k[x]\) 是主理想整环；\cref{b1:prop:pid-irreducible-prime}说明其非零素理想恰由不可约多项式生成，且这些理想都极大。把生成元取为首一，可使表示唯一。因此 \(k[x]\) 的全部素理想为 \((0)\) 和各 \((q)\)，其中 \(q\) 为首一不可约多项式。

乘法集 \(S=\{\,f^r\mid r\geq0\,\}\) 不含零，故 \((0)\) 与它不相交。对 \((q)\)，素性说明它与 \(S\) 相交当且仅当 \(q\mid f\)；这又当且仅当 \(q=x\) 或 \(q=x-1\)，因为不可约元在 \(k[x]\) 中为素元。由\cref{b1:thm:localization-prime-correspondence}，\(A\) 的素理想于是恰为
\[
(0),\qquad qA
\quad\text{其中 }q\text{ 首一不可约，且 }q\neq x,\ x-1.
\]
它们互不相同，因为收缩分别恢复 \((0)\) 和 \((q)\)。

固定一个这样的 \(q\)。在域 \(k[x]/(q)\) 中，\(f\) 的像不为零，因而可逆。由\cref{b1:thm:localization-universal}，商映射唯一延拓为
\[
A\longrightarrow k[x]/(q),\qquad
g/f^r\longmapsto\bar g\bar f^{-r}.
\]
它满射，因为每个剩余类都有 \(g/1\) 的像作为代表；其核由 \(q\mid g\) 的分式组成，正是 \(qA\)。\cref{b1:thm:ring-first-iso}给出 \(A/qA\cong k[x]/(q)\)，特别地每个非零素理想都是极大理想。至于 \((x)\) 与 \((x-1)\)，它们都含有 \(f\in S\)，在局部化中的扩张均为全环，故不能对应真素理想。
\end{solution}

\begin{exercise}\label{b1:ex:13-14}
设 \(k\) 为域，\(R=k[x,y]\)，整数 \(m\geq1\)，\(I=(x^my)\)。计算所有 \(I:x^n\)、\(n\geq0\)，求 \(I:x^\infty\)，并确定升链 \(I\subseteq I:x\subseteq I:x^2\subseteq\cdots\) 开始稳定的最小下标。再求 \(I:y^\infty\) 和 \(I:(xy)^\infty\)，解释为什么不同的分母会产生不同的饱和。
\end{exercise}
\begin{solution}
一个多项式属于 \((x^my)\)，当且仅当它的每个非零单项式都含有至少 \(m\) 个 \(x\) 因子和至少一个 \(y\) 因子：必要性由把 \(x^my\) 乘入多项式看出，充分性则对每一项同时提出这个因子。乘以 \(x^n\) 只把各项中 \(x\) 的指数统一增加 \(n\)，所以
\[
I:x^n=
\begin{cases}
(x^{m-n}y),&0\leq n<m,\\
(y),&n\geq m.
\end{cases}
\]
于是 \(I:x^\infty=(y)\)。当 \(0\leq n<m\) 时，元素 \(x^{m-n-1}y\) 属于 \(I:x^{n+1}\)，却因 \(x\) 的次数不够而不属于 \(I:x^n\)。故此前各次包含都严格，最小稳定下标恰为 \(m\)。这还说明，即使固定环 \(k[x,y]\) 和分母 \(x\)，也不能给所有理想的这类升链规定同一个稳定下标。

乘以 \(y^n\) 不会增加 \(x\) 的指数；一旦 \(n\geq1\)，又已补足所需的 \(y\) 因子。因此 \(I:y^n=(x^m)\) 对所有 \(n\geq1\) 成立，给出 \(I:y^\infty=(x^m)\)。最后，\((xy)^m=x^my^m\in I\)，所以 \(1\in I:(xy)^\infty\)，即 \(I:(xy)^\infty=R\)。在 \(R_x\) 中只有因子 \(x\) 被保证可逆，原来的理想化为 \((y)R_x\)；在 \(R_y\) 中化为 \((x^m)R_y\)；在 \(R_{xy}\) 中，\(x\) 与 \(y\) 都可逆，\(IR_{xy}=R_{xy}\)。这些计算与\eqref{b1:eq:element-saturation-contraction}逐一相符。
\end{solution}

\begin{exercise}\label{b1:ex:13-15}
设 \(I,J\) 为交换含幺环 \(R\) 的理想，约定 \(J^0=R\)。定义
\[
I:J^\infty\coloneqq\bigcup_{n\geq0}(I:J^n)
=\{\,a\in R\mid aJ^n\subseteq I\text{ 对某个整数 }n\geq0\,\}.
\]
这个构造称为\term[lixiang-baohe]{关于理想的饱和}[saturation with respect to an ideal]。
\symidx{\(I:J^\infty\)：理想 \(I\) 关于理想 \(J\) 的饱和}
\begin{enumerate}
\item 证明 \(I:J^\infty\) 是理想。
\item 若 \(J=(f_1,\ldots,f_r)\)、\(r\geq1\)，证明
\[
I:J^\infty=\bigcap_{i=1}^{r}(I:f_i^\infty).
\]
\item 在 \(k[x,y]\) 中取 \(I=(xy)\)、\(J=(x,y)\)，计算 \(I:J^\infty\)，并与 \(I\) 关于乘法集 \(S=\{\,x^ay^b\mid a,b\geq0\,\}\) 的饱和比较。
\item 若 \(J\) 有限生成，证明 \((I:J^\infty):J^\infty=I:J^\infty\)。
\item 设 \(k\) 为域，取 \(R=k[x_1,x_2,\ldots]\)、\(J=(x_1,x_2,\ldots)\)，并令
\[
I=(\,x_ix_j\ (i\ne j),\ x_i^{\,i+1}\ (i\ge1)\,).
\]
证明 \(I:J^\infty=J\)，而 \((I:J^\infty):J^\infty=R\)。
\end{enumerate}
\end{exercise}
\begin{solution}
1. 由 \(J^{n+1}\subseteq J^n\)，有 \(I:J^n\subseteq I:J^{n+1}\)。各项都是\cref{b1:prop:ideal-quotient-rules}所给的理想，因此并集也为理想：其中任意两个元素同时属于某个足够靠后的项，相减仍在该项；乘以任意环元素也留在原项。

2. 若 \(a\in I:J^\infty\)，取 \(n\geq0\) 使 \(aJ^n\subseteq I\)。因为 \(f_i^n\in J^n\)，所以 \(af_i^n\in I\)，即 \(a\in I:f_i^\infty\)，这对每个 \(i\) 都成立。

反过来，设 \(a\) 属于右端的交。对每个 \(i\)，取整数 \(n_i\geq0\) 使 \(af_i^{n_i}\in I\)。若某个 \(n_i=0\)，则 \(a\in I\)，结论已经成立。否则全部 \(n_i\geq1\)，令
\[
N=1+\sum_{i=1}^{r}(n_i-1).
\]
理想 \(J^N\) 由所有 \(f_1^{e_1}\cdots f_r^{e_r}\)、\(e_1+\cdots+e_r=N\) 生成。这来自把 \(J\) 的生成元代入 \(N\) 个因子的乘积并按分配律展开。指数 \(N\) 的选择用到了鸽巢原理：若每一类因子 \(f_i\) 至多出现 \(n_i-1\) 次，总因子数至多为 \(N-1\)。所以每个这样的单项式必有某个 \(e_i\geq n_i\)。因此 \(a\) 乘以该单项式后含有因子 \(af_i^{n_i}\)，属于 \(I\)。取这些生成元的任意线性组合，可知 \(aJ^N\subseteq I\)，即 \(a\in I:J^\infty\)。有限个生成元使这一步能够选取共同的指数 \(N\)，从分别控制各个 \(f_i\) 转为控制整个幂理想 \(J^N\)。

3. 在给定例子中，\(I:x^\infty=(y)\)、\(I:y^\infty=(x)\)：例如 \(x^na\in(xy)\) 时令 \(y=0\)，得到 \(x^n\cdot a(x,0)=0\)，整环 \(k[x]\) 中便有 \(a(x,0)=0\)，即 \(y\mid a\)；反向包含由 \(xy\in I\) 立即得到。交换变量给出另一式。于是
\[
I:J^\infty=(x)\cap(y)=(xy)=I.
\]
最后一个交的等式也可直接核对：若 \(a=xu\in(y)\)，令 \(y=0\) 得 \(x\cdot u(x,0)=0\)，所以 \(y\mid u\)，从而 \(xy\mid a\)；相反包含显然成立。然而 \(xy\in I\cap S\)，所以 \(I\) 关于 \(S\) 的饱和包含 \(1\)，等于全环。这两种饱和的条件不同：\(I:J^\infty\) 要求某次幂理想 \(J^n\) 的所有元素都把 \(a\) 送入 \(I\)，而关于 \(S\) 的饱和只要求 \(S\) 中存在一个元素做到这一点。

4.\ 若 \(J\) 有限生成，记 \(K=I:J^\infty\)。显然 \(K\subseteq K:J^\infty\)。反过来，若 \(a\in K:J^\infty\)，则存在 \(m\geq0\) 使 \(aJ^m\subseteq K\)。取 \(J^m\) 的非空有限生成组 \(u_1,\ldots,u_t\)；若 \(J^m=0\)，可取 \(u_1=0\)。对每个 \(i\)，由 \(au_i\in K\) 取 \(n_i\) 使 \(au_iJ^{n_i}\subseteq I\)。令 \(N=\max_i n_i\)，便有 \(aJ^{m+N}\subseteq I\)，所以 \(a\in K\)。因此
\[
(I:J^\infty):J^\infty=I:J^\infty\qquad(J\text{ 有限生成}).
\]

5.\ 在题目给定的无限变量环中，对每个 \(i\)，理想 \(J^i\) 的单项式生成元若含有 \(x_j\)（\(j\ne i\)），则乘以 \(x_i\) 后含有 \(x_ix_j\)；其余只有 \(x_i^i\)，乘以 \(x_i\) 后为 \(x_i^{i+1}\)。故 \(x_iJ^i\subseteq I\)，从而 \(J\subseteq I:J^\infty\)。

还须验证 \(1\notin I:J^\infty\)。显然 \(1\notin I=I:J^0\)。对任意 \(n\ge1\)，取 \(j\ge n\)，则 \(x_j^n\in J^n\)。考虑代入同态
\[
\psi_j:R\longrightarrow k[t],\qquad
x_j\longmapsto t,\quad x_i\longmapsto0\quad(i\ne j).
\]
混合生成元均映为零，纯幂生成元只有 \(x_j^{j+1}\) 映为 \(t^{j+1}\)，故 \(\psi_j(I)=(t^{j+1})\)。若 \(x_j^n\in I\)，便有 \(t^n\in(t^{j+1})\)，与 \(n\le j\) 矛盾。因此 \(J^n\nsubseteq I\)，即 \(1\notin I:J^n\)。

原点求值映射 \(R\to k\) 的核为 \(J\)，所以 \(J\) 极大。结合 \(J\subseteq I:J^\infty\subsetneq R\)，可得 \(I:J^\infty=J\)。再次饱和却有 \(J:J^\infty=R\)，因为 \(1\cdot J\subseteq J\)。因此这里关于理想的饱和不满足幂等性，有限生成假设不能从第4问中删去。失败恰在共同指数的选取：每个生成元 \(x_i\) 都有某个指数使 \(x_iJ^i\subseteq I\)，却不存在一个 \(N\) 使整个 \(J\) 都满足 \(J\cdot J^N\subseteq I\)。第2、4问中的有限生成性，正是用来把有限多个所需指数统一成一个共同指数。
\end{solution}

\begin{exercise}\label{b1:ex:13-8}
设 \(f,g\in R\) 满足 \((f,g)=R\)，\(I\) 为交换环 \(R\) 的理想。证明对每个 \(a\in R\)，
\[
a\in I\quad\Longleftrightarrow\quad
a/1\in IR_f\text{ 且 }a/1\in IR_g.
\]
说明若去掉 \((f,g)=R\)，结论可能失败。
\end{exercise}
\begin{solution}
必要性来自扩张的定义。若两个局部成员条件均成立，由\eqref{b1:eq:localized-ideal-membership}，存在 \(m,n\geq0\) 使 \(f^ma\in I\)、\(g^na\in I\)。若 \(m=0\) 或 \(n=0\)，已经得到 \(a\in I\)。否则取 \(u,v\in R\) 使 \(uf+vg=1\)，并把
\[
1=(uf+vg)^{m+n-1}
\]
按分配律展开。每个项都含有至少 \(m\) 个 \(f\) 或至少 \(n\) 个 \(g\)，因为若两者都不足，总因子数至多为 \(m+n-2\)。故存在 \(c,d\in R\) 使 \(1=cf^m+dg^n\)。两边乘以 \(a\)，得 \(a=c(f^ma)+d(g^na)\in I\)。

若 \(R=\mathbb Z\)、\(f=g=2\)、\(I=2\mathbb Z\)、\(a=1\)，两个局部化都是 \(\mathbb Z[1/2]\)，其中 \(I\) 的扩张为全环，两个局部成员条件均成立；但 \(1\notin2\mathbb Z\)。因此分母所满足的生成条件不能省去。
\end{solution}

